\documentclass[11pt,a4paper]{article}
\usepackage[utf8]{inputenc}
\usepackage[T1]{fontenc}
\usepackage[czech]{babel}
\usepackage{lmodern}
\usepackage[margin=2.2cm]{geometry}
\usepackage{amsmath,amssymb,amsthm,mathtools}
\usepackage{physics}
\usepackage[dvipsnames]{xcolor}
\usepackage{tikz}
\usetikzlibrary{arrows.meta,calc,decorations.pathmorphing,decorations.pathreplacing,patterns,positioning,babel}
\usepackage{enumitem}
\usepackage[normalem]{ulem}
\usepackage[skins,breakable,theorems]{tcolorbox}
\usepackage{titlesec}
\usepackage[colorlinks=true,linkcolor=NavyBlue,urlcolor=NavyBlue]{hyperref}

%======================================================================
% Barvy
%======================================================================
\definecolor{hlavni}{RGB}{0,70,140}      % nadpisy
\definecolor{defbarva}{RGB}{30,100,200}  % definice
\definecolor{vetabarva}{RGB}{190,40,40}  % věty, postuláty
\definecolor{prbarva}{RGB}{20,130,70}    % příklady
\definecolor{vzorbarva}{RGB}{220,120,0}  % důležité vzorce

\titleformat{\section}{\color{hlavni}\Large\bfseries}{\thesection}{1em}{}[{\color{hlavni}\titlerule[1pt]}]
\titleformat{\subsection}{\color{hlavni}\large\bfseries}{\thesubsection}{1em}{}
\titleformat{\subsubsection}{\color{hlavni!80!black}\normalsize\bfseries}{\thesubsubsection}{1em}{}

%======================================================================
% Zvýraznění
%======================================================================
% Definice (modrý box)
\newtcolorbox{definice}[1][]{enhanced,breakable,colback=defbarva!6,colframe=defbarva,
  boxrule=0.6pt,left=4pt,right=4pt,top=2pt,bottom=2pt,arc=2pt,
  fonttitle=\bfseries\small,coltitle=white,title={Definice\ifx&#1&\else: #1\fi}}
% Věta / postulát / lemma / klíčové tvrzení (červený box)
\newtcolorbox{veta}[1][Věta]{enhanced,breakable,colback=vetabarva!5,colframe=vetabarva,
  boxrule=0.6pt,left=4pt,right=4pt,top=2pt,bottom=2pt,arc=2pt,
  fonttitle=\bfseries\small,coltitle=white,title={#1}}
% Příklad (zelený proužek vlevo)
\newtcolorbox{priklad}[1][Příklad]{enhanced,breakable,colback=prbarva!4,colframe=prbarva,
  boxrule=0pt,leftrule=2.5pt,left=5pt,right=3pt,top=2pt,bottom=2pt,sharp corners,
  fonttitle=\bfseries\small,coltitle=prbarva!80!black,colbacktitle=prbarva!12,title={#1}}
% Důležitý vzorec v rámečku (v matematice):  \vzorec{E = mc^2}
\newcommand{\vzorec}[1]{\tcbhighmath[colback=vzorbarva!8,colframe=vzorbarva,boxrule=0.8pt,arc=2pt,left=3pt,right=3pt,top=2pt,bottom=2pt]{#1}}
% Klíčový pojem v textu
\newcommand{\kw}[1]{\textbf{\color{hlavni}#1}}
% Upozornění / pozor
\newcommand{\pozor}[1]{{\color{vetabarva}\bfseries #1}}

% Poznámky psané tužkou (v originále slabě, šedě)
\newcommand{\pozn}[1]{\relax\ifmmode\text{\color{gray}\small\itshape [#1]}\else{\color{gray}\small\itshape [#1]}\fi}
% Číslo strany originálních poznámek
\newcommand{\strana}[1]{\marginpar{\colorbox{hlavni!10}{\small\texttt{\color{hlavni}#1}}}}

\setlength{\parindent}{0pt}
\setlength{\parskip}{4pt}


\title{Kvantová mechanika I.\ (2014/2015)\\[2pt]\large přepis ručně psaných poznámek}
\date{}

\begin{document}
\maketitle
\tableofcontents
\newpage

%======================================================================
% QM-I-2014-1.pdf
%======================================================================
\section*{0. Organizační poznámky}
\addcontentsline{toc}{section}{0. Organizační poznámky}
\strana{QM-I-1}

\begin{tabular}{lll}
St & 9:00 -- 10:30 & P \\
   & + 10:40 -- 12:10 & Cv \\
Čt & 10:40 -- 12:10 & P
\end{tabular}
\quad $\Big\}$ většinou

\textbf{Já:} \underline{kontaktní informace} \dots\ viz web

\textbf{Zápočet:} cvičení \dots\ DÚ $= 50$\,b; koncem semestru písemka $= 50$\,b
\begin{itemize}[nosep]
  \item zápočet $\geq 50$\,b
  \item bez zkouškového příkladu $\geq 75$
  \item tričko s QM tematikou $\geq 90$\,b
\end{itemize}

\subsection*{Literatura}
\begin{enumerate}[label=\arabic*)]
  \item Pavel Cejnar: \emph{A Condensed Course of QM}; Karolinum 2013
  \item J.\ Formánek: \emph{Úvod do kvant.\ teorie} (Academia 1983, 2004)
  \item L.\ E.\ Ballentine: \emph{QM a Modern Development}; \\
        J.\ J.\ Sakurai: \emph{Modern QM} \hfill $\Big\}$ moderní přístup
  \item L.\ D.\ Landau, E.\ M.\ Lifschitz: \emph{QM Nonrelativistic Theory}; \\
        Cohen-Tannoudji, Diu, Laloe: \emph{Q.M.} \hfill $\Big\}$ klasické
  \item \emph{Cvičení:}
  \begin{itemize}[nosep]
    \item S.\ Flügge: \emph{Practical QM}
    \item \sout{J.\ Pišút, V.\ Čer} Galitsky, Karnakov, Kogan: \emph{Exploring QM} (2013 Oxford) -- Collection of 700+ Solv.\ Prob.
    \item + příklady v jiných učebnicích
  \end{itemize}
\end{enumerate}

\subsection*{Doporučené doplňkové přednášky}
\begin{itemize}[nosep]
  \item Pavel Krtouš: NTMF036: \emph{Interpretace kvant.\ mechaniky} \dots\ každé 2 roky \dots\ letos Čt 8:10 -- 10:30
  \item Karel Houfek: NTMF061: \emph{Teor.\ grup a její aplikace ve fyz.} \dots\ letos Pá 8:10 -- 11:20
\end{itemize}

\subsection*{Diskuse: \uv{Co znáte z úvodu do QM?}}
základní \uv{postuláty}, neurčitost, rovinná vlna, $\vb{L}^2, L_z$, LHO, H, jáma, SR, kontinuita (a pravd.\ interpret.), spin, \dots

\uv{Co znáte z matematiky?} \quad mat.\ statistika?

LVP + skal.\ součin, Fce $\mathbb{C}$-proměnné, zákl.\ integ.\ a difro \pozn{$\leftarrow$ předpoklad}

teorie FFT + řady, distribuce, $\mathcal{H}$-prostor, $\ell^2$, $\mathcal{L}^2$ \pozn{$\leftarrow$ udělám úvod (minimalistický)}

%----------------------------------------------------------------------
\section{Úvodní poznámky}
\strana{QM-I-2}

Doporučená interpretace QM: \\
\quad \uv{Shut up and calculate} \quad N.\ D.\ Mermin, Phys.\ Today Apr.\ 1989
\pozn{Ashcroft, Mermin: Solid State Phys.} \\
mylně připisováno Feynmanovi \pozn{+ BBT 2?.9.}

\subsection*{Plán přednášky}
\begin{itemize}[nosep]
  \item základní formalismus a mat.\ aparát QM
        \pozn{pragmatický přístup. Exaktní mat. \uv{fousy} někdy vynecháme}
  \item příklady konkrét.\ syst.\ (\sout{sada řeš})
  \begin{itemize}[nosep]
    \item částice se spinem
    \item bodová částice v 1D a 3D \dots\ odvození ze symetrie
          \pozn{příklad moderního způsobu budování teorií}
  \end{itemize}
  \item základní řešitelné modely v 1D a 3D
        \pozn{smysl práce fyzika: black-box heavy calcul $\times$ simple model; zásobárna modelů + příklady jejich aplikace}
  \item časový vývoj
  \item úvod do teorie rozptylu
        \pozn{sofistik.\ použití čas.\ záv./nezáv.\ přístupu, problémy se spoj.\ částí spektra}
  \item částice v mg.\ poli
\end{itemize}

NTMF067 \dots\ Kvantová mechanika II \pozn{příští semestr}: \\
přibližné metody, vícečásticové systémy, matice hustoty a otevřené systémy.

\subsection*{Základní aspekty QM chování}
\begin{itemize}
  \item \underline{diskrétní charakter} \pozn{některých veličin, které jsou v klas.\ mechanice spojité}
\end{itemize}

\kw{Planckova konstanta} \quad
$\vzorec{\hbar = \begin{cases} 1.05\cdot 10^{-34}\ \mathrm{J\cdot s} \\ 0.66\ \mathrm{eV\cdot fs}\end{cases}}$

$\to$ rozměr jako \quad Moment hybnosti \dots\ $\vec{L} = \vec{x}\times\vec{p}$ \\
\phantom{$\to$ rozměr jako} \quad Akce \quad $S = \int L\,\dd t$ \quad $\big(= \int p\,\dd q - Et\big)$
\pozn{čas.\ nezáv.\ V} \\
\phantom{$\to$ rozměr jako \quad Akce} \quad \underline{Hamilton} \quad Maupertuis \quad princip

\pozn{kvantování ELMG-pole -- to tady dělat nebudeme}

\begin{itemize}[nosep]
  \item 1900 Planck (záření čer.\ tělesa), \quad 1905 Einstein (fotoefekt)
  \item 1913 Bohr (model atomu) \dots \quad 1914 Franck + Hertz \quad
        $\mathrm{Hg} + e \to \mathrm{Hg}^* + e$ \quad ($4.9$\,eV)
  \item 1922 Stern--Gerlach experiment \pozn{k atomům se vrátíme příští semestr}
\end{itemize}
\pozn{prototyp kvantového měření DISKRÉTNÍ SYSTÉM -- rozebereme podrobněji:}

\strana{QM-I-3}
\pozn{dva stab.\ izotopy 107 a 109}

\begin{center}
\begin{tikzpicture}[scale=0.9]
  % pec
  \draw (0,0) rectangle (1.6,1);
  \node at (0.8,0.5) {\color{gray}Pec};
  % clona
  \draw (2.3,-0.6) -- (2.3,0.35);
  \draw (2.3,0.65) -- (2.3,1.6);
  % svazek
  \draw[-{Stealth}] (2.5,0.5) -- (6.2,0.5);
  \node[above] at (3.3,0.55) {Ag};
  \node[above] at (4.6,0.9) {${}^2S_{1/2}$ stav};
  \node[below] at (3.6,0.2) {$[\mathrm{Kr}]\,4d^{10}5s^1$};
  \draw[decorate,decoration={brace,mirror,amplitude=4pt}] (2.6,-0.5) -- (4.6,-0.5);
  \node at (3.6,-1) {$47\,e^-$};
  % magnet S (klínový) a N
  \draw (6.4,1.6) -- (6.4,1.1) -- (6.8,0.75) -- (8.4,0.75) -- (8.8,1.1) -- (8.8,1.6);
  \node at (7.6,1.3) {S};
  \draw (6.4,-0.6) -- (6.4,0.2) -- (8.8,0.2) -- (8.8,-0.6);
  \node at (7.6,-0.2) {N};
  \draw[-{Stealth}] (7.6,0.25) -- (7.8,0.7) node[right] {$\vec B$};
  % rozštěpení
  \draw[gray,dashed,-{Stealth}] (8.9,0.5) -- (10.2,1.1);
  \draw[gray,dashed,-{Stealth}] (8.9,0.5) -- (10.2,-0.1);
  % souřadnice
  \draw[-{Stealth}] (11,0) -- (11,1.3) node[right] {$z$};
  \draw[-{Stealth}] (11,0) -- (12.2,0) node[right] {$x$};
  \draw[-{Stealth}] (11,0) -- (11.6,0.45) node[right] {$y$};
\end{tikzpicture}
\end{center}

energie \dots\ dipól v MG poli \qquad $E = -\vec\mu\cdot\vec B$ \\
\phantom{energie \dots\ dipól v MG poli} \qquad
$F_z = \pdv{z}\big(\vec\mu\cdot\vec B\big) = \mu_z\cdot\pdv{B_z}{z}$
\pozn{Bohr magneton}

naměřené \quad $\mu = \pm\dfrac{\hbar}{2}\,\dfrac{e}{2m_e}$ \quad \dots\ odpovídá momentu hybnosti \quad $S_z = \pm\frac12\hbar$ \\
\pozn{nuclear spin nepřispívá $m_p \gg m_e$}

\textbf{Poznámky:}
\begin{itemize}[nosep]
  \item prototyp měření \dots\ spin \dots\ měřený systém; \\
        prostorové souřadnice atomu -- součást aparátu
  \item výsledek měření v rozporu s klas.\ představou \dots\ diskrét.\ hodnoty
  \item budeme používat jako ideal.\ příklad později
\end{itemize}

\begin{itemize}
  \item \underline{interference a difrakce}
\end{itemize}
dvouštěrbinový experiment \pozn{Feynman a $\heartsuit$ QM}

\begin{center}
\begin{tikzpicture}[scale=0.75]
  % (c) obě štěrbiny
  \node at (0,3.2) {\textcircled{c}};
  \draw (0,1.2) -- (0,1.6) -- (0.3,1.6) -- (0.3,1.2) -- cycle;
  \draw (0,1.3)--(0.3,1.3) (0,1.4)--(0.3,1.4) (0,1.5)--(0.3,1.5);
  \draw[thick] (2,0) -- (2,0.6); \draw[thick] (2,0.9) -- (2,2.1); \draw[thick] (2,2.4) -- (2,3);
  \node[gray] at (1.75,2.55) {\small A}; \node[gray] at (1.75,0.45) {\small B};
  \draw[thick] (4.5,0) -- (4.5,3);
  \draw[gray,dashed] (0.3,1.4) -- (2,2.25) -- (4.5,2);
  \draw[gray,dashed] (0.3,1.4) -- (2,0.75) -- (4.5,2);
  \node[gray] at (4.75,2) {\small x};
  \node at (2,-0.4) {$z_1$}; \node at (4.5,-0.4) {$z_2$};
  \node at (5.6,1.5) {\Large$\neq$};
  \draw (5.6,1.5) circle (0.45);
  % (b) jen A
  \node at (6.8,3.2) {\textcircled{b}};
  \draw[thick] (7.2,0) -- (7.2,2.1); \draw[thick] (7.2,2.4) -- (7.2,3);
  \node[gray] at (7.5,2.6) {\small A}; \node[gray] at (7.5,0.6) {\small B};
  \draw[thick] (9.4,0) -- (9.4,3);
  \draw[gray] (7.2,2.25) -- (9.4,2);
  \node at (10.5,1.5) {\Large$+$};
  \draw (10.5,1.5) circle (0.45);
  % (a) jen B
  \node at (11.5,3.2) {\textcircled{a}};
  \draw[thick] (12,0) -- (12,0.6); \draw[thick] (12,0.9) -- (12,3);
  \node[gray] at (12.3,2.6) {\small A}; \node[gray] at (12.3,0.6) {\small B};
  \draw[thick] (14.2,0) -- (14.2,3);
  \draw[gray] (12,0.75) -- (14.2,2);
  % osy
  \draw[-{Stealth}] (15.3,1) -- (15.3,2.2) node[right] {$x$};
  \draw[-{Stealth}] (15.3,1) -- (16.5,1) node[right] {$z$};
  \draw[-{Stealth}] (15.3,1) -- (15.8,0.4) node[below] {$y$};
\end{tikzpicture}
\end{center}

např.\ rychlé elektrony 50\,keV \dots\ rychl.\ $v$
\begin{itemize}[nosep]
  \item lze chápat jako měření polohy $x, y$; $z$ součástí měř.\ aparatury \\
        $\to$ příprava systému v $t_1 = z_1/v$, \quad měření v $t_2 = z_2/v$
\end{itemize}

experiment \textcircled{a} příprava ve stavu \quad $\psi_B(t_1) \to \psi_B(t_2)$ \quad nalezneme na stínítku \\
\phantom{experiment} \textcircled{b} \phantom{příprava ve stavu} \quad $\psi_A(t_1) \to \psi_A(t_2)$

analogie s optikou (Youngův experiment, 1805) \\
\quad $\to$ signál na stínítku \quad $|\psi_A|^2$; \quad $|\psi_B(t_2)|^2$

fáze je dána klas.\ akcí \quad $S = \displaystyle\int_{t_0}^{t_2} p(q)\,\dd q$ \quad podél dráhy částice \\
(upřesníme později \dots\ WKB aproximace; dráhový integrál)
\[
  \to\quad \psi(x,t) = \sqrt{\rho(x,t)}\;e^{\,iS(x,t)/\hbar}
  \qquad\text{pozorovatelnost interference } \Delta S \sim \hbar
\]

%----------------------------------------------------------------------
\strana{QM-I-4}
popis interference \quad signál na stínítku $\sim |\psi|^2 = |\psi_A + \psi_B|^2$ \\
\quad v rozporu s klas.\ zákony pravděpodobnosti: \quad $\neq |\psi_A|^2 + |\psi_B|^2$

vlnově-částic.\ dualismus \dots\ elektron musel v nějakém smyslu projít oběma štěrbinami, ale vždy pozorujeme celý elektron \\
$\to$ podobné úvahy vedou k zdánlivým rozporům \quad $\to$ viz MOTTO

Experimentálně prověřeno: elektrony 1961; 1970.\ léta -- individ.\ $e^-$ \\
\quad 1994--5 \dots\ $\mathrm{He}_2, \mathrm{Na}_2, \mathrm{I}_2$; \quad 1999: $\mathrm{C}_{60}$; \quad organ.\ mol.\ s 810 atomy: 2013
\pozn{$\uparrow$ i dnes jednotlivé molekuly}

\begin{minipage}{0.5\textwidth}
Podobné experimenty \dots\ difrakce na mřížce, \\
Mach--Zehnderův interferometr
\end{minipage}
\begin{minipage}{0.45\textwidth}
\begin{tikzpicture}[scale=0.8]
  % vstupní svazek
  \draw[double] (-1.5,0) -- (0,0);
  % BS dole vlevo
  \draw[thick] (-0.3,-0.3) -- (0.3,0.3); \node[below] at (0,-0.3) {\small BS};
  % zrcadlo dole vpravo
  \draw[thick] (2.2,-0.3) -- (2.8,0.3); \node[right] at (2.8,0) {\small M};
  \draw (0,0) -- (2.5,0) -- (2.5,2);
  % zrcadlo nahoře vlevo
  \draw[thick] (-0.3,1.7) -- (0.3,2.3); \node[left] at (-0.3,2.1) {\small M};
  \draw (0,0) -- (0,2) -- (2.5,2);
  % BS nahoře vpravo
  \draw[thick] (2.2,1.7) -- (2.8,2.3); \node[above right] at (2.6,2.2) {\small BS};
  \draw[-{Stealth}] (2.5,2) -- (3.5,2) node[right] {$D_1$};
  \draw[-{Stealth}] (2.5,2) -- (2.5,2.9) node[above] {$D_2$};
\end{tikzpicture}
\end{minipage}

+ další nuance: delayed choice, kvant.\ guma

\begin{itemize}
  \item \underline{statistický charakter}
        \pozn{\dots\ individ.\ výsledek nepředpověditelný -- velký soubor ano} \\
        \dots\ viz výše \dots\ srovn.\ Stern--Gerlach \\
        až na výjimky jen pravděpodobnosti, ne jistota \pozn{(integrály pohybu)}
  \item \underline{vliv měření na stav} \dots\ kolaps vln.\ funkce \\
        \dots\ měření modifikuje časový vývoj \\
        detailní rozbor viz \dots\ Pavel Krtouš: Interpretace QM, NTMF036 \\
        \pozn{(modely měření, kvantová guma, EPR, Bellovy nerovnosti, teleportace \dots)}
\end{itemize}

\textbf{DNES:} zlatý věk kvantové technologie! \pozn{(osobní výběr)}
\begin{itemize}[nosep]
  \item experimenty testující fundament.\ základy \pozn{(dříve jen thought experiments)}
  \item kvantové počítání, kryptografie, \dots
  \item mnohočásticová fyzika:
  \begin{itemize}[nosep]
    \item kompletní měření \uv{state-to-state} rozptylu \quad ($\to$ \uv{měření vln.\ fce})
          \pozn{Nobel Chem.\ 1999}
    \item časově rozlišená dynamika \dots\ femto \& atto pulzy
    \item experimenty s extrémně chladnými plyny (dnes i molekuly) \pozn{Nobel 1997, 2001} \\
          \dots\ BEC, atom-laser, optical lattices, zastav.\ světlo, metrologie
    \item transportní jevy -- formulace nerovnovážné stat.\ fyziky
    \item nanotechnologie \dots\ kvantové stroje, využití biomol., STM + manipulace
    \item graphene, nanotrubičky, Q-Hall jev \pozn{Nobel 2010}
    \item vysokoteplotní supravod.
    \item \uv{ab initio} výpočty v Q-chemii, solid state, rozptyl \pozn{Nobel 98 (Chem)}
    \item Jaderná a částic.\ fyzika \pozn{Higgs Nobel 2013, 1999; Nobel dynamika Chem.\ 2013}
  \end{itemize}
\end{itemize}

%======================================================================
% QM-I-2014-2.pdf
%======================================================================
\section{Formalismus kvantové teorie 1}
\strana{QM-I-5}

Začneme s konečnými LVP, nekonečně dim.\ Hilbert.\ p.\ později.

\subsection{Opakování LVP (komentované s ohledem na QM)}

\kw{Lineární vektor.\ prostor (LVP)} $=$ \kw{Prostor stavů} (ket space)

\begin{definice}[LVP]
LVP $= \{\text{vektorů } \phi\}$ s operacemi $+$; násob.\ $\mathbb{C}$ číslem, \\
tj.\ uzavřený vůči těmto operacím:
\[
  \vzorec{\forall\, \phi,\psi\in V \;\Rightarrow\; a\phi + b\psi \in V}
  \qquad\dots\quad \text{\kw{Princip superpozice}}
\]
\pozn{(kočky pana Schrödingera -- viz interpretace)}
\end{definice}

další algeb.\ požadavky (axiomy): asoc.\ $+$; komut.\ $+$, $\exists\, 0 + \phi = \phi$, $\exists\, {-v}$, \\
distrib.\ pro vektory i čísla, asoc.\ $\cdot\,\mathbb{C}$, $\exists\, 1\cdot\phi = \phi$

\underline{zopakujte:} lineární závislost, dimenze, báze

\begin{definice}[Skalární součin]
\kw{Skalární součin} (inner product) \dots\ amplitudy pravděpod. \\
\quad $(\psi|\phi) \in \mathbb{C}$ \quad splňující:
\begin{enumerate}[label=\arabic*)]
  \item $(\psi|\phi) = (\phi|\psi)^*$
  \item $(\phi|c_1\psi_1 + c_2\psi_2) = c_1(\phi|\psi_1) + c_2(\phi|\psi_2)$ \hfill linearita ve 2.\ argumentu \\
        1)+2) $\Rightarrow$ \quad $(c_1\phi_1 + c_2\phi_2|\psi) = c_1^*(\phi_1|\psi) + c_2^*(\phi_2|\psi)$ \hfill antilinearita v 1.\ argum.
  \item $(\phi|\phi) \geq 0$; \quad rovnost $\Leftrightarrow \phi = 0$ \hfill \dots\ posit.\ definitnost \\
        \pozn{pozn: norma $\|\phi\| = \sqrt{(\phi|\phi)}$} \hfill $\to$ důležitá pro pravděpod.\ int.
\end{enumerate}
\end{definice}

\underline{zopakujte:} ON báze \quad $\{\phi_i,\ i = 1\dots d\}$ \qquad $(\phi_i,\phi_j) = \delta_{ij}$ \\
\quad rozklad \quad $\forall\psi\ \exists c_i:\ \psi = \sum_i c_i\phi_i$ \quad přitom $c_i = \braket{\phi_i}{\psi}$ \\
\quad tj.\ $\forall$ konečný LVP je izomorfní s $\mathbb{C}^d$ \\
\quad \pozn{$\to$ připomenout zápis skalár.\ souč.\ v bázi $(\phi|\psi) = \sum f_i^* c_i$}

\begin{definice}[Duální prostor]
\kw{Duální prostor $V^*$}: $\equiv$ prostor $\forall$ lin.\ funkcionálů nad $V$ \hfill (bra space) \\
\pozn{(pozor: v $\infty$ dim budeme přísnější)}
\end{definice}

v konečné dim.\ je $V$ a $V^*$ izomorfní, tj.\ $\mathbb{C}^d$

\textbf{DK:} označme $F(\phi_i) = f_i^*$, \quad pak $\forall\, F(\psi) = f_i^* c_i \equiv (\phi|\psi)$ \\
\quad kde $\ket{\phi} = \sum f_i\ket{\phi_i}$ \quad \dots\ tj.\ libovolný lin.\ funkcionál lze reprezentovat skalárním součinem \\
\pozn{v $\infty$ dim musíme dávat bacha na konvergenci $\sum_{1}^{\infty}$ -- později}

značení (matematické): funkcionál odpovídající vektoru $\phi$:
\[ F_\phi(\psi) = (\phi,\psi) \]

%----------------------------------------------------------------------
\strana{QM-I-6}
\subsection*{Diracova notace}
ket vektory z $V$: $\ket{\psi}$ \quad \dots\ v konečné dim.\ slop.\ vekt.\ $\begin{pmatrix} c_1\\ \vdots\\ c_d\end{pmatrix}$

bra vektory z $V^*$: $\bra{\psi}$ \quad \dots\ v konečné dimenzi řád.\ vektory \\
\quad $\bra{\psi}$ vlastně odpovídá působení funkcionálu $F_\psi$, tj.\ řádkovému vektoru $(c_1^*,\dots,c_d^*)$

\underline{poznámka:} skalární součin $\braket{\phi}{\psi}$ můžeme číst jako maticové násobení řád.\ matice $1\times d$ a slop.\ matice $d\times 1$ \\
\quad přechod $\ket{\psi}\to\bra{\psi} = \ket{\psi}^\dagger$ lze chápat jako hermit.\ sdružení \\
\quad \pozn{pozn: $\big(c_1\ket{\psi_1} + c_2\ket{\psi_2}\big)^\dagger = c_1^*\bra{\psi_1} + c_2^*\bra{\psi_2} \leftrightarrow F_{(c_1\psi_1 + c_2\psi_2)}$}

\subsection*{Lineární operátory}
\hfill \underline{pozorovatelné (měřit.) veličiny} a \underline{transformace} stavů

Lineární operátor $\hat A: V \to V$ \\
\quad tj.\ $\forall\ket{\psi}\in V\ \exists\ket{\phi} \equiv \hat A\ket{\psi} \in V$ \quad \pozn{(v $\infty$ dim opatrnější)} \\
\quad linearita \quad $\hat A\big(c_1\ket{\phi_1} + c_2\ket{\phi_2}\big) = c_1\hat A\ket{\phi_1} + c_2\hat A\ket{\phi_2}$

\dots\ ON báze + linearita \dots\ $\hat A$ jednoznačně zadán čísly $a_{ij} \equiv \mel{\phi_i}{\hat A}{\phi_j}$ \\
\quad v konečné dimenzi lze ztotožnit $\hat A$ a matici $a_{ij}$ \\
$\to$ Diracova notace \quad $\mel{\phi}{\hat A}{\psi}$; \ $\hat A\ket{\psi}$; \ $\bra{\psi}\hat A$ \quad \dots\ vše maticové násobení
\pozn{např.\ $A\ket{\psi} = \sum_j a_{ij}c_j$}

\begin{itemize}
  \item rovnost operátorů \quad $\hat A = \hat B \Leftrightarrow \hat A\ket{\phi} = \hat B\ket{\phi}\ \ \forall\ket{\phi}$
  \item nulový a jednotkový operátor $\hat 0, \hat I$ \quad \pozn{(často prostě $0, 1$ $\to$ z kontextu)}
  \item sčítání operátorů: \quad $\hat A + \hat B$, \qquad $(\hat A + \hat B) + \hat C = \hat A + (\hat B + \hat C)$
  \item násobení operátorů: \quad $\hat A\hat B$ \dots\ skládání zobr., či násobení matic \\
        pozor \quad $\hat A\hat B \neq \hat B\hat A$ (obecně) \quad \dots\ def.\ $[\hat A,\hat B] = \hat A\hat B - \hat B\hat A$ \\
        obvyklý trik: \quad $\hat A\hat B = \hat B\hat A + [\hat A,\hat B]$
  \item funkce operátoru \quad $f(\hat A) = \sum_n f_n\hat A^n$ \quad \dots\ např.: $e^{\hat A} = \sum_n \frac{1}{n!}\hat A^n$
\end{itemize}

pozor: vektoru $A\ket{\psi}$ neodpovídá lin.\ funkcionál $F_{A\ket{\psi}} \neq \bra{\psi}A$ \\
\phantom{pozor: vektoru} ale \quad $F_{A\ket{\psi}} = \bra{\psi}A^\dagger$ \\
\quad kde $A^\dagger$ je hermitovsky sdružená matice $(A^\dagger)_{ij} = a_{ji}^*$

%----------------------------------------------------------------------
\strana{QM-I-7}
\subsection*{Sdružený operátor (adjoint)}
\begin{definice}[Sdružený operátor]
matematická definice: \quad $B = A^\dagger$ \quad pokud \quad $(B\phi|\psi) = (\phi, A\psi)$ \quad $\forall\phi,\psi\in V$ \\
pozn: v $\infty$ dim to komplikují def.\ obory \dots\ později
\end{definice}

v kon.\ dimenzích: dosadíme bázi:
\[
  a_{ij} = (\phi_i, A\phi_j) = (B\phi_i, \phi_j) = (\phi_j, B\phi_i)^* = B_{ji}^*
\]
tj.\ opravdu $A^\dagger$ odpovídá hermit.\ sdruž.\ matice \quad $a^\dagger_{ji} = B_{ij}$

Vlastnosti: \quad
$\left.\begin{aligned}
  (cA)^\dagger &= c^*A^\dagger \\
  (A+B)^\dagger &= A^\dagger + B^\dagger \\
  (AB)^\dagger &= B^\dagger A^\dagger
\end{aligned}\right\}$ \quad \pozn{DK viz cvičení} \quad \pozn{pozn: pro matice rovněž platí}

\subsection*{Vnější součin \ $\ket{\psi}\bra{\phi} \equiv$ operátor}
def: \quad $\big(\ket{\psi}\bra{\phi}\big)\ket{u} \equiv \big(\braket{\phi}{u}\big)\ket{\psi}$ \quad ($\braket{\phi}{u}$ \dots\ $\mathbb{C}$-číslo)

$\hookrightarrow$ v Dirac.\ notaci \quad
$\begin{pmatrix} c_1\\ \vdots\\ c_d\end{pmatrix}(f_1^*\cdots f_d^*) =
 \begin{pmatrix} c_1f_1^* & \cdots & c_1f_d^* \\ \vdots & & \vdots \\ c_df_1^* & \cdots & c_df_d^*\end{pmatrix}$

\underline{pozn:} výraz \quad $\ket{\psi}\braket{\phi}{u} = \big(\ket{\psi}\bra{\phi}\big)\ket{u} = \ket{\psi}\big(\braket{\phi}{u}\big)$ \\
\quad lze libovolně uzávorkovat (matic.\ násobení) $\to$ def.\ asociativita

platí \quad $\big(\ket{\psi}\bra{\phi}\big)^\dagger = \ket{\phi}\bra{\psi}$

\underline{pozn:} Rozvoj operátoru do báze lze psát
\[
  \hat A = \hat I\hat A\hat I = \sum_i \ket{\phi_i}\bra{\phi_i}\hat A\sum_j\ket{\phi_j}\bra{\phi_j}
  = \sum_{ij} a_{ij}\ket{\phi_i}\bra{\phi_j} = \sum_j \ket{\hat A\phi_j}\bra{\phi_j}
\]
přitom jsem užil tzv.\ \kw{rozklad jedničky} \quad $\vzorec{\hat I = \sum_n \ket{\phi_n}\bra{\phi_n}}$ \quad (relace úplnosti)

podobně lze odvozovat vyjádření jiných výrazů v bázi, např.
\[
  \braket{\phi}{\psi} = \mel{\phi}{\hat I}{\psi} = \sum_n \underbrace{\braket{\phi}{\phi_n}}_{f_n^*}\underbrace{\braket{\phi_n}{\psi}}_{c_n} = \sum_n f_n^* c_n
\]

\subsection*{(Ortogonální) projekční operátor}
\begin{definice}[Projektor]
\begin{tabular}{lll}
$\hat P^2 = \hat P$ & \pozn{(Idempotence)} & \pozn{\dots\ tj.\ $\ket{b} = P\ket{a}$, pak $P\ket{b} = \ket{b}$} \\
$P = P^\dagger$ & \pozn{(ortogonalita)} & \pozn{\dots\ tj.\ $\ket{f} = P\ket{\psi}$ a $\ket{\bar f} = \ket{\psi} - \ket{f} = (I-P)\ket{\psi}$} \\
 & & \pozn{pak $\braket{\bar f}{f} = 0 \ \Rightarrow\ \mel{\psi}{(I-P^\dagger)P}{\psi} = 0$}
\end{tabular}
\end{definice}

\textbf{PŘ:} \quad $\hat P = \ket{\phi}\bra{\phi}$ \ $\leftarrow$ pokud $\braket{\phi}{\phi} = 1$ \quad \dots\ nebo \quad
$P = \dfrac{1}{\braket{\phi}{\phi}}\ket{\phi}\bra{\phi}$ \pozn{vždy}

%----------------------------------------------------------------------
\strana{QM-I-8}
\subsection*{Samosdružené operátory \hfill \underline{pozorovatelné}}
\begin{definice}[Samosdružený (hermitovský) operátor]
v kon.\ dim \quad $A = A^\dagger$ \quad (prostě maticově) \quad $\to$ \kw{Hermitovské operátory} \\
v $\infty$ dim se tyto dva pojmy mohou lišit \quad
tj.\ $\vzorec{\mel{\phi}{A}{\psi} = \mel{\psi}{A}{\phi}^*\ \ \forall\psi,\phi}$
\end{definice}

\subsection*{Vlastní čísla a vektory}
doporučení: zopakovat lin.\ algebru, podobnostní transf.; diagonalizace

\begin{tabular}{l@{\qquad}l|l}
$\hat A\ket{\phi} = \lambda\ket{\phi}$ & hermit.\ matice & $\to \lambda\in\mathbb{R}$ \\
$\bra{\phi}\hat A = \lambda\bra{\phi}$ & & $\to$ levé a pravé vl.\ v.\ totožné \\
 & $|\hat A - \lambda\hat I| = 0$ & $\to$ vl.\ v.\ pro $\lambda\neq\lambda'$ ortogonální \\
 & {\small charakteristický pol.} &
\end{tabular}

$\to$ lze vybrat bázi z vlastních vektorů

obvyklá notace \quad $\hat A\ket{a} = a\ket{a}$ \quad \dots\ předp.\ normování $\braket{a}{a} = 1$ \\
nebo (degenerace $\equiv$ netriv.\ vlastní podprostor):
\[
  \hat A\ket{a,k} = a\ket{a,k} \qquad k = 1,\dots,d_a;\quad \textstyle\sum d_a = d
\]
\quad $\to$ ON báze \quad $\braket{a,k}{a',k'} = \delta_{aa'}\delta_{kk'}$

platí: \quad $\mel{a,k}{\hat A}{a',k'} = a\,\delta_{aa'}\delta_{kk'}$

$\to$ \kw{spektrální rozklad}
\[
  \vzorec{\hat A = \sum_a a\underbrace{\sum_k \ket{a,k}\bra{a,k}}_{\hat P_a} = \sum_a a\hat P_a}
\]
\quad projektor na vlastní podprostor \quad $\hat P_a = \sum_k\ket{a,k}\bra{a,k}$ \\
\quad (ověřte: $\hat P_a^2 = \hat P_a$; $\hat P_a^\dagger = \hat P_a$) \quad nedegenerované spektrum $\to$ prostě $\hat P_a = \ket{a}\bra{a}$

funkce operátoru: \quad $f(\hat A) = \sum_a f(a)\hat P_a$ \\
\quad \dots\ zobecnění pro případ nekonverg.\ Taylorova rozvoje

\underline{pozn:}
\begin{itemize}[nosep]
  \item projektor \dots\ Hermitovský
  \begin{itemize}[nosep]
    \item[--] vl.\ č.\ jen $0, 1$ \quad \dots\ $P\ket{\psi} = \lambda\ket{\psi} = P^2\ket{\psi} = \lambda^2\ket{\psi}\ \ \forall\ket{\psi}$ \quad $\Rightarrow \lambda^2 = \lambda$
    \item[--] působí jako $0$ na svém jádru, jinde jako $\hat I$
    \item[--] je sám sobě spektrálním rozkladem
    \item[--] $(I - P)$ je projektor na ortogonální doplněk
  \end{itemize}
  \item relace $\hat A = \sum_{ak}\ket{ak}a\bra{ak}$ je vlastně $\hat A = U\Lambda U^\dagger$, \\
        kde $U$ je matice přechodu k bázi z vl.\ v.
\end{itemize}

%======================================================================
% QM-I-2014-3.pdf
%======================================================================
\subsection{Základní principy kvantové teorie (POSTULÁTY)}
\strana{QM-I-9}

\begin{veta}[Postuláty kvantové teorie]
\begin{enumerate}[label=\textcircled{\arabic*}]
  \item Stav systému je popsán vektorem v komplexním LVP (stavový prostor $\equiv\mathcal{H}$), přesněji paprskem $\lambda\psi$, $\forall\lambda\in\mathbb{C}$.
  \item Pozorovatelným veličinám odpovídají Hermitovské lineární operátory. Přípustné hodnoty jsou dány jejich vlastními čísly. \hfill $\hat A\ket{a} = a\ket{a}$ \dots\ $A = \sum a\hat P_a$
  \item Skalární součin na stavovém prostoru definuje tzv.\ amplitudu pravděpodobnosti. Podmíněná pravděpodobnost pro naměření hodnoty $a$ pozorovatelné $\hat A$ systému připraveného ve stavu $\ket{\psi}$ je
  \[
    \vzorec{\mu(a,\psi) \equiv |\braket{a}{\psi}|^2} \qquad
    \Big(= \mel{\psi}{P_a}{\psi} = \sum_k |\braket{a,k}{\psi}|^2 \ \dots\ \text{degenerované spektrum}\Big)
  \]
  Po naměření hodnoty $a$ systém přejde do stavu
  \[
    \ket{\psi} \xrightarrow{\ \hat A\ } \ket{\bar\psi} = \ket{a}
    \qquad \Big(= \tfrac{1}{\sqrt{\mu}}P_a\ket{\psi}\ \dots\ \text{degenerované spektrum}\Big)
  \]
  (\kw{Redukce stavu}, \kw{kolaps vlnové funkce}) \pozn{$\leftarrow$ motivace měřením $x$}
\end{enumerate}
\end{veta}

\textbf{Poznámky:}
\begin{itemize}[nosep]
  \item[--] \pozn{existují různé formulace / sady postulátů}
  \item[--] \pozn{později přibereme} unitární časový vývoj \\
            \pozn{(princip koresp.)} + konkrétní tvar operátorů pro různé systémy
\end{itemize}

\underline{Ad 1)}
\begin{itemize}[nosep]
  \item[--] v $\infty$ dimenzích \dots\ požadavky úplnosti/spojitosti $\to$ Hilbertův prostor (později)
  \item[--] vlnovou funkci je výhodné normovat: $\|\psi\|^2 = \braket{\psi}{\psi} = 1$ \quad (stále není určená fáze!)
\end{itemize}

\underline{Ad 2)}
\begin{itemize}[nosep]
  \item[--] z vlastností výše plyne, že měřitelné veličiny $a\in\mathbb{R}$ \\
            a vektory tvoří bázi $\to$ měřitelná hodnota je def.\ $\forall\psi\in\mathcal{H}$
\end{itemize}

\underline{Ad 3)}
\begin{itemize}[nosep]
  \item[--] vzorec je uveden pro normovanou $\psi$. Pokud $\braket{\psi}{\psi}\neq 1$, \\
            pak \quad $\mu(a,\psi) = \dfrac{|\braket{a}{\psi}|^2}{\braket{\psi}{\psi}}$
            \quad \pozn{$\leftarrow \ket{\psi}\to\dfrac{\ket{\psi}}{\sqrt{\braket{\psi}{\psi}}}$}
  \item[--] střední hodnota pozorovatelné po mnoha měřeních:
  \[
    \langle A\rangle = \sum_a \mu(a,\psi)\,a = \sum_a \mel{\psi}{\hat P_a}{\psi}\,a = \bra{\psi}\sum_a aP_a\ket{\psi}
    = \vzorec{\mel{\psi}{A}{\psi}}
  \]
  \pozn{Pokud $\braket{\psi}{\psi}\neq 1$: \ $\langle A\rangle = \dfrac{\mel{\psi}{A}{\psi}}{\braket{\psi}{\psi}}$}
\end{itemize}

%----------------------------------------------------------------------
\strana{QM-I-10}
\subsection*{PŘÍKLAD: částice se spinem $1/2$ \ (Sakurai)}
aplikace postulátů na Stern--Gerlachův experiment \\
\quad $\to$ konstrukce modelu na základě $\uparrow$

\underline{Několik experimentů (schéma):}

\begin{center}
\begin{tikzpicture}[
    box/.style={draw,minimum width=1.4cm,minimum height=0.6cm},
    blk/.style={-{Bar[width=6pt]Bar[width=6pt]}},
    >=Stealth, x=1cm, y=1cm]
  % (0)
  \node at (-0.8,0) {(0)};
  \node[box,gray] (z0) at (0.5,0) {zdroj};
  \node[box] (s0) at (3,0) {SG\,$z$};
  \draw[->] (z0) -- (s0);
  \draw[->] (s0.east)++(0,0.15) -- ++(1.3,0.15) node[right] {$S_z = +\frac{\hbar}{2}$};
  \draw[->] (s0.east)++(0,-0.15) -- ++(1.3,-0.15) node[right] {$S_z = -\frac{\hbar}{2}$};
  % osy
  \draw[->] (10,-0.2) -- (10,0.6) node[right] {$z$};
  \draw[->] (10,-0.2) -- (11,-0.2) node[right] {$y$};
  \draw[->] (10,-0.2) -- (9.5,-0.6) node[left] {$x$};
  % (a)
  \node at (-0.8,-1.5) {(a)};
  \node[box,gray] (za) at (0.5,-1.5) {zdroj};
  \node[box] (sa) at (3,-1.5) {SG\,$z$};
  \node[box] (sa2) at (5.6,-1.5) {SG\,$z$};
  \draw[->] (za) -- (sa);
  \draw[->] (sa.east)++(0,0.15) -- (sa2.west|-0,-1.35);
  \draw[blk] (sa.east)++(0,-0.15) -- ++(0.5,0);
  \draw[->] (sa2.east)++(0,0.15) -- ++(1.1,0.15) node[right] {$S_z = \frac{\hbar}{2}$ \qquad $\mu = 1$};
  \draw[dotted] (sa2.east)++(0,-0.15) -- ++(1.1,-0.15) node[right] {$S_z = -\frac{\hbar}{2}$ \quad $\mu = 0$ nepozorováno};
  % (b)
  \node at (-0.8,-3) {(b)};
  \node[box,gray] (zb) at (0.5,-3) {zdroj};
  \node[box] (sb) at (3,-3) {SG\,$z$};
  \node[box] (sb2) at (5.6,-3) {SG\,$x$};
  \draw[->] (zb) -- (sb);
  \draw[->] (sb.east)++(0,0.15) -- (sb2.west|-0,-2.85);
  \draw[blk] (sb.east)++(0,-0.15) -- ++(0.5,0);
  \draw[->] (sb2.east)++(0,0.15) -- ++(1.1,0.15) node[right] {$S_x = \frac{\hbar}{2}$ \qquad $\mu = \frac12$};
  \draw[->] (sb2.east)++(0,-0.15) -- ++(1.1,-0.15) node[right] {$S_x = -\frac{\hbar}{2}$ \quad $\mu = \frac12$};
  % (c)
  \node at (-0.8,-4.5) {(c)};
  \node[box,gray] (zc) at (0.5,-4.5) {zdroj};
  \node[box] (sc) at (3,-4.5) {SG\,$z$};
  \node[box] (sc2) at (5.6,-4.5) {SG\,$x$};
  \node[box] (sc3) at (8.2,-4.5) {SG\,$z$};
  \draw[->] (zc) -- (sc);
  \draw[->] (sc.east)++(0,0.15) -- (sc2.west|-0,-4.35);
  \draw[blk] (sc.east)++(0,-0.15) -- ++(0.5,0);
  \draw[->] (sc2.east)++(0,0.15) -- (sc3.west|-0,-4.35);
  \draw[blk] (sc2.east)++(0,-0.15) -- ++(0.5,0);
  \draw[->] (sc3.east)++(0,0.15) -- ++(1,0.15) node[right] {$S_z = +\hbar/2$ \quad $\mu = \frac12$};
  \draw[->] (sc3.east)++(0,-0.15) -- ++(1,-0.15) node[right] {$S_z = -\hbar/2$ \quad $\mu = \frac12$};
\end{tikzpicture}
\end{center}
\hrule\medskip

\textbf{(0):} Nejjednodušší model na základě měření (0):
\begin{itemize}
  \item stav.\ prostor $\mathcal{H} = \mathcal{L}\{\ket{z+},\ket{z-}\}$
        \pozn{$\leftarrow$ jsou OG (vl.\ v.\ $\hat S_z$), bůno normalizace} \\
        $\forall$ stav \quad $\ket{\psi} = \alpha\ket{z+} + \beta\ket{z-} \equiv \begin{pmatrix}\alpha\\ \beta\end{pmatrix}$
  \item měření SG\,$z$ $\to$ hermit.\ matice
  \[
    \hat S_z = \frac{\hbar}{2}\big(\ket{z+}\bra{z+} - \ket{z-}\bra{z-}\big) = \frac{\hbar}{2}\begin{pmatrix}1&0\\0&-1\end{pmatrix}
    \quad \pozn{$\leftarrow\sigma_z$}
  \]
\end{itemize}
popis 1.\ měření: $\ket{\psi}$ neznáme, ale po měření $\ket{\bar\psi} = \ket{z+}$

\textbf{(a):} Druhé měření \dots\ před: $\ket{\psi} = \ket{z+}$ \quad \pozn{(nebo $\lambda\ket{z+}$)} \\
\quad $\mu\big(+\tfrac{\hbar}{2} = S_z, \psi\big) = |\braket{z+}{\psi}|^2 = 1$ \\
\quad po měření: \quad $\ket{\bar\psi} = \ket{z+}$

\textbf{(b):} Neznáme $\hat S_x$, ale pokusíme se jej najít

\pozn{pozn: nepotřebujeme větší stav.\ prostor: $\mathcal{L}\{\ket{z{+}x{+}}, \ket{z{-}x{+}}, \ket{z{+}x{-}}, \ket{z{-}x{-}}\}$?} \\
\pozn{měření (c) říká, že tento model nefunguje; zkusme pokračovat s $\mathcal{H} = \mathbb{C}^2$}

%----------------------------------------------------------------------
\strana{QM-I-11}
operátor \quad $\hat S_x = \dfrac{\hbar}{2}\big(\ket{x+}\bra{x+} - \ket{x-}\bra{x-}\big)$

\begin{minipage}{0.68\textwidth}
pozorujeme \quad $\mu\big(S_x = +\tfrac{\hbar}{2}, \ket{z+}\big) = |\braket{x+}{z+}|^2 = \tfrac12$ \\
\phantom{pozorujeme} \quad $\mu\big(S_x = -\tfrac{\hbar}{2}, \ket{z+}\big) = |\braket{x-}{z+}|^2 = \tfrac12$
\end{minipage}
\begin{minipage}{0.3\textwidth}
koneč.\ stav: \\ $\ket{\bar\psi} = \ket{x+}$ \\ $\ket{\bar\psi} = \ket{x-}$
\end{minipage}

\textbf{(c)} \quad $\mu\big(S_z = +\tfrac{\hbar}{2}, \ket{x+}\big) = |\braket{z+}{x+}|^2 = \tfrac12$ \\
\phantom{\textbf{(c)}} \quad $\mu\big(S_z = -\tfrac{\hbar}{2}, \ket{x+}\big) = |\braket{z-}{x+}|^2 = \tfrac12$
\begin{align*}
  \Rightarrow\quad \ket{x+} &= \tfrac{1}{\sqrt2}\ket{z+} + \tfrac{1}{\sqrt2}e^{i\delta_1}\ket{z-}
    &&\text{ortogonalita } \braket{x+}{x-} = 0 \\
  \ket{x-} &= \tfrac{1}{\sqrt2}\ket{z+} - \tfrac{1}{\sqrt2}e^{i\delta_1}\ket{z-}
    &&\swarrow\ \text{fixuje druhou fázi}
\end{align*}
\quad $\uparrow$ celková fáze libovolná \dots\ v $\hat S_x$ se neprojeví

zatím máme: \quad
$\hat S_x = \frac{\hbar}{2}\big[\ket{x+}\bra{x+} - \ket{x-}\bra{x-}\big]
 = \frac{\hbar}{2}\big[e^{-i\delta_1}\ket{z+}\bra{z-} + e^{i\delta_1}\ket{z-}\bra{z+}\big]$

\underline{IZOTROPIE PROSTORU:} \dots\ obrácení celého aparátu $\to$ stejný výsl. \pozn{$z\to y\to z$}
\[
  \Rightarrow\quad \hat S_y = \frac{\hbar}{2}\big[e^{-i\delta_2}\ket{z+}\bra{z-} + e^{i\delta_2}\ket{z-}\bra{z+}\big]
  \ \to\
  \begin{aligned}
    \ket{y+} &= \tfrac{1}{\sqrt2}\big(\ket{+} + e^{i\delta_2}\ket{-}\big) \\
    \ket{y-} &= \tfrac{1}{\sqrt2}\big(\ket{+} - e^{i\delta_2}\ket{-}\big)
  \end{aligned}
\]

\begin{minipage}{0.65\textwidth}
další otočení aparátu: \pozn{$x\to y\to x$}
\[
  \Rightarrow\quad |\braket{y\pm}{x+}|^2 = |\braket{y\pm}{x-}|^2 = \tfrac12
\]
tj.\ \quad $\frac12\big|1 \pm e^{i(\delta_1-\delta_2)}\big| = \frac{1}{\sqrt2}$
\[
  \to\quad \delta_2 - \delta_1 = \pm\frac{\pi}{2}
\]
\end{minipage}
\begin{minipage}{0.33\textwidth}
\begin{tikzpicture}[scale=1.1,gray]
  \draw[->] (-1.2,0) -- (1.6,0) node[right] {Re};
  \draw[->] (0,-1.2) -- (0,1.4) node[right] {Im};
  \draw[->,thick] (0,0) -- (1,0) node[below] {$1$};
  \draw[->,thick] (1,0) -- (1,1) node[right] {$\pm i$};
  \draw[dashed] (0,0) -- (1,1) node[midway,above left] {$\sqrt2$};
  \draw[dashed] (0,0) -- (1,-1);
  \fill (0,1) circle (1pt) (0,-1) circle (1pt);
\end{tikzpicture}
\end{minipage}

\dots\ konvence fáze v bázi $\ket{z\pm}$ \\
$\to$ dá se volit $S_x$ reálná matice \dots\ $\delta_1 = 0$ \quad $\Rightarrow\ \delta_2 = +\frac{\pi}{2}$
\pozn{(pravotočivé $\ket{x+},\ket{y+},\ket{z+}$)}

\underline{závěr:}
\[
  \begin{aligned}
    \ket{x\pm} &= \tfrac{1}{\sqrt2}\big(\ket{z+} \pm \ket{z-}\big) = \tfrac{1}{\sqrt2}\begin{pmatrix}1\\ \pm1\end{pmatrix} \\
    \ket{y\pm} &= \tfrac{1}{\sqrt2}\big(\ket{z+} \pm i\ket{z-}\big)
  \end{aligned}
  \quad\to\quad
  \begin{aligned}
    \hat S_x &= \tfrac{\hbar}{2}\big[\ket{+}\bra{-} + \ket{-}\bra{+}\big] \\
    \hat S_y &= \tfrac{\hbar}{2}\big[-i\ket{+}\bra{-} + i\ket{-}\bra{+}\big]
  \end{aligned}
\]

\kw{Pauliho matice}
\[
\vzorec{\hat{\vec\sigma} = (\hat\sigma_x,\hat\sigma_y,\hat\sigma_z) \equiv
 \left(\begin{pmatrix}0&1\\1&0\end{pmatrix}, \begin{pmatrix}0&-i\\ i&0\end{pmatrix}, \begin{pmatrix}1&0\\0&-1\end{pmatrix}\right)}, \qquad
\vzorec{\hat{\vec S} = \frac{\hbar}{2}\hat{\vec\sigma}}
\]

\underline{Pozn:} Posunovací operátory \quad
$S_+ = \hbar\ket{+}\bra{-} = \hbar\begin{pmatrix}0&1\\0&0\end{pmatrix}$, \quad
$S_- = \hbar\ket{-}\bra{+} = \hbar\begin{pmatrix}0&0\\1&0\end{pmatrix} = S_+^\dagger$

\begin{minipage}{0.5\textwidth}
platí: \quad $S_\pm = S_x \pm iS_y$ \\[2pt]
\quad $S_-\ket{+} = \hbar\ket{-}$ \qquad $S_-\ket{-} = 0$ \\
\quad $S_+\ket{-} = \hbar\ket{+}$ \qquad $S_+\ket{+} = 0$
\end{minipage}
\begin{minipage}{0.45\textwidth}
Poznáme později obecně při vyšetřování momentu hybnosti a rotač.\ symetrie
\end{minipage}

%----------------------------------------------------------------------
\strana{QM-I-12}
\underline{pozn:} později uvidíme částice s větším spinem

\begin{center}
\begin{tikzpicture}[>=Stealth, box/.style={draw,minimum width=0.9cm,minimum height=0.5cm}]
  % s = 1/2
  \node at (0,0) {$s = \frac12$};
  \node[box] (a) at (2,0) {SG};
  \draw[->] (0.6,0) -- (a.west);
  \draw[->] (a.east) -- ++(0.8,0.3);
  \draw[->] (a.east) -- ++(0.8,-0.3);
  \node at (1.6,-0.9) {$e^-, p^+, n$};
  \node at (1.6,-1.6) {$\hat S_x,\hat S_y,\hat S_z$:};
  % s = 1
  \node at (4.6,0) {$s = 1$};
  \node[box] (b) at (6.4,0) {SG};
  \draw[->] (5.1,0) -- (b.west);
  \draw[->] (b.east) -- ++(0.8,0.4);
  \draw[->] (b.east) -- ++(0.8,0);
  \draw[->] (b.east) -- ++(0.8,-0.4);
  \node at (6,-0.9) {např.\ $d$};
  \node at (6,-1.6) {matice $3\times3$};
  % s = 3/2
  \node at (9.2,0) {$s = \frac32$};
  \node[box] (c) at (11,0) {SG};
  \draw[->] (9.8,0) -- (c.west);
  \draw[->] (c.east) -- ++(0.8,0.45);
  \draw[->] (c.east) -- ++(0.8,0.15);
  \draw[->] (c.east) -- ++(0.8,-0.15);
  \draw[->] (c.east) -- ++(0.8,-0.45);
  \node at (11.4,-0.9) {atd.};
  \node at (11.2,-1.6) {matice $4\times4$};
\end{tikzpicture}
\end{center}

\subsection*{Přechod k jiné bázi}
pozorovatelná $\hat A$ \quad \dots\ katalog stavů $\ket{a_i}$, \ $i = 1,\dots,d$ \\
jiná pozorovatelná $\hat B$ \quad \dots\ $\ket{b_i}$

reprezentace $A$:
\begin{align*}
  \ket{\psi} = \underbrace{I}_{\sum_i\ket{a_i}\bra{a_i}}\ket{\psi} &= \sum_i \psi_i\ket{a_i}; \qquad \psi_i = \braket{a_i}{\psi} \\
  &= \sum_i \psi_i \underbrace{I}_{\sum_j\ket{b_j}\bra{b_j}}\ket{a_i}
   = \sum_j\underbrace{\Big[\sum_i\overbrace{\braket{b_j}{a_i}}^{U_{ji}}\psi_i\Big]}_{\tilde\psi_j}\ket{b_j}
\end{align*}
tj.\ matice přechodu k jiné bázi (reprezentaci): \quad
$\begin{pmatrix}\tilde\psi_1\\ \vdots\\ \tilde\psi_d\end{pmatrix} = U\begin{pmatrix}\psi_1\\ \vdots\\ \psi_d\end{pmatrix}$

unitární matice \quad $U^\dagger U = I$: \quad
$\sum_j (U^\dagger)_{ij}U_{ji'} = \sum_j U_{ji}^*U_{ji'} = \sum_j\braket{b_j}{a_i}^*\braket{b_j}{a_{i'}} = \delta_{ii'}$

podobně operátory:
\[
  \tilde C_{jj'} \equiv \mel{b_j}{\hat C}{b_{j'}} = \sum_{ii'}\braket{b_j}{a_i}C_{ii'}\braket{a_{i'}}{b_{j'}} = UCU^\dagger
\]

\subsection*{Nekompatibilní pozorovatelné a relace neurčitosti}
o momentech pravd.\ rozdělení: \quad $\mu_n = \sum_a \mu_a a^n = \langle\hat A^n\rangle$ \\
speciálně \quad $\mu_0 = 1$, \quad $\mu_1 = \mel{\psi}{\hat A}{\psi}$, \quad $\mu_2 = \mel{\psi}{\hat A^2}{\psi}$

rozptyl měření veličiny $A$ ve stavu $\psi$:
\[
  (\Delta A)^2 \equiv \mel{\psi}{(\hat A - \langle A\rangle)^2}{\psi} = \langle\hat A^2\rangle - \langle A\rangle^2 = \mu_2 - \mu_1^2
\]
\dots\ speciálně pro $\ket{\psi} = \ket{a}$ \quad \dots\ $\mu_a = 1$, \ $\mu_{a'\neq a} = 0$ \\
\quad tj.\ $\mu_1 = a$, $\mu_2 = a^2$ \quad a tedy \quad $\Delta A \equiv \sqrt{\mu_2 - \mu_1^2} = 0$

%----------------------------------------------------------------------
\strana{QM-I-13}
\subsection*{Relace neurčitosti}
\begin{veta}[Relace neurčitosti]
nechť \quad $[\hat A,\hat B] = i\hat C$ \quad ($\neq 0 \overset{\text{def}}{\equiv}$ \kw{Nekompatibilní}) \\
\quad \phantom{nechť} $\uparrow$ to $i$ tady je, aby $\hat C^\dagger = \hat C$
\[
  \text{potom}\qquad \vzorec{(\Delta A)\cdot(\Delta B) \geq \tfrac12\,|\mel{\psi}{\hat C}{\psi}|}
\]
\end{veta}

\underline{Důkaz:} \quad Schwarzova nerovnost \quad $\|\phi_1\|\cdot\|\phi_2\| \geq |\braket{\phi_1}{\phi_2}|$ \\
\quad pro \quad $\ket{\phi_1} = (A - \langle A\rangle)\ket{\psi}$ \quad $\to\ \|\phi_1\| = \sqrt{\mel{\psi}{(A-\langle A\rangle)^2}{\psi}} = \Delta A$ \\
\quad \phantom{pro} \quad $\ket{\phi_2} = (B - \langle B\rangle)\ket{\psi}$ \quad podob.\ $\|\phi_2\| = \Delta B$

pravá strana:
\begin{align*}
  |\braket{\phi_1}{\phi_2}| &= \big|\bra{\psi}\big(AB - A\langle B\rangle - \langle A\rangle B + \langle A\rangle\langle B\rangle\big)\ket{\psi}\big| \\
  &= \big|\bra{\psi}\big(AB - \langle A\rangle\langle B\rangle\big)\ket{\psi}\big| \\
  &= \Big|\bra{\psi}\Big(\frac{AB - BA}{2} + \frac{AB + BA}{2} - \langle A\rangle\langle B\rangle\Big)\ket{\psi}\Big| \\
  &= \Big|\underbrace{\bra{\psi}\tfrac{AB+BA}{2}\ket{\psi} - \langle A\rangle\langle B\rangle}_{\in\mathbb{R}}
     + \tfrac{i}{2}\underbrace{\mel{\psi}{C}{\psi}}_{\in\mathbb{R}}\Big|
     \geq \tfrac12|\mel{\psi}{C}{\psi}| \qquad\text{c.b.d.}
\end{align*}

\textbf{PŘ 1:}
\begin{itemize}[nosep]
  \item Heisenberg \quad $[\hat x,\hat p] = i\hbar$ \quad $\to$ \quad $\vzorec{\Delta x\cdot\Delta p \geq \frac{\hbar}{2}}$
  \item spin $1/2$ \quad $[\hat S_z,\hat S_x] = i\hbar\hat S_y$ \quad (ověřte! \quad obecně $[\hat S_a,\hat S_b] = i\hbar\,\varepsilon_{abc}\hat S_c$ \quad $\leftarrow$ Levi-Civita)
\end{itemize}
\quad tj.\ \quad $\Delta S_x\cdot\Delta S_z \geq \frac{\hbar}{2}|\langle S_y\rangle|$

důsledky:
\begin{enumerate}[label=\arabic*)]
  \item \uv{ostrá hodnota} $\hat S_z$ \dots\ např.: $\ket{\psi} = \ket{z+}$ $\to$ $\Delta S_z = 0$ \\
        pak musí být $\langle S_y\rangle = 0$, tj.\ $P_\pm = \frac12$!
  \item \uv{ostrá hodnota} $S_y$ \dots\ tj.\ $|\langle S_y\rangle| = \frac{\hbar}{2}$ \\
        $\to$ nesmí být $\Delta S_x = 0$ ani $\Delta S_z = 0$, \ tj.\ ne vlastní stavy!
\end{enumerate}

\subsection*{Kompatibilní pozorovatelné}
následující tvrzení řeší důsledky relace $[\hat A,\hat B] = 0$ pro dvě pozorovatelné (hermitovské matice) $\hat A,\hat B$

pozn: (triviální, ale důležitá) \quad $[\hat A,\hat B] = 0 \Leftrightarrow \hat A\hat B = \hat B\hat A$ \\
\quad tj.\ ve výrazech lze prohazovat $A\leftrightarrow B$

%----------------------------------------------------------------------
\strana{QM-I-14}
\begin{veta}[Lemma 1]
$[\hat A,\hat B] = 0 \ \wedge\ B\ket{\psi} = b\ket{\psi}$ \quad $\Rightarrow$ \quad $B\ket{\phi} = b\ket{\phi}$ \quad pro $\ket{\phi} = A\ket{\psi}$
\end{veta}

\underline{DK:} prostě dosadit $\phi = A\ket{\psi}$ do $B\ket{\phi}$ + užití předpokladu \\
\underline{komentář:} invariance vl.\ podprostorů $B$ vůči působení $A$

\begin{veta}[Lemma 2]
$[\hat A,\hat B] = 0$, \ $\hat A\ket{a_i} = a_i\ket{a_i}$, \ $a_1\neq a_2$ \quad $\Rightarrow$ \quad $\mel{a_1}{\hat B}{a_2} = 0$
\end{veta}

\underline{DK:} \quad
$0 = \mel{a_1}{[\hat A,\hat B]}{a_2} = \mel{a_1}{\hat A\hat B}{a_2} - \mel{a_1}{\hat B\hat A}{a_2}$ \\
\phantom{\underline{DK:}} \quad $0 = (a_1 - a_2)\mel{a_1}{\hat B}{a_2}$ \quad $\to$ \quad c.b.d.

\underline{komentář:}

\begin{center}
\begin{tikzpicture}[x=0.5cm,y=0.5cm,gray]
  % spektrální rozklad A
  \node[left] at (-0.8,-2.5) {sp.\ rozklad $\hat A$};
  \draw (0,0) -- (0,-5) (5,0) -- (5,-5);
  \node at (0.5,-0.5) {$a_1$}; \node at (1.5,-1.5) {$a_1$};
  \node at (2.5,-2.5) {$a_2$}; \node at (3.5,-3.5) {$a_2$}; \node at (4.5,-4.5) {$a_2$};
  \node at (3.5,-1) {$0$}; \node at (1,-3.8) {$0$};
  \draw[decorate,decoration={brace,amplitude=3pt}] (0,0.2) -- (2,0.2) node[midway,above=2pt] {\small $a_1$};
  \draw[decorate,decoration={brace,amplitude=3pt}] (2,0.2) -- (5,0.2) node[midway,above=2pt] {\small $a_2$};
  % šipka
  \draw[->] (6.5,-2.5) -- (10,-2.5) node[midway,above] {$[\hat A,\hat B] = 0$};
  % B
  \node at (11,-2.5) {$\hat B =$};
  \begin{scope}[xshift=6cm]
    \draw (0,0) -- (0,-5) (5,0) -- (5,-5);
    \draw[pattern=north east lines,pattern color=gray!60] (0,0) rectangle (2,-2);
    \draw[pattern=north east lines,pattern color=gray!60] (2,-2) rectangle (5,-5);
    \node at (3.5,-1) {$0$}; \node at (1,-3.5) {$0$};
    \draw[decorate,decoration={brace,amplitude=3pt}] (0,0.2) -- (2,0.2) node[midway,above=2pt] {\small $a_1$};
    \draw[decorate,decoration={brace,amplitude=3pt}] (2,0.2) -- (5,0.2) node[midway,above=2pt] {\small $a_2$};
  \end{scope}
\end{tikzpicture}
\end{center}
budeme hojně používat ve výpočtech \hfill
$\nearrow$ částečná diagonalizace $B$ -- známé spektrum $A$

\begin{veta}[Lemma 3a]
$[\hat A,\hat B] = 0$ \quad $\Rightarrow$ \quad $[\hat P_a,\hat B] = 0$ \\
\quad $\hat A = \sum_a a\hat P_a$, \qquad $\hat B = \sum_b b\hat P_b$
\end{veta}

\underline{DK:} báze vl.\ v.\ $\hat A$: \quad $\hat A\ket{a,\alpha} = a\ket{a,\alpha}$ \quad $\to$ \quad $\hat P_a = \sum_\alpha\ket{a\alpha}\bra{a\alpha}$
\begin{align*}
  \mel{a_1\alpha_1}{[\hat P_a,\hat B]}{a_2\alpha_2}
  &= \sum_\alpha\underbrace{\braket{a_1\alpha_1}{a\alpha}}_{\delta_{a_1a}\delta_{\alpha_1\alpha}}\mel{a\alpha}{\hat B}{a_2\alpha_2}
   - \sum_\alpha\mel{a_1\alpha_1}{\hat B}{a\alpha}\underbrace{\braket{a\alpha}{a_2\alpha_2}}_{\delta_{aa_2}\delta_{\alpha\alpha_2}} \\
  &= \delta_{a_1a}\mel{a_1\alpha_1}{\hat B}{a_2\alpha_2} - \delta_{aa_2}\mel{a_1\alpha_1}{\hat B}{a_2\alpha_2} \\
  &= 0 \quad \begin{cases}
       \text{pro } a_1 = a_2 & \text{díky faktoru } \delta_{a_1a} - \delta_{a_2a} \\
       \text{pro } a_1\neq a_2 & \text{díky lemmatu 2}
     \end{cases}
\end{align*}

\begin{veta}[Lemma 3]
$[\hat A,\hat B] = 0$ \quad $\Rightarrow$ \quad $[\hat P_a,\hat P_b] = 0 \quad \forall a,b$
\end{veta}

\underline{DK:} užití lemmatu 3a $\to$ $[\hat P_a,\hat B] = 0$ + opak.\ pro $\hat A\to\hat B$, $\hat B\to\hat P_a$

\underline{Důsledek:} \quad $\hat P_{ab} \equiv \hat P_a\hat P_b$ je projekční operátor \\
\quad ověření \quad $\hat P_{ab}^\dagger = \hat P_b^\dagger\hat P_a^\dagger = \hat P_b\hat P_a = \hat P_{ab}$ \\
\quad \phantom{ověření} \quad $\hat P_{ab}^2 = \hat P_a\hat P_b\hat P_a\hat P_b = \hat P_a^2\hat P_b^2 = \hat P_{ab}$

%----------------------------------------------------------------------
\strana{QM-I-15}
Dokonce platí

\begin{veta}[VĚTA]
samosdružené $\hat A,\hat B$: \quad $[\hat A,\hat B] = 0 \Leftrightarrow [\hat P_a,\hat P_b] = 0$
\end{veta}

pozn: \quad ($\Rightarrow$) jsme dokázali výše \\
\phantom{pozn:} \quad ($\Leftarrow$) je triviálním důsl.\ spekt.\ rozkladu a linearity komutátoru: \\
\phantom{pozn:} \quad\quad $[\hat A,\alpha_1\hat B_1 + \alpha_2\hat B_2] = \alpha_1[\hat A,\hat B_1] + \alpha_2[\hat A,\hat B_2]$

\begin{veta}[VĚTA]
samosdružené operátory $A, B$ mají společnou bázi vl.\ v. \quad $\Leftrightarrow$ \quad $\vzorec{[\hat A,\hat B] = 0}$
\end{veta}

\underline{DK:}
\begin{itemize}[nosep]
  \item[--] stačí vzít bázi ve vl.\ podprostoru $\hat P_{ab}\ket{\psi} = \lambda\ket{\psi}$ \\
            pro $\lambda = 1$ a $\forall a,b$, pro něž je $\hat P_{ab}\neq 0$
  \item[$\to$] tyto prostory jsou navz.\ ortogonální díky $\hat P_a\hat P_{a'} = \delta_{aa'}\hat P_a$ \\
            (podob.\ pro $b$) a navíc jejich direkt.\ součet je celý $\mathcal{H}$, \\
            neboť \quad $\hat I = \sum_{ab}\hat P_{ab}$ \quad \pozn{DK $= \hat I\cdot\hat I = \sum_a\hat P_a\cdot\sum_b\hat P_b$}
\end{itemize}
pozn \dots\ $\{(a,b)\}$ nemusí být kart.\ součin $\to$ vynecháváme $\hat P_{ab} = 0$

\underline{pozn:}
\begin{itemize}
  \item společný spektr.\ rozklad \quad
        $\hat A = \sum_{ab}a\hat P_{ab}$ \quad \pozn{$= \sum_a a\underbrace{\big(\textstyle\sum_b\hat P_{ab}\big)}_{\hat P_a}$} \\
        \phantom{společný spektr.\ rozklad} \quad $\hat B = \sum_{ab}b\hat P_{ab}$ \quad \pozn{podobně}
  \item funkce dvou proměnných \quad $f(\hat A,\hat B) = \sum_{ab}f(a,b)\hat P_{ab}$ \\
        \dots\ není jednoznačně definovaná pro nekomutující $\hat A,\hat B$
  \item dá se rozšířit na libovolný počet navzájem komut.\ oper. \\
        $\hat A^{(1)},\hat A^{(2)},\dots,\hat A^{(N)}$ navzájem komutují $\Rightarrow$ $\exists$ spol.\ báze vl.\ v.\ $\forall\hat A^{(i)}$
\end{itemize}

\subsection*{Úplný systém komutujících operátorů (ÚSKO)}
\begin{definice}[ÚSKO]
Řekneme, že $\underbrace{\hat A^{(1)},\dots,\hat A^{(N)}}_{\text{(navzájem komutující)}}$ tvoří \kw{ÚSKO}, pokud $\forall$ přípustná \\
vlastních čísel $\{a^{(1)},\dots,a^{(N)}\}$ jednoznačně definuje spol.\ vlastní vektor $\ket{\psi}$ (až na fázi)
\end{definice}

značení \quad $\ket{\psi} \equiv \ket{a^{(1)},\dots,a^{(N)}}$

\underline{pozn:}
\begin{itemize}[nosep]
  \item tj.\ spol.\ projektory $P_{a^{(1)}\dots a^{(N)}}$ už jsou jednodim.\ nebo $0$
  \item sadu $\{\hat A^{(1)},\dots,\hat A^{(N)}\}$ lze považovat za jeden operátor \\
        s nedegenerovaným spektrem vl.\ čísel \dots\ $N$-tice $\{a^{(1)},\dots,a^{(N)}\} \equiv \vec a$
\end{itemize}

%----------------------------------------------------------------------
\strana{QM-I-16}
\begin{veta}[Věta]
Pokud operátor $\hat F$ komutuje s $\forall A^{(1)},\dots,A^{(N)}$, které tvoří ÚSKO, pak lze vyjádřit jako $\hat F = f(A^{(1)},\dots,A^{(N)})$
\end{veta}

\underline{DK:} $\hat F, A^{(1)},\dots,A^{(N)}$ mají spol.\ sadu vl.\ v., ale to musí být $\ket{a^{(1)}\dots a^{(N)}}$ \\
\quad současně $\hat F\ket{a^{(1)}\dots a^{(N)}} = f\ket{a^{(1)}\dots a^{(N)}}$ \dots\ $f$ je jediné pro $\forall$ sadu $\uparrow$

\textbf{Poznámka:}
\begin{itemize}[nosep]
  \item značná část zbytku přednášky bude o hledání ÚSKO pro různé systémy a hledání jejich spektra a transf.\ mezi nimi.
  \item Pro odpovědi na fyzikální otázky potřebujeme, aby veličina, která nás zajímá, byla součástí ÚSKO, nebo jejich funkcí.
  \item Uvidíme, že pro odpověď na otázky o časovém vývoji je vhodné mít v ÚSKO Hamiltonián.
\end{itemize}

\subsection*{Pozn: DIREKTNÍ SOUČET PROSTORŮ}
\begin{minipage}{0.72\textwidth}
rozklad \quad $I = \sum_a P_a$ \quad $\longleftrightarrow$ \quad $\mathcal{H} = \bigoplus_a\mathcal{H}_a$ \\
\hspace*{4cm} \dots\ rozklad do navz.\ ortogon.\ podprostorů
\end{minipage}
\begin{minipage}{0.26\textwidth}
\pozn{ve vektorech:} \\
$\psi = \left(\begin{array}{c} f_1 \\ \hline \vdots \\ \hline f_d \end{array}\right)
 \begin{array}{l} \}\,a_1 \\ \\ \}\,a_n \end{array}$
\end{minipage}

podobně lze skládat:

\begin{definice}[Direktní součet]
$\mathcal{H} = \mathcal{H}^{(1)}\oplus\mathcal{H}^{(2)}$ \quad je LVP daný
$\Big\{\big(\ket{\psi^{(1)}},\ket{\psi^{(2)}}\big),\ \ket{\psi^{(1)}}\in\mathcal{H}^{(1)},\ \ket{\psi^{(2)}}\in\mathcal{H}^{(2)}\Big\}$ \\
\quad s operacemi:
\begin{align*}
  \ket{\psi} + \ket{\phi} &= \big(\ket{\psi^{(1)}} + \ket{\phi^{(1)}},\ \ket{\psi^{(2)}} + \ket{\phi^{(2)}}\big) \\
  \alpha\ket{\psi} &= \big(\alpha\ket{\psi^{(1)}},\ \alpha\ket{\psi^{(2)}}\big) \\
  \braket{\phi}{\psi} &= \braket{\phi^{(1)}}{\psi^{(1)}}_{\mathcal{H}^{(1)}} + \braket{\phi^{(2)}}{\psi^{(2)}}_{\mathcal{H}^{(2)}}
\end{align*}
\end{definice}

tj.\ v konečné dimenzích:
\[
  \ket{\psi^{(1)}} = \begin{pmatrix} f_1^{(1)}\\ \vdots\\ f_{d_1}^{(1)}\end{pmatrix},\quad
  \ket{\psi^{(2)}} = \begin{pmatrix} f_1^{(2)}\\ \vdots\\ f_{d_2}^{(2)}\end{pmatrix}
  \quad\longrightarrow\quad
  \ket{\psi} = \begin{pmatrix} f_1^{(1)}\\ \vdots\\ f_{d_1}^{(1)}\\ f_1^{(2)}\\ \vdots\\ f_{d_2}^{(2)}\end{pmatrix}
  \qquad\text{dimenze}
\]
tj.\ $\mathcal{H}^{(1)}\oplus\mathcal{H}^{(2)}$ je prostor, jehož bází (ortonormální) je sjednocení bází $\mathcal{H}^{(1)}$ a $\mathcal{H}^{(2)}$

\underline{Příklad:} \quad $\mathbb{C}^2 = \mathbb{C}\oplus\mathbb{C}$; \quad $\mathbb{C}^n = \mathbb{C}\oplus\mathbb{C}\oplus\dots\oplus\mathbb{C}$ \\
\phantom{\underline{Příklad:}} \quad $I = \sum_{ab}P_{ab}$ \quad \dots\ $\mathcal{H} = \bigoplus_a\bigoplus_b\mathcal{H}_{ab}$

%----------------------------------------------------------------------
\strana{QM-I-17}
fyzikálně: \dots\ přidávání dalších hodnot pozorovatelné

\begin{minipage}{0.4\textwidth}
\textbf{PŘ:} Kvantové tečky:

\begin{center}
\begin{tikzpicture}[gray]
  \foreach \x/\l in {0/A, 0.8/B, 2/C, 2.8/D}
    \node[draw,circle,inner sep=1.5pt,double] at (\x,0) {\small\l};
\end{tikzpicture}
\end{center}
\end{minipage}
\begin{minipage}{0.58\textwidth}
$\mathcal{H}^{(1)} = \mathcal{L}\{\ket{A},\ket{B}\}$ \\
\quad poloha elektronu je $\ket{A}$ nebo $\ket{B}$ \\
$\mathcal{H}^{(2)} = \mathcal{L}\{\ket{C},\ket{D}\}$ \\
\quad poloha elektronu je $\ket{C}$ nebo $\ket{D}$
\end{minipage}
\[
  \mathcal{H} = \mathcal{H}^{(1)}\oplus\mathcal{H}^{(2)} = \mathcal{L}\{\ket{A},\ket{B},\ket{C},\ket{D}\}
\]

\subsection*{DIREKTNÍ SOUČIN PROSTORŮ \ {\normalfont\small (někdy též tenzorový)}}
\begin{itemize}[nosep]
  \item[--] důležitá operace \dots\ skládání dvou kvantových systémů, respektive přidávání stupňů volnosti
  \item[--] prostor, jehož bází je kartézský součin bází
\end{itemize}

\begin{definice}[Direktní (tenzorový) součin]
nechť $\mathcal{H}^{(1)}$ \dots\ báze $\{\ket{\phi_i^{(1)}}\}_{i=1}^{d_1}$ \quad \dots\ podob.\ $\mathcal{H}^{(2)}$ \\
\quad pak def: \quad $\mathcal{H} = \mathcal{H}^{(1)}\otimes\mathcal{H}^{(2)}$ \quad je prostor s bází $\{\ket{\phi_i^{(1)}}\ket{\phi_j^{(2)}}\}_{i=1,\,j=1}^{d_1,\ d_2}$ \\
\quad tj.\ prostor s dimenzí $d = d_1\cdot d_2$ \quad (izomorfní maticím $d_1\times d_2$) \\
\quad tj.\ obecný vektor \quad $\ket{\psi} = \sum_{ij}\psi_{ij}\ket{\phi_i}\ket{\phi_j}$
\quad \dots\ někdy se vynechává horní index \dots\ automat.\ 1.\ vekt.\ $\mathcal{H}^{(1)}$ \\
\quad + součet a násobení číslem definovaný přirozeně (distrib.\ zákon)
\end{definice}

\begin{itemize}
  \item pozn: zkrácený zápis \quad $\ket{\phi_i^{(1)}}\ket{\phi_j^{(2)}} \equiv \ket{\phi_i^{(1)}}\otimes\ket{\phi_j^{(2)}} \equiv \ket{\phi_i^{(1)}\phi_j^{(2)}}$
  \item \underline{skalární součin} \quad
        $\braket{\phi^{(1)}\phi^{(2)}}{\psi^{(1)}\psi^{(2)}} = \braket{\phi^{(1)}}{\psi^{(1)}}_{\mathcal{H}^{(1)}}\braket{\phi^{(2)}}{\psi^{(2)}}_{\mathcal{H}^{(2)}}$ \\
        speciálně báze: \quad $\braket{\phi_i\phi_j}{\phi_{i'}\phi_{j'}} = \braket{\phi_i}{\phi_{i'}}\braket{\phi_j}{\phi_{j'}} = \delta_{ii'}\delta_{jj'}$
  \item \underline{faktorizované vektory:} \quad
        $\ket{\psi^{(1)}} = \sum_i f_i^{(1)}\ket{\phi_i^{(1)}}$; \quad $\ket{\psi^{(2)}} = \sum_j f_j^{(2)}\ket{\phi_j^{(2)}}$ \\
        $\to$ distrib.\ zákon \quad
        $\ket{\psi^{(1)}\psi^{(2)}} \equiv \ket{\psi^{(1)}}\ket{\psi^{(2)}} \equiv \ket{\psi^{(1)}}\otimes\ket{\psi^{(2)}} = \sum_{ij}f_i^{(1)}f_j^{(2)}\ket{\phi_i\phi_j}$
  \item \kw{entanglement:} \pozor{ne každý stav je možno napsat $\uparrow$}
  \item možno zobecnit na více prostorů \quad
        $\mathcal{H} = \mathcal{H}^{(1)}\otimes\mathcal{H}^{(2)}\otimes\mathcal{H}^{(3)}$, \quad
        $\mathcal{H} = \bigotimes_{i=1}^{N}\mathcal{H}^{(i)}$
  \item \underline{OPERÁTORY na $\mathcal{H}^{(1)}\otimes\mathcal{H}^{(2)}$:} \quad
        obecně \quad $\hat A\otimes\hat B\ket{\phi}\ket{\psi} = \hat A\ket{\phi}\otimes\hat B\ket{\psi}$ \\
        operátor $\hat A$ na $\mathcal{H}^{(1)}$ \dots\ rozšíření na $\mathcal{H}^{(1)}\otimes\mathcal{H}^{(2)}$ \dots\ $\hat A\to\hat A\otimes\hat I$ \\
        podobně $\hat B$ na $\mathcal{H}^{(2)}$ $\to$ $\hat I\otimes\hat B$
\end{itemize}
ČASTO: ÚSKO na $\mathcal{H}^{(1)}\otimes\mathcal{H}^{(2)}$ budován z ÚSKO na $\mathcal{H}^{(1)}$ a ÚSKO na $\mathcal{H}^{(2)}$

%----------------------------------------------------------------------
\strana{QM-I-18}
\subsection*{PŘÍKLADY}
\textbf{A) MATEMATICKÉ}
\begin{itemize}
  \item $\mathcal{H} = \mathbb{C}$ \quad \dots\ 1D vektory \quad $\ket{\psi} = \alpha\ket{e_1}$ \hfill $\mathcal{H}'\equiv\mathcal{H}$
  \begin{itemize}[nosep]
    \item[--] $\mathcal{H}\oplus\mathcal{H}'$ \ (izomorfní $\mathbb{C}^2$) \quad \dots\ báze $\{\ket{e_1},\ket{e_1'}\}$ \dots\ 2 komp.\ \dots\ $\begin{pmatrix}\alpha\\ \alpha'\end{pmatrix}$
    \item[--] $\mathcal{H}\otimes\mathcal{H}'$ \ (izomorfní $\mathbb{C}$) \quad \dots\ báze $\{\ket{e_1}\ket{e_1'}\}$ \dots\ 1 vektor
  \end{itemize}
  \item $\mathcal{H} = \mathbb{C}^2$
  \begin{itemize}[nosep]
    \item[--] $\mathcal{H}\oplus\mathcal{H}'$ \ ($\equiv\mathbb{C}^4$) \quad $\Big\}$ ale jiná struktura \quad
          $\nearrow \begin{pmatrix}\ \\ \ \end{pmatrix}\!\begin{smallmatrix}1\\2\end{smallmatrix}\
                    \begin{pmatrix}\ \\ \ \end{pmatrix}\!\begin{smallmatrix}1'\\2'\end{smallmatrix}$
    \item[--] $\mathcal{H}\otimes\mathcal{H}'$ \ ($\equiv\mathbb{C}^4$) \hspace{3.7cm}
          $\searrow \begin{pmatrix}\ \\ \ \\ \ \\ \ \end{pmatrix}\!\begin{smallmatrix}11'\\12'\\21'\\22'\end{smallmatrix}$
  \end{itemize}
  \item Obecně: \quad $\mathbb{C}^n\oplus\mathbb{C}^m \equiv \mathbb{C}^{n+m}$ \\
        \phantom{Obecně:} \quad $\mathbb{C}^n\otimes\mathbb{C}^m \equiv \mathbb{C}^{n\times m}$ \quad
        $\leftarrow$ izomorfní prostoru matic, ale v QM píšeme do sloupce
\end{itemize}

\textbf{B) FYZIKÁLNÍ:}
\begin{itemize}
  \item jedna částice v kvantové trojtečce
        \tikz[gray,baseline=-3pt]{\foreach \x/\l in {0/1,0.6/2,1.2/3} \node[draw,circle,inner sep=1pt] at (\x,0) {\scriptsize\l};} \\
        $\mathcal{H}^{(D)} = \mathcal{L}\{\ket{1},\ket{2},\ket{3}\}$ \ ($\equiv\mathbb{C}^3$) \\
        \quad $\ket{\psi} = \sum_d\alpha_d\ket{d}$, \quad $d = 1,2,3$ \quad $\longleftrightarrow$ \quad
        $\begin{pmatrix}\alpha_1\\ \alpha_2\\ \alpha_3\end{pmatrix}$
  \item dvě částice v trojtečce: \\
        $\mathcal{H} = \mathcal{H}^{(D)}\otimes\mathcal{H}^{(D)} = \mathcal{L}\{\ket{11},\ket{12},\ket{13},\ket{21},\dots,\ket{33}\}$ \ ($\equiv\mathbb{C}^9$) \\
        \quad $\ket{\psi} = \sum_{d_1d_2}\psi_{d_1d_2}\ket{d_1d_2}$, \qquad $\ket{d_1d_2} \equiv \ket{d_1}\otimes\ket{d_2}$ \\
        \hspace*{6.2cm} $\to$ částice 1 v $\ket{d_1}$; částice 2 v $\ket{d_2}$ \\
        \underline{pozn:} komplikace \dots\ nerozlišitelné částice $\to$ později
  \item částice se spinem $\frac12$: \quad $\mathcal{H}^{(S)} = \mathcal{L}\{\ket{+},\ket{-}\} \equiv \mathbb{C}^2$
  \item částice se spinem $\frac12$ v trojtečce: \quad $\mathcal{H} = \mathcal{H}^{(D)}\otimes\mathcal{H}^{(S)}$ \ ($\equiv\mathbb{C}^6$) \\
        \quad $\ket{\psi} = \sum_{ds}\psi_{ds}\ket{ds}$; \qquad
        $\ket{ds} \equiv \underbrace{\ket{d}}_{=1,2,3}\otimes\underbrace{\ket{s}}_{=+,-}$
        \quad $\leftarrow$ částice v místě $d$ se spinem $s$
\end{itemize}
\pozn{$\to$ na cvičení též pozorovatelné, 2 částice se spinem}

př: \underline{Entanglement} \quad např.\ $\ket{\psi} = \ket{1+} + \ket{2-}$ nelze faktorizovat \\
\quad jako \quad $\ket{\psi} = \ket{\phi}\otimes\ket{\sigma}$ \qquad \pozn{DALŠÍ PŘ NA CVIČENÍ}

%======================================================================
% QM-I-2014-4.pdf
%======================================================================
\section{Formalismus kvantové teorie 2}
\strana{QM-I-19}
\begin{itemize}[nosep]
  \item[--] případ spojitého spektra
  \item[--] pár poznámek k úvodu do funkcionální analýzy
\end{itemize}

\begin{minipage}{0.5\textwidth}
\textbf{PŘ 1:} částice v 1D \dots\ bezstrukturní

Experiment \dots\ měření polohy $x$:

\bigskip
\quad $\Rightarrow$ \quad $\mathcal{H} = \mathcal{L}\{\ket{x};\ x\in\mathbb{R}\}$
\end{minipage}
\begin{minipage}{0.48\textwidth}
\begin{tikzpicture}[gray,scale=0.8]
  % zdroj
  \draw (0,1.3) rectangle (0.3,1.6);
  \draw[dashed,->] (0.4,1.45) -- (1.4,1.6);
  % štěrbina
  \draw[thick] (1.6,1.75) -- (1.6,2.8);
  \draw[thick] (1.6,1.45) -- (1.6,0.3);
  % dvojštěrbina
  \draw[thick] (3,2.2) -- (3,3);
  \draw[thick] (3,1.4) -- (3,1.95);
  \draw[thick] (3,0) -- (3,1.15);
  % stínítko
  \draw[->] (4.6,0) -- (4.6,3.2) node[right] {$x$};
  \foreach \y in {2.7,2.3,2.22,2.15,1.6,1.4,1.3} \fill (4.6,\y) circle (1.5pt);
  \draw[->] (5,0.2) -- (5.5,0.35) node[right] {$z$};
\end{tikzpicture}
\end{minipage}

\begin{tabular}{@{}lll}
OBECNÝ VEKTOR: & $\ket{\psi} = \sum_x\psi_x\ket{x}$ \quad $\longrightarrow$ & $\ket{\psi} = \int\psi(x)\ket{x}\,\dd x$ \\[3pt]
SKALÁRNÍ SOUČIN: & $\braket{\phi}{\psi} = \sum_x\phi_x^*\psi_x$ \quad $\longrightarrow$ & $\braket{\phi}{\psi} = \int\phi(x)^*\psi(x)\,\dd x$ \\[3pt]
NORMOVÁNÍ: & $\ket{\psi}\to\ket{\psi'} = \dfrac{\ket{\psi}}{\sqrt{\braket{\psi}{\psi}}}$ & \dots\ lze jen $\int|\psi(x)|^2\dd x < \infty$
\end{tabular}

$\to$ požadujeme \quad $\mathcal{H} = L^2(\mathbb{R}) \equiv \big\{\psi(x):\ \int|\psi(x)|^2\dd x < \infty\big\}$ \\
\hspace*{5.3cm} $\uparrow$ def.\ skoro všude na $\mathbb{R}$

OPERÁTOR: \qquad $\hat x = \sum_x xP_x$ \quad $\longrightarrow$ \quad $\hat x = \int x\ket{x}\bra{x}\,\dd x$

\textbf{PŘ 2:} (polo)nekonečný řetízek kvantových teček (atomů)

\begin{center}
\begin{tikzpicture}[gray]
  \draw[dotted] (-0.8,0) -- (-0.3,0);
  \foreach \x/\l in {0/{-2},1/{-1},2/0,3/1,4/2} {
    \draw (\x,0) circle (0.2);
    \node[above] at (\x,0.25) {\scriptsize$\ket{\l}$};
  }
  \draw[dotted] (4.3,0) -- (4.8,0);
\end{tikzpicture}
\end{center}

$\mathcal{H} = \mathcal{L}\{\ket{n};\ n\in\mathbb{Z}\}$ \quad \dots\ $\ket{\psi} = \sum_n\psi_n\ket{n}$ \\
\dots\ normalizovatelnost \quad $\|\psi\|^2 = \sum_n|\psi_n|^2 < \infty$ \quad \dots\ prostor $\mathcal{H} = \ell^2(\mathbb{Z})$

zpátky k příkladu 1: \\
\quad pravděpodobnost nalezení částice v místě $x$ $\sim|\psi_x|^2$ \\
\quad vlastní stavy operátoru $\hat x$ na $L^2(\mathbb{R})$

\begin{minipage}{0.58\textwidth}
\underline{Diracova $\delta$-funkce}

můžeme chápat jako limitu $\varepsilon\to 0$
\[
  \delta(x) = \lim_{\varepsilon\to 0}\frac{1}{\varepsilon}\,\chi_{\langle 0,\varepsilon\rangle}(x)
\]
\end{minipage}
\begin{minipage}{0.4\textwidth}
\begin{tikzpicture}[scale=0.9]
  \draw[->] (-1.5,0) -- (2.2,0) node[right] {\small $x$};
  \draw[thick] (-1.5,0.02) -- (0,0.02) -- (0,1.8) -- (0.5,1.8) -- (0.5,0.02) -- (2,0.02);
  \node[below] at (0,0) {\small $0$};
  \node[below] at (0.5,0) {\small $\varepsilon$};
  \node[gray,left] at (0,1.8) {\small $1/\varepsilon$};
\end{tikzpicture}
\end{minipage}

vlastnosti: \quad $\delta(x) = 0\ \ \forall x\neq 0$; \quad $\int\delta(x)\,\dd x = 1$; \quad $\int\delta(x)f(x)\,\dd x = f(0)$ \\
$\to$ vlastní funkce operátoru $\hat x$ \quad $\hat x\psi_{x_0} = x_0\psi_{x_0}$ \quad \dots\ $\psi_{x_0}(x) = \delta(x - x_0)$

JEDNODUCHÝ NÁHLED \dots\ Dirac
\[
  \psi(x) \to \text{sloupcový vektor s spojitým indexem}
  \begin{pmatrix}\vdots\\ \psi(x_0)\\ \vdots\end{pmatrix}\!\cdots x_0
  \qquad
  \delta(x - x_0)\cdots
  \begin{pmatrix}0\\ \vdots\\ \blacktriangle\\ \vdots\\ 0\end{pmatrix}\!\cdots x_0
\]
\hfill jako konečně-dim $\mathcal{H}$, jen $\sum_n\to\int\dd x$

%----------------------------------------------------------------------
\strana{QM-I-20}
\subsection{(jemný) Úvod do pokročilých mat.\ metod QM}

Reed \& Simon: \emph{Methods of modern mathematical physics}. \quad (cca 1600 stran) \\
\quad I.\ Functional analysis \quad II.\ Fourier Analysis, Self-Adjointness \\
\quad III.\ Scattering Theory \quad IV.\ Analysis of Operators

motto k úvodní kapitole: \emph{\uv{The beginner should not be discouraged if he finds that he does not have the prerequisites for reading the prerequisites.}}

\subsubsection*{Hilbertův prostor \pozn{(ket vektory, stavy)}}
\begin{definice}[Hilbertův prostor]
LVP $\mathcal{H}$ nazveme \kw{Hilbertovým prostorem}, pokud \\
\quad $\forall$ Cauchyovská posloupnost vektorů z $\mathcal{H}$ má limitu v $\mathcal{H}$ \hfill (ÚPLNOST)
\end{definice}

\underline{Pozn:} $\{\psi_n\}_{n=1}^\infty$ je Cauchyovská, pokud:
\quad $\forall\varepsilon > 0\ \exists K:\ \forall m,n > K:\ \|\psi_m - \psi_n\| < \varepsilon$

\underline{Pozn:} úplnost normovaného prostoru $\to$ Banachův prostor \\
\quad v Hilbertovi: \quad $\|\psi\| = \sqrt{\braket{\psi}{\psi}}$ \quad norma induk.\ skal.\ součinem

\underline{Pozn:} Def.\ limity: řekneme $\psi_n\to\psi$, pokud $\|\psi_n - \psi\|\to 0$

\subsubsection*{Separabilita $\mathcal{H}$}
\begin{definice}[Separabilita]
$\mathcal{H}$ nazveme \kw{separabilní} $\Leftrightarrow$ $\exists$ spočetná hustá podmnožina \\
\hspace*{3.9cm} $\Leftrightarrow$ $\exists$ spočetná ON báze
\end{definice}

\underline{PŘ:} (Prostory)
\begin{itemize}[nosep]
  \item prostor $\ell^2$ $\equiv$ prostor posl.\ $\psi\equiv\{\psi_n\}_{n=1}^\infty$; \ $\sum|\psi_n|^2 < \infty$ \\
        \quad skalární součin \quad $\braket{\phi}{\psi} \equiv \sum_n\phi_n^*\psi_n$
  \item prostor $L^2(a,b)$; \ $(a,b)\subset\mathbb{R}$ \dots\ může být neomezený \\
        \quad $\to$ fce $\psi(x)$ na $(a,b)$: \ $\int|\psi|^2\dd x < \infty$ \\
        \quad skalární součin \quad $\braket{\phi}{\psi} \equiv \int\phi^*(x)\psi(x)\,\dd x$
  \item Banachův prostor $L^p(a,b)$: \quad $\int|\psi|^p\dd x < \infty$ \quad \dots\ $L^2$ je spec.\ př.
\end{itemize}

\underline{PŘ} (Báze)
\begin{itemize}[nosep]
  \item v prostoru $\ell^2$ je kanon.\ báze: \quad
        $\begin{aligned}[t] \{1,0,0,\dots\} &\equiv \ket{\phi_1} \\ \{0,1,0,\dots\} &\equiv \ket{\phi_2} \\ &\ \vdots \end{aligned}$
  \item v prostoru $L^2(0,L)$ \dots\ Fourierovská báze \quad
        $\ket{\phi_k} = \frac{1}{\sqrt L}e^{2ik\pi x/L}$, \quad $k = \dots,-1,0,1,\dots$ \\
        \quad nebo reálná \quad $\ket{\phi_n} = \frac{1}{\sqrt L},\ \sqrt{\frac2L}\sin\frac{2\pi x}{L},\ \sqrt{\frac2L}\cos\frac{2\pi x}{L},\ \dots$
\end{itemize}
$\forall\psi\in L^2(0,L)$: \quad $\ket{\psi} = \sum_n f_n\ket{\phi_n}$ \quad \dots\ reprezentace $\ket{\psi}\leftrightarrow\{f_n\}_{n=1}^\infty$ \\
$\to$ \pozor{$\forall$ separabilní $\mathcal{H}$ je izomorfní $\ell^2$} \hfill lin.\ operátory \dots\ \uv{matice}

%----------------------------------------------------------------------
\strana{QM-I-21}
\subsubsection*{Duální prostor \pozn{(bra vektory, spoj.\ lin.\ funkcionály)}}
\underline{Lineární funkcionály:} \quad $F:\mathcal{H}\to\mathbb{C}$ \\
\quad norma: \quad $\|F\| = \sup\dfrac{\|F(\ket{\psi})\|}{\|\psi\|}$ \quad \dots\ tj.\ $\|F(\psi)\| \leq \|\psi\|\cdot\|F\|$

Lineární funkcionály: \quad omezenost $\Leftrightarrow$ spojitost v bodě $\Leftrightarrow$ spojitost na $\mathcal{H}$

\begin{veta}[Rieszova věta o reprezentaci]
$\forall F\in\mathcal{H}^*$ ($\equiv$ prostor $\forall$ spoj.\ lin.\ funkcionálů) \ $\exists!\,\ket{\phi}\in\mathcal{H}:\ F(\psi) = \braket{\phi}{\psi}$ \\
navíc: \quad $\|F\|_{\mathcal{H}^*} = \|\phi\|_{\mathcal{H}}$ \\
tj.\ izomorfizmus: \quad $\vzorec{\mathcal{H}\equiv\mathcal{H}^*}$
\end{veta}

\underline{pozn:} $\uparrow$ lze DK díky úplnosti a $\exists$ skalár.\ součinu
\begin{itemize}[nosep]
  \item[--] nefunguje v $L^p$ prostorech: \quad $(L^p)^* \equiv L^q$, \ kde $\frac1p + \frac1q = 1$
\end{itemize}

\subsubsection*{Diracova $\delta$-funkce}
jako limita \quad $\delta = \lim_{\varepsilon\to 0^+}\delta_\varepsilon(x)$ \quad \dots\ není def.\ v $\mathcal{H}$

\underline{DK:}
\begin{itemize}[nosep]
  \item[--] jako bodová konvergence: $\delta = 0$ skoro všude, tj.\ $\int\delta = 0$
  \item[--] jako lim v $\mathcal{H}$ \dots\ není Cauchyovská posl.
\end{itemize}

\begin{minipage}{0.6\textwidth}
\pozn{podrobněji: $\delta_n \equiv \delta_{\varepsilon = \frac1n}(x)$} \\
\pozn{platí \quad $\|\delta_n - \delta_{2n}\| = \sqrt n \to\infty$}
\end{minipage}
\begin{minipage}{0.38\textwidth}
\begin{tikzpicture}[gray,scale=0.7]
  \draw[->] (-1,0) -- (2,0);
  \draw (0,0) -- (0,2) -- (0.6,2) -- (0.6,0);
  \node[left] at (0,2) {\small $n$};
  \node[below] at (0.3,0) {\small $\frac1n$};
  \node[right] at (1,1.2) {\small $\int|\cdot|^2 = n^2\cdot\frac1n = n$};
\end{tikzpicture}
\end{minipage}

\underline{$\delta$-fce jako lin.\ funkcionál:} \quad $F_\varepsilon[\psi] = \int\delta_\varepsilon(x)\psi(x)\,\dd x$ \\
\quad \dots\ není spojitý na $\mathcal{H}$ $\to$ nedefinuje spoj.\ funkcionál v lim $\varepsilon\to 0$

\begin{minipage}{0.65\textwidth}
\textbf{PŘ:} posloupnost \quad $\phi_n(x) = \dfrac{1}{1 + e^{-nx}} \xrightarrow{\ n\to\infty\ } \Theta(x)$ \\[3pt]
\quad přitom \quad $F_\varepsilon[\phi_n] \to \frac12$ \quad ($\varepsilon\to 0$, fixní $n$) \\
\quad ale \quad $F_\varepsilon[\phi_\infty] = 1$ \quad $\forall\varepsilon > 0$
\end{minipage}
\begin{minipage}{0.33\textwidth}
\begin{tikzpicture}[scale=0.8]
  \draw[gray,->] (-2,0) -- (2.2,0);
  \draw[gray] (0,-0.4) -- (0,1.5);
  \draw[thick] (-2,0.05) -- (-0.1,0.05);
  \draw (0,0.05) circle (2pt);
  \draw[thick] (0.1,1) -- (2,1);
  \draw (0,1) circle (2pt);
  \fill (0,0.5) circle (2pt);
\end{tikzpicture}
\end{minipage}

Problém $\to$ příliš velký def.\ obor. \quad $F_\varepsilon[\psi]$ je spojitý funkcionál na prostoru spojitých funkcí

\begin{definice}[$\delta$-funkce]
$\delta$ jako spojitý lin.\ funkcionál na $\psi\in\mathcal{H}$; $\psi$ spojité
\[
  \vzorec{F_\delta[\psi] \equiv \int\delta(x)\psi(x)\,\dd x \overset{\text{def}}{\equiv} \psi(0)}
\]
\end{definice}
podobně \quad $\delta(x - x_0)$ \quad \dots \quad $F_{\delta_{x_0}}[\psi] \equiv \int\delta(x - x_0)\psi(x)\,\dd x = \psi(x_0)$

%----------------------------------------------------------------------
\strana{QM-I-22}
\noindent\rule{\linewidth}{0.4pt}
\subsubsection*{ODBOČKA -- VZOREČKY PRO $\delta$-funkci}

\underline{$\delta$-limity:}

\begin{minipage}{0.7\textwidth}
\begin{itemize}
  \item $\delta_\alpha(x) = \dfrac{1}{\pi}\dfrac{\alpha}{x^2 + \alpha^2} \xrightarrow{\ \alpha\to 0\ } \delta(x)$
  \item $\delta_\alpha(x) = \dfrac{\alpha}{\sqrt\pi}e^{-\alpha^2x^2} \xrightarrow{\ \alpha\to\infty\ } \delta(x)$
\end{itemize}
\end{minipage}
\begin{minipage}{0.28\textwidth}
\begin{tikzpicture}[gray,scale=0.6]
  \draw[thin] (0,0) -- (0,1.6);
  \draw[thick,domain=-2.5:2.5,samples=80] plot (\x,{1.4*0.3^2/(\x*\x+0.3^2)});
  \begin{scope}[yshift=-2cm]
    \draw[thin] (0,0) -- (0,1.6);
    \draw[thick,domain=-2.5:2.5,samples=80] plot (\x,{1.4*exp(-6*\x*\x)});
  \end{scope}
\end{tikzpicture}
\end{minipage}

základ F.T.:
\begin{itemize}
  \item $\displaystyle\int e^{ikx}\,\dd k$ \quad přesněji \quad
        $\displaystyle\int_{-\alpha}^{\alpha}e^{ikx}\,\dd k = \frac{e^{i\alpha x} - e^{-i\alpha x}}{ix} = \frac{2\sin\alpha x}{x} \xrightarrow{\ \alpha\to\infty\ } 2\pi\delta(x)$
\end{itemize}

\begin{minipage}{0.65\textwidth}
tj. \quad $\bullet$ \ $\delta_\alpha(x) = \dfrac{\sin\alpha x}{\pi x} \xrightarrow{\ \alpha\to\infty\ } \delta(x)$

\medskip
tj. \quad $\vzorec{\displaystyle\int_{-\infty}^{\infty}e^{2\pi ikx}\,\dd k = \delta(x)}$ \quad $\delta$ je F.T.\ funkce $f(x) = 1$
\end{minipage}
\begin{minipage}{0.33\textwidth}
\begin{tikzpicture}[gray,scale=0.6]
  \draw[thin] (-3,0) -- (3,0);
  \draw[thin] (0,-0.4) -- (0,1.8);
  \draw[thick,domain=-3:3,samples=150] plot (\x,{(abs(\x)<0.001) ? 1.5 : 1.5*sin(5*\x r)/(5*\x)});
\end{tikzpicture}
\end{minipage}

\underline{Další užitečné vztahy:}
\begin{itemize}
  \item $\delta(x) = \Theta'(x)$ \quad \dots\ lze chápat z \quad $\Theta(x) = \displaystyle\int_{-\infty}^{x}\delta(y)\,\dd y$

  \begin{minipage}{0.45\textwidth}
  nebo jako limitu:
  \[ \delta_\varepsilon(x) \equiv \tfrac{1}{\varepsilon}\chi_{\langle 0,\varepsilon\rangle} = \Theta_\varepsilon'(x) \]
  \end{minipage}
  \begin{minipage}{0.5\textwidth}
  \begin{tikzpicture}[scale=0.8]
    \draw[->] (-2,0) -- (3.5,0);
    \draw[->] (0,0) -- (0,1.5);
    \node[left] at (0,1) {$1$};
    \draw[thick] (-2,0) -- (0,0) -- (0.8,1) -- (3,1) node[right] {$\Theta_\varepsilon(x)$};
    \draw[gray,thin] (0.8,0) -- (0.8,1);
    \node[below] at (0.8,0) {$\varepsilon$};
  \end{tikzpicture}
  \end{minipage}
  \item $\delta(\alpha x) = \dfrac{1}{|\alpha|}\delta(x)$ \quad \dots\ ze subst.\ v \ $\int\delta(\alpha x)\,\dd x = \frac1\alpha\int\delta(y)\,\dd y$
  \item $\delta\big(f(x)\big) = \displaystyle\sum_{x_k}\frac{\delta(x - x_k)}{|f'(x_k)|}$ \quad kde $x_k$ kořeny $f(x_k) = 0$
  \item $\delta(-x) = \delta(x)$
\end{itemize}
\noindent\rule{\linewidth}{0.4pt}

%======================================================================
% Další části (každý dokument zvlášť v parts/)
%======================================================================
\newpage
%======================================================================
% PoznQMI-4.pdf
%======================================================================
\section{Formalismus QM -- II}
\strana{QM-F-1}

Řekli jsme, že stavovým prostorem částice je $\mathcal{H} = L^2(\mathbb{R})$, tj.
\[
  \ket{\psi} = \int\psi(x')\ket{x'}\,\dd x' \qquad\dots\ \text{stav.\ prostor izomorf.\ s prost.\ } \forall\psi
\]
\[
  \|\psi\|^2 = \braket{\psi}{\psi} = \int\psi^*(x)\psi(x)\,\dd x < \infty
\]
Jak vypadají bázové vektory $\ket{x}$?

\subsection*{Diracova $\delta$-funkce}
\begin{minipage}{0.6\textwidth}
můžeme chápat jako limitu $\varepsilon\to 0$: \\
\dots\ částice lokalizovaná se stejnou amplitudou pravděpodobnosti v celém intervalu $\langle 0,\varepsilon\rangle$
\end{minipage}
\begin{minipage}{0.38\textwidth}
\begin{tikzpicture}[scale=0.9]
  \draw[->] (-2,0) -- (2.2,0) node[right] {$x$};
  \draw[thick] (-2,0.02) -- (0,0.02) -- (0,1.6) -- (0.6,1.6) -- (0.6,0.02) -- (2,0.02);
  \node[below] at (0,0) {\small $0$};
  \node[below] at (0.6,0) {\small $\varepsilon$};
  \node[left] at (0,1.6) {\small $\frac1\varepsilon$};
  \node[right] at (0.8,1.3) {$\delta_\varepsilon(x)$};
\end{tikzpicture}
\end{minipage}

\underline{problém:} \quad $\lim_{\varepsilon\to 0}\delta_\varepsilon(x)$ \ není dobře def.\ v $\mathcal{H}$

přesněji:
\begin{itemize}[nosep]
  \item[--] chápáno jako lim bod po bodu \dots\ $\delta_0 = 0$ skoro všude \quad ($\delta_0(0) = \infty$)
  \item[--] chápáno jako lim v $\mathcal{H}$ není Cauchy, neboť $\|\delta_\varepsilon\| = \sqrt{\varepsilon\cdot\frac{1}{\varepsilon^2}} = \frac{1}{\sqrt\varepsilon}\to\infty$
\end{itemize}

\begin{definice}[\uv{Rigged} Hilbertův prostor]
$\underline{\mathcal{H}}$ hustá podmn.\ v $\mathcal{H}$; \quad $\overline{\mathcal{H}} \equiv \underline{\mathcal{H}}^*$ \dots\ lin.\ spoj.\ funkcionály na $\underline{\mathcal{H}}$
\end{definice}

\textbf{PŘ:} \quad $\underline{\mathcal{H}} = \big\{\forall\psi(x)\in C^\infty(\mathbb{R});\ |\psi^{(n)}(x)(1+|x|)^m| < K\ \ \forall m,n\big\}$ \quad \kw{Schwartz space} \\
\quad v $\mathbb{R}^n$ obec.: $\sup|x^\alpha D^\beta f| < \infty$ $\forall\alpha,\beta$ \dots\ multiindexy

\textbf{PŘ:} \quad $\delta_\varepsilon(x)$ jako lin.\ funkcionál \quad $F_\varepsilon[\psi] = \int\delta_\varepsilon(x)\psi(x)\,\dd x$ \\
\quad opět $F_\varepsilon$ není limita v $\mathcal{H}^*$ \quad (nedefinuje spojitý funkcionál)

\begin{minipage}{0.65\textwidth}
\pozn{př: posl.\ fcí $\phi_n(x) = \frac{1}{1+e^{-nx}} \xrightarrow{n\to\infty} \Theta(x)$} \\
\pozn{přitom: $F_\varepsilon[\phi_n] \to \frac12$ \ ($\varepsilon\to 0$, fixní $n$)} \\
\pozn{ale \quad $F_\varepsilon[\phi_\infty] = 1$ \quad $\forall\varepsilon > 0$}
\end{minipage}
\begin{minipage}{0.33\textwidth}
\begin{tikzpicture}[gray,scale=0.8]
  \draw[->] (-2,0) -- (2.2,0);
  \draw (0,-0.6) -- (0,1.6);
  \draw[thick] (-2,0.05) -- (-0.1,0.05);
  \draw (0,0.05) circle (2pt);
  \draw[thick] (0.1,1.1) -- (2,1.1);
  \draw (0,1.1) circle (2pt);
  \fill (0,0.55) circle (2pt);
  \draw[thin,domain=-1.8:1.8,samples=60] plot (\x,{1.1/(1+exp(-3*\x))});
\end{tikzpicture}
\end{minipage}

\dots\ volba $\underline{\mathcal{H}}$ spojité funkce \quad $F_\varepsilon[\psi(x)] \to \delta[\psi(x)] = \psi(0)$

\begin{itemize}
  \item další důležitá fce, která není v $\mathcal{H}$: \quad $e^{ipx}$ \dots\ je v $\overline{\mathcal{H}}$, pokud volíme testovací fce dostatečně ubývající. \\
        většinou: \quad $\underline{\mathcal{H}}$ \dots\ $\forall$ fce z $C^\infty$ ubývající rychleji než $\forall$ mocnina $|x|^n$ v $\infty$ (i jejich derivace)
\end{itemize}
\dots\ funkce jako $\delta(x - x_0)$, $e^{ikx}$, $\sin kx$ můžeme zahrnout jako prvky $\overline{\mathcal{H}}$

\textbf{Jiný příklad:} \quad $\ell^2 = \mathcal{H}$; \quad $\underline{\mathcal{H}} = \forall$ posl.\ $c_n$: $\sum_n|c_n|^2n^m < \infty$ \ $\forall m = 0,1,2,\dots$ \\
\quad potom $\overline{\mathcal{H}}$ lze složit \dots\ prostor $\forall v_n$: \ $\sum_n v_n^*c_n < \infty$ \ $\forall c_n\in\underline{\mathcal{H}}$

%----------------------------------------------------------------------
\strana{QM-F-2}
\subsection*{Další důležité příklady $\delta$-limit}
\begin{minipage}{0.68\textwidth}
\begin{itemize}[nosep]
  \item $\delta_\alpha(x) = \dfrac1\pi\dfrac{\alpha}{x^2+\alpha^2} \xrightarrow{\ \alpha\to 0\ } \delta(x)$ \quad \pozn{Lorentz}
  \item $\delta_\alpha(x) = \dfrac{\alpha}{\sqrt\pi}e^{-\alpha^2x^2} \xrightarrow{\ \alpha\to\infty\ } \delta(x)$ \quad \pozn{Gauss}
\end{itemize}
důležitý případ: \quad $\delta_\alpha(x) = \dfrac{\sin\alpha x}{\pi x} \xrightarrow{\ \alpha\to\infty\ } \delta(x)$ \dots
\end{minipage}
\begin{minipage}{0.3\textwidth}
\begin{tikzpicture}[gray,scale=0.55]
  \draw[thin] (0,-0.2) -- (0,2.2);
  \draw[domain=-2.5:2.5,samples=80] plot (\x,{2*0.5^2/(\x*\x+0.5^2)});
  \draw[domain=-2.5:2.5,samples=80] plot (\x,{2*0.2^2/(\x*\x+0.2^2)});
  \begin{scope}[yshift=-2.6cm]
    \draw[thin] (-3,0) -- (3,0);
    \draw[thick,domain=-3:3,samples=150] plot (\x,{(abs(\x)<0.001) ? 1.8 : 1.8*sin(5*\x r)/(5*\x)});
  \end{scope}
\end{tikzpicture}
\end{minipage}

\dots\ základ Fourier.\ transformace:
\[
  \int e^{ikx}\,\dd k \ \longrightarrow\ \int_{-\alpha}^{\alpha}e^{ikx}\,\dd k = \frac{e^{i\alpha x} - e^{-i\alpha x}}{ix} = \frac{2\sin\alpha x}{x} \xrightarrow{\ \alpha\to\infty\ } 2\pi\delta(x)
\]
\noindent\rule{\linewidth}{0.3pt}
neboli \quad $\vzorec{\displaystyle\int e^{2\pi ikx}\,\dd k = \delta(x)}$

\begin{minipage}{0.6\textwidth}
užitečný vztah: \quad $\delta(x) = \Theta'(x)$ \quad \dots\ \pozn{$\lim_{\alpha\to 0}$}
\end{minipage}
\begin{minipage}{0.38\textwidth}
\begin{tikzpicture}[gray,scale=0.8]
  \draw (-2,0) -- (3,0);
  \draw (0,0) -- (0,1.2);
  \node[left] at (0,1) {\small $1$};
  \draw[thick] (-2,0) -- (0,0) -- (0.6,1) -- (2.8,1);
  \node[below] at (0,0) {\small $0$};
  \node[below] at (0.6,0) {\small $\alpha$};
\end{tikzpicture}
\end{minipage}

\noindent\rule{\linewidth}{0.3pt}

\subsection*{Použití $\delta(x)$ jako báze}
vztah \quad $\ket{\psi} = \int\psi(x')\ket{x'}\,\dd x'$ \quad \dots\ báze v $L^2$: \ $\ket{x'} = \delta(\xi - x')$
\[
  = \int\psi(x')\delta(\xi - x')\,\dd x' = \psi(\xi)\ \checkmark
\]
koef.\ rozvoje do báze: $\braket{n}{\psi} = \psi_n$ \quad \dots\ zde \quad
$\braket{x'}{\psi} = \int\delta(\xi - x')\psi(\xi)\,\dd\xi = \psi(x')\ \checkmark$

\pozor{pozor} \dots\ $\ket{x}$ nejde normovat konvenčním způsobem: \ \sout{$\braket{x}{x} = 1$} \\
místo toho \quad $\vzorec{\braket{x}{x'} = \int\dd\xi\,\delta(x - \xi)\delta(x' - \xi) = \delta(x - x')}$ \\
\quad \dots\ analogie \quad $\braket{n}{n'} = \delta_{nn'}$

Rozklad $\hat I$: \quad $I = \sum_n\ket{n}\bra{n}$ \quad \dots\ zde: \quad
$\begin{aligned}[t]\ket{x_0} &= \delta(x - x_0) \\ \bra{x_0} &= \delta(x - x_0)^* = \delta(x - x_0)\end{aligned}$
\[
  \text{tj.}\quad I = \int\dd x_0\ket{x_0}\bra{x_0} = \int\dd x_0\,\delta(\xi - x_0)\delta(\xi' - x_0) = \delta(\xi - \xi')
\]
\quad \dots\ zápis $I$ v $x$-reprezentaci \quad $\mel{x}{I}{x'} \equiv \braket{x}{x'} = \delta(x - x')$

zápis operátoru $Q$: \quad $Q\ket{x} = x\ket{x}$
\[
  \mel{x}{Q}{x'} = \int\dd\xi\,\delta(x - \xi)\,x'\,\delta(x' - \xi) = x'\delta(x - x') = x\,\delta(x - x')
\]

%----------------------------------------------------------------------
\strana{QM-F-3}
\subsection*{Více o operátorech}
$A:\ \mathcal{D}(A)\to\mathcal{R}(A)$; \quad $\overline{\mathcal{D}(A)} = \mathcal{H}$ \dots\ hustá množina \\
\quad $\mathcal{R}(A)\subset\mathcal{H}$ \dots\ obor hodnot. \quad \dots\ jinak nevíme dost dobře jak aplikovat na $\forall\ket{\psi}$

\begin{definice}[omezený operátor]
$\exists M:\ \dfrac{\|A\psi\|}{\|\psi\|} \leq M \quad \forall\psi\in\mathcal{D}(A)$ \\[2pt]
def: norma operátoru \quad $\|A\| = \sup_{\|\psi\|\leq 1}\|A\psi\|$
\end{definice}

\underline{pozn:} spojitost \dots\ v $\infty$ dim \dots\ $\forall$ lin.\ oper.\ je spojitý \\
\quad v $\mathcal{H}$: \ lin.\ operátor je spoj.\ $\Leftrightarrow$ omezený

\textbf{PŘ:}
\begin{itemize}[nosep]
  \item unitární operátor \quad $\|U\| = 1$ \quad protože $\|U\psi\| = \|\psi\|$
  \item $P$ \dots\ projektor \quad $\|P\|\leq 1$ \quad neboť
  \[
    \|P\psi\|^2 = \braket{P\psi}{P\psi} = |\braket{\psi}{P\psi}| \overset{\text{Schwarz}}{\leq} \|\psi\|\cdot\|P\psi\| \qquad\text{tj.}\quad \frac{\|P\psi\|}{\|\psi\|}\leq 1
  \]
\end{itemize}

\underline{Pozn:} díky spojitosti lze omez.\ oper.\ dodef.\ (rozšířit) na $\mathcal{H} = \overline{\mathcal{D}(A)}$ \\
\quad \dots\ úplnost \dots\ $\forall$ Cauchy $\psi_n\to\psi$ \dots\ $A\psi_n$ je rovněž Cauchy $\to\exists A\psi$

\textbf{PŘ:} F.T. \quad $\int f(x)e^{2\pi ikx}\,\dd x$ \dots\ je def.\ na $L_1$ nebo na $\mathcal{S}$ + rozšíření \\
\quad $\equiv F[f]$ \quad \dots\ $\|F\| = 1$

\pozor{Problém:} v QM se často setkáváme s neomez.\ operátory: \\
\quad $\exists\psi_n\in\mathcal{H}:\ \|\psi_n\| = 1:\ \|A\psi_n\|\to\infty$ \quad a \quad $\mathcal{D}(A)\subset\mathcal{H}$

\begin{priklad}[PŘ]
$\hat x\psi(x) = x\psi(x)$ \quad \dots\ $\mathcal{D}(A) = \big\{\psi(x)\in\mathcal{H};\ \int x^2|\psi|^2\dd x < \infty\big\}$

není omez.: \quad $\psi_n(x) = $
\tikz[baseline=2pt,scale=0.55]{
  \draw[->] (0,0) -- (4.2,0);
  \draw (0,0) -- (0,1.1);
  \node[left] at (0,1) {\scriptsize 1};
  \draw[dotted] (0,1) -- (2.5,1);
  \draw[thick] (2.5,0) -- (2.5,1) -- (3.3,1) -- (3.3,0);
  \node[below] at (2.5,0) {\scriptsize $n$};
  \node[below] at (3.4,0) {\scriptsize $n{+}1$};
}
\quad $\|\hat x\psi_n\|^2 = \int_n^{n+1}x^2\,\dd x = \frac13(3n^2 + 3n + 1)\to\infty$

podobně: \quad $\hat p = -i\hbar\dv{x}$ \quad \dots\ $\psi_n = N_ne^{inx}$ \quad \dots\ $\|\hat p\psi_n\| = \hbar n\|\psi_n\|\to\infty$
\end{priklad}

\dots\ Rigged spaces: \ $\underline{\mathcal{H}}\equiv\mathcal{D}$ \dots\ vl.\ vektory v $\overline{\mathcal{H}}$ \dots\ $\delta(x - \xi)$; $e^{\frac{i}{\hbar}p\xi}$ \\
\quad nebo nahrazení spektrálního rozkladu projektorovou mírou \dots

MATEMATICKÉ FOUSY \dots\ KOPÁČEK

%----------------------------------------------------------------------
\strana{QM-F-4}
\subsection*{Spektrum lineárního operátoru}
\underline{v $\not\infty$ dim:} \quad $A\ket{\phi} = \lambda\ket{\phi} \Leftrightarrow |A - \lambda I| = 0$ \\
\quad $\{\forall\lambda\}$ = spektrum (diskrétní) \dots\ spekt.\ rozklad \quad $A = \sum_\lambda\lambda P_\lambda$, \quad
$P_\lambda = \sum_{\ket{\phi}\text{ pro fix }\lambda}\ket{\phi}\bra{\phi}$

\underline{v $\infty$ dim:} \quad def.\ nemá smysl.

\sout{Matem.\ fousy: Symetrický $\times$ samosdružený operátor:}

\begin{definice}[spektrum lin.\ operátoru]
$\lambda\in\mathbb{C}$ nazveme \kw{bodem spektra} $\lambda\in\sigma(A)$, pokud $(A - \lambda I)$ není prosté a $\mathcal{H}$ na $\mathcal{H}$.
\end{definice}

Možnosti:
\begin{itemize}[nosep]
  \item[(i)] $(A - \lambda)$ není prosté, tj.\ $\exists\psi\neq 0:\ (A - \lambda)\psi = 0$. \\
             tj.\ $\psi\in\ker(A - \lambda)$ \quad \dots\ $\lambda\in\sigma_p(A)$ \dots\ \kw{bodové spektrum} \sout{(diskrétní)}
  \item[(ii)] $(A - \lambda)$ prosté, ale $\mathcal{R}(A - \lambda)\neq\mathcal{H}$
  \begin{itemize}[nosep]
    \item[a)] $\overline{\mathcal{R}(A - \lambda)} = \mathcal{H}$ \dots\ hustá v $\mathcal{H}$ \quad \dots\ $\lambda\in\sigma_c(A)$ \dots\ \kw{spojité spektrum}
    \item[b)] $\overline{\mathcal{R}(A - \lambda)}\neq\mathcal{H}$ \quad \dots\ $\lambda\in\sigma_R$ \dots\ \kw{reziduální spektrum}
  \end{itemize}
\end{itemize}
Zde se operátory s $\sigma_R$ zabývat nebudeme.

\begin{definice}[Normální operátor]
$[A, A^\dagger] = 0$ \quad pak \quad
$\lambda\in\sigma(A) \Leftrightarrow \exists\psi_n\in\mathcal{H},\ \|\psi_n\| = 1,\ \lim_{n\to\infty}(A - \lambda)\psi_n = 0$
\end{definice}

\begin{priklad}[PŘ: operátor $\hat x$]
\dots\ posloupnost \quad $\psi_n = \sqrt n\,\chi_{\langle x_0,x_0+\frac1n\rangle}$ \quad \dots\ $\|\psi_n\| = 1$ \\
neboť \quad $\|(\hat x - x_0)\psi_n\|^2 = \int_{x_0}^{x_0+\frac1n}(x - x_0)^2n\,\dd x = n\cdot\frac13\big(\frac1n\big)^3 = \frac13n^{-2}\to 0$ \\
ale $\psi_n$ není Cauchy (výše) \dots\ má však limitu v $\overline{\mathcal{H}}$ $\to\delta(x - x_0)$ \\
(po přenormování \dots\ norma na $\overline{\mathcal{H}}$ je jiná než na $\mathcal{H}$)
\end{priklad}

\begin{itemize}
  \item více o \underline{normálních operátorech}: \hfill \dots\ spec.\ př.\ \dots\ $A = A^\dagger$; $U^\dagger = U^{-1}$
\end{itemize}
\underline{v $\not\infty$ dim:}
\begin{enumerate}[label=\arabic*),nosep]
  \item $A\ket{\psi} = \lambda\ket{\psi} \Rightarrow A^\dagger\ket{\psi} = \lambda^*\ket{\psi}$ \hfill \dots\ unit.\ $\lambda^* = \lambda^{-1}$ \dots\ $|\lambda| = 1$
  \item $\lambda_1\neq\lambda_2$ vl.\ č.\ $\Rightarrow\ \braket{\psi_1}{\psi_2} = 0$
  \item vl.\ v.\ tvoří ON bázi
\end{enumerate}
důvod: každý operátor mohu rozložit: \quad $A = \dfrac{A + A^\dagger}{2} + i\dfrac{A - A^\dagger}{2i} = C + iD$, \\
\quad kde $C, D$ hermitovské, navíc $[A, A^\dagger] = 0 \Rightarrow [C, D] = 0$ \\
\quad tj.\ mají spol.\ bázi vl.\ v.

Normální operátor \dots\ komplex.\ zobecnění reál.\ Hermit.\ operátoru

%----------------------------------------------------------------------
\strana{QM-F-5}
\providecommand{\sumint}{\mathop{\textstyle\sum\hspace{-1.05em}\int}}
\subsection*{Další matematické fousy}
\underline{Symetrický $\times$ Samosdružený operátor:}

\begin{definice}
\kw{symetrický} (Hermitovský): \quad $(A\phi,\psi) = (\phi,A\psi)\quad\forall\phi,\psi\in\mathcal{D}(A)$ \\
\kw{samosdružený} \dots\ navíc \quad $\mathcal{D}(A) = \mathcal{D}(A^\dagger)$
\end{definice}
(pozn: relace $(\psi,A\phi) = (A^\dagger\psi,\phi)$ může definovat $A^\dagger\ket{\psi}$ i pro $\psi$ mimo $\mathcal{D}(A)$)

\pozor{pojem samosdruženosti podstatně závisí na prostoru na němž def.\ $A$!} \ (a spektra)

\begin{priklad}[Př: $P = -i\dv{x}$ \dots\ neomezený, symetrický (v záv.\ na prostoru)]
\begin{align*}
  \mel{\phi}{D^\dagger}{\psi} \equiv \braket{D\phi}{\psi} &= i\int_a^b\phi^{*\prime}(x)\psi(x)\,\dd x = \text{(per partes)} \\
  &= i\big[\psi(x)\phi^*(x)\big]_a^b - i\int\phi^*(x)\psi'(x)\,\dd x = [\ \ ] + \mel{\phi}{D}{\psi}
\end{align*}
tj.\ závisí na okraj.\ podm.: \hfill $Pe^{i\lambda x} = \lambda e^{i\lambda x}$
\begin{itemize}[nosep]
  \item[(a)] fce na $(-\infty,\infty)$ bez podm.: $P$ není sym., $\forall\lambda\in\mathbb{C}$ je vl.\ č.
  \item[(b)] fce omezené pro $|x|\to\infty$: $P$ nesym.; $\lambda\in\mathbb{R}$ jsou vl.\ č.
  \item[(c)] na $\langle 0,L\rangle$ period.\ fce: $P$ symetrický, \ $\lambda = \lambda_n = \frac{2\pi n}{L}$, $n\in\mathbb{Z}$ vl.\ č.
  \item[(d)] fce na $(-\infty,\infty)$; $|\phi|\to 0$ pro $|x|\to\infty$: $P$ symetrický, ale nemá vl.\ fce.
\end{itemize}
\end{priklad}

\begin{veta}[Samosdružený operátor]
má úplnou množinu (zobecněných) vlastních funkcí: \\
Spektrální rozklad: \quad $I = \int P_a\,\dd a + \sum_n P_n$ \quad $\leftarrow$ \kw{úplnost}
\[
  \vzorec{A = \int aP_a\,\dd a + \sum_n a_nP_n} \quad \pozn{$= \sumint_{a\in\sigma(A)} aP_a$}
\]
\end{veta}
v případě nedegener.\ spektra: \quad $P_n = \ket{a_n}\bra{a_n}$, \quad $P_a = \ket{a}\bra{a}$ \\
\phantom{v případě} degenerované: \quad $P_n = \sumint_\alpha\ket{a_n,\alpha}\bra{a_n,\alpha}$, \quad $P_nP_{n'} = P_n\delta_{nn'}$ \\
\phantom{v případě degenerované:} \quad $P_a = \sumint_\alpha\ket{a,\alpha}\bra{a,\alpha}$, \quad $P_aP_{a'} = P_a\delta(a - a')$ \quad $\leftarrow$ \underline{ortogonalita}

\begin{priklad}[PŘ: v $\mathcal{H} = L^2(\mathbb{R}^3)$ operátor $Q_1\equiv\hat x_1$]
\[
  \hat x_1 = \int\dd x_1\,x_1P_{x_1};\quad\text{kde}\quad P_{x_1} = \int\dd x_2\dd x_3\ket{x_1x_2x_3}\bra{x_1x_2x_3}
  \quad\text{tj.}\ = \int\dd x_2\dd x_3\ket{\vec x}\bra{\vec x}
\]
nebo \quad $\hat x_1 = \int\dd^3x\,x_1\ket{\vec x}\bra{\vec x}$
\end{priklad}
Příklady operátorů se spojitým i diskr.\ spektrem uvidíme později (H pro systém s vázanými stavy)

\underline{pozn:} \quad $f(\hat A) = \sumint_a f(a)\hat P_a$ \qquad
\underline{pozn:} samosdr.\ $\hat A$ omezený $\Leftrightarrow\sigma(\hat A)$ omez. \quad $\|\hat A\| = \sup\sigma(\hat A)$

%----------------------------------------------------------------------
\strana{QM-F-6}
matematické fousy? \ co je $\ket{x}\bra{x}$ za objekt? \\
$\to$ my budeme pracovat intuitivně \dots\ např.\ $\ket{x}\braket{x}{\psi} = \psi(x)\cdot\ket{x}$ \\
podrobněji:

\begin{veta}[Spektrální teorém]
$\forall$ samosdružený operátor $A$ $\exists$ funkce $E(\lambda)$; kde $\lambda\in\mathbb{R}$ a $E$ je projekční operátor na $\mathcal{H}$ taková, že:
\begin{enumerate}[label=\arabic*),nosep]
  \item $\forall\lambda_1 < \lambda_2$: \quad $E(\lambda_1)E(\lambda_2) = E(\lambda_2)E(\lambda_1) = E(\lambda_1)$
  \item $\forall\ket{\psi}$: \quad $E(\lambda+\varepsilon)\ket{\psi}\to E(\lambda)\ket{\psi}$, \quad $\varepsilon\to 0^+$
  \item \phantom{$\forall\ket{\psi}$:} \quad $E(\lambda)\ket{\psi}\to 0$; \quad $\lambda\to -\infty$
  \item \phantom{$\forall\ket{\psi}$:} \quad $E(\lambda)\ket{\psi}\to\ket{\psi}$; \quad $\lambda\to +\infty$
  \item $\int\lambda\,\dd E(\lambda) = A$
\end{enumerate}
\end{veta}
kde $\int f(\lambda)\,\dd E(\lambda)$ je definovaný jako (Stieltjesův integrál):
\[
  \int f(\lambda)\,\dd E(\lambda) \equiv \lim_{\substack{n\to\infty\\ \Delta\lambda\to 0}}\sum_{k=1}^{\infty\,(n)} f(\lambda_k)\big[E(\lambda_k) - E(\lambda_{k-1})\big]
\]

\underline{poznámka:}
\begin{itemize}
  \item Tato formulace je trik, jak napsat předchozí formulaci v $\overline{\mathcal{H}}$ do $\mathcal{H}$. Přičemž
  \[ E(\lambda) = \sum_{a_n<\lambda}P_n + \int_{a<\lambda}P_a\,\dd a \qquad (*) \]
  $\to$ ověřte, že $(*)$ splňuje 1)--5)
\end{itemize}

Ad 5):
\begin{center}
\begin{tikzpicture}[x=1cm,y=1cm]
  \draw[->] (-0.5,0) -- (9.5,0) node[right] {$a$};
  \draw[->] (0,-0.3) -- (0,4.2) node[right] {$E$};
  \draw[thick] (0.8,0) -- (0.8,1.2) -- (2,1.2) -- (2,2.2) -- (3.8,2.2);
  \fill (1.2,1.2) circle (2pt);
  \node[above] at (1.2,1.25) {$E = P_1$};
  \node[above] at (2.7,2.25) {$E = P_1{+}P_2$};
  \draw[thick,domain=3.8:8,samples=40] plot (\x,{2.2 + 0.6*sqrt(\x-3.8)*1.3});
  \fill (5,{2.2+0.6*sqrt(1.2)*1.3}) circle (2pt);
  \draw[->] (5,{2.2+0.6*sqrt(1.2)*1.3}) to[bend right=20] (5.8,2.9) node[right] {$E = P_1 + P_2 + \int_{-\infty}^{\lambda}P_a\,\dd a$};
  \draw[ultra thick,gray] (3.8,0) -- (9,0);
  \draw[dashed,gray] (3.8,0) -- (3.8,2.2);
  \node[below] at (0.8,0) {$a_1$};
  \node[below] at (2,0) {$a_2$};
  \draw[decorate,decoration={brace,mirror,amplitude=4pt}] (0.6,-0.5) -- (2.2,-0.5) node[midway,below=4pt] {$\sigma_b(A)$};
  \draw[decorate,decoration={brace,mirror,amplitude=4pt}] (3.8,-0.5) -- (9,-0.5) node[midway,below=4pt] {$\sigma_c(A)$};
\end{tikzpicture}
\end{center}

tj.\ pro spojité spektrum:
\[
  \int\lambda\,\dd E(\lambda) \equiv \sum_{k=1}^{\infty}\underbrace{f(\lambda_k)}_{\lambda_k}\underbrace{\big[E(\lambda_k) - E(\lambda_{k-1})\big]}_{\Delta\lambda_k\to 0\ \dots\ E\text{ spojité: }\Delta E\doteq\frac{\dd E}{\dd\lambda}\cdot\Delta\lambda}
  = \sum\lambda_k\frac{\dd E}{\dd\lambda}(\lambda_k)\cdot\Delta\lambda
\]
\[
  \xrightarrow{\ \Delta\lambda\to 0\ } \int\lambda\frac{\dd E(\lambda)}{\dd\lambda}\,\dd\lambda = \int\lambda P_\lambda\,\dd\lambda
  \qquad\text{neboť}\quad \dv{\lambda}\int_{-\infty}^{\lambda}P_a\,\dd a = P_\lambda
  \quad\text{tj.}\quad \vzorec{\frac{\dd E(\lambda)}{\dd\lambda} = P_\lambda}
\]
ale $E(\lambda)$ má smysl i $\mathcal{H}$; \ $P_\lambda$ jen v $\overline{\mathcal{H}}$

%----------------------------------------------------------------------
\strana{QM-F-7}
\underline{diskrétní spektrum:} \\
stačí si uvědomit, že $E(\lambda_k) - E(\lambda_{k-1})\neq 0$ jen v okolí bodů spektra a tam je
$E(\lambda_k) - E(\lambda_{k-1}) = P_n$ \quad \dots\ $\lambda_k\geq a_n > \lambda_{k-1}$ \\
tj.: \quad $\int\lambda\,\dd E(\lambda) = \sum_n a_nP_n$
\begin{itemize}
  \item lze také chápat jako $E(\lambda) = \sum_n\Theta(\lambda - a_n)P_n$, potom:
  \[
    \int\lambda\,\dd E(\lambda) = \int\lambda\frac{\dd E}{\dd\lambda}\cdot\dd\lambda
    = \int\lambda\sum_n\underbrace{\dv{\lambda}\Theta(\lambda - a_n)}_{\delta(\lambda - a_n)}P_n\,\dd\lambda = \sum_n a_nP_n
  \]
\end{itemize}
\noindent\rule{\linewidth}{0.3pt}

\subsection*{Doplnění k axiomům QM}
\begin{veta}[Axiomy QM]
\begin{itemize}[nosep]
  \item[\fbox{A1}] Stav systému je dán paprskem v $\mathcal{H}$.
  \item[\fbox{A2}] Každé dynamické proměnné (měřitelné veličině, pozorovatelné) odpovídá samosdružený lineární operátor, jehož spektrum udává možné výsledky měření.
  \item[\fbox{A3}] Výsledky měření předpovídáme pomocí spektrálního rozkladu operátoru: \\
  hustota pravděpodobnosti nalezení hodnoty $a$ ve stavu $\ket{\psi}$
  \[
    \mu_\psi(a)\,\dd a = \mel{\psi}{\hat P(a)}{\psi}\,\dd a \qquad \big(= \mel{\psi}{\dd E(a)}{\psi}\big)
  \]
  pravděpodobnost nalezení hodnoty $a\in\mathcal{M}$:
  \[
    \mu = \int_{\mathcal{M}}\mu_\psi(a)\,\dd a
  \]
  stav po měření: \quad $\ket{\psi'} = \int_{\mathcal{M}}P(a)\ket{\psi}\,\dd a$ \quad $\big(= \int_{\mathcal{M}}\dd E(a)\ket{\psi}\big)$
\end{itemize}
\end{veta}

\underline{pozn:}
\begin{itemize}[nosep]
  \item[--] dále platí \quad $\mu = \braket{\psi'}{\psi'} \equiv \|\psi'\|^2$ \quad
            \pozn{\dots\ $\braket{\psi'}{\psi'} = \int_{\mathcal{M}}\int_{\mathcal{M}}\mel{\psi}{\underbrace{P(a')P(a)}_{\delta(a,a')I}}{\psi}\,\dd a\,\dd a'$}
  \item[--] $\mu$ nalezení vůbec nějaké hodn.\ $= 1 = \int_{\mathbb{R}}\mu_\psi(a)\,\dd a = \mel{\psi}{\hat I}{\psi}$ \quad \pozn{$\hat I = \int_{\mathbb{R}}P_a\,\dd a$}
\end{itemize}

\begin{priklad}[PŘ: Pravděpodobnost nalezení částice v místě $x\in\langle a,b\rangle$ (1D)]
\dots\ částice ve stavu $\ket{\psi} = \int\psi(x)\ket{x}\,\dd x$ \\
\dots\ měřitelná veličina \quad $\hat Q\equiv\hat x = \int\ket{x}\bra{x}\,\dd x$ \quad \dots\ spektr.\ rozklad $P_x = \ket{x}\bra{x}$
\[
  \text{tj.}\quad \mu_{\langle a,b\rangle} = \int_a^b\braket{\psi}{x}\braket{x}{\psi}\,\dd x = \int_a^b|\psi|^2\,\dd x
\]
po měření ve stavu \quad $\ket{\psi'} = \int_a^bP_x\ket{\psi}\,\dd x = \int_a^b\psi(x)\ket{x}\,\dd x$ \\
\phantom{po měření ve stavu} \quad tj.\ $\psi'(x) = \chi_{\langle a,b\rangle}(x)\psi(x)$
\end{priklad}
\underline{pozn:} \quad $\rho(x)\equiv|\psi(x)|^2$ \quad \dots\ \kw{hustota pravd.\ výskytu} částice v místě $x$

%----------------------------------------------------------------------
\strana{QM-F-8}
\subsection*{Další statistické veličiny}
z matematického hlediska měření $A$ ve stavu $\ket{\psi}$ \\
$\to$ náhodná proměnná: opakované měření dá pokaždé jinou hodnotu $a\in\sigma(A)$ s rozdělovací funkcí $\mu_\psi(a)$

střední hodnota náhod.\ prom.: \quad $\langle a\rangle \equiv \int a\,\mu_\psi(a)\,\dd a = \int a\mel{\psi}{P_a}{\psi}\,\dd a$
\[
  \text{tj.}\quad \vzorec{\langle a\rangle = \mel{\psi}{A}{\psi}} \qquad (\text{normované } \ket{\psi}!)
\]
vyšší momenty náhod.\ prom.: \quad $\langle a^n\rangle \equiv \int a^n\mu_\psi(a)\,\dd a = \mel{\psi}{A^n}{\psi}$ \\
obec.\ $\to$ fce náhod.\ prom. \quad $\langle f(a)\rangle \equiv \int f(a)\mu_\psi(a)\,\dd a = \mel{\psi}{f(A)}{\psi}$

důležitý případ \underline{variance} (střední kvadr.\ odchylka, rozptyl): \\
\quad značení $\sigma_a^2$ nebo \quad $(\Delta a)^2 \equiv \big\langle(a - \langle a\rangle)^2\big\rangle = \langle a^2\rangle - \langle a\rangle^2$

\underline{pozn:} ve vl.\ stavu: $A\ket{\psi} = a\ket{\psi}$ \quad \dots\
$\left.\begin{aligned}\langle A\rangle &= a\\ \langle A^2\rangle &= a^2\end{aligned}\right\}\to(\Delta A)^2 = 0$

\subsection*{Relace neurčitosti}
\begin{veta}[Relace neurčitosti]
nechť $[\hat A,\hat B] = i\hat C$ \quad potom \quad $\vzorec{\Delta a\cdot\Delta b \geq \tfrac12|\mel{\psi}{\hat C}{\psi}|}$
\end{veta}
\underline{pozn:} im.\ jednotka zaručuje: pokud $A = A^\dagger$, $B = B^\dagger\Rightarrow C = C^\dagger$

\underline{DK:} \quad Schwarzova nerovnost: \quad $\|\phi_1\|\cdot\|\phi_2\| \geq |\braket{\phi_1}{\phi_2}|$ \quad $(*)$ \\
\quad pro \quad $\ket{\phi_1} = (A - \langle a\rangle)\ket{\psi}$ \quad \dots\ $\|\phi_1\| = \sqrt{\mel{\psi}{(A - \langle a\rangle)^2}{\psi}} = \Delta a$ \\
\quad podob.\ \quad $\ket{\phi_2} = (B - \langle b\rangle)\ket{\psi}$ \quad \dots\ $\|\phi_2\| = \Delta b$ \\
pravá str.\ $(*)$:
\begin{align*}
  |\braket{\phi_1}{\phi_2}| &= |\mel{\psi}{(A - \langle a\rangle)(B - \langle b\rangle)}{\psi}| = |\mel{\psi}{AB - \langle a\rangle\langle b\rangle}{\psi}| \\
  &= \Big|\bra{\psi}\frac{AB - BA}{2} + \frac{AB + BA}{2} - \langle a\rangle\langle b\rangle\ket{\psi}\Big| \\
  &= \Big|\underbrace{\tfrac{i}{2}\mel{\psi}{C}{\psi}}_{\text{ryze imag}} + \underbrace{\bra{\psi}\tfrac{AB+BA}{2}\ket{\psi} - \langle a\rangle\langle b\rangle}_{\text{reálné}}\Big|
  \geq \tfrac12|\mel{\psi}{C}{\psi}| \qquad\text{c.b.d.}
\end{align*}

nejznámější příklad: \quad $[x_1,p_1] = i\hbar$ \quad $\to$ \quad $\vzorec{\Delta x\cdot\Delta p \geq \tfrac{\hbar}{2}}$

%----------------------------------------------------------------------
\strana{QM-F-9}
\subsection*{Kompatibilní pozorovatelné}
\begin{veta}
Pokud $[A,B] = 0$, pak vl.\ prostory $A$ jsou invariantní vůči $B$ a naopak:
\[ A\ket{\psi} = a\ket{\psi} \ \Rightarrow\ A\ket{\phi} = a\ket{\phi} \quad\text{kde}\quad \ket{\phi} = B\ket{\psi} \]
\end{veta}

\begin{veta}
Nechť $A, B$ jsou samosdružené operátory. \\
Potom existuje úplná množina společných vl.\ v.\ $\Leftrightarrow$ $[A,B] = 0$
\end{veta}

\underline{DK:} \hfill {\color{red}Dělal jsem jinak -- viz dodatek} \\
\quad ($\Rightarrow$) je triviální důsledek spektrálního rozkladu \\
\quad ($\Leftarrow$) konstruktivní dk \dots\ sada vektorů $P_b\ket{a_n}$ $\forall b, a_n$
\begin{itemize}[nosep]
  \item[\dots] některé mohou být nulové, ale je úplná: \quad $\ket{\psi} = \sum_n\psi_n\ket{a_n} = \sum_{nb}\psi_nP_b\ket{a_n}$
  \item[\dots] jsou vl.\ v.\ $B$ \dots\ $b$
  \item[\dots] jsou vl.\ v.\ $A$ \dots\ a) $(A - a_n)\hat P_b\ket{a_n}$ je také vl.\ v.\ $B$ přísl.\ $b$ \dots\ $\psi_{nb}$ \\
        neboť \quad $B(A - a_n)P_b\ket{a_n} = (A - a_n)BP_b\ket{a_n} = b(A - a_n)P_b\ket{a_n}$ \\
        b): \quad $(A - a_n)\ket{a_n} = \sum_b(A - a_n)P_b\ket{a_n} = 0 = \sum_b\psi_{nb}$ \\
        $\to$ jednotlivé členy sumy jsou OG (různé $b$) a tedy $\psi_{nb} = 0$ $\forall b, n$
\end{itemize}

\underline{pozn:} důsl.\ \quad $[A,B] = 0 \Leftrightarrow [P_a,P_b] = 0\quad\forall a,b$ \\
\quad \pozn{DK: $\Leftarrow$ spektr.\ rozklad; $\Rightarrow$ rozepíšu $P_a, P_b$ pomocí spol.\ báze}

\underline{pozn:} některé $P_b\ket{a_n}$ mohli být nulové $\to$ množina $\forall\ket{ba_n}$ nemusí být kartézský součin \dots\ později např.\ $\ket{\ell m}$

+ Rozšíření na více operátorů
\subsection*{Úplný systém komutujících operátorů (ÚSKO)}
$\hat A^{(1)},\dots,\hat A^{(N)}$ navzájem komutují $\Rightarrow$ $\exists$ spol.\ báze vl.\ v.\ $\forall$ operátorů

\begin{definice}[ÚSKO]
$\forall$ sada vl.\ č.\ $\{a^{(1)},\dots,a^{(N)}\}$ (každá přípustná) jednoznačně def.\ vl.\ v.\ $\ket{\psi}$ (až na fázi) \\
\dots\ \underline{značení} \quad $\ket{\psi} = \ket{a^{(1)}\dots a^{(N)}}$
\end{definice}

\begin{veta}
operátor $\hat F$ komutuje se $\forall A^{(1)},\dots,A^{(N)}$ prvky ÚSKO $\Rightarrow$ $\hat F = f(A^{(1)},\dots,A^{(N)})$
\end{veta}
\underline{DK:} $\hat F, A^{(1)},\dots,A^{(N)}$ mají spol.\ sadu vl.\ v., ale to musí být $\ket{a^{(1)}\dots a^{(N)}}$ \\
\quad spektr.\ rozklad $\hat F$ $\to$ $\forall a^{(1)},\dots,a^{(N)}$ $\exists$ jediné $f = f(a^{(1)},\dots,a^{(N)})$

\underline{pozn:} celou množinu $\hat A^{(1)},\dots,\hat A^{(N)}$ lze považovat za jeden oper.\ s nedegenerovaným spektrem \dots\ vl.\ č.\ vektor

\textbf{PŘ:} \quad $Q_1, Q_2, Q_3$ \quad $\to$ spol.\ báze $\ket{\vec x}$

%----------------------------------------------------------------------
\strana{QM-F-10}
\begin{priklad}[typický příklad]
$\mathcal{H} = \mathcal{H}^{(A)}\otimes\mathcal{H}^{(B)}$; \quad
$\begin{aligned}[t] &A^{(1)}\dots A^{(N)} \ \dots\ \text{ÚSKO NA } \mathcal{H}^{(A)} \\ &B^{(1)}\dots B^{(M)} \ \dots\ \text{ÚSKO NA } \mathcal{H}^{(B)}\end{aligned}$

potom \quad $A^{(1)}\otimes I,\dots,A^{(N)}\otimes I,\ I\otimes B^{(1)},\dots,I\otimes B^{(M)}$ \ je ÚSKO na $\mathcal{H}$
\end{priklad}
\underline{pozn:} operátor $C = A + B$ \dots\ diagonalizovat \dots\ $[A,B] = 0 \Leftrightarrow$ najít spol.\ bázi $\to$ častý případ v QM

\subsection*{Různé ekvivalentní reprezentace QM}
\begin{itemize}
  \item \underline{volba ÚSKO} $\to$ báze + vyjádření vln.\ fce a oper.\ v bázi
\end{itemize}
\begin{itemize}[nosep]
  \item[--] diskrétní případ: \quad $I = \sum_n\ket{n}\bra{n}$ \quad \dots\ stavy \quad
  $\ket{\psi} = \sum_n\ket{n}\underbrace{\braket{n}{\psi}}_{\psi_n} \equiv \ket{\psi}$ \quad (izomorfizm.) ztotožnění \\
  operátory \quad $A\ket{\psi} = \sum_{nn'}\ket{n}\underbrace{\mel{n}{A}{n'}\braket{n'}{\psi}}_{A_{nn'}\psi_{n'}}$ \quad \dots\ $A\equiv A_{nn'}$
  \item[--] spojitý případ: \quad $I = \int\ket{\nu}\bra{\nu}\,\dd\nu$ \quad
  $\ket{\psi} = \int\dd\nu\,\psi(\nu)\ket{\nu}$ \quad ($\braket{\nu}{\psi} \equiv \ket{\psi}$)
  \[
    A\ket{\psi} = \int\dd\nu\,\dd\nu'\ket{\nu}\underbrace{\mel{\nu}{A}{\nu'}}_{A_{\nu\nu'}\equiv A(\nu,\nu')}\braket{\nu'}{\psi}
    \quad\dots\ \text{jádro integrálního operátoru}
  \]
\end{itemize}

\begin{priklad}[PŘ]
\begin{tabular}{lll}
$\hat x$ & $\leftrightarrow\ x\,\delta(x - x')$ & \dots\ $\int\dd x'\,x\,\delta(x - x')\psi(x') = x\psi(x)$ \\[2pt]
$V(\hat x)$ & $\phantom{\leftrightarrow}\ V(x)\delta(x - x')$ & $A\ket{\psi}$ \dots\ $V(x)\psi(x)$ \\[2pt]
$\hat p$ & $\leftrightarrow\ {-i\hbar}\,\delta'(x - x')$ & $= \int\dd x'\,i\hbar\,\delta'(x - x')\psi(x')\,\dd x' = -i\hbar\dv{x}\psi(x)$
\end{tabular}

\quad $\uparrow$ jiný zápis $-i\hbar\,\delta(x - x')\dv{x'}$

\pozn{podrobněji \ $-i\hbar\,\delta'(x - x') \equiv -i\hbar\dv{x}\delta(x - x') = +i\hbar\dv{x'}\delta(x - x')$ + per partes}
\end{priklad}

%----------------------------------------------------------------------
\strana{QM-F-11}
\subsection*{DODATEK}
\pozn{celá strana psána tužkou}


\textbf{Lemma 1:} \quad $[A,B] = 0\ \wedge\ B\ket{\psi} = b\ket{\psi}$ \quad $\Rightarrow$ \quad $B\ket{\phi} = b\ket{\phi}$; kde $\ket{\phi} = A\ket{\psi}$ \\
\quad DK: zjevné

\textbf{Lemma 2:} \quad $[A,B] = 0$, \ $a_1\neq a_2$ \quad $\Rightarrow$ \quad $\mel{a_1}{B}{a_2} = 0$ \\
\quad DK: \quad $0 = \mel{a_1}{[A,B]}{a_2} = \mel{a_1}{AB}{a_2} - \mel{a_1}{BA}{a_2}$ \\
\quad\phantom{DK:} tj. \ $0 = (a_1 - a_2)\mel{a_1}{B}{a_2}$ \quad $\Rightarrow$ c.b.d.

\textbf{Lemma 3a:} \quad $[A,B] = 0$ \quad $\Rightarrow$ \quad $[P_a,B] = 0$ \\
\quad DK: \ v bázi vl.\ v.\ $A$: \quad $A\ket{a,\alpha} = a\ket{a,\alpha}$ $\to$ $P_a = \sum_\alpha\ket{a\alpha}\bra{a\alpha}$
\begin{align*}
  \mel{a_1\alpha_1}{[P_a,B]}{a_2\alpha_2} &= \bra{a_1\alpha_1}\sum_\alpha\ket{a\alpha}\bra{a\alpha}B\ket{a_2\alpha_2}
   - \underbrace{\bra{a_1\alpha_1}B\sum_\alpha\ket{a\alpha}}_{\delta_{a_1a}\mel{a}{B}{a}}\underbrace{\braket{a\alpha}{a_2\alpha_2}}_{\delta_{aa_2}\delta_{\alpha\alpha_2}} \\
  \text{podle lemm.\ 1, 2!}\qquad &= \delta_{aa_1}\delta_{aa_2}\big\{\mel{a\alpha_1}{B}{a\alpha_2} - \mel{a\alpha_1}{B}{a\alpha_2}\big\} \\
  &= 0
\end{align*}

\textbf{Lemma 3:} \quad $[A,B] = 0$ \quad $\Rightarrow$ \quad $[P_a,P_b] = 0\quad\forall a,b$ \\
\quad DK: uplatnění lemm.\ 3a na $[P_a,B]$

$\Rightarrow$ $\exists$ společná báze \dots\ projektory $P_aP_b = P_bP_a$


\newpage
%======================================================================
% PoznQMI-5.pdf
%======================================================================
\section{Částice v 1D}
\strana{QM-C-1}

\begin{itemize}
  \item ÚSKO $= \hat x$ \quad \dots\quad $\hat p = -i\hbar\dv{x}$ \qquad
        $H = \dfrac{p^2}{2m} = -\dfrac{\hbar^2}{2m}\dv[2]{x}$
\end{itemize}

\underline{více o hybnosti:} \hfill De Broglie \quad $\lambda = \dfrac{2\pi\hbar}{p} = \dfrac{h}{p}$
\quad \pozn{$\leftarrow$ $x$ změna fáze $= 2\pi$}

\quad vl.\ vektory \dots\ $\psi(x) = Ne^{\frac{i}{\hbar}px} \equiv \braket{x}{p}$ \quad \dots\ pro $p\in\mathbb{R}$ \quad (jinak není z $\bar{\mathcal{H}}$)

normalizace:
\[
  \braket{p}{p'} = |N|^2\int e^{\frac{i}{\hbar}(p'-p)x}\,\dd x = 2\pi\hbar\,|N|^2\,\delta(p - p')
  \quad\dots\quad N = \frac{1}{\sqrt{2\pi\hbar}}
\]
\quad $\nwarrow$ obě substituci z \ $\delta(x) = \int e^{2\pi ikx}\,\dd k$

\kw{relace úplnosti:} \quad $\vzorec{I = \int\ket{p}\bra{p}\,\dd p}$ \\
\quad $\swarrow$ v $x$ reprez.: \quad $\delta(x - x') = \dfrac{1}{2\pi\hbar}\displaystyle\int\dd p\; e^{ip(x-x')/\hbar}$

\underline{\kw{hybnostní reprezentace}:}
\[
  \braket{p}{\psi} = \psi(p) = \int\dd x\,\braket{p}{x}\braket{x}{\psi}
  = \vzorec{\frac{1}{\sqrt{2\pi\hbar}}\int\dd x\; e^{-\frac{i}{\hbar}px}\psi(x)}
  \qquad\text{Fourierova transf.}
\]
operátor hybnosti \quad $\hat p\ket{\psi} = p\,\psi(p)$ \quad \dots\ věta o derivaci F.T.

operátor souřadnice:
\[
  \mel{p}{\hat x}{\psi} = \int\dd x\,\braket{p}{x}\,x\,\braket{x}{\psi}
  = \int\dd x\Big[i\hbar\pdv{p}\braket{p}{x}\Big]\braket{x}{\psi} = i\hbar\pdv{p}\psi(p)
\]
$\to$ někdy se hodí pro řešení příkladů (částice v konst.\ poli)

\begin{priklad}[Př: částice v homogenním poli]
$\hat H = \dfrac{p^2}{2m} + V(x)$, \qquad kde $V(x) = -Fx$
\begin{itemize}[nosep]
  \item[--] stacionární stavy: \quad $H\ket{\psi} = E\ket{\psi}$ \quad \dots\ jednodušší v $p$-reprez.
  \[
    H\psi(p) = \frac{p^2}{2m}\psi(p) - i\hbar F\psi'(p) = E\psi(p)
    \quad \pozn{$\to \frac{\psi'}{\psi} = \frac{-E}{i\hbar F} + \frac{p^2}{2m}\cdot\frac{1}{i\hbar F}$}
  \]
  \[
    \to\quad \psi(p) = A\exp\Big\{\frac{i}{\hbar}\Big(\frac{E}{F}p - \frac{p^3}{6mF}\Big)\Big\}
  \]
  \item[--] řešení v $x$-reprez.\ pomocí F.T.:
  \[
    \psi(x) \equiv \braket{x}{\psi} = \int\dd p\,\braket{x}{p}\braket{p}{\psi}
    = \frac{1}{\sqrt{2\pi\hbar}}\int_{-\infty}^{\infty}\dd p\,\exp\Big\{\frac{i}{\hbar}\Big(px + p\frac{E}{F} - \frac{p^3}{6mF}\Big)\Big\}
  \]
  \pozn{jen Re část; Im -- lichá $\int = 0$} \\
  dá se vyjádřit pomocí spec.\ fce \quad $\mathrm{Ai}(x) \equiv \dfrac{1}{\pi}\displaystyle\int_0^\infty\cos\Big(\frac{t^3}{3} + xt\Big)\dd t$
  \[
    \to\quad \vzorec{\psi(x) = C\cdot\mathrm{Ai}\Big[\sqrt[3]{\tfrac{2mF}{\hbar^2}}\Big(x + \frac{E}{F}\Big)\Big]}
  \]
\end{itemize}
\end{priklad}

\underline{\kw{energetická reprezentace}}

$H_0 = \dfrac{p^2}{2m}$ \quad \dots\ $\ket{p}$ jsou vl.\ fce \quad $H_0\ket{p} = \dfrac{p^2}{2m}\ket{p}$ \\
\dots\ pro dané $E$ dvě vl.\ fce \quad $p = \pm\sqrt{2mE}$ \quad \dots\ $P_E = \ket{p}\bra{p} + \ket{-p}\bra{-p}$ \\
$H_0$ není ÚSKO \dots\ dá se přidat \kw{parita} \quad $\mathcal{P}\psi(x) \equiv \psi(-x)$ \\
snadno ověříte \dots\ $\mathcal{P}^\dagger = \mathcal{P}$, \ $\mathcal{P}^2 = I$ \quad $\to$ vl.\ č.\ $\mathcal{P}\ket{\psi} = \lambda\ket{\psi}$, \ $\lambda^2 = 1$ \\
\hspace*{6cm} tj.\ $\lambda = \pm1$ \quad \dots\ parita sudá/lichá

%----------------------------------------------------------------------
\strana{QM-C-2}
Platí $[\mathcal{P},H_0] = 0$ \quad \dots\ společná báze \quad
$\ket{E\lambda} = \dfrac{N}{\sqrt2}\big(\ket{p} \pm \ket{-p}\big)$

normalizace:
\begin{align*}
  \braket{E\lambda}{E'\lambda'} &= \delta_{\lambda\lambda'}\frac{N^2}{2}\big(\bra{p} \pm \bra{-p}\big)\big(\ket{p'} \pm \ket{-p'}\big)
    \qquad p,p' > 0 \quad \pozn{$\delta(p+p')$ nepřisp.} \\
  &= \delta_{\lambda\lambda'}N^2\,\delta(p - p') = \delta_{\lambda\lambda'}N^2\frac{p}{m}\,\delta(E - E')
\end{align*}

\begin{veta}[Důležitá pozn.: substituce v $\delta$-funkci]
\[
  \vzorec{\delta\big(f(x)\big) = \sum_{x_0}\frac{\delta(x - x_0)}{|f'(x_0)|}}
\]
tj.\ \quad $\delta\big(\sqrt{2mE} - \sqrt{2mE'}\big) = \sqrt{\dfrac{2E}{m}}\,\delta(E - E') = \dfrac{p}{m}\,\delta(p - p')$
\end{veta}

tj.: \quad $N = \sqrt{\dfrac{m}{p}}$ \quad a tedy \quad
$\ket{E+} = \sqrt{\dfrac{m}{\pi\hbar p}}\cos(px)$, \qquad
$\ket{E-} = i\sqrt{\dfrac{m}{\pi\hbar p}}\sin(px)$

spektrální rozklad: \quad
$I = \displaystyle\sum_{\lambda=\pm1}\int_0^\infty\dd E\,\ket{E\lambda}\bra{E\lambda}$, \qquad
$H_0 = \displaystyle\sum_{\lambda=\pm1}\int_0^\infty\dd E\,E\,\ket{E\lambda}\bra{E\lambda}$

\subsection*{Gaussovské vlnové balíky}
\[
  \psi(p) = \frac{1}{\sqrt[4]{2\pi(\Delta p)^2}}\,e^{-\frac{(p-p_0)^2}{4\Delta p^2}}
  \quad\xleftrightarrow{\ \text{F.T.}\ }\quad
  \psi(x) = \frac{1}{\sqrt[4]{2\pi(\Delta x)^2}}\,e^{-\left(\frac{x}{2\Delta x}\right)^2}e^{\frac{i}{\hbar}p_0x}
\]
kde \quad $\vzorec{\Delta p\,\Delta x = \frac{\hbar}{2}}$ \quad \dots\
$\langle p\rangle_\psi = p_0$; \ $\langle x\rangle_\psi = x_0$; \ $\langle(p-p_0)^2\rangle = \Delta p^2$, \ $\langle(x-x_0)^2\rangle = \Delta x^2$

\underline{DK:}
\begin{itemize}
  \item Normalizace: \quad G-integrál: $\int e^{-x^2}\dd x = \sqrt\pi$
  \[
    \braket{\psi}{\psi} = |N|^2\int e^{-2\left(\frac{p-p_0}{2\Delta p}\right)^2}\dd p
    = |N|^2\int e^{-\left(\frac{p-p_0}{\sqrt2\Delta p}\right)^2}\frac{\dd(p-p_0)}{\sqrt2\Delta p}\sqrt2\Delta p
    = |N|^2\cdot\sqrt\pi\cdot\sqrt2\Delta p = 1 \ \checkmark
  \]
  \item Střední hodnota:
  \[
    \langle p\rangle_\psi = N^2\int p\,e^{-\frac{(p-p_0)^2}{2\Delta p^2}}\dd p
    = N^2\underbrace{\int(p-p_0)e^{-\frac{(p-p_0)^2}{2\Delta p^2}}\dd p}_{0} + p_0\underbrace{\braket{\psi}{\psi}}_{1} \ \checkmark
  \]
  \item variance:
  \[
    \langle(p-p_0)^2\rangle_\psi = N^2\int(p-p_0)^2e^{-\left(\frac{p-p_0}{\sqrt2\Delta p}\right)^2}\dd p
    = N^2\big(\sqrt{2\Delta p^2}\big)^3\int\underbrace{\Big(\frac{p-p_0}{\Delta p\sqrt2}\Big)^2}_{\xi^2}e^{-\xi^2}\dd\xi
  \]
  přitom \quad $\int\xi^2e^{-\xi^2}\dd\xi = -\dv{I}{\alpha}\Big|_{\alpha=1} = +\frac12\sqrt\pi$, \quad
  kde $I(\alpha) = \int e^{-\alpha\xi^2}\dd\xi = \sqrt{\pi/\alpha}$ \\
  tj.\ \quad $\langle(p-p_0)^2\rangle_\psi = \dfrac{1}{\sqrt{2\pi\Delta p^2}}\cdot\big(\sqrt{2\Delta p^2}\big)^3\cdot\frac12\sqrt\pi = (\Delta p)^2$ \ $\checkmark$
  \item ověření
  \begin{align*}
    \psi(x) &= \int\braket{x}{p}\psi(p)\,\dd p
      = \frac{1}{\sqrt[4]{2\pi}}\frac{1}{\sqrt{\Delta p}}\frac{1}{\sqrt{2\pi\hbar}}\int e^{\frac{i}{\hbar}px}e^{-\frac{(p-p_0)^2}{4\Delta p^2}}\dd p \\
    &= \frac{1}{\sqrt[4]{2\pi}}\cdot\frac{1}{\sqrt{\Delta p\cdot 2\pi\hbar}}\,e^{\frac{i}{\hbar}p_0x}
       \int\exp\Big[-\underbrace{\Big(\Big(\frac{p-p_0}{2\Delta p}\Big)^2 - \frac{i}{\hbar}(p-p_0)x\Big)}_{\left(\frac{p-p_0}{2\Delta p} - \frac{i}{\hbar}x\Delta p\right)^2 + \frac{x^2}{\hbar^2}(\Delta p)^2}\Big]\dd x \\
    &= \frac{1}{\sqrt[4]{2\pi}}\cdot\frac{1}{\sqrt{\Delta p\cdot2\pi\hbar}}\cdot 2\Delta p\sqrt\pi\cdot
       e^{-\frac{x^2}{\hbar^2}\Delta p^2}e^{\frac{i}{\hbar}p_0x}
     = \frac{1}{\sqrt[4]{2\pi}}\cdot\sqrt{\frac{2\Delta p}{\hbar}}\cdot e^{\frac{i}{\hbar}p_0x}\,
       e^{-\left(\frac{x}{2}\frac{2\Delta p}{\hbar}\right)^2}
  \end{align*}
  \pozn{$\frac{2\Delta p}{\hbar} = \frac{1}{\Delta x}$} \ $\checkmark$
\end{itemize}

%----------------------------------------------------------------------
\strana{QM-C-3}
\subsection{Částice v potenciálovém poli \quad $H = \frac{p^2}{2m} + V(x)$}

\subsubsection*{$\bullet$ Lineární harmonický oscilátor}

\underline{Klasická mechanika} \quad
$H = \dfrac{p^2}{2m} + \dfrac12kx^2 = \dfrac{p^2}{2m} + \dfrac{m\omega^2}{2}x^2$, \qquad
$\omega = \sqrt{\dfrac{k}{m}}$, \ tj.\ $T = \dfrac{2\pi}{\omega}$

Hamilton: \quad
$\left.\begin{aligned}
  \dot p &= -\pdv{H}{x} = -kx \\
  \dot x &= \pdv{H}{p} = \frac{p}{m}
\end{aligned}\right\}
\ \longrightarrow\
\begin{aligned}
  x &= x_0\cos[\omega(t-\tau)] \\
  p &= -m\omega x_0\sin[\omega(t-\tau)]
\end{aligned}$

ještě hezčí \dots\ \kw{bezrozměrné veličiny}: \\
(problém: v klasické mech.\ neexistuje typ.\ délková škála) \\
(pro LHO $\to$ půjčíme si $\hbar$ z QM \quad \fbox{$x_0 = \sqrt{\dfrac{\hbar}{m\omega}}$}, \quad $p_0 = \dfrac{\hbar}{x_0}$ \pozn{$=\sqrt{\hbar\omega m}$})

\begin{minipage}{0.4\textwidth}
$p = \sqrt{\dfrac{1}{m\hbar\omega}}\,P$ \qquad $q = \sqrt{\dfrac{m\omega}{\hbar}}\,x = x/x_0$ \\[4pt]
\quad \dots\ $p_0 = \dfrac{\hbar}{x_0}$
\[
  \to\quad
  \begin{aligned}
    q &= q_0\cos[\omega(t-\tau)] \\
    p &= -q_0\sin[\omega(t-\tau)]
  \end{aligned}
\]
\end{minipage}
\begin{minipage}{0.25\textwidth}
\begin{tikzpicture}[gray,scale=0.8]
  \draw[->] (-1.4,0) -- (1.6,0) node[right] {$q$};
  \draw[->] (0,-1.4) -- (0,1.6) node[above] {$p$};
  \draw (0,0) circle (1.1);
  \draw[->] (0,0) -- (0.65,0.89) node[midway,left] {\small $q_0$};
\end{tikzpicture}
\end{minipage}
\begin{minipage}{0.33\textwidth}
mapování do $\mathbb{C}$:
\[
  a \equiv \frac{1}{\sqrt2}(q + ip) = a_0e^{-i\omega t}
\]
\end{minipage}

\noindent\rule{\linewidth}{0.3pt}

\underline{Kvantová mechanika}
\[
  H = \frac{p^2}{2m} + \frac{m\omega^2}{2}x^2
  \quad\Big(\overset{x\text{-repr.}}{=} -\frac{\hbar^2}{2m}\dv[2]{x} + \frac{m\omega^2}{2}x^2
  \ \Big|\ \overset{p\text{-repr.}}{=} \frac{1}{2m}p^2 - \frac{\hbar^2m\omega^2}{2}\dv[2]{p}\Big)
\]
opět přejdeme k bezrozm.: \quad (OVĚŘTE: $p = -i\pdv{q}$, \ $q = i\pdv{p}$)
\[
  \vzorec{H = \frac{\hbar\omega}{2}(p^2 + q^2)}
  \qquad
  \Big[\overset{q\text{-repr.}}{=} \frac{\hbar\omega}{2}\Big(-\pdv[2]{q} + q^2\Big)
  \ \Big|\ \overset{p\text{-repr.}}{=} \frac{\hbar\omega}{2}\Big(p^2 - \pdv[2]{p}\Big)\Big]
\]

Podobně jako v klasickém případě zavedeme:
\[
  \vzorec{\hat a = \tfrac{1}{\sqrt2}(\hat q + i\hat p)}
  \quad\longleftrightarrow\quad
  \begin{aligned}
    \hat q &= \tfrac{1}{\sqrt2}(\hat a + \hat a^\dagger) \\
    \hat p &= \tfrac{1}{i\sqrt2}(\hat a - \hat a^\dagger)
  \end{aligned}
  \qquad\text{navíc: } [\hat q,\hat p] = \tfrac{1}{\hbar}[\hat X,\hat P] = i
\]
\[ \to\quad \hat a^\dagger = \tfrac{1}{\sqrt2}(\hat q - i\hat p) \]

\subsubsection*{VLASTNÍ VEKTORY $H$}
pozn: připomenout $\to$ stacionární stavy; význam pro měření $E$

Lze řešit přímo \pozn{přeškrtnuto: $\big[-\frac{\hbar^2}{2m}\dv[2]{x} + \frac{m\omega^2}{2}x^2\big]\psi(x) = E\psi(x)$, $q = x/x_0$} \\
\quad $\to$ \ $\dfrac{\hbar\omega}{2}\Big(-\pdv[2]{q} + q^2\Big)\psi(q) = E\psi(q)$ \\
\quad + asymptotika $q\to\infty$ ($q \gg \frac{2E}{\hbar\omega}$) \dots\ $\psi\sim e^{\pm q^2/2}$
+ hledání řeš.\ ve tvaru $p(q)e^{-q^2/2}$ \\
\hspace*{6cm} $\to$ UDĚLÁME JINAK

pozn: stejná rovnice vyjde v $p$-reprezentaci (srovn.\ klasické řešení)

%----------------------------------------------------------------------
\strana{QM-C-4}
\underline{nalezení spektra:}
\[
  [a,a^\dagger] = \tfrac12\big(-i[\hat q,\hat p] + i[\hat p,\hat q]\big) = 1
\]
\begin{definice}
$\hat N = \hat a^\dagger\hat a$ \quad \dots\ samosdružený operátor, fyz.\ význam?
\[
  \hat N = \tfrac12\big(\hat q^2 + \hat p^2 + i[\hat q,\hat p]\big) = \frac{\hat H}{\hbar\omega} - \frac12
  \qquad\text{tj.}\quad \vzorec{\hat H = \hbar\omega\big(\hat N + \tfrac12\big)}
\]
\end{definice}
tj.\ $[\hat H,\hat N] = 0$ \quad \dots\ spol.\ báze \quad \dots\ def.\ $\hat N\ket{n} = n\ket{n}$ \\
\quad potom \quad $\hat H\ket{n} = \hbar\omega\big(n + \frac12\big)\ket{n}$

\underline{spektrum $\hat N$?}

\begin{minipage}{0.68\textwidth}
důležité komut.\ relace \quad
$\begin{aligned}
  [\hat N,\hat a] &= \hat a^\dagger[a,a] + [a^\dagger,a]a = -a \\
  [N,a^\dagger] &= -\big([N,a]\big)^\dagger = a^\dagger
\end{aligned}$
\end{minipage}
\begin{minipage}{0.3\textwidth}
\fbox{$\begin{aligned} [N,a] &= -a \\ [N,a^\dagger] &= a^\dagger \end{aligned}$}
\end{minipage}

důsledek: \quad
$\begin{aligned}[t]
  Na^\dagger\ket{n} &= \big\{a^\dagger N + [N,a^\dagger]\big\}\ket{n} = (n+1)a^\dagger\ket{n} \\
  Na\ket{n} &= \big\{aN + [N,a]\big\}\ket{n} = (n-1)a\ket{n}
\end{aligned}$

\dots\ interpretace: \quad $a^\dagger, a$ \dots\ \kw{kreační, anihilační operátor}

normalizace: \quad $\mel{n}{a^\dagger a}{n} = n$ \quad $\to$ \quad $\vzorec{a\ket{n} = c\ket{n-1} = \sqrt n\ket{n-1}}$ \\
\quad \pozor{pozor} $\to$ obsahuje fázovou konvenci, smíme volit jen fázi $\ket{0}$, ostatní fixované

podob.: \quad $\mel{n}{aa^\dagger}{n} = \mel{n}{a^\dagger a + [a,a^\dagger]}{n} = (n+1)$ \quad
$\to$ \quad $\vzorec{a^\dagger\ket{n} = \sqrt{n+1}\ket{n+1}}$

s každou vl.\ hodnotou $n$ spektrum obsahuje $n\pm1, n\pm2, \dots$ \\
\quad + $N$ je \underline{pozitivně definitní}: $\mel{\psi}{a^\dagger a}{\psi} = \|a\psi\|^2 \geq 0$ \dots\ nelze sestupovat nekonečně dolů \\
\quad $\to$ nejnižší vl.\ č.: $n = 0$: \quad $a\ket{0} = 0$ \quad \dots\ tj.\ $\sigma = \{n = 0,1,2,\dots\}$
\begin{itemize}[nosep]
  \item vlastní čísla $\nearrow$
  \item formálně i vlastní vektory! \quad $\vzorec{\ket{n} = \frac{(a^\dagger)^n}{\sqrt{n!}}\ket{0}}$
\end{itemize}

Vlastní vektory v $x$ repr.\ (nebo $p$) neznáme, ale umíme počítat maticové elementy operátorů:
\[
  \hat x = \sqrt{\frac{\hbar}{2m\omega}}\,(a + a^\dagger) \qquad
  \hat P = i\sqrt{\frac{m\hbar\omega}{2}}\,(a^\dagger - a)
\]
\pozn{$\sqrt{\frac{\hbar}{2m\omega}} = \frac{x_0}{\sqrt2}$, \quad $\sqrt{\frac{m\hbar\omega}{2}} = \frac{\hbar}{x_0\sqrt2}$}

přitom: \quad $\mel{n'}{a}{n} = \sqrt n\,\delta_{n',n-1}$ \hfill $\mel{n'}{a^\dagger}{n} = \sqrt{n+1}\,\delta_{n',n+1}$

{\color{gray}
\[
  a = \begin{array}{c|cccc}
       & 0 & 1 & 2 & \cdots \\ \hline
     0 & 0 & \sqrt1 & 0 & \\
     1 & 0 & 0 & \sqrt2 & \\
     2 & 0 & 0 & 0 & \sqrt3 \\
     \vdots & & & & \ddots
  \end{array}
  \ (\text{sloupce } n,\ \text{řádky } n')
  \qquad\qquad
  a^\dagger = \begin{pmatrix}
     0 & & & \\
     \sqrt1 & 0 & & \\
     & \sqrt2 & 0 & \\
     & & \sqrt3 & \ddots
  \end{pmatrix}
\]}

%----------------------------------------------------------------------
\strana{QM-C-5}
\[
  \to\quad \mel{n'}{x}{n} = \frac{x_0}{\sqrt2}\begin{pmatrix} 0 & \sqrt1 & & \\ \sqrt1 & 0 & \sqrt2 & \\ & \sqrt2 & \ddots & \ddots \\ & & \ddots & \end{pmatrix}
  \qquad
  \mel{n'}{p}{n} = \frac{i\hbar}{x_0\sqrt2}\begin{pmatrix} 0 & \sqrt1 & & \\ -\sqrt1 & 0 & \sqrt2 & \\ & -\sqrt2 & \ddots & \ddots \\ & & \ddots & \end{pmatrix}
\]
$\to$ každý operátor $f(p,q)$ lze vyjádřit pomocí $a, a^\dagger$

\underline{\kw{vlnová funkce v $x$-reprezentaci}:}

$a\ket{0} = 0$ \quad \dots\ $\psi(x) = \braket{x}{0}$ \quad \dots\ $\psi(x) = \phi\big(q = \frac{x}{x_0}\big)$ \\
\quad $\hookrightarrow$ \ $(q + ip)\phi(q) = \Big(q + \dv{q}\Big)\phi(q) = 0$ \quad $\to$ \underline{řešení} \ $\phi_0(q) = Ne^{-\frac12q^2}$
\[
  \text{tj.}\quad \vzorec{\psi_0(x) = \frac{1}{\sqrt[4]{\pi x_0^2}}\,e^{-\frac12(x/x_0)^2}}
  \quad\longleftarrow\quad \int e^{-(x/x_0)^2}\dd x = x_0\sqrt\pi = \braket{\phi_0}{\phi_0} \ \ \text{(normování)}
\]

\underline{excitované stavy:} \quad
$a^\dagger = \frac{1}{\sqrt2}(q - ip) = \frac{1}{\sqrt2}\Big(\frac{x}{x_0} - \frac{ix_0}{\hbar}(-i\hbar)\dv{x}\Big) = \frac{1}{x_0\sqrt2}\Big(x - x_0^2\dv{x}\Big)$
\[
  \to\quad \psi_n(x) \equiv \braket{x}{n} = \frac{\mel{x}{(a^\dagger)^n}{0}}{\sqrt{n!}}
  = \frac{1}{\sqrt[4]\pi}\frac{1}{\sqrt{2^nn!}}\frac{1}{x_0^{n+\frac12}}\Big[x - x_0^2\dv{x}\Big]^ne^{-\frac12(x/x_0)^2}
\]
\[
  \text{nebo}\quad \vzorec{\psi_n(x) = \sqrt{\frac{1}{\sqrt\pi\,x_0\,n!\,2^n}}\;H_n(x/x_0)\,e^{-\frac12(x/x_0)^2}} \qquad (*)
\]
kde \quad $H_n(q) \equiv e^{q^2/2}\Big[q - \dv{q}\Big]^ne^{-q^2/2}$ \quad \dots\ \kw{Hermitův polynom}

\underline{pozn:}
\begin{itemize}
  \item Hermitovy polynomy jako OG-poly s vahou $w(x) = e^{-x^2}$ \\
        $\to$ \ $\int H_n(q)H_{n'}(q)w(q)\,\dd q = 0$ \quad \dots\ obecná teorie OG-poly $\to$ Gaussova kvadratura \\
        \pozn{Gram-Schmidt ortogonalizace monomů}
  \item ortogonální polynomy obecně splňují rekurentní relace, které získáme násobením $\hat x$: \\
        $\to$ v energetické reprezentaci:
        \[
          \mel{x}{\hat x}{n} = x\psi_n(x) = \frac{x_0}{\sqrt2}\big(\mel{x}{a}{n} + \mel{x}{a^\dagger}{n}\big)
          = \frac{x_0}{\sqrt2}\big(\sqrt n\,\psi_{n-1}(x) + \sqrt{n+1}\,\psi_{n+1}(x)\big)
        \]
        + dosazení $(*)$: \pozn{faktor $\frac{1}{\sqrt{2^nn!}}$} \quad
        $\dfrac{x}{x_0}H_n\Big(\dfrac{x}{x_0}\Big) = nH_{n-1}\Big(\dfrac{x}{x_0}\Big) + \dfrac12H_{n+1}\Big(\dfrac{x}{x_0}\Big)$
  \item pozn: velmi obecný systém \dots\ model stavů poblíž minima \\
        $\leftrightarrow$ izomorfní s mnoha částicemi v 1hladinovém systému \\
        \dots\ spektrum $= N\epsilon_0$ \quad \dots\ bosony (později)
  \item pozn: \pozn{vytvářející funkce: $e^{\xi^2 - (\xi-\eta)^2}$ \dots\
        $H_n(\xi) = \dv[n]{\eta}e^{\xi^2-(\xi-\eta)^2}\Big|_{\eta=0}$ \quad
        $\to\ = \sum_{n=0}^\infty\frac{H_n(\xi)}{n!}\eta^n$}
\end{itemize}

%----------------------------------------------------------------------
\strana{QM-C-6}
\subsection{Obecné vlastnosti stacionárních stavů částice v 1D skalárním potenciálu}
\[
  H\psi = E\psi \quad\xrightarrow[x\text{-repr.}]{}\quad
  \vzorec{-\frac{\hbar^2}{2m}\psi'' + V(x)\psi(x) = E\psi(x)} \qquad (*)
\]

\underline{OBECNĚ:} $\bullet$ \ OBYČEJNÁ LINEÁRNÍ DIF.\ ROVNICE 2.\ ř.
\begin{itemize}[nosep]
  \item[$\to$] existují dvě lin.\ nezávislá řešení $\psi(x) = \phi_1(x), \phi_2(x)$
  \item[$\to$] Lineární nezávislost \dots\ def.\ \
        $W(x) = \begin{pmatrix}\phi_1(x) & \phi_2(x)\\ \phi_1'(x) & \phi_2'(x)\end{pmatrix}$ \\
        $\Leftrightarrow \det W \neq 0$ \quad přitom \quad $\dv{x}|W| = \phi_1\phi_2'' - \phi_2\phi_1'' = 0$ \quad \pozn{dosazení za $\phi''$}
  \item[$\to$] obecné řešení $(*)$: \quad
        $\begin{pmatrix}\psi(x)\\ \psi'(x)\end{pmatrix} = W(x)\cdot A$; \quad
        $A = \begin{pmatrix}A_1\\ A_2\end{pmatrix}$ \dots\ vektor koeficientů
\end{itemize}

\begin{itemize}
  \item Nespojitosti potenciálu -- \kw{napojování řešení}:
  \begin{itemize}[nosep]
    \item[$\to$] vlnová funkce musí být spojitá -- jinak generuje členy $\delta'(x)$ v potenciálu
    \item[$\to$] derivace $\psi'(x)$ musí být spojitá -- jinak členy $\delta(x)$ v potenciálu
  \end{itemize}
  pozn: někdy se toho využívá a konstruují se modelové systémy s $\delta$-funkčním potenciálem \\
  \quad $V(x) = \underbrace{v(x)}_{\text{regulární člen}} + \lambda\delta(x - x_0)$
  \[
    \text{potom:}\quad \psi'(x_0+\varepsilon) - \psi'(x_0-\varepsilon) \equiv \Delta\psi' = \int_{x_0-\varepsilon}^{x_0+\varepsilon}\psi''\,\dd x
    \overset{(*)}{=} \frac{2m}{\hbar^2}\int_{x_0-\varepsilon}^{x_0+\varepsilon}\big[-E + v(x) + \lambda\delta(x - x_0)\big]\psi(x)\,\dd x
  \]
  \[ \text{tj.}\quad \vzorec{\Delta\psi' = \frac{2m}{\hbar^2}\lambda\,\psi(x_0)} \]
  $\to$ v bodech, kde je nanejvýš nespojitost typu skoku, tj.\ $\psi,\psi'$ spojité, \\
  můžeme podmínky napojování psát: \quad
  $\begin{pmatrix}\psi(x_0)\\ \psi'(x_0)\end{pmatrix} = W_-\cdot A_- = W_+A_+$ \\
  kde $W_\pm$ jsou matice Wronskiánu pod a nad bodem $x_0$ \\
  a $A_\pm = \begin{pmatrix}A_1^\pm\\ A_2^\pm\end{pmatrix}$ jsou koeficienty rozvoje řešení \dots\ pokud známe \\
  $A_+$, můžeme dopočíst $A_-$ a naopak \dots\ 2 rovnice, dvě neznámé (regulární \dots\ $\det W\neq0$)

  $\to$ lze redukovat na jednu rovnici \dots\ spojitost $\psi'/\psi = \dv{x}\ln\psi$ \\
  \dots\ \kw{Logaritmická derivace} \quad \dots\ důležitý jen poměr $A_1/A_2$ \quad (zbytek = normování)
  \item z poznámek výše: možno zkonstruovat dvě globální LNZ řešení \dots\ zbytek \dots\ okraj.\ podmínky
\end{itemize}

%----------------------------------------------------------------------
\strana{QM-C-7}
\begin{itemize}
  \item \kw{okrajové podmínky}

  napišme $(*)$ jako: \quad $\dfrac{\psi''}{\psi} = -\dfrac{2m}{\hbar^2}\big(E - V(x)\big)$
  \ $= \Big(-\dfrac{p(x)^2}{\hbar^2}\Big)$ \quad \dots\ znaménko \\
  $\to$ terminologie z klas.\ fyziky: body obratu, klasicky povolená/zakázaná oblast

  $\to$ obvyklý případ: $V(|x|\to\infty)\to$ konst \dots\ bůno $= 0$ \quad (fázová konvence/volba počátku $E$)

  \underline{pro $E < 0$:} dvě lnz.\ řešení: \quad $\psi = Ae^{\varkappa x} + Be^{-\varkappa x}$ \quad $x\to\infty$, \quad
  $\varkappa = \sqrt{\frac{2m}{\hbar^2}|E - V|}$ \\
  \quad (konst.\ pot $V = 0$) \quad \dots\ vyhovuje jen $A = 0$ $\uparrow$ \qquad $B = 0$ pro $x\to-\infty$

  pozn: dá se DK silnější tvrzení (Formánkova učebnice): \\
  \quad pokud $\frac{2m}{\hbar^2}(V - E) \geq \varkappa^2$ pro $x\geq a$, pak $\exists$ řešení
  $\psi_1 \geq Ae^{\varkappa x}$ a $\psi_2 \leq Be^{-\varkappa x}$ \\
  \quad \pozn{stručně princip důkazu \dots}

  \underline{pro $E > 0$:} dvě lnz.\ řešení: \quad
  $\psi = A\underbrace{e^{ikx}}_{\text{částice}\to} + B\underbrace{e^{-ikx}}_{\text{částice}\leftarrow}$ \quad $x\to\pm\infty$, \quad
  $k^2 = \frac{2m}{\hbar^2}(E - V) = p^2/\hbar^2$ \\
  \quad $\to$ srovn.\ s volnou částicí

  pozn: opět se dá DK silnější (Formánek): \\
  \quad pokud $\frac{2m}{\hbar^2}\big(E - V(x)\big) = k^2 + \Delta(x) \xrightarrow{x\to\infty} k^2$, \quad
  kde $\lim_{x\to\infty}x^\alpha\Delta(x) = 0$ pro nějaké $\alpha > 1$ \\
  \quad pak $\exists$ dvě lnz řešení $\sin(kx + \delta)$ a $\cos(kx + \delta)$ \\
  \quad \pozn{$\to$ princip -- fázová rovnice -- konverguje ke konst.}
\end{itemize}

\begin{minipage}{0.68\textwidth}
$\to$ shrnutí \dots\ \kw{kvantová jáma} (potenciálová) (angl.\ studna = well) \\
\pozn{celé $V > 0$ neobsahuje případ $E < 0$}
\end{minipage}
\begin{minipage}{0.3\textwidth}
\begin{tikzpicture}[gray,scale=0.8]
  \draw[thin] (-2,0) -- (2.2,0);
  \draw[thick] plot[smooth] coordinates {(-2,0.05) (-1.2,0) (-0.5,-0.6) (0.3,-0.3) (1.1,0.45) (1.6,0.1) (2.2,0.05)};
  \draw[->] (1.6,-0.6) -- (1.15,0.3);
  \node[below right] at (1.4,-0.5) {\scriptsize nebo};
\end{tikzpicture}
\end{minipage}

klasická částice $\to$ povolené energie $E > V_0 = \min V(x)$ \\
\hspace*{3.2cm} $\to$ pro $E < 0$ \dots\ vázané stavy \\
\hspace*{3.2cm} $\to$ pro $E > 0$ \dots\ rozptylové (nebo za bariérou chycené)

kvantová částice:
\begin{veta}[Vázané stavy]
\fbox{$E < 0$} $\to$ \kw{vázané stavy} \dots\ $\psi\in L^2$, ale obě okrajové podmínky
($A = 0$ v $\infty$, $B = 0$ v $-\infty$) \\
lze splnit jen pro diskrétní $E = E_B$ \\
navíc \quad $E_B = \langle H\rangle = \underbrace{\langle T\rangle}_{\geq0} + \underbrace{\langle V\rangle}_{\geq V_0} \geq V_0$
\quad \dots\ tj.\ $E_B\in\langle V_0, 0)$
\end{veta}

\begin{veta}[Oscilační věta]
dk.\ rozborem chování $\frac{\psi''}{\psi} = -\frac{p^2(x)}{\hbar^2}$ \pozn{nesnadné}

$E_0 < E_1 < E_2 \dots$ energie váz.\ stavů \dots\ $\ket{\psi_\nu}$; $\nu = 0,1,2,\dots$ přísl.\ stavy \\
potom $\psi_\nu(x)$ má právě $\nu$ nulových bodů $\psi_\nu(x_0) = 0$ \\
a navíc mezi $\forall$ dvěma nul.\ body $\psi_\nu$ leží nulový bod $\psi_{\nu+1}$
\end{veta}
pozn: \dots\ semiklasický pohled $\hbar$ \dots\ plocha v $px$-prostoru \\
pozn: pokud $V(x) < 0\ \forall x$ \dots\ dá se dk, že $\exists$ alesp.\ jeden váz.\ stav (Var.\ princip)

%----------------------------------------------------------------------
\strana{QM-C-8}
\fbox{$E > 0$}
\begin{itemize}[nosep]
  \item[$\to$] jen zobecněné vl.\ stavy $\psi(x)$ omezená fce \\
        pro $\forall$ energii $E$ existují dva stavy \dots\ klasifikujeme později $\to$ teorie rozptylu
\end{itemize}

\begin{center}
\begin{tikzpicture}[>=Stealth]
  \node[left,align=right] at (-0.3,0.3) {okraj.\ podm.: \\ {\small lze DK, že jsou} \\ {\small ortogonální} $\to$};
  % levý box: částice nalétá zleva
  \draw (0,-0.6) rectangle (5,1.2);
  \node at (1.2,0.85) {$e^{ikx} + Re^{-ikx}$};
  \node at (3.9,0.85) {$Te^{ikx}$};
  \draw[->] (0.6,0.4) -- (1.6,0.4) arc (90:-90:0.12) -- (0.8,0.16);
  \draw[->] (3.4,0.35) -- (4.3,0.35);
  \node at (2.5,-0.3) {částice nalétá zleva};
  % pravý box: částice nalétá zprava
  \draw (5.3,-0.6) rectangle (10.3,1.2);
  \node at (6.2,0.85) {$Te^{-ikx}$};
  \node at (8.9,0.85) {$e^{-ikx} + Re^{ikx}$};
  \draw[->] (6.6,0.35) -- (5.7,0.35);
  \draw[->] (9.5,0.4) -- (8.5,0.4) arc (90:270:0.12) -- (9.3,0.16);
  \node at (7.8,-0.3) {částice nalétá zprava};
\end{tikzpicture}
\end{center}

$\to$ v případě symetrického potenciálu $V(x) = V(-x)$ \\
\quad tj.\ $[\mathcal{P},V] = 0$, \ $\mathcal{P}$ \dots\ parita, lze volit řešení s paritou $\lambda = \pm1$ \\
\hspace*{5cm} \dots\ $\psi(|x|\to\infty)\to\cos(k|x| + \delta), \ \pm\sin(k|x| + \delta)$ \\
pozn: $\hookrightarrow$ pro váz.\ stavy platí $\psi_\nu(-x) = (-1)^\nu\psi(x)$ \quad (srovn.\ oscilační věta)

$\to$ to, že $V(x)\in\mathbb{R}$, implikuje, že pokud $\psi$ řeší $(*)$, pak $\psi^*$ řeší \\
\quad tj.\ také $\frac12(\psi + \psi^*) = \mathrm{Re}\,\psi$ a $\frac{1}{2i}(\psi - \psi^*) = \mathrm{Im}\,\psi$ řeší $(*)$ \\
\quad \dots\ lze vždy vybrat reál.\ řešení \quad \pozn{zmínit se o $\mathrm{Im}\,V\neq0$ -- otevřené systémy}

\fbox{$V_{+\infty}\neq V_{-\infty}$} \dots\ hodí se jako model:
\begin{itemize}[nosep]
  \item[] $E < \min(V_+,V_-)$ \dots\ váz.\ stavy
  \item[] $\min(V_\pm) < E < \max(V_+,V_-)$ \dots\ $\forall E$ existuje jeden stav \dots\ kontinuum
  \item[] $E > \max(V_+,V_-)$ \dots\ dvě kontinua
\end{itemize}
pozn: v klas.\ fyzice \dots\ stavy v kontinuu (trapping) \\
\phantom{pozn:} v kvant.\ fyzice \dots\ není díky tunelování, ale později \dots\ rezonance

\fbox{$V = \infty$} \dots\ někdy je užitečný model \\
\quad \dots\ dá se chápat jako $\lim V\to\infty$ \\
\quad vyjde \dots\ $\psi = 0$ v oblasti, kde $V = \infty$ \\
\quad\quad $\psi'$ nemusí být spojitá na hranici oblasti \quad + pozn $\blacktriangle$ \dots\ na druhé str.

\subsection*{PŘESNĚ ŘEŠITELNÉ MODELY}
\begin{itemize}
  \item napojování a $\delta$-potenciály (včetně $V = \infty$)
  \item konst.\ potenciál, Lineární pot.\ a LHO (obecné řeš.\ definují spec.\ fci)
  \item \kw{Woods--Saxon} \quad $V(x) = -\dfrac{V_0}{1 + e^{\alpha x}}$ \quad $\psi$ \dots\ Hypergeometrická funkce

  \begin{minipage}{0.45\textwidth}
  \begin{tikzpicture}[gray,scale=0.9]
    \draw[thick,domain=-3:3,samples=60] plot (\x,{1.4/(1+exp(-2.5*\x))});
    \node[left] at (-3,0) {$-V_0$};
    \node[right] at (3,1.4) {$V = 0$};
  \end{tikzpicture}
  \end{minipage}
  \begin{minipage}{0.5\textwidth}
  pro $E > 0$:
  \[
    |R(E)|^2 = \frac{\sinh^2\big(\pi(k - k')/\alpha\big)}{\sinh^2\big(\pi(k + k')/\alpha\big)}
  \]
  \end{minipage}

  \kw{kvantový odraz} (nemá klas.\ obdobu)
\end{itemize}

%----------------------------------------------------------------------
\strana{QM-C-9}
\begin{itemize}
  \item pot.\ jáma a bariéra \quad $V(x) = \pm V_0\cosh(\alpha x)$ \quad \dots\ spec.\ fce.
  \item \kw{Morseho potenciál}: \quad $V(x) = D\big(e^{-2a(x-x_0)} - 2e^{-a(x-x_0)}\big)$

  \begin{minipage}{0.6\textwidth}
  váz.\ stavy: \quad $E_\nu = -\big(\lambda - n - \tfrac12\big)^2\cdot\dfrac{a^2\hbar^2}{2m}$ \\[4pt]
  \hspace*{2.2cm} $\lambda = \dfrac{1}{a\hbar}\sqrt{2mD}$

  \medskip
  $\psi_n$ \dots\ dá se vyjádřit pomocí Laguerre poly.
  \end{minipage}
  \begin{minipage}{0.35\textwidth}
  \begin{tikzpicture}[gray,scale=0.8]
    \draw[->] (0,-1.3) -- (0,1.5);
    \draw[->] (0,0) -- (3.5,0);
    \draw[thick,domain=0.55:3.4,samples=80] plot (\x,{1.2*(exp(-2*1.6*(\x-1.1)) - 2*exp(-1.6*(\x-1.1)))*0.9});
    \draw[thin] (1.1,-1.08) -- (1.1,-1.3) node[below] {\small $x_0$};
    \draw[<->] (3.1,0) -- (3.1,-1.08) node[midway,right] {\small $D$};
    \draw[dotted] (1.1,-1.08) -- (3.1,-1.08);
  \end{tikzpicture}
  \end{minipage}
\end{itemize}

\underline{pozn $\blacktriangle$:} (Ballentine) \\
v ošklivých potenciálech mohou existovat váz.\ st.\ v kont. \\
\quad př: \quad $V(r)\xrightarrow{r\to\infty}\dfrac{\sin(2kr)}{r}$ \quad a \quad
$\psi(r) = f(r)\dfrac{\sin kr}{kr}$, \quad $E = \frac12k^2$ \\
\quad \dots\ $\psi\in L^2$ a přitom $E > 0$ \quad \dots\ důvod: odrazy od bariér, \\
\hspace*{5cm} jejichž výška jde pomalu k $0$ pro $x\to\infty$

\underline{pozn:} v 1D se dá DK, že pro $V(x) < 0 = V_\infty$ $\exists$ alesp.\ jeden váz.\ stav. \\
\hspace*{3.5cm} $\hookrightarrow$ později z variač.\ principu


\newpage
%======================================================================
% PoznQMI-6.pdf
%======================================================================
\section{Částice ve 3D}
\strana{QM-C3-1}

\kw{Stavový prostor} \quad $\mathcal{H} = L^2(\mathbb{R}^3) = \mathcal{H}_1\otimes\mathcal{H}_1\otimes\mathcal{H}_1$ \\
\underline{ÚSKO:} \quad (úplné syst.\ komutujících oper.\ $=$ různé reprezentace)

\subsection*{1) \underline{Operátory polohy}: \ $\hat{\vec X} = (\hat x_1,\hat x_2,\hat x_3)$ \ \pozn{$=$ ÚSKO}}
$\to$ souřadnicová reprezentace: \quad $\braket{x_1,x_2,x_3}{\psi} \equiv \psi(x_1,x_2,x_3) \equiv \psi(\vec x)$
\[
  \braket{\psi}{\psi'} = \int\psi^*(\vec x)\,\psi'(\vec x)\,\dd^3x
\]

\subsection*{2) \underline{Operátor hybnosti (impulzu)} \ \dots\ $\hat{\vec P} = (\hat p_1,\hat p_2,\hat p_3)$ \ \pozn{$=$ ÚSKO}}
transformace $x\leftrightarrow p$:
\[
  \braket{\vec x}{\vec p} \equiv \braket{x_1x_2x_3}{p_1p_2p_3} = \vzorec{\frac{1}{\sqrt{(2\pi\hbar)^3}}\,e^{\frac{i}{\hbar}\vec p\cdot\vec x}}
\]
\hfill $\uparrow$ vl.\ fce odp.\ vl.\ č.\ $(p_1,p_2,p_3) = \vec p$

$\to$ v $x$-reprezentaci \quad $\vec p\,\psi(\vec x) = -i\hbar\vec\nabla\psi(\vec x)$

$\to$ normalizace k $\delta$: \quad $\braket{\vec p}{\vec p\,'} = \delta(\vec p - \vec p\,')$ \quad
$\big[= \delta(p_1 - p_1')\,\delta(p_2 - p_2')\,\delta(p_3 - p_3')\big]$

\pozn{ověření: $\displaystyle\int\dd^3\vec x\,\frac{1}{(2\pi\hbar)^3}e^{\frac{i}{\hbar}\vec x(\vec p\,' - \vec p)}
 = \frac{1}{(2\pi)^3}\int\frac{\dd x_1}{\hbar}e^{i\frac{x_1}{\hbar}(p_1'-p_1)}\int\frac{\dd x_2}{\hbar}e^{i\frac{x_2}{\hbar}(p_2'-p_2)}\int\frac{\dd x_3}{\hbar}e^{i\frac{x_3}{\hbar}(p_3'-p_3)}$} \\
\hspace*{3cm}\pozn{$\displaystyle = \frac{1}{(2\pi)^3}\,2\pi\delta(p_1'-p_1)\,2\pi\delta(p_2'-p_2)\,2\pi\delta(p_3'-p_3) = \delta(\vec p\,' - \vec p)$}

\subsection*{3) \underline{Energie a směr hybnosti} \ \pozn{$=$ ÚSKO}}
\dots\ volná částice \quad $\hat H_0 = \dfrac{p^2}{2m} = \dfrac{1}{2m}\sum_\alpha\hat p_\alpha\cdot\hat p_\alpha$; \quad $\vec n = \dfrac{\vec p}{p}$

platí \quad $[\hat p_\alpha,\hat p_\beta] = 0$ \quad $\to$ \quad $[\hat p_\alpha,\hat H_0] = 0$ \qquad $H_0 = f(p_\alpha)$; \ $\vec n = \vec g(p_\beta)$ \\
$\to$ vl.\ funkce $\ket{\vec p}$ jsou i vl.\ fcemi $\ket{E,\vec n}$ až na normování:
\[
  \braket{E,\vec n}{E',\vec n\,'} = \delta_2(\vec n - \vec n\,')\,\delta(E - E')
\]
\hspace*{3.5cm}\pozn{$\uparrow$ $\delta$-fce na 1 sféře} \quad \pozn{$\nwarrow$ $\delta$-fce na $\langle 0,\infty)$}

tj.\ \quad $\ket{E,\vec n} = N\ket{\vec p}$ \\
tj.\ \quad $\braket{E,\vec n}{E',\vec n\,'} = |N|^2\delta_3(\vec p - \vec p\,') = |N|^2\delta_1(p_1 - p_1')\,\delta_1(p_2 - p_2')\,\delta_1(p_3 - p_3')$ \\
\hspace*{4.5cm}\pozn{$\vec p = \vec p\,' + \Delta\vec p = \vec p\,' + (|p| - |p'|)\vec e_r + (\Delta n_\theta\vec e_\theta + \Delta n_\varphi\vec e_\varphi)\,p$}

\begin{minipage}{0.36\textwidth}
\begin{tikzpicture}[scale=0.95,gray]
  \coordinate (O) at (0,0);
  \draw[->] (O) -- (0,3) node[left] {$p_3$};
  \draw[->] (O) -- (2.6,0.1) node[below right] {$p_2$};
  \draw[->] (O) -- (-1.4,-1.2) node[below] {$p_1$};
  % vektor p'
  \draw[->,thick,black] (O) -- (1,1.5) node[midway,right] {$\vec p\,'$};
  \draw[thin] (0,0.9) arc (90:57:0.9); \node at (0.2,1.1) {\small$\theta$};
  \draw[thin] (-0.4,-0.35) arc (220:330:0.3); \node at (-0.2,-0.5) {\small$\varphi$};
  \draw[dashed,thin] (1,1.5) -- (1,-0.25) -- (O);
  % plošný element
  \draw[thin] (0.7,1.75) -- (1.3,1.9) -- (1.4,1.3) -- (0.75,1.15) -- cycle;
  \draw[->,black] (1,1.5) -- (1.35,2.1) node[above] {$\vec e_r$};
  \draw[->,black] (1,1.5) -- (1.7,1.75) node[right] {$\vec e_\varphi$};
  \draw[->,black] (1,1.5) -- (1.3,1.0) node[right] {$\vec e_\theta$};
  \node[black] at (0.9,-1.3) {\small$\Delta\vec p = \Delta|p|\,\vec e_r + |p|\,\Delta\vec n$};
\end{tikzpicture}
\end{minipage}
\begin{minipage}{0.62\textwidth}
\begin{align*}
  &= \frac{|N|^2}{p^2}\,\delta_2(\vec n - \vec n\,')\,\delta(|p| - |p'|) \\
  &\qquad \Big\lfloor\ p = \sqrt{2mE},\ \ \frac{\dd p}{\dd E} = \sqrt{\frac{m}{2E}} = \frac{m}{p} \\
  &= \frac{|N|^2}{p^2}\,\frac{p}{m}\,\delta_2(\vec n - \vec n\,')\,\delta(E - E')
  \quad\to\quad N = \sqrt{mp}
\end{align*}
\end{minipage}
\[
  \to\quad \vzorec{\ket{E\vec n} = \sqrt{mp}\,\ket{\vec p} = \sqrt{\frac{pm}{(2\pi\hbar)^3}}\;e^{\frac{i}{\hbar}\vec x\cdot\overbrace{\vec n\sqrt{2mE}}^{\vec p_{E,\vec n}}}}
\]

\pozn{Lamé $h_\theta, h_\varphi$ pro $r = 1$} \quad $\delta(\vec n - \vec n\,') = \delta(\theta - \theta')\,\delta\big(\sin\theta(\varphi - \varphi')\big)$

\underline{mimochodem:} tj.: \quad
$\displaystyle\int f(\vec n)\,\delta(\vec n - \vec n\,')\,\dd\Omega = f(\vec n\,')
 \overset{\text{na sféře}}{=} \int_0^\pi\!\!\int_0^{2\pi} f(\theta,\varphi)\,\delta(\theta - \theta')\frac{\delta(\varphi - \varphi')}{|\sin\theta|}\sin\theta\,\dd\varphi\,\dd\theta = f(\theta',\varphi')$

%----------------------------------------------------------------------
\strana{QM-C3-2}
\subsection*{4) \underline{Orbitální moment hybnosti} \ \pozn{do ÚSKO lze přibrat např.\ $E$ nebo $|\vec p|$}}
\[
  \hat L_\alpha = \varepsilon_{\alpha\beta\gamma}\underbrace{\hat Q_\beta}_{=\hat x}\hat P_\gamma
  \qquad\text{tj.}\quad \vec L = \vec X\times\vec P
\]
+ již dříve jsme ukázali: \quad $[L_\alpha,L_\beta] = i\hbar\,\varepsilon_{\alpha\beta\gamma}L_\gamma$ \quad \dots\ nelze hledat spol.\ vl.\ v.\ $L_1,L_2,L_3$

ale \quad $[L^2,L_1] = \underbrace{[L_1^2,L_1]}_{0} + [L_2^2,L_1] + [L_3^2,L_1]$
\[
  = -i\hbar\,(L_2L_3 + L_3L_2) + i\hbar\,(L_3L_2 + L_2L_3) = 0
\]
podobně \quad $[L^2,L_\alpha] = 0 \quad \forall\alpha = 1,2,3$

$\to$ \underline{standardní volby ÚSKO}: \quad $\{\hat A,\hat L^2,\underbrace{\hat L_z}_{=L_3}\}$ \quad kde $\hat A = |\hat p|$ nebo $\hat H_0$

\sout{operátory $L_\alpha, L^2$ v $x$ reprezentaci}

{\color{prbarva}\rule{\linewidth}{0.8pt}}
\subsection{Delší odbočka: Kvantová teorie momentu hybnosti}

obecně: sada operátorů splňující \quad $[J_\alpha,J_\beta] = i\hbar\,\varepsilon_{\alpha\beta\gamma}J_\gamma$ \\
\quad $\to$ reprezentace grupy rotací na $\mathcal{H}$; generátory $\uparrow$

\begin{itemize}
  \item def: \quad $J^2 \equiv J_\alpha J_\alpha$ \quad $\to$ \quad $[J^2,J_\alpha] = 0$
  \item $\Rightarrow$ $\exists$ sada vl.\ v.
  \[
    J^2\ket{\beta,m} = \hbar^2\beta\ket{\beta,m}, \qquad J_3\ket{\beta,m} = \hbar m\ket{\beta,m}
  \]
  \dots\ $\beta, m$ \dots\ bezrozměrné \quad \pozn{ještě může být degener.\ \dots\ další indexy $\alpha_1,\alpha_2$ \dots\ potlačíme}
  \item platí: \quad
  $\beta = \bra{\beta,m}\frac{J^2}{\hbar^2}\ket{\beta,m} = \bra{\beta,m}\big(\tfrac{J_1}{\hbar}\big)^2 + \big(\tfrac{J_2}{\hbar}\big)^2 + m^2\ket{\beta,m} \geq 0$ \\
  \hspace*{6cm} $\uparrow\quad\uparrow\ \geq 0$ \dots\ skal.\ součin $\to$ axiom $\braket{\psi}{\psi}\geq 0$ pro $\ket{\psi} = J_\alpha\ket{\beta m}$ \\
  tj.\ \quad $\beta \geq m^2 \geq 0$ \\
  \quad $\to$ spektrum $J_3$ je omezené
  \item \underline{posunovací operátory} a jejich vlastnosti:
  \[
    \text{def:}\ \ \boxed{\begin{aligned} J_+ &\equiv J_1 + iJ_2 \\ J_- &\equiv J_1 - iJ_2 \end{aligned}} = J_+^\dagger
    \quad\Rightarrow\quad
    \begin{aligned}
      {}[J_3,J_+] &= i\hbar J_2 - i(i\hbar)J_1 = \hbar J_+ &&(+)\\
      [J_3,J_-] &= -\hbar J_- &&(-)\\
      [J_+,J_-] &= i[J_2,J_1] - i[J_1,J_2] = 2\hbar J_3 &&(\pm)
    \end{aligned}
  \]
\end{itemize}

-- \underline{význam $J_+$:}
\[
  J_3\big(J_+\ket{\beta m}\big) = \big(J_+J_3 + \underbrace{[J_3,J_+]}_{(+)}\big)\ket{\beta m} = (\hbar m + \hbar)\,J_+\ket{\beta m}
\]
\quad $\to$ $J_+\ket{\beta m}$ je vl.\ v.\ $J_3$ odp.\ hodnotě $(m+1)$ nebo $0$-vý vektor

\quad {\color{prbarva}\underline{pozn:} $J_\pm\ket{\beta m}$ zůstávají vl.\ v.\ $J^2$ odpov.\ vl.\ č.\ $\hbar^2\beta$, neboť $[J^2,J_\pm] = 0$!}

-- \underline{význam $J_-$:}
\[
  J_3\big(J_-\ket{\beta m}\big) = \big(J_-J_3 + \underbrace{[J_3,J_-]}_{(-)}\big)\ket{\beta m} = \hbar(m-1)\,J_-\ket{\beta m}
\]
\quad $\to$ vl.\ v.\ $J_3$ pro $(m-1)$ nebo $0$

%----------------------------------------------------------------------
\strana{QM-C3-3}
\begin{itemize}
  \item \underline{maximální hodnota $m$}: $\overset{\text{def}}{\equiv} j$ \quad (existuje \dots\ spektrum $J_3$ omezené) \\
  musí \quad $\|J_+\ket{\beta,j}\|^2 = 0$ \\
  tj.\ \quad $\bra{\beta j}J_-J_+\ket{\beta j} = \bra{\beta j}\underbrace{J_1^2 + J_2^2 + i(J_1J_2 - J_2J_1)}\ket{\beta j} = \hbar^2(\beta - j^2 - j) = 0$ \\
  obecně: \quad $\hookrightarrow$ \quad $\boxed{J_-J_+ = J^2 - J_3^2 - \hbar J_3}\ \ (-+)$
  \quad $\Rightarrow$ \quad \underline{$\beta = j(j+1)$}

  podobně:
  \item \underline{minimální hodnota $m$}: $\equiv k$ \\
  \quad $J_-\ket{\beta k} = 0$ \\
  $\to$ $\bra{\beta k}J_+J_-\ket{\beta k} = \bra{\beta k}\underbrace{J_-J_+ + [J_+,J_-]}\ket{\beta k} = \hbar^2(\beta - k^2 + k) = 0$ \\
  \quad $\hookrightarrow$ \quad $\boxed{J_+J_- = J^2 - J_3^2 + \hbar J_3}\ \ (+-)$ \quad tj.\ $\beta = j(j+1) = k^2 - k$ \\
  \hspace*{5cm} $\Rightarrow$ \quad \underline{$k = -j$} \quad (nebo $(j+1) > j$, ale nemůže být min $>$ max)

  \item \underline{závěr:}
  \begin{veta}[Spektrum momentu hybnosti]
    spektrum $J_3$: ($\hbar\,\times$) \quad $m = -j, -j+1,\dots,j$ \\
    \quad přičemž rozdíl $j - k = 2j$ = celé číslo $\Rightarrow$ \quad $j = 0,\tfrac12,1,\tfrac32,2,\dots$ \\
    tj.\ pro $\forall j$ $\exists$ $2j+1$ vektorů s různým $m$

    \underline{změna značení:}
    \[
      \vzorec{\begin{aligned}
        J_3\ket{jm} &= \hbar m\ket{jm} \qquad\dots\ m = -j,\dots,j \\
        J^2\ket{jm} &= \hbar^2 j(j+1)\ket{jm}
      \end{aligned}} \qquad (*)
    \]
  \end{veta}
  která $j$ se realizují záleží na $\mathcal{H}$ a konkr.\ tvaru $\hat J_\alpha$
\end{itemize}
\rule{0.75\linewidth}{0.3pt}

\begin{priklad}[PŘ:]
\begin{itemize}[nosep]
  \item částice se spinem $\frac12$: \quad
        $\vec J = \vec S = \frac{\hbar}{2}\vec\sigma = \frac{\hbar}{2}\left(\begin{pmatrix}0&1\\1&0\end{pmatrix},\begin{pmatrix}0&-i\\ i&0\end{pmatrix},\begin{pmatrix}1&0\\0&-1\end{pmatrix}\right)$ \\
        \dots\ víme, že \quad $[\sigma_\alpha,\sigma_\beta] = 2i\,\varepsilon_{\alpha\beta\gamma}\sigma_\gamma$ \quad $\to$ \quad $[J_\alpha,J_\beta] = i\hbar\,\varepsilon_{\alpha\beta\gamma}J_\gamma$ \\
        \dots\ $\sigma_z$ má vl.\ č.\ $\pm 1$ $\Rightarrow$ $S_z$ vl.\ č.\ $\pm\hbar/2$ \quad \dots\ odpovídá $j = \frac12$
  \item za chvíli ukážeme, že $L_\alpha = \varepsilon_{\alpha\beta\gamma}Q_\beta P_\gamma$ odpovídá ($j =$) $l = 0,1,2,\dots$
\end{itemize}
\end{priklad}

\begin{itemize}
  \item \underline{normalizace:} \\
  \quad $J_+\ket{jm} = C\ket{j\,m+1}$ \quad \dots\ už víme $\bra{jm}J_-J_+\ket{jm} \overset{(-+)}{=} \hbar^2\big[j(j+1) - m(m+1)\big]$
  \[
    \to\quad \vzorec{J_+\ket{jm} = \hbar\sqrt{j(j+1) - m(m+1)}\,\ket{j\,m+1} = \hbar\sqrt{(j-m)(j+m+1)}\,\ket{j\,m+1}}
  \]
  podobně z $(+-)$:
  \[
    \vzorec{J_-\ket{jm} = \hbar\sqrt{j(j+1) - m(m-1)}\,\ket{j\,m-1} = \hbar\sqrt{(j+m)(j-m+1)}\,\ket{j\,m-1}}
  \]
  \dots\ podobně jako u LHO \dots\ obsahuje fázovou konvenci \\
  vektor o němž víme jen $(*)$ $\to$ $e^{i\alpha}\ket{jm}$! \quad \pozn{OBECNĚ}
\end{itemize}

%----------------------------------------------------------------------
\strana{QM-C3-4}
\begin{itemize}
  \item \underline{Maticová reprezentace $J_1,J_2,J_3$ ve vl.\ podpr.\ $J^2$} \\
  \pozn{souvisí s ireducibilními reprezent.\ grupy rotací \dots\ viz teor.\ grup} \\
  $\to$ $J^2, J_z$ se doplní na ÚSKO: $J^2, J_z, A$ \ $\to$ $\ket{\alpha,jm}$ \hfill $[A,J^2] = 0$, \ $[A,J_z] = 0$ \\
  \hspace*{4.3cm} $\hookrightarrow$ diagonální v $\circlearrowleft$ \hfill \pozn{předp.: $[A,J_\alpha] = 0\ \forall\alpha$}
  \begin{align*}
    \text{tj.}\quad \bra{\alpha,jm}J^2\ket{\alpha',j'm'} &= \hbar^2 j(j+1)\,\delta_{jj'}\delta_{mm'}\delta_{\alpha\alpha'} \\
    \bra{\alpha,jm}J_z\ket{\alpha',j'm'} &= \hbar m\,\delta_{jj'}\delta_{mm'}\delta_{\alpha\alpha'} \\
    \bra{\alpha,jm}J_1\ket{\alpha',j'm'} &= (J_1)^{(j)}_{mm'}\,\delta_{jj'}\delta_{\alpha\alpha'} \qquad\text{podob.\ } J_2
  \end{align*}
  $(J_1)_{mm'}$ lze najít explicitně: \quad $J_1 = \frac12(J_+ + J_-)$, \quad $J_2 = \dfrac{J_+ - J_-}{2i}$ \\
  + už víme \quad $\bra{\alpha,jm}J_+\ket{\alpha',j'm'} = \delta_{\alpha\alpha'}\delta_{jj'}\underbrace{\delta_{m\,m'+1}\,\hbar\sqrt{j'(j'+1) - m'(m'+1)}}_{(J_+)_{mm'}}$ \\
  \hspace*{4cm} $J_- = J_+^\dagger$
\end{itemize}

\textbf{ČÁSTICE SE SPINEM} $1/2, 1, 3/2,\dots$ atd.; pro celé $j$ $\to$ orb.\ moment

\pozn{TABULKA:}
\begin{center}
\renewcommand{\arraystretch}{1.3}
\begin{tabular}{c|c|c|c|c|c}
 & \fbox{$j=0$} & \fbox{$j=\frac12$} & \fbox{$j=1$} & \fbox{$j=\frac32$} & \\ \hline
$J_+ =$ & $(0)$ &
$\hbar\begin{pmatrix}0&1\\0&0\end{pmatrix}$ &
$\hbar\begin{pmatrix}0&\sqrt2&0\\0&0&\sqrt2\\0&0&0\end{pmatrix}$ &
$\hbar\begin{pmatrix}0&\sqrt3&0&0\\0&0&\sqrt4&0\\0&0&0&\sqrt3\\0&0&0&0\end{pmatrix}$ & $\cdots$ \\ \hline
$J_x =$ & &
$\frac{\hbar}{2}\begin{pmatrix}0&1\\1&0\end{pmatrix}$ &
$\frac{\hbar}{\sqrt2}\begin{pmatrix}0&1&0\\1&0&1\\0&1&0\end{pmatrix}$ &
$\frac{\hbar}{2}\begin{pmatrix}0&\sqrt3&0&0\\ \sqrt3&0&\sqrt4&0\\0&\sqrt4&0&\sqrt3\\0&0&\sqrt3&0\end{pmatrix}$ & $\cdots$ \\ \hline
$J_y =$ & &
$\frac{\hbar}{2}\begin{pmatrix}0&-i\\ i&0\end{pmatrix}$ &
$\frac{\hbar}{\sqrt2}\begin{pmatrix}0&-i&0\\ i&0&-i\\0&i&0\end{pmatrix}$ &
$\frac{\hbar}{2}\begin{pmatrix}0&-i\sqrt3&0&0\\ i\sqrt3&0&-i\sqrt4&0\\0&i\sqrt4&0&-i\sqrt3\\0&0&i\sqrt3&0\end{pmatrix}$ & $\cdots$
\end{tabular}
\end{center}
{\small (řádky/sloupce indexovány $m, m' = j, j-1, \dots, -j$; \ např.\ pro $j = \frac32$: $\frac32, \frac12, -\frac12, -\frac32$)}

{\color{teal}\rule{\linewidth}{0.8pt}}\\[-4pt]
\hfill{\color{teal}konec odbočky}

\underline{zpět orbitální moment hybnosti:} \\
\quad $L_\alpha$ splňují komut.\ relace generátorů rotací \\
\quad $\Rightarrow$ $\exists$ báze $\ket{\alpha,lm}$: \quad
$\begin{aligned}[t] L^2\ket{\alpha,lm} &= \hbar^2 l(l+1)\ket{\alpha,lm} \\ L_z\ket{\alpha,lm} &= \hbar m\ket{\alpha,lm}\end{aligned}$ \\
\quad kde $m = -l,\dots,l$, ale nevíme, jaká $l$ jsou dovolená \\
\quad ($\sigma(L^2)$ je podmnožinou $\{0,\frac12,1,\frac32,2,\dots\}$)

%----------------------------------------------------------------------
\strana{QM-C3-5}
\begin{itemize}
  \item \underline{operátor $\vec L$ v $x$-reprezentaci} \\
  samosdruženost: \quad $L_\alpha^\dagger = \varepsilon_{\alpha\beta\gamma}P_\gamma^\dagger Q_\beta^\dagger = \varepsilon_{\alpha\beta\gamma}Q_\beta P_\gamma = L_\alpha$ \\
  \hspace*{5cm} $\uparrow$ $Q^\dagger = Q$, $P^\dagger = P$; \ $[Q_\beta,P_\gamma] = 0$ pro $\gamma\neq\beta$
  \[
    \vec L = \vec Q\times(-i\hbar\vec\nabla_Q) \quad\dots\quad
    \begin{aligned}
      L_x &= -i\hbar\big(y\partial_z - z\partial_y\big), \qquad L_y = -i\hbar\big(z\partial_x - x\partial_z\big) \\
      L_z &= -i\hbar\big(x\partial_y - y\partial_x\big)
    \end{aligned}
  \]
\end{itemize}

\underline{pozn:} \underline{$p$-reprezentace:}
\[
  \vec L = (i\hbar\vec\nabla_p)\times\vec p = -i\hbar\,\vec p\times\vec\nabla_p
  \quad\dots\ \text{např.:}\ L_x = -i\hbar\Big(p_y\pdv{p_z} - p_z\pdv{p_y}\Big)
\]
\dots\ totožný tvar $\Rightarrow$ $\braket{\vec x}{lm}$; $\braket{\vec p}{lm}$ budou mít totožné funkční tvary (až na fázi) \\
\hspace*{2cm} (jen úhlová závislost)

\underline{vyjádření ve sférických souřadnicích:}

\begin{minipage}{0.32\textwidth}
\begin{tikzpicture}[scale=1.1,gray]
  \coordinate (O) at (0,0);
  \draw[->] (O) -- (0,2.6) node[right] {$z$};
  \draw[->] (O) -- (2.3,0) node[right] {$y$};
  \draw[->] (O) -- (-1.2,-1.3) node[left] {$x$};
  \draw[->,thick,black] (O) -- (1.3,0.9);
  \draw[thin] (0,0.9) arc (90:35:0.9); \node at (0.25,1.15) {\small$\theta$};
  \draw[thin] (-0.4,-0.43) arc (227:300:0.5); \node at (-0.05,-0.75) {\small$\varphi$};
  \draw[dashed,thin] (1.3,0.9) -- (1.3,-0.3) -- (O);
  \draw[->,black] (1.3,0.9) -- (1.85,1.25) node[right] {$\vec e_r$};
  \draw[->,black] (1.3,0.9) -- (1.1,1.55) node[above] {$\vec e_\varphi$};
  \draw[->,black] (1.3,0.9) -- (1.6,0.35) node[right] {$\vec e_\theta$};
\end{tikzpicture}
\end{minipage}
\begin{minipage}{0.66\textwidth}
\begin{tabular}{l@{\qquad}ll}
$x = r\cos\varphi\sin\theta$ & Lamé: & $h_r = 1$ \\
$y = r\sin\varphi\sin\theta$ & & $h_\theta = r$ \qquad $h = \big\|\tfrac{\partial}{\partial\theta}(x,y,z)\big\|$ \\
$z = r\cos\theta$ & & $h_\varphi = r\sin\theta$
\end{tabular}
\end{minipage}
\[
  \boxed{\begin{aligned}
    \vec e_r &= \tfrac{\partial}{\partial r}(x,y,z) = \cos\varphi\sin\theta\,\vec e_x + \sin\varphi\sin\theta\,\vec e_y + \cos\theta\,\vec e_z \\
    h_\theta^{-1}\vec e_\theta &= \tfrac{\partial}{\partial\theta}(x,y,z) = \cos\varphi\cos\theta\,\vec e_x + \sin\varphi\cos\theta\,\vec e_y - \sin\theta\,\vec e_z \\
    h_\varphi^{-1}\vec e_\varphi &= \tfrac{\partial}{\partial\varphi}(x,y,z) = -\sin\varphi\,\vec e_x + \cos\varphi\,\vec e_y
  \end{aligned}}
\]
\underline{pozn:} \quad
$\begin{pmatrix}e_r\\ e_\theta\\ e_\varphi\end{pmatrix} = M\begin{pmatrix}e_x\\ e_y\\ e_z\end{pmatrix}$
\quad kde $M$ je OG matice: $M^{-1} = M^T$ \dots\ zpětná transf.

\underline{operátor $\vec L$:} \quad
$\boxed{\vec L = -i\hbar\,\vec Q\times\nabla = -i\hbar\,(r\vec e_r)\times\Big(\vec e_r\pdv{r} + \frac{\vec e_\theta}{h_\theta}\pdv{\theta} + \frac{\vec e_\varphi}{h_\varphi}\pdv{\varphi}\Big)}$ \\
\hspace*{5cm} přitom: \quad $\vec e_r\times\vec e_r = 0$, \quad $\vec e_r\times\vec e_\theta = \vec e_\varphi$, \quad $\vec e_r\times\vec e_\varphi = -\vec e_\theta$
\[
  \to\quad \vec L = -i\hbar\Big(\vec e_\varphi\pdv{\theta} - \frac{\vec e_\theta}{\sin\theta}\pdv{\varphi}\Big)
  = \begin{pmatrix}
      -i\hbar\big(-\sin\varphi\,\pdv{\theta} - \cos\varphi\,\frac{\cos\theta}{\sin\theta}\pdv{\varphi}\big) \\
      -i\hbar\big(\cos\varphi\,\pdv{\theta} - \sin\varphi\,\frac{\cos\theta}{\sin\theta}\pdv{\varphi}\big) \\
      -i\hbar\,\pdv{\varphi}
    \end{pmatrix}
  \begin{matrix} = L_x \\[8pt] = L_y \\[8pt] = L_z \end{matrix}
\]

\underline{pozn:} vyjádření $\vec L$ operátorů neobsahuje $r$! \dots\ působí jen na úhl.\ část. \\
\quad $\ket{lm}$ tvoří bázi v úhlových prom., tj.\ na 1-sféře \\
v ÚSKO $(L^2, L_z, A)$ \dots\ operátor $A$ dělá kvantování v $r$ \\
\hspace*{3cm} $A$ \dots\ komutuje s $\vec L$ $\to$ sfér.\ sym.\ \dots\ tj.\ skalár, \\
\hspace*{6cm} např.\ délka vektoru, stopa tenzoru
\[
  \vzorec{L_\pm = (L_x \pm iL_y) = \hbar e^{\pm i\varphi}\Big[\pm\pdv{\theta} + i\frac{\cos\theta}{\sin\theta}\pdv{\varphi}\Big]}
\]

%----------------------------------------------------------------------
\strana{QM-C3-6}
\begin{itemize}
  \item \underline{vyjádření $L^2$}
\end{itemize}
1) \quad $L^2 = -\hbar^2\Big[\vec e_\varphi\pdv{\theta} - \dfrac{\vec e_\theta}{\sin\theta}\pdv{\varphi}\Big]\Big[\vec e_\varphi\cdot\pdv{\theta} - \dfrac{\vec e_\theta}{\sin\theta}\cdot\pdv{\varphi}\Big]$

{\color{gray}
+ užití: \quad $\dfrac{\partial\vec e_\varphi}{\partial\theta} = 0$, \quad
$\dfrac{\partial\vec e_\varphi}{\partial\varphi} = -(\cos\varphi,\sin\varphi,0) = -\vec e_r\sin\theta - \vec e_\theta\cos\theta$ \\
\hspace*{1.4cm} $\dfrac{\partial\vec e_\theta}{\partial\theta} = -\vec e_r$, \quad $\dfrac{\partial\vec e_\theta}{\partial\varphi} = \cos\theta\cdot\vec e_\varphi$
\[
  \vzorec{L^2 = -\hbar^2\Big[\frac{\partial^2}{\partial\theta^2} + \frac{\cos\theta}{\sin\theta}\pdv{\theta} + \frac{1}{\sin^2\theta}\frac{\partial^2}{\partial\varphi^2}\Big]}
\]
\rule{0.6\linewidth}{0.3pt}

srovn.\ \dots\ $\Delta$ ve sfér.\ souř.
\begin{align*}
  \Delta &= \frac{1}{h_rh_\theta h_\varphi}\Big[\pdv{r}\frac{h_\theta h_\varphi}{h_r}\pdv{r} + \pdv{\theta}\frac{h_rh_\varphi}{h_\theta}\pdv{\theta} + \pdv{\varphi}\frac{h_rh_\theta}{h_\varphi}\pdv{\varphi}\Big] \\
  &= \underbrace{\frac{1}{r^2}\pdv{r}r^2\pdv{r}}_{\frac{\partial^2}{\partial r^2} + \frac2r\pdv{r}}
   + \frac{1}{r^2}\Big\{\underbrace{\frac{1}{\sin\theta}\pdv{\theta}\sin\theta\pdv{\theta}}_{\frac{\partial^2}{\partial\theta^2} + \frac{\cos\theta}{\sin\theta}\pdv{\theta}} + \frac{1}{\sin^2\theta}\frac{\partial^2}{\partial\varphi^2}\Big\}
   \quad \boxed{\{\ \} = -\frac{L^2}{\hbar^2r^2}}
\end{align*}
}

jiné odvození:
\[
  L^2 = -\hbar^2(\vec Q\times\vec\nabla)\cdot(\vec Q\times\vec\nabla) = -\hbar^2\big[(\vec Q\times\vec\nabla)\times\vec Q\big]\cdot\vec\nabla
\]
\[
  \overset{\text{BAC-CAB}}{=} -\hbar^2\big[(\vec Q\cdot\vec Q)\vec\nabla - \vec Q(\vec\nabla\cdot\vec Q)\big]\cdot\vec\nabla
  = -\hbar^2\Big\{r^2\Delta + r\pdv{r} - 3r\pdv{r} - \frac{\partial^2}{\partial r^2}\Big\}
\]
\hspace*{5.5cm} $\vec Q\cdot\nabla = r\pdv{r}$ \quad $\vec\nabla\vec Q = I$ \quad $\vec\nabla\cdot\vec Q = 3$
\[
  \text{tj.}\quad L^2 = -\hbar^2\Big(r^2\Delta - \frac{\partial^2}{\partial r^2} - 2r\pdv{r}\Big)
  \qquad \pozn{úhlová část $\Delta$-operátoru}
\]

\rule{\linewidth}{0.3pt}
\begin{itemize}
  \item \underline{vlastní podprostory $L^2, L_z$} \dots\ obecné pozn: \\
  \kw{homogenní polynomy stupně $l$} \dots\ $p_l(\lambda x,\lambda y,\lambda z) = \lambda^l p_l(x,y,z)$

  \pozn{PŘ: $x^2 + y^2 - \frac13z^2$ \dots\ $p_2(x,y,z)$; \quad $(x+iy)^3 = x^3 + 3ix^2y - 3xy^2 - iy^3$ \dots\ $p_3(x,y,z)$}

  platí:
  \begin{enumerate}[label=\arabic*)]
    \item $p_l(x,y,z) = r^l f(\theta,\varphi)$
    \item $L_\alpha\,p_l(x,y,z) = \tilde p_l(x,y,z)$ \quad $\swarrow$ opět homogenní st.\ $l$ \dots\ je vidět z $x$-repr.\ $L_\alpha$
    \item $\Delta\,p_l(x,y,z) = f(\theta,\varphi)\underbrace{\frac{1}{r^2}\pdv{r}r^2\pdv{r}r^l}_{l(l+1)r^{l-2}} - \frac{L^2}{\hbar^2}r^{l-2}f(\theta,\varphi)$ \\
    tj.\ pokud $\Delta p_l(x,y,z) = 0$ \quad $\Rightarrow$ \quad $L^2f(\theta,\varphi) = \hbar^2 l(l+1)f(\theta,\varphi)$
    \item pokud $\Delta p_l = 0$ \quad $\Rightarrow$ \quad $\Delta\tilde p_l = 0$, \quad kde $\tilde p_l = L_\alpha p_l$ \\
    \dots\ plyne z $[L^2,L_\alpha] = 0$ a vyjádření $\Delta$ pomocí $L^2$
  \end{enumerate}
\end{itemize}

%----------------------------------------------------------------------
\strana{QM-C3-7}
\underline{závěr:} $\ket{lm}$ lze hledat v prostoru polynomů \\
\hspace*{1.6cm} \dots\ homogenních stupně $l$ splňujících $\Delta p = 0$

-- počet homog.\ pol.\ stupně $l$: \\
\quad \pozn{$=$ počet koef.\ u $x^ay^bz^c$; $a+b+c = l$ $\to$ $(l+1) + l + \dots + 1 = \frac{(l+2)(l+1)}{2} = N(l)$}

počet $p_l$ splňujících $\Delta p_l = 0$: \\
\quad $\Delta$ mění stupeň $p_l$ o dva \dots\ $\Delta p = 0$ je $N(l-2)$ podmínek \\
\quad $\hookrightarrow$ $= N(l) - N(l-2) = \frac12\big[(l+2)(l+1) - l(l-1)\big] = 2l + 1$ \quad \dots\ konzistentní

\dots\ obvyklé značení \quad $f(\theta,\varphi) \equiv Y_{lm}(\theta,\varphi)$ \quad \dots \quad
{\color{gray}\fbox{$\begin{aligned} L^2Y_{lm} &= \hbar^2 l(l+1)Y_{lm} \\ L_zY_{lm} &= \hbar m\,Y_{lm}\end{aligned}$}}

$\to$ nalezení $Y_{lm}$: \quad \pozn{$-i\partial_\varphi Y = mY$} \\
{\color{gray}
\quad $\Rightarrow$ \quad $Y_{ll} = e^{im\varphi} = (\cos\varphi + i\sin\varphi)^m = \frac{1}{r^m}(x+iy)^m$ \quad $\leftarrow$ \\
\quad\quad $\to$ funguje jen pro $m = l$ \dots\ jinak netriv.\ $\theta$-část, ale ostatní pomocí $L_-$ \\
\quad nebo \quad $Y_{l,-m} = e^{-im\varphi} = (\cos\varphi - i\sin\varphi)^m = \frac{1}{r^m}(x-iy)^m$ \quad $\leftarrow$ \quad fáze?
}

\underline{tradiční způsob nalezení $Y_{lm}(\theta,\varphi)$} \dots\ separace proměnných
\[
  Y_{lm}(\theta,\varphi) = f(\theta)\,g(\varphi)
\]
\begin{itemize}
  \item $\to$ \quad $L_zg(\varphi) = -i\hbar\pdv{\varphi}g(\varphi) = \hbar m\,g(\varphi)$ \quad s period.\ podm.: $g(\varphi + 2\pi) = g(\varphi)$ \\
  $\to$ \quad $g(\varphi) = e^{im\varphi}$; \quad $m = 0,\pm1,\pm2,\dots$ \quad (+ z dřívějška víme $|m|\leq l$) \\
  \hspace*{6cm} $\to$ jen celočíselná $l$
\end{itemize}
+ dosazení $Y_{lm}(\theta,\varphi) = f(\theta)e^{im\varphi}$ do $L^2Y_{lm} = \hbar^2 l(l+1)Y_{lm}$:
\[
  -\Big[\frac{1}{\sin\theta}\pdv{\theta}\sin\theta\pdv{\theta} - \frac{m^2}{(\sin\theta)^2}\Big]f(\theta) = l(l+1)f(\theta)
\]
\hspace*{2cm} + substituce $z = \cos\theta$ \quad \dots\ $\dd z = -\sin\theta\,\dd\theta$ \hfill $\lambda\equiv l(l+1)$
\[
  \to\quad \Big\{\dv{z}\big[(1-z^2)\big]\dv{z} + \lambda - \frac{m^2}{1-z^2}\Big\}f(\theta) = 0
\]
\hspace*{6cm} hledáme řešení nesingulární v $\pm 1$

$\to$ \kw{přidružené Legendre funkce}:
\[
  \vzorec{Y_{lm}(\theta,\phi) = (-1)^m\sqrt{\frac{(2l+1)(l-m)!}{4\pi\,(l+m)!}}\;P_l^m(\cos\theta)\,e^{im\varphi}}
\]
\quad kde \quad $P_l^m(\xi) = \dfrac{1}{2^l l!}(1-\xi^2)^{m/2}\dfrac{\dd^{l+m}}{\dd\xi^{l+m}}(\xi^2 - 1)^l$

\begin{minipage}[t]{0.48\textwidth}
\underline{Relace ortogonality:}
\[ \int Y_{lm}^*(\vec n)\,Y_{l'm'}(\vec n)\,\dd\Omega_n = \delta_{ll'}\delta_{mm'} \]
\end{minipage}
\begin{minipage}[t]{0.5\textwidth}
\underline{Relace úplnosti:}
\[ \sum_{l=0}^\infty\sum_{m=-l}^{l} Y_{lm}(\vec n)\,Y_{lm}^*(\vec n\,') = \delta_2(\vec n - \vec n\,') \]
\pozn{$= \dfrac{\delta(\theta-\theta')\delta(\varphi-\varphi')}{|\sin\theta|} = \delta(\cos\theta - \cos\theta')\,\delta(\varphi - \varphi')$}
\end{minipage}

\pozn{rozebrat \dots\ $Y(\vec n) = Y(\theta,\varphi) = Y(\cos\theta,\varphi)$ \dots\ záměna diferenciálu; také v $\nearrow$}

%----------------------------------------------------------------------
\strana{QM-C3-8}
\subsubsection*{Detailnější odvození $Y_{lm}(\theta,\varphi)$}
\[
  0 = L_-Y_{l,-l}(\theta,\varphi) = \hbar e^{-i\varphi}\Big[-\pdv{\theta} + i\frac{\cos\theta}{\sin\theta}\pdv{\varphi}\Big]f(\theta)e^{-il\varphi}
  = \hbar e^{-i(l+1)\varphi}\Big[-\pdv{\theta} + l\frac{\cos\theta}{\sin\theta}\Big]f(\theta) = 0
\]
\quad tj.\ \quad $\dfrac{f'}{f} = l\dfrac{\cos\theta}{\sin\theta}$ \quad $\to$ \quad $f(\theta) = c_l(\sin\theta)^l$

$c_l$ z normalizace: \quad $\int\dd\varphi\int\sin\theta\,\dd\theta\,|c_l|^2\sin^{2l}\theta = 1$ \quad $\to$ \quad $c_l = \dfrac{1}{\sqrt{2\pi I}}$ \\
\quad kde \quad $I_l = \displaystyle\int_0^\pi\sin\theta\,\dd\theta\,\sin^{2l}\theta = \frac{2^{2l+1}(l!)^2}{(2l+1)!}$
\quad $\to$ \quad $c_l = \dfrac{\text{\sout{$(-1)^l$}}}{2^l l!}\sqrt{\dfrac{(2l+1)!}{4\pi}}$
\hfill \begin{tabular}[t]{l}\underline{fázová konvence} \\ $\to P_{l0}(\theta = 0) > 0$ \\ \dots\ $Y_{l0}(\theta,\varphi) > 0$\end{tabular}

\underline{trošku hraní s $L_+$:} \dots\ platí: \quad $\dfrac{\dd}{\dd(\cos\theta)} = -\dfrac{1}{\sin\theta}\cdot\dfrac{\dd}{\dd\theta}$
\begin{align*}
  \Rightarrow\quad \frac{\dd}{\dd\cos\theta}\big((\sin\theta)^{-m}f\big) &= -\frac{1}{\sin\theta}\big[(\sin\theta)^{-m}f' - m(\sin\theta)^{-m-1}\cos\theta\,f\big] \\
  &= -\sin^{-m-1}\theta\,\Big[f' - m\frac{\cos\theta}{\sin\theta}f\Big]
\end{align*}
\[
  \to\quad L_+fe^{im\varphi} = \hbar e^{i\varphi}\Big[\pdv{\theta} - m\frac{\cos\theta}{\sin\theta}\Big]f(\theta)e^{im\varphi}
  = -\hbar e^{i(m+1)\varphi}(\sin\theta)^{m+1}\frac{\dd}{\dd\cos\theta}\big[(\sin\theta)^{-m}f(\theta)\big]
\]
\[
  \Rightarrow\quad L_+^p fe^{im\varphi} = (-\hbar)^p e^{i(m+p)\varphi}(\sin\theta)^{m+p}\frac{\dd^p}{(\dd\cos\theta)^p}\big[(\sin\theta)^{-m}f\big]
\]
+ normování: \quad $L_+\ket{lm} = \hbar\sqrt{(l-m)(l+m+1)}\,\ket{l\,m+1}$
\[
  N\ket{lm} = (L_+)^{l+m}\ket{l,-l} = \hbar^{l+m}\sqrt{2l}\cdot\sqrt{2l-1}\cdots\sqrt{l-m+1}\cdot\sqrt1\cdot\sqrt2\cdots\sqrt{l+m}\;\ket{lm}
\]
\[
  \to\quad N = \frac{1}{\hbar^{l+m}}\sqrt{\frac{(l-m)!}{(l+m)!}\frac{1}{2l!}}
\]
vše dohromady:
\[
  \vzorec{Y_{lm}(\theta,\varphi) = \frac{(-1)^{l+m}}{2^l l!}\sqrt{\frac{2l+1}{4\pi}\frac{(l-m)!}{(l+m)!}}\;e^{im\varphi}(\sin\theta)^m\frac{\dd^{l+m}}{(\dd\cos\theta)^{l+m}}(\sin\theta)^{2l}}
\]
\quad + substituce \quad $\xi = \cos\theta$, \ $\sqrt{1-\xi^2} = \sin\theta$ \quad $\to$ vzorec na předch.\ str.

\underline{pozn:} $Y_{l0}(\theta)$ \dots\ Legendreův polynom (v prom.\ $\cos\theta$) \\
podrobněji: \quad
$Y_{l0}(\theta) = \dfrac{(-1)^l}{2^l l!}\sqrt{\dfrac{2l+1}{4\pi}}\;\underbrace{\dfrac{\dd^l}{\dd\xi^l}(1-\xi^2)^l\Big|_{\xi=\cos\theta}}_{\text{polynom st.\ } 2l-l}
 \equiv \sqrt{\dfrac{2l+1}{4\pi}}\,P_l(\cos\theta)$

relace ortogonality:
\[
  \delta_{ll'} = \int Y_{l0}(n)\,Y_{l'0}^*(n)\,\dd^2n = \int_0^{2\pi}\dd\varphi\int_0^\pi\underbrace{\sin\theta\,\dd\theta}_{\dd\xi}\sqrt{\frac{2l+1}{4\pi}}\sqrt{\frac{2l'+1}{4\pi}}\,P_l(\xi)P_{l'}(\xi)
\]
\[
  \to\quad \vzorec{\int_{-1}^{1}P_l(\xi)P_{l'}(\xi)\,\dd\xi = \frac{2}{2l+1}\delta_{ll'}}
  \qquad \begin{tabular}{l}\dots\ ortogonální polynomy \\ s vahou $w(\xi) = 1$ \\ na intervalu $\langle -1,1\rangle$ \\ $\to$ souvislost \dots\ Gauss\end{tabular}
\]

\underline{pozn:} symetrie: \quad $Y_{lm}(-\vec n) = Y_{lm}(\vec n)(-1)^l$ \\
\hspace*{3.25cm} $Y_{lm}^*(\vec n) = (-1)^m\,Y_{l,-m}(\vec n)$ \quad
\pozn{$\nearrow$ faktor $(-1)^m$ není vidět, dá se odv.\ podobně z $L_-$ operátoru (Cohen-Tannoudji)}

%----------------------------------------------------------------------
\strana{QM-C3-9}
\underline{pozn:} \kw{reálné harmoniky}: \quad
$\begin{aligned}[t]
  C_{lm} &= \tfrac{1}{\sqrt2}\big(Y_{lm} + Y_{lm}^*\big) \\
  S_{lm} &= \tfrac{1}{\sqrt2}\big(Y_{lm} - Y_{lm}^*\big) \quad\leftarrow \text{někdy se čistí faktorem } (-1)^{m+\dots}
\end{aligned}$ \hfill \fbox{$m\geq 0$} \\
\quad + $Y_{l0}$ beze změny

$\to$ $Y_{lm}^* \sim Y_{l,-m}$ \dots\ tj.\ jen unit.\ transf.\ v rámci $\pm m$ \\
$\to$ fázový faktor $e^{im\varphi}$ na 1-kouli $r = 1$: \\
\quad $\to$ je vždy doplněn faktorem $(\sin\theta)^m$: \quad $(r\sin\theta)^m e^{im\varphi} = (x + iy)^m$ \\
tj.\ \quad $C_{lm}\sim\cos m\varphi$, \quad $S_{lm}\sim\sin m\varphi$ \quad $\mathrm{Re}/\mathrm{Im}$ \dots\ $(x+iy)^m \pm (x-iy)^m$ \\
navíc $P_{lm}$ je vyjádř.\ $\sim\sin m\varphi\cdot\mathrm{poly}_{l-m}(\cos\theta)$ \dots\ homog.\ poly st.\ $m$

{\color{prbarva}\rule{\linewidth}{0.6pt}}\\[-4pt]
\hfill{\color{prbarva}konec odbočky}

\subsection*{4) \underline{ÚSKO obsahující $L^2, L_z$} \dots\ přidání $A$}
\dots\ většinou máme $[A,L_\alpha] = 0$ \quad $\to$ \quad $[A,L_\pm] = 0$ \\
\dots\ obecně \dots\ ve vl.\ podpr.\ $l,m$ diagonalizujeme $A$ $\to$ vl.\ č.\ $a_0,a_1,\dots,a_\alpha$ \\
\dots\ v principu $a$ závisí na $l,m$, ale \quad $AL_\pm\ket{a} = L_\pm A\ket{a}$ $\Rightarrow$ nezáv.\ na $m$
\[
  \to\quad
  \begin{aligned}
    L^2\ket{\alpha;lm} &= \hbar^2 l(l+1)\ket{lm,\alpha} && l = 0,1,2,\dots \\
    L_z\ket{lm,\alpha} &= \hbar m\ket{lm,\alpha} && m = -l,-l+1,\dots,l \\
    A\ket{lm,\alpha} &= a_\alpha^{(l)}\ket{lm,\alpha}
  \end{aligned}
\]
$\to$ spektrum $a$ závisí na konkr.\ volbě \dots\ kvantuje radiální prom.

\subsubsection*{4a) \underline{$A = H_0 \equiv p^2/2m$} -- řešení v $p$-reprezentaci}
$\to$ ihned \quad $\braket{\vec p}{Elm} = N\,\delta(E_p - E)\,Y_{lm}(\hat p)$ \qquad $\hat p \equiv \dfrac{\vec p}{|p|}$

normování: \quad $\braket{Elm}{E'l'm'} = \displaystyle\int\dd^3p\,|N|^2\delta\Big(E - \frac{p^2}{2m}\Big)\delta\Big(E' - \frac{p^2}{2m}\Big)Y_{lm}^*(\hat p)Y_{l'm'}(\hat p)$
\begin{align*}
  \hspace*{3cm} &= \delta_{ll'}\delta_{mm'}|N|^2\int p^2\,\dd p\,\frac{\delta(p - p_E)}{p/m}\,\delta\Big(E' - \frac{p^2}{2m}\Big) \\
  &= pm\,|N|^2\,\delta_{ll'}\delta_{mm'}\,\delta(E - E')
\end{align*}
\quad $N = \dfrac{1}{\sqrt{pm}}$ \pozn{$= (2Em^3)^{-1/4}$} \quad $\nwarrow$

\underline{závěr:} \quad $\vzorec{\braket{\vec p}{Elm} = \frac{1}{\sqrt{pm}}\,\delta\Big(E - \frac{p^2}{2m}\Big)\,Y_{lm}(\hat p)}$

\subsubsection*{4b) \underline{$A = H_0 \equiv p^2/2m$} -- řešení v $x$-reprezentaci}
$\to$ v principu Fourier.-transformace předchozího, ale def.\ spec.\ fce
\[
  \psi_{Elm}(\vec Q) = \int\dd^3p\,\braket{\vec x}{\vec p}\braket{\vec p}{Elm}
  \quad\dots\ \text{dá se odvodit z rozkladu:}
\]
\[
  \vzorec{e^{i\vec k\cdot\vec Q} = 4\pi\sum_l\sum_m i^l\,j_l(kr)\,Y_{lm}^*(\hat k)\,Y_{lm}(\hat Q)}
  \qquad \vec k \equiv \frac{\vec p}{\hbar}
\]
následující lze chápat jako odvození $\uparrow$ a definice $j_l(kr)$

%----------------------------------------------------------------------
\strana{QM-C3-10}
hledáme \quad $\psi_{klm}(\vec Q) \equiv \braket{\vec Q}{klm} = R_{kl}(r)\,Y_{lm}(\theta,\varphi)$ \\
\quad kde \quad $E = \dfrac{p^2}{2m}$, \quad $k = \dfrac{p}{\hbar} = \dfrac{\sqrt{2mE}}{\hbar}$ \quad \dots\ $H_0\ket{klm} = E\ket{klm}$

normování a skal.\ součin: \quad $\braket{klm}{k'l'm'} = \delta_{ll'}\delta_{mm'}\int\dd r\,r^2R_{k'l}^*R_{kl}$ \\
$\to$ je výhodné zavést \dots\ $\chi_{kl} \equiv r\,R_{kl}(r)$ \dots\ radiální skal.\ souč.: $\braket{\chi}{\chi'} = \int\chi^*\chi'\,\dd r$
\[
  \to\quad \psi_{klm}(\vec Q) = \frac{1}{r}\,\chi_{kl}(r)\,Y_{lm}(\theta,\varphi)
\]
\[
  H_0\psi_{klm}(\vec Q) = -\frac{\hbar^2}{2m}\Delta\psi_{klm} = -\frac{\hbar^2}{2m}\frac{1}{r^2}\Big[\pdv{r}r^2\pdv{r} - l(l+1)\Big]\psi = E\psi
\]
\[
  \to\quad \underbrace{R_{kl}''(r) + \frac2r R_{kl}'}_{\pozn{$\frac1r(rR)'' = \frac1r(R + rR')' = \frac1r(2R' + rR'')$}}
  + \Big[k^2 - \frac{l(l+1)}{r^2}\Big]R_{kl} = 0
\]
\hspace*{7cm} $\uparrow$ centrifugální člen (jako v klas.\ fyzice) \quad $\dfrac{L^2}{2mr^2}$
\[
  \to\quad \vzorec{\chi''(r) + \Big[k^2 - \frac{l(l+1)}{r^2}\Big]\chi(r) = 0}
  \quad\swarrow\ \text{jako 1D (SR) s } V_l(r) = \frac{\hbar^2 l(l+1)}{2mr^2}
\]

\underline{řešení pro $l = 0$:} \quad $\chi(r) = A\sin kr + B\cos(kr)$

\underline{okrajové podmínky:} \dots\ spojitost $\psi$ a $\psi'$ ($\nabla\psi$) v počátku \\
\hspace*{3.5cm} tj.\ $R(r = 0) = 0$ pro $l > 0$; \ $\chi_l(0) = 0$ $\forall l$

$\to$ bezrozm.\ jednotky: \quad $z = kr$ \quad $\to$ \quad $R(r) = j_l(kr) = j_l(z)$ \\
tj.\ \quad $\Big[\dfrac{\dd^2}{\dd z^2} + \dfrac2z\dfrac{\dd}{\dd z} + 1 - \dfrac{l(l+1)}{z^2}\Big]j_l(z) = 0$ \quad \dots\ univerz.\ rovnice nezáv.\ na $E$

\sout{nalezení asympt.\ chování \dots\ $j_l(z) = z^{l+1}$} \quad + def.\ $\hat j_l(z) = z\,j_l(z)$
\[
  \to\quad \hat j''(z) + \Big[1 - \frac{l(l+1)}{z^2}\Big]\hat j(z) = 0 \qquad\dots\ \text{asymptotika } \hat j \approx z^\alpha,\ z\to 0
\]
\quad \dots\ $\alpha(\alpha-1)z^{\alpha-2} - l(l+1)z^{\alpha-2} = 0$ \quad $\to$ \quad $\alpha = l+1$ \quad
{\color{gray}\fbox{\sout{$\alpha = -l$}} $\to$ nevyhovuje poč.\ podm.}

vytažení asymptotiky: \quad
$\begin{aligned}[t]
  \hat j(z) &= z^{l+1}\varphi(z) \\
  \hat j'(z) &= z^{l+1}\Big(\varphi' + \frac{l+1}{z}\varphi\Big) \\
  \hat j''(z) &= z^{l+1}\Big(\varphi'' + \frac2z(l+1)\varphi' + \frac{l(l+1)}{z^2}\varphi\Big)
\end{aligned}$

%----------------------------------------------------------------------
\strana{QM-C3-11a}
{\small\hfill (verze a)}
\[
  \to\quad \varphi'' + \frac2r(l+1)\varphi' + \varphi = 0
\]
\hspace*{5cm} \dots\ řešení $l = 0$: \quad $\varphi_0(z) = \dfrac{\sin z}{z}$ \quad (z rovn.\ pro $\hat j$)

+ \underline{pozorujte:} funkce $\psi = \dfrac{\varphi'}{r}$ splňuje: \quad $\psi'' + \dfrac{2(l+2)}{r}\psi' + \psi = 0$ \\
$\Rightarrow$ \quad $\varphi_{l+1}(z) = \psi(z) = \dfrac1z\dfrac{\dd}{\dd z}\varphi_l(z)$
\quad \dots\ přitom: $R_{kl}(r) = j(kr) = \dfrac{\hat j(kr)}{kr} = (kr)^l\varphi(kr)$

tj.\ \quad $R_{kl}(r) = C\,j(kr)$, \quad kde:
\[
  \vzorec{j_l(z) \equiv (-1)^l z^l\Big(\frac1z\frac{\dd}{\dd z}\Big)^l\frac{\sin z}{z}}
\]
\quad druhé lin.\ nez.\ řešení: \quad
$n_l(z) \equiv (-1)^{l+1}z^l\Big(\dfrac1z\dfrac{\dd}{\dd z}\Big)^l\dfrac{\cos z}{z}$

$\to$ \kw{sférické cylindrické funkce} \dots\ Besselova a Neumannova

\underline{asymptotické chování:} \quad $z\to 0$ \quad
$j_l(z) = \dfrac{z^l}{(2l+1)!!}$, \quad $n_l(z) = \dfrac{(2l+1)!!}{z^{l+1}}$ \quad \pozn{DK nebude}

$z\to\infty$ \quad \dots\ jen nejvyšší mocnina $z$: derivujeme jen $\sin/\cos$: \\
\quad přičemž \quad $\dfrac{\dd}{\dd z}\sin(z + \delta) = \sin\big(z - \tfrac{\pi}{2} + \delta\big)$ \quad $\to$
\[
  j_l(z) \xrightarrow{\ z\to\infty\ } \frac1z\sin\Big(z - \frac{\pi}{2}l\Big)
  \qquad n_l(z) = -\frac1z\cos\Big(z - \frac{\pi}{2}l\Big)
\]

\underline{Normování:} \quad platí \quad $\displaystyle\int_0^\infty j_l(kr)\,j_l(k'r)\,r^2\dd r = \frac{\pi}{2k^2}\delta(k - k')$

\underline{Důkaz:} z toho, že $H_0$ je samosdružený víme, že $\int j(kr)j(k'r)\sim\delta$ \\
\quad $\to$ zajímá nás jen singulární ($\delta$-chování) \\
posun o kvadraticky integ.\ fci neovlivní: \quad $\int(j_l - f)(j_l - f)\sim$ regulární $F(k,k') + \delta$

\dots\ stačí asymptotika:
\begin{align*}
  \int_0^\infty\sin(kr+\delta)\sin(k'r+\delta)\frac{1}{kk'}\,\dd r
  &= \frac{1}{kk'}\Big(-\frac14\Big)\int_0^\infty\underbrace{(e^{iA} - e^{-iA})(e^{iA'} - e^{-iA'})}_{e^{i(A+A')} + e^{-i(A+A')} - e^{i(A-A')} - e^{-i(A-A')}}\dd r \\
  &\hspace*{4cm}\hookrightarrow\ \sim\delta(k + k')\ \dots\ \text{nepřisp., } k,k' > 0 \\
  &= \frac{1}{2kk'}\int_{-\infty}^\infty\cos(k-k')r\,\dd r = \frac{1}{4kk'}\int_{-\infty}^\infty e^{i(k-k')r}\,\dd r \\
  &= \frac{2\pi}{4kk'}\,\delta(k - k')
  \quad\to\quad \text{\underline{normalizované}:}\ \vzorec{R_{kl}(r) = \sqrt{\frac{2}{\pi}}\,k\,j_l(kr)}
\end{align*}
\hfill ($\leftarrow\ \delta(k - k')$)

%----------------------------------------------------------------------
\strana{QM-C3-11b}
\subsubsection*{Algebraické řešení dle Cohena-Tannoudji {\normalfont\small(verze b)}}
\dots\ ukázali jsme: $\forall$ vektor: \quad $[L_\alpha,P_\beta] = i\hbar\,\varepsilon_{\alpha\beta\gamma}P_\gamma$
\begin{enumerate}[label=\arabic*)]
  \item[$\Rightarrow$ 1)] $[L_z,P_+] = [L_z,P_x] + i[L_z,P_y] = \hbar(iP_y + P_x) = \hbar P_+$ \\
        \hspace*{1.6cm} $\hookrightarrow$ def.\ \fbox{$P_+ = P_x + iP_y$}
  \item[2)] $[L_+,P_z] = [L_x,P_z] + i[L_y,P_z] = \hbar(-iP_y - P_x) = -\hbar P_+$
  \item[3)] $[L^2,P_+] = [L_y^2 + L_z^2,P_x] + i[L_x^2 + L_z^2,P_y] = \hbar\big[i(-L_yP_z - P_zL_y + L_zP_y + P_yL_z)$ \\
        \hspace*{7.4cm} $- L_xP_z - P_zL_x + L_zP_x + P_xL_z\big]$ \\
        $= \hbar\{L_zP_+ + P_+L_z - P_zL_+ - L_+P_z\} = 2\hbar(P_+L_z - P_zL_+) + \hbar\big(\underbrace{[L_z,P_+]}_{\hbar P_+} - \underbrace{[L_+,P_z]}_{-\hbar P_+}\big)$ \\
        $= 2\hbar(P_+L_z - P_zL_+) + 2\hbar^2P_+$
\end{enumerate}

\begin{veta}[Lemma]
\[ P_+\ket{klm}\Big|_{m=l} \sim \ket{k\,l{+}1,l{+}1} \]
\end{veta}
\underline{DK:}
\begin{itemize}[nosep]
  \item $[P_\alpha,H_0]$ $\Rightarrow$ $P_+\ket{klm}$ je vl.\ v.\ $H_0$ s vl.\ č.\ $E = \dfrac{\hbar^2k^2}{2m}$
  \item viz 2) \ $[L_z,P_+]$ $\Rightarrow$ $L_zP_+\ket{klm} = P_+L_z\ket{klm} + [L_z,P_+]\ket{klm} = \hbar(m+1)P_+\ket{klm}$
  \item 3) $\Rightarrow$ $L^2P_+\ket{klm}\big|_{m=l} = P_+L^2\ket{klm} + [L^2,P_+]\ket{kll} = \big[\hbar^2 l(l+1) + 2\hbar^2l - 0 + 2\hbar^2\big]P_+\ket{kll}$ \\
        \hspace*{3.4cm} $= \hbar^2(l+1)(l+2)\,P_+\ket{kll}$ \quad \dots\ + jednoznačnost
\end{itemize}
4) \quad $[P_+,x + iy] = [P_x,x] + (i)^2[P_y,y] = -i\hbar + i\hbar = 0$

\underline{Postup:}
\begin{center}
\begin{tikzpicture}[>=Stealth, node distance=1.8cm]
  \node (a) at (0,0) {$\ket{k00}$};
  \node (b) at (2.4,0) {$\ket{k11}$};
  \node (c) at (4.8,0) {$\ket{k22}$};
  \node (d) at (7,0) {$\dots$};
  \draw[->] (a) -- node[above] {$P_+$} (b);
  \draw[->] (b) -- node[above] {$P_+$} (c);
  \draw[->] (c) -- node[above] {$P_+$} (d);
  \node (b1) at (2.4,-1.2) {$\ket{k1m}$};
  \node (c1) at (4.8,-1.2) {$\ket{k2m}$};
  \node (d1) at (7,-1.2) {$\dots$};
  \draw[->] (b) -- node[right] {$L_-$} (b1);
  \draw[->] (c) -- node[right] {$L_-$} (c1);
  \draw[gray,double,->] (a) -- (0,-1.9);
  \draw[gray] (5.8,-0.6) -- (8.8,-0.6) -- (8.8,-2.1) -- (7,-2.1);
  \node[gray,right] at (8.9,-1.2) {\small$L_-$ působí jen na $\theta,\varphi$};
\end{tikzpicture}
\end{center}
přitom \quad $\braket{\vec Q}{klm} = N\,j_l(kr)\,Y_{lm}(\theta,\varphi)$ \\
$\to$ jak působí $P_+$? \quad \dots \quad $Y_{ll}(\theta,\varphi) = A(\sin\theta)^le^{il\varphi} = A\Big(\dfrac{x+iy}{r}\Big)^l$
\[
  P_+\frac{\sin kr}{kr}\ \dots\ -i\hbar(\partial_x + i\partial_y)f(r) \sim \underbrace{\Big(\pdv{r}{x} + i\pdv{r}{y}\Big)}_{\frac{x+iy}{r}}\pdv{f}{r}
  \sim \underbrace{(x+iy)}_{\sim\,Y_{11}(\theta,\varphi)}\frac1r\pdv{f}{r}
\]
indukce: \quad
$P_+(x+iy)^l\underbrace{\Big(\frac1r\pdv{r}\Big)^lf(r)}_{g(r)} = (x+iy)^lP_+g(r) \sim (x+iy)^{l+1}\frac1r\pdv{r}g(r)$ \\
\hspace*{8cm} $= (x+iy)^{l+1}\Big(\dfrac1r\pdv{r}\Big)^{l+1}f(r)$
\[
  \Rightarrow\ (\text{def})\quad j_l(z) \equiv (-1)^lz^l\Big(\frac1z\frac{\dd}{\dd z}\Big)^l\frac{\sin z}{z} \qquad\dots\ z = kr
\]
druhé LNZ řešení \quad $n_l(z) \equiv (-1)^{l+1}z^l\Big(\dfrac1z\dfrac{\dd}{\dd z}\Big)^l\dfrac{\cos z}{z}$

%----------------------------------------------------------------------
\strana{QM-C3-12}
\underline{chování pro malé $z$:} \\
\quad Taylorova řada pro $\sin z/\cos z$ \quad $\to$ \quad
$j_l(z) \underset{z\to0}{\approx} \dfrac{z^l}{(2l+1)!!}$, \qquad $n_l(z) \approx \dfrac{(2l+1)!!}{z^{l+1}}$

\underline{Asymptotické chování pro velké $z$:} \\
\dots\ největší mocnina $z$ \dots\ derivujeme jen $\sin/\cos$:
\[
  \frac{\dd}{\dd z}\sin(z+\delta) = \sin\Big(z - \frac{\pi}{2} + \delta\Big)
  \ \underset{z\to\infty}{\to}\
  \vzorec{j_l(z) \simeq \frac1z\sin\Big(z - \frac{\pi}{2}l\Big)},\quad n_l \simeq -\frac1z\cos\Big(z - \frac{\pi}{2}l\Big)
\]

\underline{Prvních pár funkcí:}
\begin{align*}
  j_0(z) &= \frac{\sin z}{z} & n_0(z) &= -\frac{\cos z}{z} \\
  j_1(z) &= \frac{\sin z}{z^2} - \frac{\cos z}{z} & n_1(z) &= -\frac{\cos z}{z^2} - \frac{\sin z}{z} \\
  j_2(z) &= \Big(\frac{3}{z^3} - \frac1z\Big)\sin z - \frac{3}{z^2}\cos z & n_2(z) &= -\Big(\frac{3}{z^3} - \frac1z\Big)\cos z - \frac{3}{z^2}\sin z
\end{align*}
\pozn{v originále u $j_1$ psáno $\frac{\sin z}{z^3}$; správně $\frac{\sin z}{z^2}$}

\begin{center}
\begin{tikzpicture}[gray,x=0.33cm,y=1.6cm]
  % j0
  \draw[->] (-0.5,0) -- (17,0);
  \draw[->] (0,-0.4) -- (0,1.3);
  \node[left] at (0,1) {$1$};
  \node at (5,1.15) {$l = 0$};
  \node at (12,0.9) {$j_0$};
  \draw[thick,black,domain=0.01:16.5,samples=150] plot (\x,{sin(\x r)/\x});
  % jl, l velké
  \begin{scope}[xshift=7.2cm]
    \draw[->] (-0.5,0) -- (19,0);
    \draw[->] (0,-0.4) -- (0,1.3);
    \node at (3,1.15) {$l$ velké};
    \node at (12.5,0.3) {$j_l$};
    \node[left] at (1.5,0.35) {$z^l$};
    \draw[thick,black,domain=0:5.5,samples=60] plot (\x,{0.75*exp(-0.3*(5.5-\x)*(5.5-\x))});
    \draw[thick,black,domain=5.5:18.5,samples=120] plot (\x,{0.75*cos(deg(1.1*(\x-5.5)))*exp(-0.06*(\x-5.5))});
    \draw[dashed] (4.4,0) -- (4.4,0.52);
    \node[below] at (4.4,0) {$R_i$};
    \draw[<-] (4.6,0.55) -- (6.5,0.95) node[right] {\small inflex.\ bod \dots\ klas.\ bod obratu};
  \end{scope}
\end{tikzpicture}
\end{center}

+ \underline{Normování:} \quad $\displaystyle\int_0^\infty j_l(kr)\,j_l(k'r)\,r^2\dd r = \frac{\pi}{2k^2}\delta(k - k')$ \\
tj.\ \quad $\begin{cases}
  \braket{k'l'm'}{klm} = \displaystyle\int_0^\infty r^2\dd r\int\dd^2\vec n\,\psi_{k'l'm'}^*(r\vec n)\,\psi_{klm}(r\vec n) = \delta(k - k')\delta_{ll'}\delta_{mm'} \\[8pt]
  \displaystyle\sum_{lm}\int_0^\infty\dd k\,\psi_{klm}(\vec Q)\,\psi_{klm}^*(\vec Q\,') = \delta(\vec Q - \vec Q\,')
\end{cases}$ \\
\quad kde \quad $\vzorec{\psi_{klm}(r\vec n) = \sqrt{\frac{2}{\pi}}\,k\,j_l(kr)\,Y_{lm}(\vec n)}$

Energetická normalizace: \quad $\braket{E'l'm'}{Elm} = \delta(E - E')\delta_{ll'}\delta_{mm'}$
\[
  \psi_{Elm}(r\vec n) = \boxed{\frac1\hbar\sqrt{\frac{m}{k}}}\sqrt{\frac2\pi}\,k\,j_l(kr)\,Y_{lm}(\vec n) \qquad k = \sqrt{2mE}/\hbar
\]
\hspace*{3cm}\pozn{$\hookrightarrow\ \big[\frac{\dd E}{\dd k}\big]^{-1/2} = 1\big/\sqrt{\frac{\hbar^2k}{m}}$} \qquad \pozn{$\hookrightarrow E = \frac{\hbar^2k^2}{2m}$}

\underline{pozn:} Po částech konstantní potenciál ve 3D: \\
\quad napojování \quad $R_{kl}(r) = a\,j_l(\tilde kr) + b\,n_l(\tilde kr)$, \quad kde $\tilde k = \frac1\hbar\sqrt{2m(E - V)}$ \quad + spojité $R, R'$ \\
\quad + asymptotika \quad $h_l^{(\pm)}(z) \equiv j_l \pm in_l \xrightarrow{\ z\to\infty\ } \dfrac1z\underbrace{e^{\mp i(l+1)\frac{\pi}{2}}}_{(\mp i)^{l+1}}e^{\pm iz}$ \\
\quad\quad $\hookrightarrow$ \dots\ pro $\tilde k = i\varkappa$ \dots\ jen $h^{(+)}$ povolena v $+\infty$

\sout{$\delta$-chování v počátku:}

%----------------------------------------------------------------------
\strana{QM-C3-13}
\subsection*{5) \underline{Částice v radiálním poli} $\to$ ÚSKO \dots\ $H, L^2, L_z$}
\[
  \braket{\vec Q}{lmE} = \psi(x,y,z) = \psi(r,\theta,\varphi) = R_{El}(r)\,Y_{lm}(\theta,\varphi)
\]
\[
  H\psi = E\psi \ \Rightarrow\ -\boxed{\frac{\hbar^2}{2m}}\frac{1}{r^2}\Big[\pdv{r}r^2\pdv{r} - l(l+1)\Big]R_{kl}(r) + V(r)R_{kl} = ER_{kl}
  \quad\dots\ E = \frac{\hbar^2k^2}{2m}
\]
\[
  \to\quad R_{kl}''(r) + \frac2r R_{kl}' + \Big[k^2\underbrace{- \frac{l(l+1)}{r^2} - U(r)}_{(\star)}\Big]R_{kl}(r) = 0
\]
\[
  (\star) = -V_{\text{eff}}^{(l)}(r)\cdot\frac{2m}{\hbar^2} = -\frac{2m}{\hbar^2}\Big(V(r) + \frac{\overbrace{l(l+1)\hbar^2}^{L^2}}{2mr^2}\Big)
\]
kde: \quad $U(r) = \dfrac{2m}{\hbar^2}V(r)$ \quad nebo $\nearrow$
\[
  \text{nebo}\quad \vzorec{\chi''(r) + \Big[k^2 - \frac{l(l+1)}{r^2} - U(r)\Big]\chi(r) = 0}
\]

\underline{některé řešitelné případy:}
\subsubsection*{5a) \underline{Po částech konstantní potenciál:}}
zavedeme na každé části \quad $\varkappa^2 = k^2 - U$, \quad řešení známe:
\[
  R_{kl}(r) = a\,j_l(\varkappa r) + b\,n_l(\varkappa r)
\]
nebo \quad $\chi(r) = a\,r\,j_l(\varkappa r) + b\,r\,n_l(\varkappa r)$ \quad $\leftarrow$ hodí se hlavně pro $l = 0$, kde $j_l = \frac{\sin z}{z}$, $n_l = -\frac{\cos z}{z}$

+ \underline{okrajové podmínky:} \quad $\chi(r)\big|_{r=0} = 0$ \quad $\to$ \quad \underline{$b = 0$ na prvním úseku}

pro vázané stavy: na úsecích, kde $E < V$ \dots\ tj.\ $k^2 < U$ \dots\ $\varkappa$ ryze imag.\ $= i\lambda$ \\
\quad hodí se zavést: \quad $h_l^{(\pm)}(z) = j_l(z) \pm in_l(z) \xrightarrow{\ z\to\infty\ } (\mp i)^{l+1}\frac1z e^{\pm iz}$

\quad př: pro $h^{(+)}\approx\frac1z\big(\sin(kr - \frac l2\pi) - i\cos(kr - \frac l2\pi)\big) = \frac{-i}{z}e^{i(kr - \frac l2\pi)}$

\quad \dots\ \underline{na posledním úseku} \quad $R_{kl}(r) = a\,h_l^{(+)}(\varkappa r)$ + \sout{$b\,h_l^{(-)}(\varkappa r)$} \hfill {\color{gray}$b = 0$} \\
\hspace*{6.6cm} $\downarrow$ \hspace*{1.7cm} $\downarrow$ \\
\hspace*{6.3cm} $e^{-\lambda r}$ \hspace*{1.4cm} $e^{+\lambda r}$ \qquad $r\to\infty$

\underline{napojování:} spojité $R, R'$ nebo $\chi,\chi'$ \dots\ dá se dělat pomocí $(\ln\chi)'$, $(\ln R)'$ \\
\hspace*{5.5cm} \dots\ nezávisle na norm.

\subsubsection*{5b) \underline{Částice v Coulombickém poli} \quad $V(r) = \frac{\gamma}{r}$ \dots\ $\gamma < 0$ \dots\ přitažlivý}
\hfill\pozn{rozptyl. stavy později}

\underline{pozn:}
\begin{itemize}[nosep]
  \item důležitý případ $\to$ atom vodíku + ionty \\
        \hspace*{2.75cm} $\to$ muonium, protonium, pozitronium \dots \\
        \hspace*{2.75cm} $\to$ báze se využívá při výpočtech
  \item vyšší symetrie (dynamická) -- Runge-Lenzův vektor \\
        $\to$ \uv{náhodná} degenerace \\
        -- separabilní ve sférických, parabolických, eliptických souř.\ (což se někdy hodí)
  \item kvalitativní diskuse v termínech $V_{\text{eff}}^{(l)}(r)$
\end{itemize}

%----------------------------------------------------------------------
\strana{QM-C3-14}
\underline{Radiální (SR):} \quad \dots\ opět se dá řešit algebraicky pomocí posun.\ oper. \\
\hspace*{2.5cm} $\to$ zde \dots\ stručné shrnutí řešení dif.\ rovnice
\[
  \chi''(r) + \frac{2m}{\hbar^2}\Big[E \pm \frac{|\gamma|}{r} - \frac{\hbar^2 l(l+1)}{2mr^2}\Big]\chi(r) = 0
  \qquad\dots\ \begin{tabular}{l} $+|\gamma|$ \dots\ přitažlivý \\ $-|\gamma|$ \dots\ odpudivý\end{tabular}
\]
\hspace*{6cm}{\color{vetabarva}\dots\ diskutovat obrázkem}

přepis do bezrozm.: \quad $\rho = \beta r$ \dots\ $\beta$ \dots\ 1/délka \quad \dots\ $\chi(r) = u(\rho)$
\[
  \to\quad u'' + \underbrace{\boxed{\frac{2mE}{\hbar^2\beta^2}}}_{\text{bezrozm.}}u \pm \boxed{\frac{\gamma 2m}{\hbar^2\beta}}\cdot\frac{u}{\rho} - \frac{l(l+1)}{\rho^2}u = 0
\]
\quad $\hookrightarrow$ bezrozm.\ \dots\ standard $= \frac14$ \dots\ tj.\ $\beta = \sqrt{\dfrac{8m|E|}{\hbar^2}}$ \\
\quad\quad $\hookrightarrow$ \quad $\lambda = \dfrac{2m|\gamma|}{\beta\hbar^2} = \dfrac{|\gamma|}{\hbar}\sqrt{\dfrac{m}{2|E|}}$
\[
  \to\quad \vzorec{u'' + \Big[-\frac14 + \frac\lambda\rho - \frac{l(l+1)}{\rho^2}\Big]u = 0}
\]
\hspace*{3cm}\pozn{$\uparrow$ pro $E < 0$, jinak $(+)$} \quad \pozn{$\uparrow$ pro $\gamma < 0$, jinak $(-)$}

\underline{asymptotika}: \quad $\rho\to 0$ \dots\ stejně jako volná č.\ ($\frac\lambda\rho$ méně singul.\ než $\frac{l(l+1)}{\rho^2}$) \\
\hspace*{3cm} \dots\ tj.\ $u(\rho)\sim\rho^{l+1}$ \dots\ regulární řešení \\
$\rho\to\infty$ \quad \dots\ $u'' = \frac14u$ \quad $\to$ \quad $u(\rho) = e^{\pm\rho/2}$, \quad $\to$ přežije $\ominus$

\underline{položíme}: \quad $u(\rho) = \rho^{l+1}e^{-\rho/2}w(\rho)$:
\[
  \rho w'' + (2l + 2 - \rho)w' + \big[\lambda - (l+1)\big]w = 0
\]
\rule{\linewidth}{0.3pt}

\begin{minipage}[t]{0.55\textwidth}
\underline{ODBOČKA:} \dots\ široká třída ODR 2.\ ř.\ lze převést na (Formánek):

\underline{Gaussova rovnice:}
\[ z(1-z)\bar w'' + \big[c - (1+a+b)z\big]\bar w' - ab\,\bar w = 0 \]
definuje \kw{hypergeometrickou funkci}: \\
\quad $F(a,b,c,z)$ \dots\ lze řešit řadou
\end{minipage}
\vrule\hspace{6pt}
\begin{minipage}[t]{0.42\textwidth}
\underline{Degenerovaný případ:}
\[ \rho w'' + (c - \rho)w' - aw = 0 \]
\[ \to\quad F(a,c,\rho) = \lim_{b\to\infty}F\Big(a,b,c,\frac\rho b\Big) \]
\end{minipage}

\rule{\linewidth}{0.3pt}
\dots\ řešení řadou dá: \quad $w_1(\rho) = F(a,c,\rho) = 1 + \dfrac{1}{1!}\dfrac{a}{c}\rho + \dfrac{1}{2!}\dfrac{a(a+1)}{c(c+1)}\rho^2 + \dots$ \\
\hspace*{2.5cm} druhé řeš.: \quad $w_2(\rho) = \rho^{1-c}F(1+a-c,\,2-c,\,\rho)$ \dots\ singul.\ v poč.

\underline{náš problém:} \quad $c = 2(l+1)$, \quad $a = l + 1 - \lambda$

řada se pro $a = 0,-1,-2,\dots$ redukuje na polynom. \\
jinak se dá DK: \quad $F(a,c,\rho)\sim e^\rho$ \dots\ tj.\ $u = \rho^{l+1}we^{-\rho/2}\sim e^{+\rho/2}$ \ \pozor{špatné}

$\Rightarrow$ \kw{kvantovací podmínka}: \quad $a = l + 1 - \lambda = -\nu$ \quad \dots\ $\nu = 0,1,2,\dots$ \\
\hspace*{6.5cm} $\uparrow$ lze splnit jen pro $\gamma < 0$, jinak je tady $\oplus$
\[
  \hookrightarrow\ \text{tj.}\ \lambda^2 = \frac{\gamma^2}{\hbar^2}\frac{m}{2|E|}
  \ \to\ \vzorec{E_n = -\frac{m\gamma^2}{2\hbar^2n^2} = -\frac{\mathrm{Ry}}{n^2}}
  \quad\nwarrow\ \text{vodík } \gamma = -\frac{e^2}{4\pi\varepsilon_0}\ (\mathrm{SI})
\]
\hspace*{11cm} $\to 1\,\mathrm{Ry} = 13.6\,\mathrm{eV} = \frac12\,\mathrm{Ha}$

\quad kde \quad $n = \nu + l + 1$ \quad \dots\ \kw{hlavní kvantové číslo} \\
\quad\quad \dots\ význam $\nu$ \dots\ řád polynomu $=$ počet nul $R_{nl}(r)$ \quad {\color{vetabarva}radiální kvantové číslo}

%----------------------------------------------------------------------
\strana{QM-C3-15}
\underline{systematizace hladin:} \quad $n = l+1, l+2,\dots$ \\
\quad nebo \quad $n = 1,2,3,\dots$; \quad $l = 0,1,\dots,n-1$; \quad $m = -l,-l+1,\dots,l$

\begin{minipage}{0.75\textwidth}
stupeň degenerace pro dané $E_n$ \dots\ $\vzorec{N = \sum_{l=0}^{n-1}(2l+1) = n^2}$ \\
$\to$ \uv{Náhodná degenerace} \\
\quad \dots\ ve skutečném H \dots\ další st.\ volnosti: \quad
$\begin{aligned}[t] &\text{spin } e^- \dots\ 2n^2 \\ &\text{spin } p^+ \dots\ 4n^2\end{aligned}$
\end{minipage}
\begin{minipage}{0.23\textwidth}
$\Bigg\}$ \small sejmutí degenerace po započt.\ korekcí \dots\ Dirac, pole
\end{minipage}

\underline{typické rozměry atomu:} \quad
$\rho = \beta r = \dfrac{\sqrt{8m|E_n|}}{\hbar}r = \dfrac{\sqrt{4m\frac{m\gamma^2}{\hbar^2n^2}}}{\hbar} = \dfrac{2m|\gamma|}{\hbar^2n}\cdot r$
\[
  \rho = \frac{2r}{na} \quad\dots\ \text{kde}\quad
  \boxed{\begin{aligned} a &= \Big[\frac{m|\gamma|}{\hbar^2}\Big]^{-1} \doteq 0.53\,\text{Å} \\ &\text{\kw{Bohrův poloměr}}\end{aligned}}
\]

\underline{Vlnové funkce:}
\[
  R_{nl}(r) = N_{nl}\underbrace{\Big(\frac{2r}{na}\Big)^l}_{\rho}\underbrace{F\big(-\nu,\,2(l+1),\,\rho\big)}_{\hookrightarrow\ \text{(zobecněný) Laguerův polynom}}e^{-\rho/2}
\]
def: \quad $L_n^\alpha(x) \equiv \dbinom{n+\alpha}{n}F(-n,\alpha+1,x)$
\[
  \to\quad \vzorec{R_{nl}(r) = \frac2n\sqrt{\frac{(n-l-1)!}{(n+l)!}}\;a^{-\frac32}\Big(\frac{2r}{na}\Big)^lL_\nu^{2l+1}\Big(\frac{2r}{na}\Big)e^{-\frac{r}{na}}}
  \qquad \nu = n - l - 1
\]

\underline{více o Laguer.\ polynomech}
\begin{itemize}
  \item altern.\ def: \quad $L_n^\alpha(x) = \displaystyle\sum_{m=0}^{n}(-1)^m\binom{n+\alpha}{n-m}\frac{x^m}{m!} = \frac{1}{n!}e^xx^{-\alpha}\frac{\dd^n}{\dd x^n}\big(e^{-x}x^{n+\alpha}\big)$
  \item ortogonální systém vůči míře $e^{-z}z^\alpha$ \\
  tj.\ \quad $\displaystyle\int_0^\infty L_n^\alpha(x)\,L_{n'}^\alpha(x)\,x^\alpha e^{-x}\,\dd x = \delta_{nn'}\frac{\Gamma(\alpha+n+1)}{n!}$
\end{itemize}

\underline{Pár prvních stavů:}
\begin{align*}
  R_{10}(r) &= 2\sqrt{\tfrac{1}{a^3}}\;e^{-r/a} \qquad\leftarrow\ \text{základní} \\
  R_{20}(r) &= \sqrt{\tfrac{1}{2a^3}}\Big[1 - \frac{r}{2a}\Big]e^{-r/2a} \\
  R_{21}(r) &= \tfrac12\sqrt{\tfrac{1}{6a^3}}\;\frac ra\,e^{-r/2a} \\
  &\ \ \vdots
\end{align*}

%----------------------------------------------------------------------
\strana{QM-C3-16}
\subsubsection*{5c) \underline{Izotropní 3D LHO}}
\[
  V(r) = \tfrac12 m\omega^2r^2 \qquad \big(= \tfrac12m(\omega^2x^2 + \omega^2y^2 + \omega^2z^2)\big)
\]
$\to$ \quad $\psi_{Elm}(x) = R_{nl}(r)\,Y_{lm}(\theta,\varphi)$ \quad \dots\ $\chi = rR$
\[
  \to\quad \chi''(r) + \frac{2m}{\hbar^2}\Big(E - \frac12m\omega^2r^2 - \frac{\hbar^2l(l+1)}{2mr^2}\Big)\chi = 0
\]
\quad \dots\ škálování \quad $x_0 = \sqrt{\dfrac{\hbar}{m\omega}}$ \quad \dots\ $r = \rho x_0$, \quad $\chi(r) = u(\rho)$
\[
  \to\quad u'' + \Big(\frac{2E}{\hbar\omega} - \rho^2 - \frac{l(l+1)}{\rho^2}\Big)u = 0
  \qquad\dots\ \lambda = \frac{2E}{\hbar\omega}\ \dots\ \text{bezrozm.\ energie}
\]
\quad \dots\ asymptotika \quad $\to$ \quad $u(\rho) = \rho^{l+1}e^{-\frac12\rho^2}w(\rho)$
\[
  \to\quad w''(\rho) + \frac2\rho\big((l+1) - \rho^2\big)w'(\rho) + (\lambda - 2l - 3)\,w(\rho) = 0
\]
\quad subst: \quad $z = \rho^2$ \quad \dots\ $w(\rho) = v(z)$
\[
  \to\quad z\,v'' + \big(\underbrace{l + \tfrac32}_{c} - z\big)v' + \big(\underbrace{\tfrac\lambda4 - \tfrac l2 - \tfrac34}_{-a}\big)v = 0
\]
\quad \dots\ opět rovnice pro degener.\ hypgeom.\ fci. \\
\quad\quad $\to$ řešení \dots\ Laguerovy polynomy \quad $v(z) = F(a,c,z)$

kvantovací podmínka: \quad $\frac34 + \frac l2 - \frac\lambda4 = a = -\nu$
\[
  \Rightarrow\quad \vzorec{E = 4\cdot\frac{\hbar\omega}{2}\Big(\nu + \frac l2 + \frac34\Big) = \hbar\omega\Big(\underbrace{2\nu + l}_{\Lambda} + \frac32\Big)}
  \qquad \nu = 0,1,\dots\quad l = 0,1,\dots
\]
\hspace*{7cm} $\Lambda$ \dots\ hlavní kvantové číslo

pro fixní $\Lambda$: \dots\ sudé: $l = 0,2,\dots,\Lambda$ \quad \dots\ $N_{\deg} = \sum_{l=0,2}^{\Lambda}(2l+1) = \frac12(\Lambda+1)(\Lambda+2)$ \quad \pozn{$= \frac12(\frac\Lambda2+1)(2\Lambda+2)$} \\
\hspace*{2.75cm} lichá: $l = 1,3,\dots,\Lambda$ \quad \dots\ $N_{\deg} = \sum_{l=1,3}^{\Lambda}(2l+1) = \frac12(\Lambda+1)(\Lambda+2)$ \quad \pozn{$= \frac12(\frac{\Lambda+1}{2})(2\Lambda+4)$}

\subsection*{6) \underline{Separabilní systém}}
\dots\ na Hilb.\ prostoru \quad $\mathcal{H} = \mathcal{H}^{(1)}\otimes\mathcal{H}^{(2)}$ \quad \dots\ $H = H_1 + H_2$ \\
\hspace*{4.5cm} $\hookrightarrow$ ÚSKO \dots\ $(H_1,A)$, \ $(H_2,B)$

$\Rightarrow$ $H_1, A, H_2, B$ \dots\ ÚSKO na $\mathcal{H}$ \\
\quad navíc \quad $\ket{\psi} = \ket{\psi^{(1)}}\otimes\ket{\psi^{(2)}} = \ket{E_1\alpha}\otimes\ket{E_2\beta}$
\[
  \text{a}\quad H\ket{\psi} = \big(H_1\ket{E_1\alpha}\big)\otimes\ket{E_2\beta} + \ket{E_1\alpha}\otimes\big(H_2\ket{E_2\beta}\big) = \vzorec{(E_1 + E_2)\ket{\psi}}
\]

\begin{priklad}[PŘ: Izotropní LHO]
\[
  E = E_x + E_y + E_z = \hbar\omega\Big(\frac32 + \underbrace{n_x + n_y + n_z}_{\Lambda}\Big)
\]
$\psi(x,y,z) = \psi_{n_x}(x)\,\psi_{n_y}(y)\,\psi_{n_z}(z)$

stupeň degenerace?
\begin{center}
\begin{tikzpicture}[x=0.55cm]
  \foreach \x in {0,1.5,2.7,3.9,5.1,6.3,7.5} \fill (\x,0) circle (2pt);
  \draw (0.75,-0.25) -- (0.75,0.25);
  \draw (4.5,-0.25) -- (4.5,0.25);
  \draw[decorate,decoration={brace,amplitude=4pt}] (-0.2,0.35) -- (7.7,0.35) node[midway,above=3pt] {$\Lambda$};
  \draw[gray,decorate,decoration={brace,mirror,amplitude=3pt}] (-0.2,-0.35) -- (0.6,-0.35) node[midway,below=3pt] {\small $n_x$};
  \draw[gray,decorate,decoration={brace,mirror,amplitude=3pt}] (0.9,-0.35) -- (4.35,-0.35) node[midway,below=3pt] {\small $n_y$};
  \draw[gray,decorate,decoration={brace,mirror,amplitude=3pt}] (4.65,-0.35) -- (7.7,-0.35) node[midway,below=3pt] {\small $n_z$};
  \node[gray,right] at (8,-0.6) {\small $\in(0,\Lambda)$};
\end{tikzpicture}
\end{center}
\dots\ volím 2 přihrádky z $\Lambda + 1$ \dots\ kombinace s opakováním
\[
  N_{\deg} = \binom{\Lambda+1}{2} + \Lambda + 1 = \frac12(\Lambda+1)(\Lambda+2)
\]
\end{priklad}

%----------------------------------------------------------------------
\strana{QM-C3-17}
\begin{priklad}[Jiný příklad: Dvě částice interagující prostřednictvím $V(Q^{(1)} - Q^{(2)})$]
\underline{v klasické mechanice:} \quad
$L(\vec Q^{(1)},\vec Q^{(2)},\dot{\vec Q}^{(1)},\dot{\vec Q}^{(2)}) = \frac12M^{(1)}\dot{\vec Q}^{(1)2} + \frac12M^{(2)}\dot{\vec Q}^{(2)2} + V$

nové proměnné:
\[
  \begin{aligned}
    \vec X &= \frac{M^{(1)}\vec Q^{(1)} + M^{(2)}\vec Q^{(2)}}{M^{(1)} + M^{(2)}} &&\dots\ \text{CMS} \\
    \vec x &= \vec Q^{(1)} - \vec Q^{(2)} &&\dots\ \text{relat.\ poloha}
  \end{aligned}
  \quad\Rightarrow\quad
  \begin{aligned}
    \vec P &= \pdv{L}{\dot X} = \vec p^{(1)} + \vec p^{(2)} \\
    \vec\pi &= \pdv{L}{\dot x} = \frac{M^{(2)}\vec p^{(1)} - M^{(1)}\vec p^{(2)}}{M^{(1)} + M^{(2)}}
  \end{aligned}
\]
+ dá s DK: \quad
$H = \underbrace{\dfrac{\vec P^2}{2M}}_{\text{CMS}} + \boxed{\underbrace{\dfrac{\vec\pi^2}{2\mu} + V(\vec x)}_{\text{rel}}}$ \quad kde \quad $\vzorec{\mu = \frac{M_1M_2}{M_1 + M_2}}$

\underline{v kvantové mechanice:}
\begin{align*}
  H &= -\frac{\hbar^2}{2M^{(1)}}\Delta^{(1)} - \frac{\hbar^2}{2M^{(2)}}\Delta^{(2)} + V(Q^{(1)} - Q^{(2)}) \\
  &= -\frac{\hbar^2}{2M}\Delta_{\text{CMS}} - \frac{\hbar^2}{2\mu}\Delta_{\text{rel}} + V(\vec x) = H_{0\,\text{CMS}} + H_{\text{rel}}
\end{align*}
\rule{\linewidth}{0.3pt}
\underline{DK:}
\begin{align*}
  \pdv{Q^{(1)}} &= \pdv{X}{Q_1}\nabla_X + \pdv{x}{Q_1}\nabla_x = \frac{M_1}{M}\nabla_X + \nabla_x
  &&\to\ \Delta_Q^{(1)} = \frac{M_1^2}{M^2}\Delta_X + \Delta_x + \frac{2M_1}{M}\nabla_X\cdot\nabla_x \\
  \nabla_{Q^{(2)}} &= \pdv{X}{Q_2}\nabla_X + \pdv{x}{Q_2}\cdot\nabla_x = \frac{M_2}{M}\nabla_X - \nabla_x
  &&\to\ \Delta_Q^{(2)} = \frac{M_2^2}{M^2}\Delta_X + \Delta_x - \frac{2M_2}{M}\nabla_X\cdot\nabla_x
\end{align*}
\[
  \frac{1}{M_1}\Delta^{(1)} + \frac{1}{M_2}\Delta^{(2)} = \frac1M\Delta_X + \frac{M_1 + M_2}{M_1M_2}\Delta_x
\]
\rule{\linewidth}{0.3pt}

\underline{Důsledek:} \quad $\mathcal{H} = \mathcal{H}^{(1)}\otimes\mathcal{H}^{(2)} = \mathcal{H}^{(\text{CMS})}\otimes\mathcal{H}^{(\text{rel})}$ \\
\hspace*{2cm} $H = H_{0\,\text{CMS}} + H_{\text{rel}}$ \hfill \pozn{pro $V(\vec x) = V(|\vec x|)$ \dots\ sfér.\ symetrie} \\
$\to$ stacionární stavy \quad $\ket{\psi} = \ket{\vec P_{\text{CMS}}}\otimes\ket{Elm}$ \\
\quad \dots\ tj.\ např.\ atom vodíku v souř.\ reprezentaci:
\[
  \psi(\vec X,\underbrace{\vec x}_{=r\vec n}) = \frac{1}{\sqrt{(2\pi\hbar)^3}}\,e^{\frac{i}{\hbar}\vec P_{\text{CMS}}\cdot\vec X}\cdot R_{nl}(r)\,Y_{lm}(\vec n)
\]
$\to$ jediné, čím se pozná, že 2-částicový syst.\ \dots\ $M\to\mu = \dfrac{M_1M_2}{M_1 + M_2}$

\dots\ stejné jako $H = pe$ je $e^+e^-$ (pozitronium), $p^+\mu^-$, $p^+\bar p$ (protonium), $de^-$ (deuterium) \\
\hspace*{2cm} $\to$ liší se jen $\mu$ \quad \dots\ př.\ \dots\ $m_e$ od $\mu$ pro H se liší o méně než $1\,$\textperthousand
\end{priklad}


\newpage
%======================================================================
% PoznQMI-7.pdf
%======================================================================
\section{Unitární časový vývoj}
\strana{QM-T-1}

\kw{časová Schrödingerova rovnice}:
\[
  \vzorec{i\hbar\pdv{t}\ket{\psi(t)} = H\ket{\psi(t)}} \qquad (\text{tSR})
\]
\dots\ odvodili jsme pro bezstrukt.\ částici, ale \underline{POSTULUJEME OBECNĚ} \\
\hspace*{6cm} $\to$ existence $H$ pro $\forall$ systém

\underline{formální řešení:} \quad
$\ket{\psi(t-s)} = J_s\ket{\psi(t)} = e^{\frac{i}{\hbar}H\cdot s}\ket{\psi(t)}$ \\
\hspace*{6.2cm} $\uparrow$ čas.\ posun v Galilei transf.

obvyklejší zápis: \quad $\ket{\psi(t+s)} = e^{-\frac{i}{\hbar}Hs}\ket{\psi(t)} \equiv U(s)\ket{\psi(t)}$ \\
\phantom{obvyklejší zápis:} \quad nebo \quad
$\ket{\psi(t')} = U(t',t)\ket{\psi(t)} = e^{-\frac{i}{\hbar}(t'-t)H}\ket{\psi(t)}$ \\
\hspace*{3cm} \sout{OPERÁ} \kw{EVOLUČNÍ OPERÁTOR} $\nearrow$

\underline{pozn:} (tSR) lze zobecnit rovněž na interakci, tj.\ $H$ explicitně závislé na čase. Řešení (tSR) pak \pozor{NELZE} psát jako exp! \\
\quad $\to$ příští semestr \dots\ Dysonova řada jako zákl.\ poruch.\ počtu \\
\quad\quad \dots\ nebo numericky jako difro v $t$ \dots \quad
$\to\ \ket{\psi(t+\Delta t)} \simeq \ket{\psi(t)} + \frac{H}{i\hbar}\Delta t\ket{\psi(t)}$ \\
\quad NAPŘ: vnější pole $A(Q)$, $V(Q)$ záv.\ na čase

\subsection{Stac.\ stavy a int.\ pohybu}
pro čas.\ nezáv.\ $H$: \quad \dots\ \kw{stacionární stavy} \quad
\fbox{$H\ket{E_n,\alpha} = E_n\ket{E_n,\alpha}$ \ (SR)}

časový vývoj: \quad $\ket{E_n,\alpha(t)} = e^{-\frac{i}{\hbar}E_nt}\ket{E_n,\alpha}\big|_{t=0}$

\begin{minipage}{0.75\textwidth}
\underline{libovolné měření:} \ (veličina $\hat A$)
\begin{align*}
  \langle A\rangle &\equiv \mel{E_n\alpha(t)}{A}{E_n\alpha(t)} = \langle A\rangle_{t=0} \\
  p_a &= \mel{E_n\alpha(t)}{P_a}{E_n\alpha(t)} = p_a\big|_{t=0}
\end{align*}
\end{minipage}
\begin{minipage}{0.22\textwidth}
$\Bigg\}$ \ \tikz[baseline=-3pt]{\node[draw=prbarva,thick,rounded corners=10pt,align=center,font=\bfseries\small,text=hlavni]{NEZÁVISÍ\\NA\\ČASE};}
\end{minipage}

\begin{veta}[Stacionární stavy]
\begin{itemize}[nosep]
  \item Pozorovatelné ve stacionárních stavech nezávisí na čase
\end{itemize}
\end{veta}

\underline{obecný stav:} \quad $\ket{\psi} = \sum c_{n\alpha}\ket{E_n\alpha}$ \quad $\to$ \quad
$\ket{\psi(t)} = \sum e^{-\frac{i}{\hbar}E_nt}c_{n\alpha}\ket{E_n\alpha}$

obecně \dots\ měření závisí na čase: \quad \dots\ př: $\ket{\psi} = e^{-i\omega_1t}c_1\ket{1} + e^{-i\omega_2t}c_2\ket{2}$
\[
  \dots\ \langle A\rangle_t = \underbrace{|c_1|^2\mel{1}{A}{1} + |c_2|^2\mel{2}{A}{2}}_{\text{nezáv.\ na čase}}
  + \underbrace{c_1^*c_2\mel{1}{A}{2}e^{i(\omega_1-\omega_2)t} + \text{c.c.}}_{\text{interferenční člen}}
\]

\underline{\kw{integrály pohybu}:} \quad $[H,A] = 0$ \quad \dots\ tj.\ $\forall$ vl.\ hodnoty $a$ \dots\ $[H,P_a] = 0$ \\
$\Rightarrow$ měření nezávisí na čase v libovolném stavu
\[
  \mel{\psi(t)}{A}{\psi(t)} = \bra{\psi}e^{\frac{i}{\hbar}Ht}\underset{\curvearrowleft}{A\,e^{-\frac{i}{\hbar}Ht}}\ket{\psi} = \mel{\psi}{A}{\psi}\big|_{t=0}
\]
podobně pro $P_a$ $\to$ pravděpodobnost $p_a$ nezáv.\ na čase \\
speciální případ \dots\ $H$ je integrálem pohybu

%----------------------------------------------------------------------
\strana{QM-T-2}
\begin{priklad}[PŘ]
pro sfér.\ sym.\ \sout{případ je} potenciál je $L_\alpha$ (orbitální mom.\ hybnosti) integrálem pohybu \\
$\to$ částice připravená ve stavu $\ket{lm}$ v něm i zůstane \quad ($p_{l=1,m=1} = 1$ nezáv.\ na čase)
\end{priklad}

\underline{poznámka:} tvrzení lze obrátit: \quad $\langle P_a\rangle$ nezávisí na čase $\forall a,\forall\psi$ \quad $\Rightarrow$ \quad $[H,A] = 0$

\underline{poznámka:} standardní symetrie: \sout{trans}
\begin{center}
\begin{tabular}{lll}
translační & $\longleftrightarrow\ [H,P] = 0$ & \dots\ zachování hybnosti \\
rotační & $\longleftrightarrow\ [H,L_\alpha] = 0$ & \dots\ zach.\ mom.\ hybn. \\
časová & $\longleftrightarrow\ [H(t_1),H(t_2)] = 0$ & \dots\ zach.\ energie
\end{tabular}
\end{center}

\subsection{Rovnice kontinuity pro pravděpodobnost výskytu}
\begin{itemize}
  \item pravděpodobnost výskytu částice v místě $x$: \quad $p_x = \braket{\psi}{x}\braket{x}{\psi}$ \\
        \hspace*{3cm} $\hookrightarrow$ spoj.\ spektrum \dots\ pravd.\ hustota $\equiv\rho(x) = |\psi(x)|^2$ \\
        \dots\ nalezení částice v intervalu $\int\rho(x)\,\dd x$ \dots\ spec.\ $\int_{-\infty}^{\infty}\rho(x)\,\dd x = 1 \equiv \braket{\psi}{\psi}$
  \item Během unitárního vývoje se $\braket{\psi}{\psi}$ zachovává \\
        $\to$ \kw{rovnice kontinuity}: \quad $\vzorec{\partial_t\rho(x) + \operatorname{div}\vec\jmath(x,t) = 0}$ \\
        \dots\ dá se definovat \underline{hustota toku pravděpodobnosti} $\nearrow$
\end{itemize}
\begin{align*}
  \pdv{t}\psi^*(x)\psi(x) &= \psi^*(x)\pdv{t}\psi(x) + \psi(x)\pdv{t}\psi^*(x) \qquad + \text{použití SR} \\
  &= \frac{1}{i\hbar}\Big(-\frac{\hbar^2}{2m}\Big)\big[\psi^*\Delta\psi - \psi\Delta\psi^*\big]
   = \nabla\cdot\Big(\frac{\hbar i}{2m}\big[\psi^*\nabla\psi - \psi\nabla\psi^*\big]\Big)
\end{align*}
\dots\ srovn.\ s rovn.\ kont.\ $\Rightarrow$ \quad
$\vzorec{\vec\jmath = -\frac{i\hbar}{2m}\big[\psi^*\nabla\psi - \psi\nabla\psi^*\big] = \frac{\hbar}{m}\operatorname{Im}\psi^*\nabla\psi}$

\begin{itemize}
  \item \underline{pozn:} srovn.\ \dots\ operátor rychlosti \quad $\hat{\vec v} = \frac{\vec p}{m} = -\frac{i\hbar}{m}\nabla$ \\
        \dots\ tj.\ $\vec\jmath = \operatorname{Re}\psi^*\hat{\vec v}\psi$ \quad \dots\ pro $\psi\in\mathcal{H}$ (tj.\ $L^2$ funkce): \quad
        $\mel{\psi}{\hat{\vec v}}{\psi} = \int\vec\jmath\,\dd^3x$
  \item \underline{pozn:} pokud $\psi(x,t) = R(x,t)e^{\frac{i}{\hbar}S(x,t)}$ \quad $\to$ \quad
        $\vec\jmath = R^2\frac1m\vec\nabla S = \rho\cdot\frac{\vec\nabla S}{m}$
  \begin{itemize}[nosep]
    \item[--] pro reál.\ fci $\psi$ \dots\ $\vec\jmath = 0$
    \item[--] pro stac.\ stavy $\pdv{\rho}{t} = 0$ $\to$ $\operatorname{div}\vec\jmath = 0$
    \item[--] pro rovin.\ vlnu \quad $\psi = Ne^{i\vec p\cdot\vec x/\hbar}$ $\to$ $\vec\jmath = |N|^2\cdot\frac{\vec p}{m}$
    \item[--] pro součet rov.\ vln v 1D \quad $\psi = A_+e^{i\vec p\cdot\vec x/\hbar} + A_-e^{-i\vec p\cdot\vec x/\hbar}$ $\to$
          $\vec\jmath = \frac{\vec p}{m}\big\{|A_+|^2 - |A_-|^2\big\}$ \\
          ale obecně ne aditivní
  \end{itemize}
\end{itemize}

%----------------------------------------------------------------------
\strana{QM-T-3}
\subsection{Evoluční operátor v souřadnicové repr.}
(opět $H$ nezáv.\ na čase; navíc volná částice)
\begin{itemize}[nosep]
  \item[--] Budeme se zabývat vývojem vlnového balíku pro volnou částici
  \item[--] Hodí se zvlášť rozlišit případ, kdy děláme evoluci do budoucna: \hfill $(t_2 > t_1)$
\end{itemize}
\[
  \ket{\psi(t_2)} = G^{(+)}(t_2,t_1)\ket{\psi(t_1)} \equiv \Theta(t_2-t_1)\,U(t_2,t_1)\ket{\psi(t_1)}
\]
\noindent\rule{\linewidth}{0.3pt}
\begin{itemize}
  \item důvod \dots\ $G^{(+)}$ má lepší Fourier.\ transf.\ než $U$ \hfill \kw{Retardovaný $G$-operátor}
  \item $G^{(+)}$ je \kw{Greenovým operátorem} (SRt) pro pravou stranu:
  \[
    \Big(i\hbar\pdv{t_2} - H\Big)\hat G^{(+)}(t_2,t_1) = i\hbar\,\delta(t_2-t_1)
    \qquad \text{\pozn{pozn.\ \dots\ někdy se $i\hbar$ zahrnuje do def.\ $G^{(+)}$}}
  \]
  (později rozdělíme $H = H_0 + V$ a $H_0\to G_0$; $V$ na pravou stranu)
  \item podobně se def \quad $\hat G^{(-)}(t_2,t_1) = -\Theta(t_1-t_2)\,\hat U(t_2,t_1)$ \quad \dots\ \kw{ADVANCOVANÝ $G$-operátor} \\
        $\to$ stejná rovnice \dots\ jiná okraj.\ podm: $G^{(-)} = 0$ pro $t_1 < t_2$
\end{itemize}
\noindent\rule{\linewidth}{0.3pt}

vyjádření v souřadnicové reprezentaci:
\[
  \braket{x_2}{\psi(t_2)} = \bra{x_2}G^{(+)}(t_2,t_1)\int\dd x_1\ket{x_1}\braket{x_1}{\psi(t_1)}
\]
tj.
\[
  \psi(x_2,t_2) = \int\dd x_1\,\underbrace{G^{(+)}(x_2,t_2,x_1,t_1)}_{\text{\kw{propagátor}}}\,\psi(x_1,t_1) \qquad \pozn{($*$)}
\]
pozn: lze chápat jako $G$-funkci pro počát.\ podmínku, tj.: \\
\quad časový vývoj $\psi(x_2,t_2) = G(x_2,t_2,x_1,t_1)$ pro vlnovou funkci \\
\quad s počát.\ podm.\ $\delta$-funkce \quad $\psi(x,t)\big|_{t=t_1} = \delta(x-x_1)$

\pozn{($*$)} je důsledkem linearity (tSR)

\subsubsection*{Propagátor volné částice}
\begin{align*}
  G_0^{(+)}(\vec x_2,t_2,\vec x_1,t_1) &\equiv \Theta(t_2-t_1)\bra{\vec x_2}e^{-\frac{i}{\hbar}\frac{p^2}{2m}(t_2-t_1)}
     \underbrace{\phantom{I}}_{I = \int\ket{\vec p}\bra{\vec p}\dd^3p}\ket{\vec x_1} \\
  &= \Theta(t_2-t_1)\int\frac{\dd^3p}{(2\pi\hbar)^3}\exp\Big\{\frac{i}{\hbar}\Big[\vec p\cdot\underbrace{(\vec x_2-\vec x_1)}_{\vec x}
     - \frac{\vec p^{\,2}}{2m}\underbrace{(t_2-t_1)}_{t}\Big]\Big\}
\end{align*}
\noindent\rule{\linewidth}{0.3pt}
\dots\ součin 3 integrálů: \quad $\displaystyle\int\frac{\dd p_x}{2\pi\hbar}e^{\frac{i}{\hbar}[p_x\cdot x - \frac{p_x^2}{2m}t]}$
\quad \dots\ osciluje pro $p\to\infty$ \\
\hspace*{7cm} \dots\ konverguje $\lim_{\varepsilon\to 0^+}\int\dots\big|_{t\to t-i\varepsilon}$

\noindent\rule{0.6\linewidth}{0.3pt} \\
$p^2\frac{t}{2m} - px = \frac{t}{2m}\big(p - \frac{x}{2}\cdot\frac{2m}{t}\big)^2 - \frac{x^2}{4}\cdot\frac{2m}{t}$ \\
\noindent\rule{0.6\linewidth}{0.3pt}
\[
  = \int\frac{\dd p_x}{2\pi\hbar}e^{-\frac{i}{\hbar}\cdot\frac{t}{2m}(p - \frac{xm}{t})^2}\cdot e^{\frac{i}{\hbar}\cdot\frac{mx^2}{2t}}
  = \underbrace{\frac{1}{2\pi\hbar}\cdot\sqrt{\frac{\pi}{\frac{i}{\hbar}\cdot\frac{t}{2m}}}}_{\to\ \sqrt{\frac{m}{2\pi i\hbar t}}}\;e^{\frac{i}{\hbar}\frac{mx^2}{2t}}
\]

%----------------------------------------------------------------------
\strana{QM-T-4}
tj. \hspace{3cm} \pozn{pozn: dimenze jinde se neprojeví $\to$ ${}^{3/2}$}
\[
  \vzorec{G_0^{(+)}(\vec x_2,t_2,\vec x_1,t_1) = \Theta(t_2-t_1)\Big(\frac{m}{2\pi\hbar i(t_2-t_1)}\Big)^{3/2}
  \exp\Big\{\frac{im}{2\hbar}\frac{(\vec x_2-\vec x_1)^2}{(t_2-t_1)}\Big\}}
\]

\subsection{Časový vývoj vlnového balíku pro vol.\ částici}
Gauss.\ balík $\ket{\psi}$:
\[
  \psi(p) = \frac{1}{\sqrt[4]{2\pi(\Delta p)^2}}e^{-\left(\frac{p-p_0}{2\Delta p}\right)^2}{\color{prbarva}e^{-\frac{i}{\hbar}px_0}}
  \ \underset{\int\frac{\dd x}{\sqrt{2\pi\hbar}}e^{-\frac{i}{\hbar}px}}{\overset{\int\frac{\dd p}{\sqrt{2\pi\hbar}}e^{\frac{i}{\hbar}px}}{\rightleftarrows}}\
  \psi(x) = \frac{1}{\sqrt[4]{2\pi(\Delta x)^2}}e^{-\left(\frac{x-x_0}{2\Delta x}\right)^2}e^{+\frac{i}{\hbar}p_0x}
\]
kde: \quad $\braket{\psi}{\psi} = 1$ \qquad $\mel{\psi}{\hat p}{\psi} = p_0$ \qquad $\mel{\psi}{\hat x}{\psi} = x_0$ \\[3pt]
\fbox{$\Delta x\cdot\Delta p = \frac{\hbar}{2}$} \qquad
$\mel{\psi}{(\hat p-p_0)^2}{\psi} = (\Delta p)^2$ \quad $\mel{\psi}{(\hat x-x_0)^2}{\psi} = (\Delta x)^2$

\underline{časový vývoj:} \quad $\ket{\psi(t)} = e^{-\frac{i}{\hbar}Ht}\ket{\psi}$ \\
(1). buď pomocí $G_0^+$: \quad $\braket{x}{\psi(t)} \equiv \psi(x,t) = \int\dd^3x'\,G_0^{(+)}(\vec x,t,x',0)\,\psi(x')$ \\
(2). nebo $p$-repre: \quad $\psi(p,t) = e^{-\frac{i}{\hbar}\frac{p^2}{2m}t}\psi(p)$ \quad + F.T.\ do $x$-repre. \\
\pozn{každopádně vede na úpravu na čtverec + Gauss integrál}

AD (2): uděláme pro $x_0 = 0$; \ def \ \fbox{$\mu = \frac{p}{2\Delta p}$} \ \fbox{$\mu_0 = \frac{p_0}{2\Delta p}$} \quad \dots\ tj.\ $\big(\frac{p-p_0}{2\Delta p}\big)^2 = (\mu-\mu_0)^2$

dále \quad $e^{\frac{i}{\hbar}px} = e^{2i\overbrace{\scriptstyle\frac{2}{\hbar}\Delta p}^{1/\Delta x}\cdot\frac{p}{2\Delta p}\cdot\frac{x}{2}} = \exp\{2i\mu\varkappa\}$ \quad \dots\ \fbox{$\varkappa \equiv \frac{x}{2\Delta x}$}
\[
  \bullet\ -\frac{i}{\hbar}\frac{p^2}{2m}t = -i\underbrace{\Big(\frac{2}{\hbar}\Delta p\Big)}_{1/\Delta x}\Big(\frac{p}{2\Delta p}\Big)^2\cdot\frac{\Delta p}{m}\cdot t
  = -i\mu^2\,\boxed{\frac{\Delta v}{\Delta x}\cdot t} \equiv -i\mu^2\cdot\alpha
\]

nyní: \quad
$\displaystyle\psi(x,t) = \int\frac{\dd p}{\sqrt{2\pi\hbar}}e^{\frac{i}{\hbar}px}\psi(p,t)
 = \frac{1}{\sqrt[4]{2\pi(\Delta p)^2}}\int\frac{\dd p}{\sqrt{2\pi\hbar}}e^{-(\mu-\mu_0)^2 - i\alpha\mu^2 + 2i\mu\varkappa}$

\noindent\rule{\linewidth}{0.3pt}
úprava na čtverec: \quad $-\big[\mu^2(1+i\alpha) - 2\mu(\mu_0+i\varkappa) + \mu_0^2\big] = -a(\mu-q)^2 + b$ \\
\hspace*{7.5cm} kde: \ $a = (1+i\alpha)$ \\
+ záměna $\int\dd p = 2\Delta p\int\dd\mu$ \hspace{3.3cm} $q = (\mu_0+i\varkappa)/a$ \\
\hspace*{8.2cm} $b = aq^2 - \mu_0^2$ \\
\noindent\rule{\linewidth}{0.3pt}
\[
  = \frac{1}{\sqrt[4]{2\pi}}\underbrace{\frac{2\Delta p}{\sqrt{2\pi\hbar\Delta p}}}_{1/\sqrt{\Delta x}}
  \underbrace{\int\frac{\dd\mu}{\sqrt\pi}e^{-a(\mu-q)^2}}_{\to\ 1/\sqrt a}\cdot
  \exp\Big\{\underbrace{\frac{(\mu_0+i\varkappa)^2}{1+i\alpha} - \mu_0^2}_{\text{extrakce abs.\ hodnoty a fáze}}\Big\}
\]
\noindent\rule{\linewidth}{0.3pt}
$b = -\dfrac{(\varkappa-i\mu_0)^2}{1+\alpha^2}(1-i\alpha) - \mu_0^2\dfrac{1+\alpha^2}{1+\alpha^2}
 = -\dfrac{1}{1+\alpha^2}\big\{\varkappa^2 - 2\mu_0\varkappa\alpha + \mu_0^2\alpha^2 - i(2\varkappa\mu_0 + \alpha\varkappa^2 - \alpha\mu_0^2)\big\}$ \\
\noindent\rule{0.75\linewidth}{0.3pt}
\[
  = \frac{1}{\sqrt[4]{2\pi(\Delta x)^2}}\cdot\frac{1}{\sqrt{1+i\alpha}}\exp\Big\{-\frac{(\varkappa-\mu_0\alpha)^2}{1+\alpha^2}\Big\}
  \exp\Big\{i\frac{\alpha(\varkappa^2-\mu_0^2) + 2\varkappa\mu_0}{1+\alpha^2}\Big\}
\]

%----------------------------------------------------------------------
\strana{QM-T-5}
\underline{interpretace:}
\[
  |\psi(x,t)|^2 = \frac{1}{\sqrt{2\pi}}\frac{1}{\Delta x\sqrt{1+\alpha^2}}\exp\Big\{-\frac{(\varkappa-\mu_0\alpha)^2}{1+\alpha^2}\Big\}
\]
kde: \quad $\varkappa - \mu_0\alpha = \dfrac{x}{2\Delta x} - \dfrac{p_0}{2\Delta p}\cdot\dfrac{\Delta p}{m}\cdot\dfrac{t}{\Delta x} = \dfrac{x - v_0t}{2\Delta x}$ \quad \dots\ pohyb vln.\ balíku \\[3pt]
\quad $\Delta x^2(1+\alpha^2) = \Delta x^2 + \Delta v^2t^2 \equiv \sigma(t)^2$ \quad \dots\ rozplývání bal. \\
\hspace*{9.5cm} $\uparrow$ semi-klas.\ interpretace
\[
  \text{tj.}\quad \vzorec{|\psi(x,t)|^2 = \frac{1}{\sqrt{2\pi}}\frac{1}{\sigma(t)}\exp\Big\{-\Big(\frac{x - v_0t}{2\sigma(t)}\Big)^2\Big\}}
\]

\underline{pozn:} \sout{v} 3D \dots\ prostě součin tří takových faktorů \\
\underline{pozn:} odvození lze provést pro $p_0 = 0$ + udělat Galilei transf. \\
\quad $\to$ faktor $(x - v_0t)$

\begin{minipage}{0.45\textwidth}
\subsubsection*{Rozplývání vlnového balíku:}
\vspace{1.2cm}
{\color{red}minimální vln.\ balík, který jsme zvolili jako poč.\ podm.}
\end{minipage}
\begin{minipage}{0.53\textwidth}
\begin{tikzpicture}[scale=1]
  \draw[->] (0,0) -- (5,0) node[right] {$t$};
  \draw[->] (0,0) -- (0,3.6) node[left] {$\Delta x,\ \sigma$};
  \draw[gray] (0,0) -- (4.6,3.4) node[right,font=\small] {$\frac{\Delta p}{m}\cdot t = \alpha\cdot\Delta x$};
  \draw[red,thick,domain=0:4.2,samples=60] plot (\x,{sqrt(0.81+0.55*\x*\x)});
  \fill[red] (0,0.9) circle (2.5pt);
  \node[gray,left] at (0,0.9) {$\Delta x$};
  \node[red] at (2.2,2.1) {$\sigma(t)$};
  \draw[red,dashed,->] (-1.6,2) to[bend left=10] (-0.1,1);
\end{tikzpicture}
\end{minipage}

\underline{pozn:} Gauss balíky = hustá množina v $L^2$; \ Pravd.\ výskytu v kon.\ oblasti $\sim\frac1t$ (1D) a $\frac{1}{t^3}$ (3D) \ (díky rozplývání)

\subsection{Časový vývoj v Heisenbergově obraze}
doposud: \underline{\kw{Schrödingerův obraz}:} \\
\quad $\ket{\psi_S} \equiv \ket{\psi_S(t)} = U(t)\ket{\psi_S(0)}$ \hfill \pozn{obecně $\ket{\psi(t)} = U(t-t_0)\ket{\psi(t_0)}$, BÚNO $t_0 = 0$}

fyzikální předpovědi:
\[
  \langle A\rangle_t = \mel{\psi_S(t)}{A}{\psi_S(t)} = \bra{\psi_S(0)}\underbrace{U^\dagger(t)AU(t)}_{A_H(t)}\ket{\psi_S(0)}
\]
\dots\ podobně \quad $p_a(t) = \mel{\psi_S(t)}{P_a}{\psi_S(t)} = \mel{\psi_S(0)}{P_a^{(H)}(t)}{\psi_S(0)}$

$\to$ \underline{\kw{Heisenbergův obraz}:} \quad \dots\ ztotožníme s $\ket{\psi_S} = \ket{\psi_H}$ v referenčním čase $t = 0$

stav.\ vektory čas.\ nezávislé: \quad $\ket{\psi_H} = \ket{\psi_S(0)}$ \\
měřitelné veličiny záv.\ na $t$: \quad $A_H(t) = U^\dagger(t)AU(t)$
\[
  \frac{\dd A}{\dd t} = \frac{\dd}{\dd t}\big\{U^\dagger(t)AU(t)\big\} \underset{\uparrow\ i\hbar\frac{\dd}{\dd t}U = HU}{=}
  \frac{1}{i\hbar}\big\{U^\dagger AHU - HU^\dagger AU\big\} = \frac{1}{i\hbar}U^\dagger[A,H]U
\]
\[
  \to\quad \vzorec{\frac{\dd A_H(t)}{\dd t} = \frac{1}{i\hbar}[A_H(t),H]\ \Big({+}\,\pdv{A}{t}\Big)}
  \qquad \text{(samozřejmě $H_H = H$ nezáv.\ na čase)}
\]
\kw{Heisenbergova pohyb.\ rovnice} \hfill {\color{prbarva}$\nwarrow$ příklad potenc.\ závislost: $V(\hat x,t)$}

%----------------------------------------------------------------------
\strana{QM-T-6}
\begin{itemize}
  \item \underline{pozn:} \quad $(AB)_H = A_HB_H$ \\
        \phantom{\underline{pozn:}} \quad $\big(f(A)\big)_H = f(A_H)$ \qquad
        $A = \sum_a P_a\cdot a$ \ $\to$ \ $A_H = \sum_a aP_a^{(H)}$
\end{itemize}

\begin{priklad}[PŘ: částice v potenc.\ poli $V(\hat x)$ \dots\ $\hat H = \frac{\hat p^2}{2m} + V(\hat x)$]
Heisenbergovy pohyb.\ rovnice:
\begin{align*}
  \dot x_H &= \frac{1}{i\hbar}[x_H,H] = \pozn{$[x,p] = i\hbar$; $[x,f(p)] = i\hbar f'(p)$} = \frac{p_H}{m} \\
  \dot p_H &= \frac{1}{i\hbar}[p_H,H] = -\Big[\frac{\dd V(x)}{\dd x}\Big]_H
\end{align*}
\end{priklad}

\begin{veta}[Ehrenfestův teorém]
\[
  \vzorec{m\frac{\dd}{\dd t}\langle x\rangle = \langle p\rangle \qquad
  \frac{\dd}{\dd t}\langle p\rangle = -\Big\langle\pdv{V}{x}\Big\rangle}
\]
\end{veta}
\dots\ od klasických pohyb.\ rovnic se liší \quad $\hookleftarrow$ \ $\pdv{x}V\big|_{x=\langle x\rangle} \equiv V'(\langle x\rangle)$ \\
\dots\ tj.\ OK pro \quad -- dobře lokalizovanou částici \\
\hspace*{3.05cm} -- hladké $V(x)$ \quad \pozn{$\leftarrow$ $V$ se málo mění na vzdál.\ srovn.\ s $\Delta x$}

\begin{priklad}[PŘ: Volný Gauss.\ balík \dots\ viz.\ str 4--5]
\[
  \dot x_H = \frac{1}{i\hbar}[x_H,H] = \frac{p_H}{m} \ \to\ x_H = x + \frac{p}{m}t \pozn{s}, \qquad
  \dot p_H = 0 \ \to\ p_H = p_S = p
\]
čas.\ vývoj.\ středních hodnot \hfill (bez ohledu na tvar $\psi$!)
\begin{align*}
  \langle x\rangle_t &= \langle x_H\rangle_0 = \langle x\rangle_0 + \frac{\langle p\rangle_0}{m}t \\
  \Delta x^2(t) &= \langle(x - \langle x\rangle)^2\rangle_t = \langle x_H^2\rangle_0 - \langle x_H\rangle_0^2 \\
  &= \Delta x^2(t=0) + \frac{\Delta p(t=0)^2}{m^2}t^2 + \underbrace{\frac1m\big(\langle xp+px\rangle_0 - 2\langle x\rangle_0\langle p\rangle_0\big)t}_{\substack{\text{obecně může způsobit dočas.}\\ \text{zužování balíku ($=0$ pro Gausse)}}}
\end{align*}
\begin{center}
\begin{tikzpicture}[gray,scale=0.8]
  \node at (-2,1.5) {\small Obecně};
  \draw[->] (0,0) -- (0,2.2) node[left] {$(\Delta x)^2$};
  \draw[->] (0,0) -- (3.5,0) node[right] {$t$};
  \draw[thick] (0,1.9) .. controls (0.6,0.4) and (1.4,0.3) .. (3,2);
  \node[right,align=left,font=\small] at (3.3,1.3) {kladné\\a pro $t\to\infty$ rostoucí};
\end{tikzpicture}
\end{center}
DK: $\langle xp + px\rangle$ je $0$, dokud je $\psi(x)$ nebo $\psi(p)$ reálné \\
\quad např.\ $\psi(x)$ \dots \quad
$\langle xp\rangle = \braket{\phi_x}{\phi_p} = i\cdot$reál.\ číslo $= -i\hbar\int\psi(x)\cdot x\,\psi'(x)$ \\
\quad $\ket{\phi_x} \equiv x\ket{\psi}$ \quad\ \
$\langle px\rangle = \braket{\phi_p}{\phi_x} = -i\cdot$stejné číslo $= \big(\braket{\phi_x}{\phi_p}\big)^*$ \\
\quad $\ket{\phi_p} \equiv p\ket{\psi}$ \quad pokud je $p_0\neq 0$ i $x_0\neq 0$, pak $\psi(x)$ i $\psi(p)$ jsou komplexní, \\
\hspace*{2cm} ale $\langle xp+px\rangle$ bylo $0$ s \ $\psi(x) = \frac{1}{\sqrt{\cdots}}e^{-\frac{(x)^2}{(2\Delta x)^2}}e^{\frac{i}{\hbar}p_0x}$ \\
$\to$ záměna $x\to x - x_0$ \quad $\langle(x-x_0)p + p(x-x_0)\rangle = 0 = \langle xp+px\rangle - 2x_0p_0$
\end{priklad}

%----------------------------------------------------------------------
\strana{QM-T-7}
\begin{priklad}[PŘ: Lineární harmonický oscilátor v Heisenberg.\ repr.]
\[
  \left.\begin{aligned}
    \dot x_H &= \frac{1}{i\hbar}[x_H,H] = \frac{p_H}{m} \\
    \dot p_H &= \frac{1}{i\hbar}[p_H,H] = -m\omega^2x_H = -kx_H
  \end{aligned}\right\}\ \text{jako klas.\ rovnice, ale pro operátory}
\]
\[
  \text{\underline{důsl}:}\quad
  \left.\begin{aligned}
    \frac{\dd}{\dd t}\langle x\rangle &= \frac1m\langle p\rangle \\
    \frac{\dd}{\dd t}\langle p\rangle &= -k\langle x\rangle
  \end{aligned}\right\}\
  \begin{minipage}{0.45\textwidth}
  střední hodnoty splňují stejné rce jako klas, bez ohledu na stav \\ (může být zcela neklasický)
  \end{minipage}
\]
ale \pozor{pozor}: $x,p$ nekompatibilní pozorovatelné $\nearrow$
\end{priklad}

\subsubsection*{Koherentní (kvaziklasické) stavy LHO}
\underline{cíl:}
\begin{itemize}[nosep]
  \item[--] najít kvantový stav $\ket{\psi}$ LHO, který se v limitě $E\gg\hbar\omega$ bude podobat klasické částici.
  \item[--] vedlejší produkt: najdeme vývoj Gauss.\ balíku pro LHO
\end{itemize}
vidíme, že $\langle x\rangle$ a $\langle p\rangle$ automaticky klas.\ pohyb. \\
$\to$ další požadavky: $\Delta x,\Delta p$ malé a $E = H_{\text{klas}}(\langle x\rangle,\langle p\rangle)$; $\Delta E$ malé \\
$\to$ vyjde: důležité jsou vl.\ st.\ $\hat a = (\hat q + i\hat p)/\sqrt2$

v klasické fyzice: \quad $a = a_0e^{-i\omega t}$ \quad \dots\ může zkusit, že Heisenbergova rovnice \\
\hspace*{3cm} $\hookrightarrow\ \equiv\frac{1}{\sqrt2}(q_0 + ip_0)$ \hspace{1.5cm} $\Rightarrow$ \quad $\hat a_H = \hat a_S\cdot e^{-i\omega t}$

vyjdeme z:
\begin{itemize}
  \item \underline{požadavek:} \quad
        $\langle E\rangle \equiv \langle\hbar\omega(\hat N + \tfrac12)\rangle = \hbar\omega\langle a^\dagger a\rangle + \underbrace{\tfrac{\hbar\omega}{2}}_{\substack{\text{zanedb}\\ \text{pro } E\gg\hbar\omega}}
         = \hbar\omega\cdot|a_0|^2 = \frac{\hbar\omega}{2}(p^2 + q^2)$
\end{itemize}
\quad tj.\ současně: \ (1) $\langle a\rangle = a_0 \equiv \frac{1}{\sqrt2}(q_0 + ip_0)$ \\
\quad $E\gg\hbar\omega$ \qquad\quad (2) $\langle a^\dagger a\rangle = |a_0|^2$
\[
  \Rightarrow\ \text{platí jen ve stavu}\quad a\ket{\psi} = a_0\ket{\psi} \overset{\text{def}}{\equiv} \alpha\ket{\psi}
  \quad\to\ \text{def: \kw{koherentní stav}}
\]
\underline{pozn:} základní stav $\ket{0}$ je koh.\ stavem pro $\alpha = 0$, ale semiklasický jen pro $\alpha\gg1$

\noindent\rule{0.7\linewidth}{0.3pt} \\
\underline{DK:} \quad def \ $\ket{v} = a\ket{\psi} - a_0\ket{\psi}$
\[
  \text{pak}\quad 0 \leq \braket{v}{v} = \underbrace{\mel{\psi}{a^\dagger a}{\psi}}_{|a_0|^2} - a_0^*\underbrace{\mel{\psi}{a}{\psi}}_{=a_0}
  - a_0\underbrace{\mel{\psi}{a^\dagger}{\psi}}_{=a_0^*} + |a_0|^2\underbrace{\braket{\psi}{\psi}}_{=1}
\]
\quad tj.\ $0 = \braket{v}{v}$, ale rovnost $\curvearrowleft$ platí jen pro $\ket{v} = 0$ \quad (def.\ skalár.\ souč.\ v LVP) \\
\noindent\rule{0.7\linewidth}{0.3pt}

dále budeme vyšetřovat vlastnosti koh.\ st.\ a ukážeme, že splňuje $\forall$ podmínky hledané na klas.\ lim pro $\alpha\gg1$

\sout{$\bullet$ $\Delta x,\Delta p$ je malé} \ \underline{vyjádření $\ket{\psi}\equiv\ket{\alpha}$ ve stac.\ stavech pro LHO}
\[
  \ket{\alpha} = \sum_n c_n\ket{n} \ \to\ \hat a\ket{\alpha} = \alpha\ket{\alpha} \ \Rightarrow\ c_{n+1} = \frac{\alpha}{\sqrt{n+1}}c_n
\]
tj.\ \quad $\ket{\alpha} = N\sum_n\frac{\alpha^n}{\sqrt{n!}}\ket{n}$ \quad + normování \quad
$\braket{\alpha}{\alpha} = \sum_n\frac{|\alpha|^{2n}}{n!} = e^{|\alpha|^2}$ \quad $\Rightarrow$ \quad $N = e^{-|\alpha|^2/2}$

%----------------------------------------------------------------------
\strana{QM-T-8}
\underline{Měření energie:} \quad $p_n = |c_n|^2 = \dfrac{|\alpha|^{2n}}{n!}e^{-|\alpha|^2}$ \quad $\to$ \kw{Poissonovo rozdělení}

maximum? \quad $p_n = \frac{|\alpha|^2}{n}p_{n-1}$ \quad $\to$ max pro $|\alpha|^2\simeq n$
\[
  \langle E\rangle \equiv \langle H\rangle = \sum p_n\underbrace{E_n}_{\hbar\omega(n+\frac12)} = \text{(dá se spočíst analyt.)}
  \underset{\text{ale níže}}{=} \hbar\omega\big(\langle a^\dagger a\rangle + \tfrac12\big) = \hbar\omega\big(|\alpha|^2 + \tfrac12\big)
\]
rozptyl $\Delta E$? (variance)
\begin{align*}
  \langle H^2\rangle &= \hbar^2\omega^2\mel{\alpha}{(a^\dagger a + \tfrac12)^2}{\alpha}
   = \hbar^2\omega^2\bra{\alpha}\big(a^\dagger\underbrace{aa^\dagger}_{=a^\dagger a + [a,a^\dagger]}a + a^\dagger a + \tfrac14\big)\ket{\alpha} \\
  &= \hbar^2\omega^2\big(|\alpha|^4 + 2|\alpha|^2 + \tfrac14\big)
\end{align*}
\[
  \to\quad \Delta E^2 = \langle H^2\rangle - \langle H\rangle^2 = \hbar^2\omega^2|\alpha|^2
  \quad\to\quad \vzorec{\frac{\Delta E}{\langle E\rangle} \simeq \frac{1}{|\alpha|} \ll 1}\ \text{ v lim } n\to\infty
\]

\underline{hodnoty $\Delta x,\Delta p$?} \quad \dots\ automaticky \quad
$\left.\begin{aligned}
  \langle x\rangle &= x_0\sqrt2\cdot\operatorname{Re}\alpha = x_c \\
  \langle p\rangle &= p_0\sqrt2\cdot\operatorname{Im}\alpha = p_c
\end{aligned}\right\}$ \dots\ klasické \\
\hspace*{9cm} \pozn{$x_0 = \sqrt{\frac{\hbar}{m\omega}}$, \ $p_0 = \frac{\hbar}{x_0}$}
\[
  \left.\begin{aligned}
    \langle x^2\rangle_\alpha &= \frac{x_0^2}{2}\big[(\alpha+\alpha^*)^2 + 1\big] \\
    \langle p^2\rangle_\alpha &= \frac{p_0^2}{2}\big[1 - (\alpha-\alpha^*)^2\big]
  \end{aligned}\right\}
  \Rightarrow\
  \begin{aligned}
    \Delta x &= \frac{x_0}{\sqrt2} \\
    \Delta p &= \frac{p_0}{\sqrt2}
  \end{aligned}
  \ \Rightarrow\ \vzorec{\Delta x\cdot\Delta p = \frac{\hbar}{2}}
\]
\hfill \dots\ \underline{minimalizuje neurč.!}

\underline{zápis $\ket{\alpha}$ v $x$-reprezentaci:}
\[
  \text{platí:}\quad \psi_\alpha(x) = e^{i\theta_\alpha}\sqrt[4]{\frac{1}{2\pi(\Delta x)^2}}\cdot
  \exp\Big\{-\Big(\frac{x - \langle x\rangle}{2\Delta x}\Big)^2 + \frac{i}{\hbar}\langle p\rangle x\Big\}
\]
\dots\ tj.\ \kw{minimalizující Gauss balík}

\pozn{
DK: 1) všimněte si, že zákl.\ stav $\braket{x}{0} = \phi_0\big(\frac{x}{x_0}\big)\frac{1}{\sqrt{x_0}}$, \ $\phi_0 = \frac{1}{\sqrt[4]{\pi}}e^{-\frac12q^2}$ \\
\quad splňuje $\uparrow$ s $\Delta x = \frac{x_0}{\sqrt2}$ a $\langle x\rangle = 0$ \\[2pt]
2) Najdeme unit.\ operátor $D(\alpha)$: $\ket{0}\to\ket{\alpha}$
\[
  \ket{\alpha} = e^{-|\alpha|^2/2}\sum_n\frac{\alpha^n}{\sqrt{n!}}\ket{n} = e^{-|\alpha|^2/2}\sum_n\frac{(\alpha a^\dagger)^n}{n!}\ket{0} = e^{-|\alpha|^2/2}e^{\alpha a^\dagger}\ket{0}
\]
\dots\ není unit., ale \ $D(\alpha) = e^{\alpha a^\dagger - \alpha^*a} = e^{-|\alpha|^2/2}e^{\alpha a^\dagger}e^{-\alpha^*a}$ \\
je unitární, a $e^{-\alpha^*a}\ket{0} = \ket{0}$ \hfill $\hookrightarrow$ užití $e^Ae^B = e^{A+B}\cdot e^{\frac12[A,B]}$ \\[2pt]
tj.\ $\psi_\alpha(x) = \mel{x}{D(\alpha)}{0} =$
\[
  D(\alpha) = \exp\Big\{\underbrace{\tfrac{\alpha-\alpha^*}{\sqrt2}}_{ip_c}\hat q - i\underbrace{\tfrac{\alpha+\alpha^*}{\sqrt2}}_{q_c}\hat p\Big\}
  = \exp\{ip_c\hat q - iq_c\hat p\} = e^{ip_c\hat q}e^{-iq_c\hat p}
  \underbrace{e^{\frac{\alpha^{*2}-\alpha^2}{4}}}_{e^{-\frac12[\hat q,\hat p]p_cq_c} = e^{\frac{i}{2}p_cq_c} = e^{i\theta_\alpha}}
\]
\[
  \text{tj.}\quad \psi_\alpha(x) = \bra{x}e^{ip_c\hat q}e^{-iq_c\hat p}\ket{0}
  = e^{\frac{i}{\hbar}\langle p\rangle\cdot x}\underbrace{\bra{x}e^{-\frac{i}{\hbar}\langle x\rangle\hat p}\ket{0}}_{\substack{\text{posunutí o }\langle x\rangle\\ \to\ \braket{x}{0}|_{x = x-\langle x\rangle}}}
  \qquad \hat q = \frac{\hat x}{x_0} = \frac{p_0}{\hbar}\hat x
\]
$\Rightarrow$ c.b.d. \qquad pozn: dá se alternativně DK z vyjádření $\ket{\alpha} = \sum_n c_n\ket{n}$ v $x$-repr.\ pomocí Hermite + rozgenerující fce $H_n$
}

%----------------------------------------------------------------------
\strana{QM-T-9}
\underline{časový vývoj?}
\[
  \text{platí}\quad \ket{\alpha(t)} = \sum_n c_ne^{-\frac{i}{\hbar}E_nt}\ket{n} = e^{-i\omega t/2}\sum_n\frac{(\alpha e^{-i\omega t})^n}{\sqrt{n!}}\ket{n}\cdot e^{-|\alpha|^2/2}
\]
\[
  \text{tj.}\quad \vzorec{\ket{\alpha(t)} = e^{-i\omega t/2}\ket{\alpha = \alpha_0e^{-i\omega t}}}
\]
$\Rightarrow$ zůstává koher.\ stavem s $\alpha$ měnícím se podle klas.\ fyziky \\
$\Rightarrow$ Gaussův balík zůstává, jen se posouvá $\langle x\rangle,\langle p\rangle$ \dots\ jako klas.\ fyz. \\
\quad navíc zůstává $\Delta x\Delta p = \frac{\hbar}{2}$ a $\langle E\rangle$ se samozřejmě zachovává

%----------------------------------------------------------------------
\strana{QM-T-10}
\subsection{Evoluční operátor pro $H(t)$ \ {\normalfont\small -- až nakonec kapitoly}}
\pozn{v originále stránka označená QM-T-\sout{8} 10}
\[
  \ket{\psi(t)} = U(t,t_0)\ket{\psi(t_0)}
\]
(SRt) $\to$ \quad $i\hbar\pdv{t}U(t,t_0) = H(t)U(t,t_0)$ \quad \dots\ s poč.\ podm.\ $U(t_0,t_0) = I$

integrací: \quad $U(t,t_0) = I + \dfrac{1}{i\hbar}\displaystyle\int_{t_0}^{t}H(t_1)U(t_1,t_0)\,\dd t_1 = \mathcal{F}[U]$

\quad iterování \quad $U^{(0)} = I$: \quad $U^{(n+1)} = \mathcal{F}[U^{(n)}]$
\begin{align*}
  \to\ U(t,t_0) = I &+ \frac{1}{i\hbar}\int_{t_0}^{t}H(t_1)\,\dd t_1 + \Big(\frac{1}{i\hbar}\Big)^2\int_{t_0}^{t}\int_{t_0}^{t_1}H(t_1)H(t_2)\,\dd t_2\dd t_1 + \dots \\
  \text{($*$) \kw{Dysonova řada}} \quad &+ \Big(\frac{1}{i\hbar}\Big)^n\int_{t_0}^{t}\int_{t_0}^{t_1}\!\!\cdots\int_{t_0}^{t_{n-1}}H(t_1)H(t_2)\cdots H(t_n)\,\dd t_n\cdots\dd t_2\dd t_1 + \dots
\end{align*}

\begin{itemize}
  \item dá se zkonstruovat $\{H(t)\}$: \quad $[H(t_1),H(t_2)] = 0$ \quad \dots\ ne moc častý příp.
\end{itemize}
\begin{minipage}{0.7\textwidth}
pak
\begin{align*}
 \int_{t_0}^{t}\int_{t_0}^{t_1}H(t_1)H(t_2)\,\dd t_2\dd t_1
 &= \frac12\Big[\int_{t_0}^{t}\int_{t_0}^{t}H(t_1)H(t_2)\,\dd t_1\dd t_2\Big] \\
 &= \frac12\Big[\int_{t_0}^{t}H(t')\,\dd t'\Big]^2
\end{align*}
\end{minipage}\hfill
\begin{minipage}{0.25\textwidth}
\begin{tikzpicture}[gray,scale=0.9]
  \draw[->] (0,0) -- (0,2) node[left] {$t_2$};
  \draw[->] (0,0) -- (2.2,0) node[below] {$t_1$};
  \fill[pattern=north east lines,pattern color=gray] (0,0) -- (1.7,1.7) -- (1.7,0) -- cycle;
  \draw[dashed] (0,0) -- (1.8,1.8);
  \draw[dotted] (1.7,0) -- (1.7,1.7);
\end{tikzpicture}
\end{minipage}

podob. \quad
$\displaystyle\int_{t_0}^{t}\int_{t_0}^{t_1}\!\!\cdots\int_{t_0}^{t_{n-1}}H(t_1)\cdots H(t_n)\,\dd t_n\cdots\dd t_1 = \frac{1}{n!}\Big[\int_{t_0}^{t}H(t')\,\dd t'\Big]^n$
\[
  \to\quad U(t,t_0) = \exp\Big\{-\frac{i}{\hbar}\int_{t_0}^{t}H(t')\,\dd t'\Big\}
\]

\begin{itemize}
  \item \underline{obecně:} zavedeme \kw{čas.\ uspořádání} \quad $T\big(H(t_i)H(t_j)\big) = H(t_<)H(t_>)$ \\
        \hspace*{6cm} $t_< \equiv \min(t_1,t_2)$, \ $t_> = \max(t_1,t_2)$
\end{itemize}
pak \quad $T\big([H(t_1),H(t_2)]\big) = 0$, a tedy
\begin{align*}
  \int_{t_0}^{t}\cdots\int_{t_0}^{t_{n-1}}H(t_1)\cdots H(t_n)\,\dd t_n\cdots\dd t_1
  &= \frac{1}{n!}\int_{t_0}^{t}\dd t_1\cdots\int_{t_0}^{t}\dd t_n\,T\big(H(t_1)\cdots H(t_n)\big) \\
  &= \frac{1}{n!}T\Big(\Big[\int_{t_0}^{t}\dd t'\,H(t')\Big]^n\Big)
\end{align*}
\[
  \to\quad \vzorec{U(t,t_0) = T\exp\Big\{-\frac{i}{\hbar}\int_{t_0}^{t}H(t')\,\dd t'\Big\}}
\]
\dots\ vypadá zběsile, ale s úspěchem se používá \\
\quad $\to$ výhodné pro poruch.\ rozvoje v mocninách operátorů

\newpage
%======================================================================
% Poznamky_1_stacionarni_poruchy.pdf
%======================================================================
\section{Stacionární poruchová teorie}
\strana{PT-1}
\hfill (Rayleigh--Schrödinger)

\underline{cíl:}
\begin{itemize}[nosep]
  \item přesně řešitelný problém
  \item[$\hookrightarrow$] najít, jak se změní energie a vln.\ fce vázaných stavů
\end{itemize}

\begin{tabular}{@{}ll}
problém: & $H = H_0 + \lambda H_1$ \qquad\qquad $\lambda H_1$ \uv{\kw{malá porucha}} \\
+ známe & $H_0\ket{n} = E_n^{(0)}\ket{n}$ \hspace{2.6cm} upřesníme později
\end{tabular}

$\lambda$ \dots\ v principu může být kontrolovatelné vnější pole, \\
\quad ale většinou se vkládá pro usnadnění matem.\ analýzy a nakonec $\lambda = 1$
\quad ($\nearrow$ umožňuje rozlišit řády)

+ Formální rozvoj do řady: \qquad $(H_0 + \lambda H_1)\ket{\psi_n} = E_n\ket{\psi_n}$ \hfill $(*)$
\begin{align*}
  E_n &= E_n^{(0)} + \lambda E_n^{(1)} + \lambda^2E_n^{(2)} + \dots && (RE) \\
  \ket{\psi_n} &= \ket{\psi_n^{(0)}} + \lambda\ket{\psi_n^{(1)}} + \lambda^2\ket{\psi_n^{(2)}} + \dots && (R\psi)
\end{align*}
+ dosazení do $(*)$, srovnání stejných řádů $\lambda$:
\begin{align*}
  (P0):&\quad \big[H_0 - E_n^{(0)}\big]\ket{\psi_n^{(0)}} = 0 \\
  (P1):&\quad \big[H_0 - E_n^{(0)}\big]\ket{\psi_n^{(1)}} = \big(E_n^{(1)} - H_1\big)\ket{\psi_n^{(0)}} \\
  (P2):&\quad \big[H_0 - E_n^{(0)}\big]\ket{\psi_n^{(2)}} = \big(E_n^{(1)} - H_1\big)\ket{\psi_n^{(1)}} + E_n^{(2)}\ket{\psi_n^{(0)}} \\
  &\qquad\vdots \\
  (PS):&\quad \big[H_0 - E_n^{(0)}\big]\ket{\psi_n^{(s)}} = \big(E_n^{(1)} - H_1\big)\ket{\psi_n^{(s-1)}} + E_n^{(2)}\ket{\psi_n^{(s-2)}} + \dots + E_n^{(s)}\ket{\psi_n^{(0)}}
\end{align*}
\hfill $\swarrow$ má formu rekur.\ relací: $\ket{\psi_n^{(s)}}$ pomocí $\ket{\psi_n^{(s')}}$, \ $s' < s$

rovnice $(*)$ neurčuje $\ket{\psi}$ jednoznačně $\to$ normalizace a fáze

je výhodné volit \kw{intermediální normalizaci}: \qquad $\braket{n}{\psi_n} = 1$ \hfill $(**)$ \\
\quad (to je povoleno, pokud $\braket{n}{\psi_n}\neq 0$; dá se předp., že pro $H_1$ malé nebude $\ket{\psi_n}$ ortogonální k $\ket{n}$), navíc fixuje i fázi

$(**)$ + dosazení $(R\psi)$ \quad $\Rightarrow$ \quad $\braket{n}{\psi_n^{(0)}} = 1$; \quad $\braket{n}{\psi_n^{(s)}} = 0 \quad \forall s > 0$

řešení hledáme ve tvaru rozvoje \quad $\ket{\psi_n} = \sum_{n'}\ket{n'}\braket{n'}{\psi_n}$ \\
\quad tj. \quad $\ket{\psi_n^{(s)}} = \sum_{n'}\ket{n'}\braket{n'}{\psi_n^{(s)}}$ \quad \dots\ $n'\neq n$ pro $s > 0$

(výhoda $(**)$ \dots\ projekce $(PS)$ na $\bra{n}$ \dots\ spousta členů zmizí)
\[
  \to\quad \bra{n}(PS):\quad
  \underbrace{\mel{n}{H_0 - E_n^{(0)}}{\psi_n^{(s)}}}_{0} = -\mel{n}{H_1}{\psi_n^{(s-1)}} + 0 + \dots + E_n^{(s)}\underbrace{\braket{n}{\psi_n^{(0)}}}_{1}
\]
\quad $\hookrightarrow$ neboť $(P0)$: řešení \quad $E_n^{(0)} = \varepsilon_n$, \quad $\ket{\psi_n^{(0)}} = \ket{n}$ \quad $\leftarrow$ kvůli $(**)$

tj.\ obecně
\[
  \vzorec{E_n^{(s)} = \mel{n}{H_1}{\psi_n^{(s-1)}}}
\]
\quad \dots\ stačí vln.\ fci o řád nižší

%----------------------------------------------------------------------
\strana{PT-2}
\underline{vlnové funkce:} \quad \tikz[baseline=(c.base)]\node[draw,rounded corners=5pt,inner sep=1.5pt](c){\small 0.ř.}; \quad $\ket{\psi_n^{(0)}} = \ket{n}$

\noindent\rule{\linewidth}{0.4pt}

\tikz[baseline=(c.base)]\node[draw,rounded corners=5pt,inner sep=1.5pt](c){\small 1.ř.}; \quad $E_n^{(1)} = \mel{n}{H_1}{\psi_n^{(0)}} = \mel{n}{H_1}{n}$
\begin{align*}
  \bra{m}(P1):&\quad (\varepsilon_m - \varepsilon_n)\braket{m}{\psi_n^{(1)}} = \mel{m}{E_n^{(1)} - H_1}{n} = -\mel{m}{H_1}{n}; \quad m\neq n \\
  \to&\quad \braket{m}{\psi_n^{(1)}} = -\frac{\mel{m}{H_1}{n}}{\varepsilon_m - \varepsilon_n} \quad\text{pro } m\neq n; \qquad \braket{n}{\psi_n^{(1)}} = 0 \text{ podle } (**)
\end{align*}
\[
  \text{tj.}\quad \vzorec{\ket{\psi_n^{(1)}} = \sum_{m\neq n}\frac{\ket{m}\mel{m}{H_1}{n}}{\varepsilon_n - \varepsilon_m}} \qquad (\psi 1)
\]
\underline{diskuse:} výraz $\uparrow$ umožňuje odhadnout podmínky platnosti rozvoje \\
\quad \dots\ aby $\ket{\psi_n^{(1)}}$ malé, musí být \quad $|\mel{m}{H_1}{n}| \ll |\varepsilon_n - \varepsilon_m| \quad \forall m,n$

\noindent\rule{\linewidth}{0.4pt}

\tikz[baseline=(c.base)]\node[draw,rounded corners=5pt,inner sep=1.5pt](c){\small 2.ř.}; \quad
\[
  \vzorec{E_n^{(2)} = \mel{n}{H_1}{\psi_n^{(1)}} = \sum_{m\neq n}\frac{\mel{n}{H_1}{m}\mel{m}{H_1}{n}}{\varepsilon_n - \varepsilon_m}} \qquad (E2)
\]
\noindent\rule{\linewidth}{0.4pt}

\subsection*{Případ degenerovaného spektra}
formule $(\psi 1)$ a $(E2)$ \pozor{divergují!} \\
Problém souvisí s tím, že neplatí $\ket{\psi_n^{(0)}} = \ket{n}$, ale lin.\ kombinace $\ket{n}$-ek se stejnou energií (ale nevíme jaká lin.\ kombinace). \\
Je výhodné změnit značení: \qquad $H_0\ket{m,r} = \varepsilon_m\ket{m,r}$ \hfill $(\bullet)$

\underline{Problém:} i \tikz[baseline=(c.base)]\node[draw,rounded corners=5pt,inner sep=1.5pt](c){\small 0-tý řád}; v jistém smyslu závisí na poruše; konkrétně \\
\quad plné řešení pro $H = H_0 + \lambda H_1$ je \quad
$\ket{\psi_{m,t}(\lambda)} \xrightarrow{\ \lambda\to 0\ } \ket{\psi_{m,t}^{(0)}}$, \quad $E_{m,t}\xrightarrow{\ \lambda\to 0\ }\varepsilon_m$ \\
podle $(P0)$ je $\ket{\psi_{m,t}^{(0)}}$ řešením $(\bullet)$, ale nyní to neznamená $\ket{\psi_{m,t}^{(0)}} = \ket{m,r}$! \\
(NA PŘEDNÁŠCE ZMÍNIT, JAK VYPADÁ ZOBEC.\ $(Pi)$ \dots\ V PODST.\ TOTOŽNÉ) \\
může obecně být: \quad $\ket{\psi_{m,t}^{(0)}} = \sum_r C_r^{(t)}\ket{m,r}$ \\
\quad rovnice $(\bullet)$ nedává žádnou restrikci na $C_r^{(t)}$ \dots\ (podobná situace v časové poruch.\ teorii (ME)) \\
\quad $\to$ nutno z vyšších řádů \dots\ pochopitelné

\tikz[baseline=(c.base)]\node[draw,rounded corners=5pt,inner sep=1.5pt](c){\small 1.\ řád}; \quad $(P1)\ \to\ \big[H_0 - \underbrace{E_m^{(0)}}_{\varepsilon_m}\big]\ket{\psi_{m,t}^{(1)}} = \big[E_{m,t}^{(1)} - H_1\big]\underbrace{\ket{\psi_{m,t}^{(0)}}}_{\sum_r C_r^{(t)}\ket{m,r}}$ \quad (+ dosazení)
\[
  \bra{m,r'}(P1):\quad \underbrace{[\varepsilon_m - \varepsilon_m]}_{=0}\braket{m,r'}{\psi_{m,t}^{(1)}} = \sum_r\Big\{E_{m,t}^{(1)}\delta_{r'r} - \mel{m,r'}{H_1}{m,r}\Big\}C_r^{(t)}
\]
\begin{veta}[Zobecnění formule pro 1.\ řád opravy energie -- současně udává 0.\ ř.\ vln.\ fce]
\[
  \text{tj.}\qquad \vzorec{\sum_r\mel{m,r'}{H_1}{m,r}\,C_r^{(t)} = E_{m,t}^{(1)}\,C_{r'}^{(t)}}
\]
\end{veta}
umožňuje definovat \uv{správné} vektory 0.\ řádu: \quad $\ket{m,t} \equiv \ket{\psi_{m,t}^{(0)}} = \sum_r C_r^{(t)}\ket{m,r}$

zobecnění normalizační podm.\ $(**)$: \quad $\braket{m,t}{\psi_{m,t}(\lambda)} = 1$ \\
\quad $\to$ v 0.ř. \quad $\braket{m,t'}{\psi_{m,t}^{(0)}} = \delta_{tt'}$ \qquad
$\hookrightarrow$ implikuje $\braket{m,t}{\psi_{m,t}^{(s)}} = 0 \quad \forall s > 0$ \\
\hspace*{6.3cm} neříká nic o $\braket{m,t'}{\psi_{m,t}^{(s)}}$; $s > 0$; $t\neq t'$

%----------------------------------------------------------------------
\strana{PT-3}
\underline{1.\ řád vlnové funkce:} \\
Nyní již budeme používat bázi $\ket{m,t}$, v níž \quad $\mel{m,t}{H_1}{m,t'} = E_{m,t}^{(1)}\delta_{tt'}$ \\
\hspace*{6cm} \pozor{pozor} $\mel{m,t}{H_1}{m',t'}$ není omezené $\lrcorner$
\[
  \begin{array}{l}\bra{m',t'}(P1):\\ (m'\neq m)\end{array}\quad
  (\varepsilon_{m'} - \varepsilon_m)\braket{m',t'}{\psi_{m,t}^{(1)}} = E_m^{(1)}\underbrace{\braket{m',t'}{m,t}}_{0} - \underbrace{\mel{m',t'}{H_1}{m,t}}_{\neq 0\ \dots\ \text{viz}\ \uparrow}
\]
\[
  \text{tj.}\quad \vzorec{\braket{m',t'}{\psi_{m,t}^{(1)}} = \frac{\mel{m',t'}{H_1}{m,t}}{\varepsilon_m - \varepsilon_{m'}}} \qquad m'\neq m
\]
\pozor{pozor:} $\nearrow$ nefixuje $\braket{m,t'}{\psi_{m,t}^{(1)}}$ \textcircled{?} \ (dá se určit z (P1) projekt.\ $\bra{m,t'}$, $t'\neq t$) \\
\hspace*{5.6cm} ale může přispět i (P2)

obecně pro korekce energie vyšších řádů platí stále: \quad $E_{m,t}^{(s)} = \mel{m,t}{H_1}{\psi_{m,t}^{(s-1)}}$ \\
\quad díky tomu, že $\braket{m,t}{\psi_{m,t}^{(s)}} = 0 \quad \forall s > 0$; speciálně

\underline{2.\ řád energie:} \pozn{přeškrtnutý první pokus o výpočet}
\begin{align*}
  E_{m,t}^{(2)} &= \mel{m,t}{H_1}{\psi_{m,t}^{(1)}} \\
  &= \sum_{t'}\underbrace{\mel{m,t}{H_1}{m,t'}}_{E_{m,t}^{(1)}\delta_{tt'}}\underbrace{\braket{m,t'}{\psi_{m,t}^{(1)}}}_{=0\text{ pro } t = t'}
   + \sum_{m'\neq m}\sum_{t'}\frac{\mel{m,t}{H_1}{m',t'}\mel{m',t'}{H_1}{m,t}}{\varepsilon_m - \varepsilon_{m'}}
\end{align*}
\hfill tj.\ členy, které neznáme \textcircled{?}, nepřispívají (totéž ve vyšších řádech)
\[
  \text{tj.\ opět platí}\qquad \vzorec{E_{m,t}^{(2)} = \sum_{t',\,m'\neq m}\frac{|\mel{m,t}{H_1}{m',t'}|^2}{\varepsilon_m - \varepsilon_{m'}}}
\]
\noindent\rule{\linewidth}{0.4pt}

\textbf{ZÁVĚR:} všechny problémy degenerovaného spektra se odstr.

\underline{pozn.\ na okraj:} \quad $\bra{m,t'}(P1)$, \ $t'\neq t$:
\begin{align*}
  E_{m,t}^{(1)}\braket{m,t'}{\psi_{m,t}^{(1)}} &= \mel{m,t'}{H_1}{\psi_{m,t}^{(1)}} \\
  &= \sum_{t''}\underbrace{\mel{m,t'}{H_1}{m,t''}}_{E_{m,t'}^{(1)}\delta_{t't''}}\braket{m,t''}{\psi_{m,t}^{(1)}}
   + \sum_{m''\neq m}\sum_{t''}\mel{m,t'}{H_1}{m'',t''}\frac{\mel{m'',t''}{H_1}{m,t}}{\varepsilon_m - \varepsilon_{m''}}
\end{align*}
\[
  \to\quad \big(E_{m,t}^{(1)} - E_{m,t'}^{(1)}\big)\braket{m,t'}{\psi_{m,t}^{(1)}} = \sum_{m''\neq m,\,t''}\frac{\mel{m,t'}{H_1}{m'',t''}\mel{m'',t''}{H_1}{m,t}}{\varepsilon_m - \varepsilon_{m''}}
\]
tato rovnice určuje $\braket{m,t'}{\psi_{m,t}^{(1)}}$, pokud $H_1$ snímá degeneraci \\
\quad (není automatické, např.\ vodík \dots\ $E_n$; degenerace dle $l, m$; \\
\quad porucha může sejmout jen $l$-degeneraci, např.\ korekce na jaderný potenciál) \\
tento problém mohou vyřešit vyšší řády, nebo jde o princip.\ problém, že i $H_0 + H_1$ má degeneraci, a tudíž $\ket{\psi_{m,t}}$ nejsou určeny jednoznačně; je např.\ možné položit $\braket{m,t'}{\psi_{m,t}^{(1)}} = 0$, pokud $t$ je dobré kvantové číslo pro $H = H_0 + H_1$

$\uparrow$ VŠE V TÉTO POZN.\ PODSTATNÉ JEN, POKUD NÁS ZAJÍMÁ: VLN.\ FCE \\
vln.\ fci potřebujeme i do vyšších řádů korekcí $E_{m,t}^{(s)}$ \\
\quad -- obecně stále komplikovanější výrazy, lepší $\lrcorner$

%----------------------------------------------------------------------
\subsection{Brillouin--Wignerova teorie poruch}
\strana{PT-4}
až dosud: Rayleigh--Schrödinger \dots\ poruchová teorie \underline{explicitní} \\
\hspace*{5cm} \dots\ výrazy obsahující neporušené $E_n^{(0)} = \varepsilon_n$ \\
nyní \underline{implicitní} výrazy obsahující $E_n$, ale jednodušší struktura vyšších ř.

\begin{definice}[projekční operátory]
\[ Q_n = \sum_{n'\neq n}\ket{n'}\bra{n'} = 1 - \ket{n}\bra{n} \]
\end{definice}
+ stále používáme normalizaci $\braket{n}{\psi_n} = 1$, \quad tj.\ $\ket{\psi_n} = \ket{n} + Q_n\ket{\psi_n}$

rovnici pro přesné vl.\ v.\ píšeme ($\lambda = 1$):
\begin{align*}
  (E_n - H_0)\ket{\psi_n} &= H_1\ket{\psi_n} \qquad\qquad (I) \\
  \underbrace{Q_n(E_n - H_0)}_{(E_n - H_0)Q_n}\ket{\psi_n} &= Q_nH_1\ket{\psi_n}
  \qquad\to\ [H_0,Q_n] = 0 \text{ díky tomu, že mají stejné vl.\ v.}
\end{align*}
\[
  \Rightarrow\quad Q_n\ket{\psi_n} = \underbrace{(E_n - H_0)^{-1}}_{\text{rozumíme ve vl.\ prostoru } Q_n \text{ odp.\ } \lambda = 1}Q_nH_1\ket{\psi_n}
  \qquad \text{\pozor{!} nutno dobře def.\ inverzi}
\]
def \quad $R_n \equiv (E_n - H_0)^{-1}Q_n = Q_n(E_n - H_0)^{-1} \equiv Q_n(E_n - H_0)^{-1}Q_n$

\dots\ lze psát \quad $\vzorec{\ket{\psi_n} = \ket{n} + R_nH_1\ket{\psi_n}}$ \quad \dots\ form.\ řeš.\ $= (I - R_nH_1)^{-1}\ket{n}$
\pozn{geom.\ řada -- alternativní odvození} \\
\dots\ podobnost s Lippmann--Schwingerovou rovnicí není náhodná \\
$\to$ lze řešit iteracemi (podob.\ Born.\ řadě) \dots\ pro \uv{$H_1 \ll H_0$}
\[
  \ket{\psi_n} = \ket{n} + R_nH_1\ket{n} + (R_nH_1)^2\ket{n} + \dots
\]
\begin{itemize}
  \item $\bra{n}(I)$: \quad $E_n = \varepsilon_n + \mel{n}{H_1}{\psi_n} = \varepsilon_n + \mel{n}{H_1}{n} + \mel{n}{H_1(R_nH_1)}{n} + \mel{n}{H_1(R_nH_1)^2}{n} + \dots$
\end{itemize}
\begin{veta}[Brillouin--Wignerova řada]
\begin{align*}
  \text{tj.}\quad
  E_n = \underbrace{\varepsilon_n}_{(a)} &+ \underbrace{\mel{n}{H_1}{n}}_{(b)}
  + \underbrace{\sum_{n'\neq n}\frac{\mel{n}{H_1}{n'}\mel{n'}{H_1}{n}}{E_n - \varepsilon_{n'}}}_{(c)} \\
  &+ \underbrace{\sum_{n'\neq n}\sum_{n''\neq n}\frac{\mel{n}{H_1}{n'}\mel{n'}{H_1}{n''}\mel{n''}{H_1}{n}}{(E_n - \varepsilon_{n'})(E_n - \varepsilon_{n''})}}_{(d)} + \dots
\end{align*}
\end{veta}

\underline{Poznámky:}
\begin{enumerate}[label=\arabic*)]
  \item implicitní rovnice \dots\ obsahuje $E_n$ na pravé straně. \\
        použití: \quad A) iterace \\
        \phantom{použití:} \quad B) dosazení $E_n$ přísl.\ řádu \dots\ např.\ ve III.\ ř.\ $E_n = \varepsilon_n + \langle H_1\rangle_n$ v (c) \\
        \hspace*{8.6cm} a $E_n = \varepsilon_n$ v (d)
  \item po dosazení $E_n = \sum E_n^{(s)}\lambda^s$ a rozvoji pravé strany do řady v $\lambda$ rekonstruujeme Rayleigh--Schrödingera.
\end{enumerate}

%----------------------------------------------------------------------
\subsection{Stacionární poruchová teorie -- řešené příklady}
\strana{PT-5}

\begin{priklad}[\textcircled{1} Lineární harm.\ oscilátor s lin.\ poruchou]
\[
  H_0 = \frac{p^2}{2m} + \frac12 m\omega^2x^2; \qquad H_1 = \lambda x
\]
\underline{přesné řešení:} \quad $H = H_0 + H_1 = \frac12 m\omega^2(x + x_0)^2 + \bar\varepsilon_0$
\qquad $x_0 = \dfrac{\lambda}{m\omega^2}$, \quad $\bar\varepsilon_0 = -\dfrac12\dfrac{\lambda^2}{m\omega^2}$

tj.\ spektrum \quad $\varepsilon_n = \hbar\omega(n + \frac12)$ \quad $\longrightarrow$ \quad $E_n = \hbar\omega(n + \frac12) + \bar\varepsilon_0$

\underline{stacionární p.t.:} \quad platí \quad $H_1 = \lambda x = \lambda\sqrt{\dfrac{\hbar}{2m\omega}}\,(a + a^\dagger) \equiv A(a + a^\dagger)$
\begin{align*}
  E_n^{(1)} &= \mel{n}{H_1}{n} = A\mel{n}{a + a^\dagger}{n} = 0 \\
  E_n^{(2)} &= \sum_{n'\neq n}\frac{|\mel{n}{H_1}{n'}|^2}{\varepsilon_n - \varepsilon_{n'}}
   = A^2\sum_{n'\neq n}\frac{|\mel{n}{a + a^\dagger}{n'}|^2}{\varepsilon_n - \varepsilon_{n'}}
   = \frac{A^2}{\hbar\omega}\Big\{|\mel{n}{a^\dagger}{n-1}|^2 - |\mel{n}{a}{n+1}|^2\Big\} \\
  &= \frac{A^2}{\hbar\omega}\big\{n - (n+1)\big\} = -\frac{\lambda^2}{2m\omega^2}
\end{align*}
\[
  \text{tj.}\quad \vzorec{E_n = \varepsilon_n + E_n^{(1)} + E_n^{(2)} + \mathcal{O}(\lambda^3) = \hbar\omega(n + \tfrac12) - \frac{\lambda^2}{2m\omega^2}}
  \quad\dots\ \text{přesný výsledek}
\]
+ zdůvodnit proč \ + otázka do publika: vlnové funkce?

vln.\ fce do 1.ř.:
\[
  \ket{\psi_n} = \ket{n} + \sum_{n'\neq n}\frac{\ket{n'}\mel{n'}{H_1}{n}}{\varepsilon_n - \varepsilon_{n'}}
  = \ket{n} + \frac{A}{\hbar\omega}\Big\{\sqrt n\,\ket{n-1} - \sqrt{n+1}\,\ket{n+1}\Big\}
\]
\end{priklad}

\begin{priklad}[\textcircled{2} Matice $2\times2$ (diagonální s mimodiag.\ poruchou)]
\[
  H_0 = \begin{pmatrix}\varepsilon_1 & 0\\ 0 & \varepsilon_2\end{pmatrix} \qquad
  H_1 = \begin{pmatrix}0 & v\\ v^* & 0\end{pmatrix}
\]
$\to$ přesné řešení \quad $E = \frac12(\varepsilon_1 + \varepsilon_2) \pm \frac12\sqrt{(\varepsilon_1 - \varepsilon_2)^2 + 4|v|^2}$ \\
\quad degenerace: $\varepsilon_1 = \varepsilon_2 \equiv \varepsilon$ \quad \dots \quad $E = \varepsilon \pm |v|$

B-W teorie: \quad
$\left.\begin{aligned}
  E_1 &= \varepsilon_1 + \frac{|v|^2}{E_1 - \varepsilon_2} \\
  E_2 &= \varepsilon_2 + \frac{|v|^2}{E_2 - \varepsilon_1}
\end{aligned}\right\}$
\quad $\Longleftrightarrow$ \quad $(E - \varepsilon_1)(E - \varepsilon_2) = |v|^2$ \quad $\to$ přes.\ řešení $\uparrow$

R-S teorie: \quad $E_1 = \varepsilon_1 + \dfrac{|v|^2}{\varepsilon_1 - \varepsilon_2}$ \qquad OK za předp.\ $\dfrac{|v|}{|\varepsilon_1 - \varepsilon_2|} \ll 1$

\dots\ na CV ověřit rozvojem přesného výrazu
\end{priklad}

\dots\ tj.\ \kw{B-W teorie nemá problém s degenerací ani blízkostí degenerace} \\
\quad (formalismus nezávisí na degeneraci)

\dots\ alternativa \dots\ viz Formánkova učebnice \dots\ R-S teorie formulovaná pomocí proj.\ operátorů.

%----------------------------------------------------------------------
\strana{PT-6}
\subsubsection*{\textcircled{3} Indukovaný el.\ dipól atomu v E-poli}

\underline{Elektrický dipól:} \dots\ rozložení náboje \quad $\rho(\vec r) = -|\psi(\vec r)|^2\cdot e$
\[
  \vec d \equiv \int\rho(\vec r)\cdot\vec r\,\dd^3r = -\int\psi^*(\vec r)\,e\vec r\,\psi(\vec r)\,\dd^3r \equiv -\langle e\hat{\vec r}\rangle_\psi
\]
\quad operátor el.\ dipólu: \quad $\hat{\vec d} = -e\hat{\vec r}$ \quad (vektor)

\begin{definice}[Polarizovatelnost atomu $\alpha$]
\begin{tabular}{@{}p{0.55\linewidth}l}
$\langle|\hat{\vec d}|\rangle = \alpha|\vec E|$ \quad \dots\ konst.\ úměrnosti v záv.\ vel.\ indukov.\ dipólu na E-poli & (D) \\[4pt]
energie dipólu v E-poli: \quad $-\frac12\alpha|\vec E|^2$ & (E)
\end{tabular}
\end{definice}

\begin{priklad}[Příklad: H(1s) v poli E]
\dots\ $H_0$ \dots\ Hamilt.\ atomu vodíku
\begin{align*}
  H_1 &= -\vec d\cdot\vec E = e\vec E\cdot\vec r \qquad + \text{volba } \vec E \text{ ve směru osy } z \\
  &= e|E|\,r\cos\theta
\end{align*}
0-tý řád: \dots\ $H_0\ket{nlm} = E_n\ket{nlm}$ \\
\quad \dots\ degenerované spektrum, ale pro začátek se budeme zajímat o H v zákl.\ stavu \dots\ nedegenerovaná 1s hladina $\ket{100}$

1.ř.: \quad $E_{100}^{(1)} = \mel{100}{H_1}{100} = 0$ \quad ze symetrie, neboť $|\psi|^2$ je invar.\ vůči $r\to -r$, \\
\hspace*{5.5cm} ale $H_1$ není invariantní při $\nearrow$ \\
\quad \dots\ platí dokonce ve $\forall$ stavech: \quad $\mel{nlm}{H_1}{nlm} = 0$

\begin{itemize}
  \item \underline{polarizovatelnost lze zjistit ze vzorce (E):}
  \[
    E_{100}^{(2)} = \sum_{n'lm}\frac{|\mel{n'lm}{H_1}{100}|^2}{\varepsilon_1 - \varepsilon_{n'}}
    \quad\Rightarrow\quad
    \alpha = 2e\sum_{n'lm}\frac{|\mel{n'lm}{r\cos\theta}{100}|^2}{-\varepsilon_1 + \varepsilon_{n'}}
  \]
  \pozn{znaménko $-$ ve vzorci (E) zakroužkováno a spojeno čarou se jmenovatelem}

  později uvidíme, že výběrová pravidla dovolí jen $\mel{n'10}{r\cos\theta}{100}$; ostatní jsou 0. \\
  (Zde teď je vidět, že $m = m'$ a v podstatě i $l = 1$), neboť $r\cos\theta\ket{100}$ je úměrné $Y_{10}(\theta,\varphi)$
  \[
    \text{tj.}\quad \alpha = 2e\sum_{n'}\frac{|\mel{n'10}{r\cos\theta}{100}|^2}{\varepsilon_1 - \varepsilon_{n'}}
  \]
  \dots\ \pozor{problém} \dots\ obsahuje $\forall$ vázané $+$ kontinuum \\
  \quad (na předn.\ napsat i vzorec s $\int\dd E$)

  \item \underline{polariz.\ z (D):}
  \[
    \ket{\psi_{100}} \doteq \ket{100} + \ket{\psi_{100}^{(1)}} \quad\dots\quad \braket{\psi_{100}}{\psi_{100}} = 1 + \mathcal{O}(E^2) \quad\text{díky } \braket{n}{\psi_n^{(1)}} = 0
  \]
  \[
    \langle d\rangle = \frac{\mel{\psi_{100}}{d}{\psi_{100}}}{\braket{\psi_{100}}{\psi_{100}}}
    = \mel{100}{d}{\psi_{100}^{(1)}} + \mel{\psi_{100}^{(1)}}{d}{100} + \mathcal{O}(E^2) \equiv \alpha|E|
  \]
  \underline{pozn:} vzorec pro $\ket{\psi_{100}^{(1)}}$ pomocí báze $\ket{nlm}$ vede ke stejnému výrazu ($\uparrow$ problém), \\
  tj.\ problémy s kontinuem, ale nikdo nás nenutí hledat $\ket{\psi_n^{(1)}}$ v bázi $\ket{n}$ \dots\ můžeme přímo řešit (P1)

  \item \underline{Přímé řešení (P1)} \\
  \dots\ hodí se obecně pro výpočet \ $E_n^{(2)} = \mel{n}{H_1}{\psi_n^{(1)}}$ \ v systémech obsahujících kontinuum.
\end{itemize}
\end{priklad}

%----------------------------------------------------------------------
\strana{PT-7}
$(P1)$: \quad $\big[H_0 - E_n^{(0)}\big]\ket{\psi_n^{(1)}} = \big(E_n^{(1)} - H_1\big)\ket{\psi_n^{(0)}}$ \\
v našem případě: \quad $E_{100}^{(1)} = 0$
\[
  \text{tj.}\quad \Big[-\frac{\hbar^2}{2m}\nabla^2 + V(r) - \varepsilon_1\Big]\psi^{(1)}(\vec r) = -e|E|\,r\cos\theta\;\underbrace{\psi^{(0)}(\vec r)}_{\text{nezávisí na } \theta,\varphi}
\]
levou stranu známe z řešení atomu vodíku. \quad Předp.: \quad $\psi^{(1)}(\vec r) = f(r)\,Y_{lm}(\theta,\varphi)$
\[
  \to\quad Y_{lm}(\theta,\varphi)\Big\{-\frac{\hbar^2}{2\mu}\frac{1}{r^2}\dv{r}r^2\dv{r}f(r) + \Big[\frac{\hbar^2(l+1)l}{2\mu r^2} + V(r) - \varepsilon_1\Big]f(r)\Big\}
\]
\dots\ na pravé str.\ je výraz úměrný $Y_{10}\sim\cos\theta$ \dots\ tj.\ $l = 1$, $m = 0$
\[
  \text{tj.}\quad -\frac{\hbar^2}{2\mu}\frac{1}{r^2}\dv{r}r^2\dv{r}f(r) + \Big[\frac{\hbar^2}{\mu r^2} + V(r) - \varepsilon_1\Big]f(r) = -e|E|\,r\,\psi^{(0)}(r)
\]
pro atom vodíku: \quad $V(r) = -\dfrac{e^2}{r}$; \quad $\varepsilon_1 = -\dfrac{e^2}{2a_0}$; \quad
$\psi^{(0)}(r) = (\pi a_0^3)^{-1/2}e^{-r/a_0}$; \quad $a_0 = \dfrac{\hbar^2}{\mu e^2}$

+ hledání řešení ve tvaru \quad $f(r) = p(r)\,e^{-r/a_0}$ \qquad $p(r)$ -- polynom \\
vede na \quad $p(r) = -(\pi a_0^3)^{-1/2}\Big(\dfrac{|E|}{e}\Big)\big(a_0r + \tfrac12 r^2\big)$
\quad (druhé řešení $\infty$ řada, která nevede na $L^2$ řeš.)

dosazení do
\begin{align*}
  E_{100}^{(2)} &= \mel{100}{H_1}{\psi_{100}^{(1)}} = -|E|^2(\pi a_0^3)^{-1}\int(\cos\theta)^2(a_0r^2 + r^3)\,e^{-2r/a_0}\,\dd^3r \\
  &= -\frac94|E|^2a_0^3 = -\frac12\alpha|E|^2
\end{align*}
\[
  \text{tj.}\quad \vzorec{\alpha = \frac92 a_0^3} \quad\dots\ \text{polarizovatelnost zákl.\ stavu H-atomu}
\]

\begin{priklad}[\textcircled{4} Lineární Starkův efekt v H-atomu]
stejný problém, ale $n = 2$ \dots\ $p$ stavy H-atomu:
\[
  \ket{\psi^{(0)}} = c_1\ket{200} + c_2\ket{211} + c_3\ket{210} + c_4\ket{21{-1}}
\]
$(c_1,c_2,c_3,c_4)$ \dots\ vl.\ stavy poruchy $H_1 = e|E|\,r\cos\theta$ v bázi $\uparrow$

$\mel{2l'm'}{H_1}{2lm}$ nenulové jen pro $m = m'$ \quad $\Rightarrow$ \quad matice
\[
  \begin{pmatrix} 0 & 0 & -v & 0\\ 0 & 0 & 0 & 0\\ -v & 0 & 0 & 0\\ 0 & 0 & 0 & 0\end{pmatrix}
  \qquad\text{kde}\quad \mel{210}{H_1}{200} = -3e|E|a_0 \equiv -v
\]
\begin{center}
\renewcommand{\arraystretch}{1.4}
\begin{tabular}{cc}
vl.\ v. & vl.\ č. \\ \hline
$\tfrac{1}{\sqrt2}(1,0,-1,0)$ & $v$ \\
$\tfrac{1}{\sqrt2}(1,0,1,0)$ & $-v$ \\
$(0,1,0,0)$ & $0$ \\
$(0,0,0,1)$ & $0$
\end{tabular}
\qquad
\begin{minipage}{0.4\textwidth}
tj.\ \kw{přetrvává degenerace} pro $\pm m = \pm1$, \\
\kw{snímá se} pro $l = 0,1$; $m = 0$
\end{minipage}
\end{center}
\end{priklad}

%----------------------------------------------------------------------
\strana{PT-8}
\begin{priklad}[\textcircled{5} Van der Waalsova interakce]
\begin{center}
\begin{tikzpicture}[>=Stealth]
  \draw[->] (-0.5,0) -- (5.5,0) node[below right] {$z$};
  \node[draw,circle,inner sep=0.5pt] (p1) at (0,0) {\scriptsize$+$};
  \node[draw,circle,inner sep=0.5pt] (p2) at (3.8,0) {\scriptsize$+$};
  \node[below] at (1.9,0) {$r$};
  % elektron 1
  \fill (1.2,0.8) circle (1.5pt) node[above right=-1pt] {\scriptsize$-$};
  \draw[->] (p1) -- (1.15,0.77) node[midway,above left=-2pt] {$\vec r_1$};
  % elektron 2
  \fill (3.0,0.9) circle (1.5pt) node[above=1pt] {\scriptsize$-$};
  \draw[->] (p2) -- (3.03,0.86) node[midway,right=1pt] {$\vec r_2$};
\end{tikzpicture}
\end{center}
\[
  H = H_0 + V \qquad H_0 = -\frac{\hbar^2}{2m}\big(\nabla_1^2 + \nabla_2^2\big) - \frac{e^2}{r_1} - \frac{e^2}{r_2}
\]
\[
  V = \frac{e^2}{r} + \frac{e^2}{|\vec r + \vec r_2 - \vec r_1|} - \frac{e^2}{|\vec r + \vec r_2|} - \frac{e^2}{|\vec r - \vec r_1|}
\]
\quad $\hookrightarrow$ rozvinout \quad $(1 + x)^\alpha = 1 + \alpha x + \frac12\alpha(\alpha - 1)x^2 + \dots$
\[
  V = \frac{e^2}{r^3}\big(x_1x_2 + y_1y_2 - 2z_1z_2\big) + \mathcal{O}\Big(\frac{1}{r^4}\Big) \qquad r\to\infty
\]

1.\ řád \quad $\mel{\psi_0}{V}{\psi_0}$ \ ?

$\ket{\psi_0} = \ket{nlm}_1\ket{nlm}_2$ \quad \dots\ zákl.\ stav $n = 1$, $l = 0$, $m = 0$ \\
\dots\ první řád vymizí \dots\ $\sim \underbrace{\mel{nlm}{x_1}{nlm}_1}_{0\text{ v zákl.\ stavu}}\underbrace{\mel{nlm}{x_2}{nlm}_2}_{0}$
\quad ($x\sim Y_{1m}$)

korekce 2.\ řádu:
\[
  \vzorec{E^{(2)}(r) = \frac{e^4}{r^6}\sum_{k\neq 0}\frac{|\mel{\psi_k}{x_1x_2 + y_1y_2 - 2z_1z_2}{\psi_0}|^2}{E_0^{(0)} - E_k^{(0)}} = -\frac{C}{r^6} < 0}
\]
\end{priklad}


\newpage
%======================================================================
% Poznamky_2_variacni_metody.pdf
%======================================================================
\section{Ritzův variační princip}
\strana{VP-1}

Střední hodnota energie ve stavu $\ket{\psi}$ je
\[
  \vzorec{E[\psi] \equiv \langle H\rangle_\psi = \frac{\mel{\psi}{H}{\psi}}{\braket{\psi}{\psi}}} \qquad\text{(RVP)}
\]
Podívejme se na změnu $E$ při záměně $\ket{\psi}\to\ket{\psi} + \ket{\delta\psi}$
\begin{align*}
  \delta E \equiv E[\psi + \delta\psi] - E[\psi]
  &= \frac{\mel{\psi}{H}{\psi} + \mel{\delta\psi}{H}{\psi} + \mel{\psi}{H}{\delta\psi}}{\braket{\psi}{\psi} + \braket{\delta\psi}{\psi} + \braket{\psi}{\delta\psi}} - E + O(\|\delta\psi\|^2) \\
  &= \frac{\mel{\delta\psi}{H - E}{\psi} + \mel{\psi}{H - E}{\delta\psi}}{\braket{\psi}{\psi} + \braket{\delta\psi}{\psi} + \braket{\psi}{\delta\psi}} + O(\|\delta\psi\|^2)
\end{align*}
Stacionární hodnoty funkcionálu (RVP) $\Leftrightarrow$
\[
  \mel{\delta\psi}{H - E}{\psi} + \mel{\psi}{H - E}{\delta\psi} = 0 \qquad \forall\ket{\delta\psi}
\]
splněno $\forall\ket{\delta\psi}$, tj.\ i pro $i\ket{\delta\psi}$:
\[
  \times(\pm i)\,\big| \qquad -i\mel{\delta\psi}{H - E}{\psi} + \mel{\psi}{H - E}{\delta\psi}\,i = 0
\]
\[
  \Rightarrow\qquad \mel{\delta\psi}{H - E}{\psi} = \mel{\psi}{H - E}{\delta\psi} = 0 \qquad \forall\ket{\delta\psi}
\]

\begin{veta}[Tvrzení 1]
Funkcionál (RVP) nabývá stacionárních hodnot pro stacionární stavy $H\ket{\psi_n} = E_n\ket{\psi_n}$. \quad \dots\ $E[\ket{\psi_n}]$ \dots\ stac.\ hodnota.
\end{veta}

\underline{Důsledek:} Pokud nějak odhadneme $\ket{\tilde\psi_n} = \ket{\psi_n} + \ket{\delta\psi}$, tak \kw{chyba ve výpočtu energie je druhého řádu} v $\|\delta\psi\|$.

\underline{Poznámka:} Striktně řečeno výroky
\begin{align*}
  \mel{\delta\psi}{H - E}{\psi} &= 0 \quad \forall\ket{\delta\psi} && \text{(V1)} \\
  \text{a}\quad (H - E)\ket{\psi} &= 0 && \text{(V2)}
\end{align*}
nejsou ekvivalentní, respektive $\mel{\delta\psi}{H - E}{\psi}$ lze vyčíslit pro větší třídu funkcí. (V1) se nazývá \kw{\uv{SLABOU FORMULACÍ}} úlohy (v tomto případě S.R.). Kinetický člen (pro částici v potenc.\ poli):
\begin{align*}
  \mel{\delta\psi}{H}{\psi} &= -\frac{\hbar^2}{2m}\int\delta\psi(\vec r)^*\,\Delta\psi(\vec r)\,\dd^3\vec r + \int\delta\psi(\vec r)^*V(\vec r)\psi(\vec r)\,\dd^3\vec r \\
  &\overset{\text{per partes}}{=} \frac{\hbar^2}{2m}\int(\nabla\delta\psi)^*\cdot\big(\nabla\psi(\vec r)\big)\,\dd^3\vec r + \int\delta\psi(\vec r)^*V(\vec r)\psi(\vec r)\,\dd^3\vec r
\end{align*}
$\uparrow$ tato forma je matematicky výhodná \dots\ nemusí $\exists$ druhé derivace $\psi$ \\
\dots\ Sobolevovy prostory \dots

%----------------------------------------------------------------------
\strana{VP-2}
\begin{veta}[Tvrzení 2]
Nechť $E_0$ je energie základního stavu pro $\hat H$. Potom
\[
  \vzorec{E_0 \leq E[\psi] \equiv \frac{\mel{\psi}{\hat H}{\psi}}{\braket{\psi}{\psi}}} \qquad \forall\ket{\psi}
\]
\end{veta}

\underline{Důkaz:}
\[
  E[\psi] = \frac{\sum_n E_n|\braket{n}{\psi}|^2}{\sum_n|\braket{n}{\psi}|^2} \geq \frac{E_0\sum_n|\braket{n}{\psi}|^2}{\sum_n|\braket{n}{\psi}|^2} = E_0
\]
+ je vidět \dots\ rovnost $\Leftrightarrow$ $\ket{\psi}$ patří do vl.\ podprostoru $\hat H$ příslušného $E_0$.

\underline{Použití:} Určitá třída \kw{testovacích funkcí} $\ket{\psi(\alpha_1,\alpha_2,\dots,\alpha_n)}$ \\
\quad $E[\psi(\alpha_1\dots\alpha_n)]$ je funkce parametrů $\alpha_1\dots\alpha_n$ + hledáme min. \\
\quad tj.\ $\dfrac{\partial E}{\partial\alpha_i} = 0\ \ \forall\alpha_i$ \quad \dots\ minimum v $\vec\alpha = \vec\alpha^{(0)} \equiv (\alpha_1^{(0)}\dots\alpha_n^{(0)})$ \\
\quad pak $\ket{\psi(\alpha^{(0)})}$ je aproximace vln.\ fce zákl.\ stavu a $E[\psi(\alpha^{(0)})]$ je \kw{horní odhad energie}.

\underline{Pozn:} lze zkonstruovat i dolní odhad (viz Ballentine)

\underline{Speciální případ:} třída testovacích funkcí tvoří podprostor:
\[
  \ket{\psi} = \alpha_1\ket{\phi_1} + \alpha_2\ket{\phi_2} + \dots + \alpha_n\ket{\phi_n}
\]
\[
  \text{(V1)}\ \Rightarrow\ \sum_{ij}\delta\alpha_i^*\mel{\phi_i}{H}{\phi_j}\alpha_j - E\sum_{ij}\delta\alpha_i^*\braket{\phi_i}{\phi_j}\alpha_j = 0 \qquad \forall\delta\alpha_i
\]
\[
  \text{tj.}\qquad \vzorec{\sum_j\Big[\mel{\phi_i}{H}{\phi_j} - E\braket{\phi_i}{\phi_j}\Big]\alpha_j = 0} \qquad \forall i
\]
zobecněná rovnice pro vl.\ č.\ matice \quad \pozn{vl.\ č.\ $\det\{\mathcal{H} - E\mathcal{S}\} = 0$}
\begin{itemize}
  \item speciálně pro ON-bázi $\braket{\phi_i}{\phi_j} = \delta_{ij}$ \quad $\to\ \sum_j\mathcal{H}_{ij}\alpha_j = E\alpha_i$ \\
        $\to$ vlastní čísla a vl.\ vektory matice $\mathcal{H}_{ij} \equiv \mel{\phi_i}{\hat H}{\phi_j}$
  \item v obecném případě je třeba nejdříve diagonalizovat matici $\mathcal{S}_{ij} = \braket{\phi_i}{\phi_j} = U^\dagger\Lambda U$ \\
        \dots\ \kw{Löwdinova symetrická ortogonalizace}
\end{itemize}
\pozn{$H\ket{\psi} = E\,S\ket{\psi}$ \ $\to$ \ $\underbrace{\tfrac{1}{\sqrt S}H\tfrac{1}{\sqrt S}}_{\tilde H}\cdot\underbrace{\sqrt S\ket{\psi}}_{\ket{\tilde\phi}} = E\cdot\sqrt S\ket{\psi}$, \quad $\tilde H\ket{\tilde\phi} = E\ket{\tilde\phi}$}

%----------------------------------------------------------------------
\strana{VP-3}
\begin{veta}[Tvrzení 3 (Hylleraas--Undheim theorem)]
Nechť $E_0\leq E_1\leq E_2\leq\dots$ jsou přesné vl.\ č.\ $\hat H$ (každé tolikrát, kolik je jeho násobnost) \\
a nechť $\varepsilon_0\leq\varepsilon_1\leq\varepsilon_2\leq\dots$ jsou vl.\ č.\ problému $\det\{\mathcal{H} - \varepsilon\mathcal{S}\} = 0$.

Potom: \qquad $\vzorec{E_i \leq \varepsilon_i} \qquad \forall i = 1,\dots N$ \quad \dots\ (dimenze podprostoru test.\ fcí)
\end{veta}

\underline{Náznak důkazu:}

\begin{minipage}{0.55\textwidth}
Postupně přidáváme řádky a sloupce matice:
\end{minipage}
\begin{minipage}{0.4\textwidth}
\begin{tikzpicture}[scale=0.8]
  \draw (0,0) rectangle (2.2,1.7);
  \draw (0,0.4) -- (2.2,0.4);
  \draw (1.8,0) -- (1.8,1.7);
  \node at (0.9,1.05) {$M'$};
  \node[right] at (2.4,0.85) {$= M$};
\end{tikzpicture}
\end{minipage}

\begin{minipage}{0.3\textwidth}
dá se DK, že
\end{minipage}
\begin{minipage}{0.4\textwidth}
\begin{tikzpicture}[x=1cm,y=0.45cm,>=Stealth]
  % levý sloupec: lambda' (M'), pravý sloupec: lambda (M)
  \foreach \y in {4.5,3,1.5,0} \draw (0,\y) -- (0.5,\y);
  \foreach \y in {5.2,3.8,2.3,0.8,-0.6} \draw (2,\y) -- (2.5,\y);
  \node[left] at (0,4.5) {$\lambda'_{N-1}$};
  \node[left] at (0,0) {$\lambda'_1$};
  \node[left] at (0,2.25) {$\vdots$};
  \node[right] at (2.5,5.2) {$\lambda_N$};
  \node[right] at (2.5,-0.6) {$\lambda_1$};
  \node[right] at (2.5,2.3) {$\vdots$};
  \draw[->] (0.5,3) -- (2.6,2.0);
  \draw[->] (0.5,1.5) -- (2.6,0.5);
  \draw[->] (0.5,0) -- (2.6,-0.9);
\end{tikzpicture}
\end{minipage}
\begin{minipage}{0.25\textwidth}
\dots\ navzájem se dělí
\end{minipage}

\noindent\rule{3cm}{0.4pt}

DK, že se navzájem dělí \dots\ najdeme bázi, v níž $M'$ diagonální
\[
  \to\quad M = \left(\begin{array}{ccc|c}
    \lambda'_1 & & & M_{1N} \\
    & \ddots & & \vdots \\
    & & \lambda'_{N-1} & \\ \hline
    M_{N1} & \cdots & & M_{NN}
  \end{array}\right)
  \quad\dots\quad
  f(\lambda) \equiv \det(M - \lambda\hat I)
\]
\[
  f(\lambda) = (M_{NN} - \lambda)\prod_{n=1}^{N-1}(\lambda'_n - \lambda) - \sum_{m=1}^{N-1}|M_{mN}|^2\prod_{n\neq m}(\lambda'_n - \lambda)
\]
$f(\lambda)$ je char.\ polynom. Vyčíslený v $\lambda = \lambda'_k$:
\[
  f(\lambda'_k) = -|M_{kN}|^2\prod_{n\neq k}(\lambda'_n - \lambda'_k)
  \quad\to\quad \text{střídá znaménko pro } k = 1,2,\dots
\]
\hfill $+\ f(\lambda\to-\infty)\to+\infty$, \quad $f(\lambda\to+\infty)\to(-1)^N\infty$

\noindent\rule{3cm}{0.4pt}

\textbf{Poznámky:}
\begin{itemize}
  \item Základní způsob použití $\to$ základní stav jako minimum (i nelineární test.\ prostor)
  \item Excitované stavy, ale základní pro danou symetrii \dots\ jako $\uparrow$ \\
  \begin{minipage}[t]{0.5\textwidth}
    PŘ: sudé a liché fce: \quad
    $H = \left(\begin{array}{c|c} H_S & 0 \\ \hline 0 & H_L\end{array}\right)$
  \end{minipage}
  \begin{minipage}[t]{0.45\textwidth}\raggedright
    sfér.\ sym.\ potenciál \\
    $H = \begin{pmatrix} \boxed{H_0} & & 0 \\ & \boxed{H_1} & \\ 0 & & \ddots\ \boxed{H_l}\end{pmatrix}$
  \end{minipage} \\
  + složitější \dots\ sym.\ molekul, krystalů \dots\ teorie grup.
  \item Netriviální excitace $\to$ Hylleraas--Undheim \dots\ max; + stac.\ stav.
  \item souvislost s poruchovou teorií: \quad $H = H_0 + H_1$ \\
        \dots\ pro hledání $E_0$ použijeme test.\ prostor $\ket{\phi_k} = \ket{n=0,k}$ \\
        $\to$ korekce 1.\ ř.\ $=$ vl.\ č.\ $\mel{0,m}{H_1}{0,k}$ \dots\ totožné s variačním principem. \\
        (konzist.\ s tím, že korekce 2.\ ř.\ je vždy $< 0$) \quad $\Rightarrow$ horní odhad!
\end{itemize}

%----------------------------------------------------------------------
\strana{VP-4}
\subsection*{Dolní odhad?}
nechť $\ket{\psi}$ je odhad vl.\ fce pro $E_n$ (normovaná) \\
\dots\ def \quad $\varepsilon \equiv \mel{\psi}{H}{\psi}$ \quad a zbytek \quad $\ket{R} \equiv (H - \varepsilon)\ket{\psi}$ \\
pak
\[
  \braket{R}{R} = \mel{\psi}{(H - \varepsilon)^2}{\psi} = \sum_{n'}|\braket{n'}{\psi}|^2(E_{n'} - \varepsilon)^2
  \underset{\uparrow}{\geq} (E_n - \varepsilon)^2\underbrace{\sum_{n'}|\braket{n'}{\psi}|^2}_{=1}
\]
odhad, ale musíme něco vědět! (nebo doufat) $\to$ \pozor{za předp.\ $E_n$ je nejbližší vl.\ č.\ k hodnotě $\varepsilon$!}

pak \quad $\to$ \quad $\vzorec{\varepsilon - \Delta \leq E_n \leq \varepsilon + \Delta}$ \quad kde \quad $\Delta \equiv \sqrt{\braket{R}{R}}$

lepší odhady: je třeba vědět více o spektru $\to$ viz Ballentine

\begin{priklad}[PŘÍKLADY: Atom vodíku (Ballentine)]
\[
  H = \underbrace{\frac{p^2}{2\mu}}_{K} \underbrace{-\ \frac{e^2}{r}}_{+P}; \qquad
  \psi(r) = e^{-r/a} \qquad N = \braket{\psi}{\psi} = \pi a^3 = \int|\psi|^2\dd^3r
\]
\pozn{přesné $\psi = \frac{1}{\sqrt{\pi a_0^3}}e^{-r/a}$ $\leftarrow$ normované}
\begin{align*}
  \to\quad K &= \frac{\hbar^2\pi a}{2\mu}
     = \frac{\hbar^2}{2\mu}\cdot\frac{4\pi}{a^2}\int_0^\infty e^{-2r/a}r^2\,\dd r
     = \frac{\hbar^2}{2\mu}\int\Big|\frac{\partial\psi}{\partial r}\Big|^2\dd^3r \\
  P &= -\pi e^2a^2 = -e^2\,4\pi\int_0^\infty e^{-2r/a}r\,\dd r = -\int|\psi|^2\frac{e^2}{r}\,\dd^3r
\end{align*}
tj.\ \quad $\langle H\rangle = \dfrac{K + P}{N} = \dfrac{\hbar^2}{2\mu a^2} - \dfrac{e^2}{a}$
\quad $\to$ \quad min.\ pro \quad $\vzorec{a = \frac{\hbar^2}{\mu e^2}}$

\underline{Srovn:}
\[
  \begin{array}{c|ccc}
    \psi = & Ce^{-r/a} & \dfrac{C}{r^2 + a^2} & Cre^{-r/a} \\[6pt] \hline
    \rule{0pt}{14pt}\dfrac{\langle H\rangle}{E_0} = & -1 & -0.81 & -0.75 \\[6pt]
    1 - |\braket{\psi}{100}|^2 & 0 & 0.21 & 0.05
  \end{array}
\]
\dots\ tj.\ \pozor{lepší vln.\ fce neznamená lepší energii!}
\end{priklad}

\newpage
%======================================================================
% Poznamky_4_skladani_ang_mom.pdf
%======================================================================
\section{Skládání momentu hybnosti}
\strana{AA-1}

\textbf{Opakování z minulého semestru:}

\begin{definice}[Moment hybnosti]
\kw{Moment hybnosti} \dots\ libovolný operátor (vektorový) \quad
$[\hat J_i,\hat J_j] = i\hbar\,\varepsilon_{ijk}\hat J_k$, \quad $i = x,y,z$
\end{definice}

$\Rightarrow$ spektrum $J^2, J_z$ \dots\ spol.\ vl.\ vektory $\ket{j,m,\alpha}$

\begin{minipage}{0.6\textwidth}
\begin{align*}
  J^2\ket{j,m,\alpha} &= j(j+1)\hbar^2\ket{j,m,\alpha} \\
  J_z\ket{j,m,\alpha} &= m\hbar\ket{j,m,\alpha}
\end{align*}
\end{minipage}
\begin{minipage}{0.38\textwidth}
$j = 0,\frac12,1,\frac32,\dots$ (nebo podmnožina) \\
$m = -j,-j+1,\dots,j$ \\
($\forall j$; $2j+1$ hodnot + různá degenerace)
\end{minipage}

přitom
\[
  \vzorec{\begin{aligned}
    J_+\ket{j,m} &= \sqrt{(j-m)(j+m+1)}\,\hbar\ket{j,m+1} \\
    J_-\ket{j,m} &= \sqrt{(j+m)(j-m+1)}\,\hbar\ket{j,m-1}
  \end{aligned}}
\]
\quad kde \quad $J_\pm \equiv J_x \pm iJ_y$

\begin{priklad}[příklady]
\begin{itemize}[nosep]
  \item spin $\frac12$: \quad $\hat{\vec S} = \frac{\hbar}{2}(\sigma_x,\sigma_y,\sigma_z)$ \quad
        $\leftarrow \left(\begin{pmatrix}0&1\\1&0\end{pmatrix},\begin{pmatrix}0&-i\\ i&0\end{pmatrix},\begin{pmatrix}1&0\\0&-1\end{pmatrix}\right)$
  \item orbit.\ moment: \quad
        $\hat{\vec L} = \hat{\vec x}\times\hat{\vec p} = -i\hbar\Big(y\pdv{z} - z\pdv{y},\ z\pdv{x} - x\pdv{z},\ x\pdv{y} - y\pdv{x}\Big)$
        \pozn{v originále složky zapsány v pořadí $z$, $y$, $x$ se šipkou na prohození}
\end{itemize}
\end{priklad}

\kw{Skládání systémů:} \quad $\mathcal{H} = \mathcal{H}_1\otimes\mathcal{H}_2$ \quad
\dots\ $\ket{\psi} = \ket{\phi_1}\otimes\ket{\phi_2}$ (nebo lin.\ kombinace) \\
\hspace*{6cm} báze $\ket{\phi_{mn}} = \ket{\phi_m}_1\otimes\ket{\phi_n}_2$ \\
\quad operátory \quad $\hat A_1 \to \hat A_1\otimes\hat I$ \qquad
$\hat A_1\otimes\hat I\big(\ket{\phi_1}\otimes\ket{\phi_2}\big) = \hat A_1\ket{\phi_1}\otimes\ket{\phi_2}$

částice se spinem $\frac12$: \dots\ báze $\ket{\vec Q,s} \equiv \ket{x,y,z}\otimes\ket{\pm}$ \quad \dots\ $\mathcal{H} = \mathcal{L}^2(\mathbb{R}^3)\otimes\mathbb{R}^2$

\quad vln.\ fce má teď komponenty \quad
$\psi_\pm(x,y,z) \equiv \underbrace{\bra{x,y,z}\otimes\bra{\pm}}_{\in\mathcal{H}^*}\underbrace{\ket{\psi}}_{\mathcal{H}}$

\quad tj.\ $\ket{\psi} \equiv \begin{pmatrix}\psi_+(\vec Q)\\ \psi_-(\vec Q)\end{pmatrix}$ \quad \dots\ dvoukomponentová vln.\ fce

celkový moment hybnosti \quad $\hat{\vec J} = \hat{\vec L} + \hat{\vec S} \equiv \hat{\vec L}\otimes\hat I + \hat I\otimes\hat{\vec S}$

tj.\ např.\ $J_z\ket{\psi}$ \dots\ v souřad.\ reprez.\
\[
  = -i\hbar\begin{pmatrix} x\partial_y\psi_+ - y\partial_x\psi_+ \\ x\partial_y\psi_- - y\partial_x\psi_-\end{pmatrix}
  + \frac{i\hbar}{2}\underbrace{\sigma_z\begin{pmatrix}\psi_+\\ \psi_-\end{pmatrix}}_{\begin{pmatrix}1&0\\0&-1\end{pmatrix}\begin{pmatrix}\psi_+\\ \psi_-\end{pmatrix}}
\]

v případě, že $H$ je sfér.\ sym.\ nezáv.\ na spinu \\
\dots\ místo $Q$-reprezentace \quad
$\left.\begin{aligned} &\ket{n\,l\,m}\ket{+} \\ &\ket{n\,l\,m}\ket{-}\end{aligned}\right\}$
vl.\ v.\ $H, L^2, L_z, S_z, S^2$ \ ($S^2 \hookrightarrow$ vždy $\frac34\hbar^2$) \\
\hspace*{5.8cm} $\to \ket{n,l,m,s_z}$

$\to$ ukážeme, že \sout{dob} ÚSKO je též \quad $J^2, J_z, L^2, S^2$ \quad \dots\ vl.\ v.\ později

\begin{priklad}[příklad 2: dvě částice se spinem $\frac12$]
$\mathcal{H} = \mathcal{H}_1\otimes\mathcal{H}_2 = \mathbb{R}^2\otimes\mathbb{R}^2$

celk.\ spin $\to$ \quad $\hat{\vec S} = \hat{\vec S}_1 + \hat{\vec S}_2 \equiv \hat{\vec S}_1\otimes\hat I + \hat I\otimes\hat{\vec S}_2$
\quad \dots\ ztotožňujeme $S_1 \equiv S_1\otimes I$ \dots\ v značení nebudeme rozlišovat.

přitom \quad $[S_{1i},S_{2j}] = 0$
\end{priklad}

%----------------------------------------------------------------------
\strana{AA-2}
\begin{priklad}[příklad 2 -- pokračování]
Přitom \quad $[S_{1i},S_{1j}] = i\hbar\,\varepsilon_{ijk}S_{1k}$ \quad \dots\ podobně $\hat{\vec S}_2$
\[
  \Rightarrow\quad [S_i,S_j] = [S_{1i} + S_{2i},\,S_{1j} + S_{2j}]
  = i\hbar\,\varepsilon_{ijk}\big(\underset{\substack{\uparrow\\ [1,1]}}{S_{1k}} + \underset{\substack{\uparrow\\ [2,2]}}{S_{2k}}\big)
  = i\hbar\,\varepsilon_{ijk}S_k
  \qquad [1,2] = [2,1] = 0
\]
$\Rightarrow$ spektrum $S$ musí být spektrum oper.\ mom.\ hybnosti.

ÚSKO: $S_1^2, S_{1z}, S_2^2, S_{2z}$ \dots\ přiroz.\ báze v $\mathcal{H}_1\otimes\mathcal{H}_2$: \quad
$\ket{s_z}_1\otimes\ket{s_z}_2 \equiv \ket{s_{z1},s_{z2}}$ \pozn{často vynech.} \\
\quad ($s_z$ resp.\ $m_1$, $m_2$ $= \pm\frac12$) \\
\quad $\to$ dimenze prost.\ 4: \quad báze $\big\{\ket{+}\ket{+},\ \ket{+}\ket{-},\ \ket{-}\ket{+},\ \ket{-}\ket{-}\big\}$

Jiná možnost \dots\ ÚSKO: $S^2, S_z$, \sout{\dots} \quad \dots\ báze $\ket{S,m}$ \quad -- není jasné, jestli další index?

potřebujeme 4 báz.\ vektory:
\begin{enumerate}[label=\textcircled{\arabic*}]
  \item $\ket{S,m} \equiv \ket{1,1} = \ket{++}$ \quad \dots\ $s = 1$, $m = 1$ \\
        $S_z = (S_{1z}\ket{+})\ket{+} + \ket{+}(S_{2z}\ket{+}) = \big(\frac{\hbar}{2} + \frac{\hbar}{2}\big)\ket{+}\ket{+} = \hbar\ket{++}$
        \quad $\Rightarrow$ $s \geq 1$ \ (\,$\uparrow$)
  \item $\ket{s,m} = \ket{1\,0}$: \quad $\sqrt{(1+1)(1-1+1)}\,\ket{1\,0} = J_-\ket{1\,1}$
        \[
          \hat S_-\ket{1\,1} = \big(\hat S_{1-} + \hat S_{2-}\big)\ket{++}
          = \sqrt{\big(\tfrac12+\tfrac12\big)\big(\tfrac12-\tfrac12+1\big)}\,\big(\ket{-+} + \ket{+-}\big)
        \]
        \quad $\nwarrow$ (v případě spinu $\frac12$ vždy 1 nebo 0)
        \[ \Rightarrow\quad \ket{1\,0} = \tfrac{1}{\sqrt2}\big(\ket{-+} + \ket{+-}\big) \]
  \item $\ket{s\,m} = \ket{1\,{-1}} = \ket{--}$ \quad \dots\ opět je vidět, že $m = -1$, tj.\ $S \geq 1$ \\
        nebo \quad $J_-\ket{1\,0} = \sqrt{(1+0)(1-0+1)}\,\ket{1\,{-1}} = \big(\ket{--} + \ket{--}\big)\cdot\frac{1}{\sqrt2}$
  \item zatím máme jen 3.\ dimenze v $\mathcal{L}\{\ket{++},\ket{--},\ket{+-},\ket{-+}\}$ \\
        (předchozí vektory jsou očividně kolmé \dots\ např.\
        $\braket{++}{-+} = \underbrace{\braket{+}{-}}_{0}\underbrace{\braket{+}{+}}_{1}$) \\
        další LNZ (nejlépe kolmý) vektor \dots\ ? $\to$ \quad
        $\frac{1}{\sqrt2}\big(\ket{+-} - \ket{-+}\big) = \ket{s,m} \equiv \ket{0,0}$
\end{enumerate}

ověření hodnoty $S^2$: relace
\[
  \vec S^{\,2} = (\vec S_1 + \vec S_2)^2 = \vec S_1^{\,2} + 2\vec S_1\cdot\vec S_2 + \vec S_2^{\,2}
  = S_1^2 + S_2^2 + 2S_{1z}S_{2z} + 2(S_{1x}S_{2x} + S_{1y}S_{2y})
\]
přitom \quad $S_x = \frac12(S_+ + S_-)$, \qquad $S_y = \frac{1}{2i}(S_+ - S_-)$
\begin{align*}
  \Rightarrow\ 2(S_{1x}S_{2x} + S_{1y}S_{2y})
  &= 2\cdot\Big(\tfrac14(S_{1+} + S_{1-})(S_{2+} + S_{2-}) - \tfrac14(S_{1+} - S_{1-})(S_{2+} - S_{2-})\Big) \\
  &= S_{1+}S_{2-} + S_{1-}S_{2+}
\end{align*}
\[
  \to\quad \vzorec{\vec S^{\,2} = S_1^2 + S_2^2 + 2S_{1z}S_{2z} + (S_{1+}S_{2-} + S_{1-}S_{2+})}
\]
\begin{align*}
  \Rightarrow\ \vec S^{\,2}\ket{++} &= \Big[\hbar\big(\tfrac34 + \tfrac34 + 2\cdot\tfrac14\big) + 0\Big]\ket{++} = 2\hbar\ket{++} = 1(1+1)\hbar\ket{++} &&\dots\ s = 1 \\
  \vec S^{\,2}\big(\ket{+-} + \ket{-+}\big) &= \Big\{\hbar\big(\tfrac34 + \tfrac34 - 2\tfrac14\big) + \hbar\Big\}\big(\ket{+-} + \ket{-+}\big) = 2\hbar\big(\ket{+-} + \ket{-+}\big) &&\dots\ s = 1 \\
  \vec S^{\,2}\big(\ket{+-} - \ket{-+}\big) &= \Big\{\hbar\big(\tfrac34 + \tfrac34 - 2\tfrac14\big) - \hbar\Big\}\big(\ket{+-} - \ket{-+}\big) = 0 &&\dots\ s = 0
\end{align*}
\end{priklad}

%----------------------------------------------------------------------
\strana{AA-3}
\subsection*{ZOBECNĚNÍ: Q-teorie skládání mom.\ hyb.\ \ {\normalfont\color{defbarva}(dle Sakuraie)}}

\begin{enumerate}[label=\textcircled{\arabic*}]
  \item $[\hat J_{1i},\hat J_{1j}] = i\hbar\,\varepsilon_{ijk}\hat J_{1k}$ \quad \dots\ podob.\ $\hat{\vec J}_2$ \\
        $[\hat J_{1i},\hat J_{2j}] = 0 \quad \forall i,j$
\end{enumerate}

\textcircled{4}\ $\left\{\ \begin{minipage}{0.88\textwidth}
\kw{možné ÚSKO:} \quad $\vec J_1^{\,2}, J_{1z}, \vec J_2^{\,2}, J_{2z}$ \quad \dots\
$\ket{j_1j_2;m_1m_2} \equiv \ket{j_1m_1}\otimes\ket{j_2m_2}$

přitom \quad $\vec J_1^{\,2}\ket{j_1j_2;m_1m_2} = \hbar^2 j_1(j_1+1)\ket{j_1j_2m_1m_2}$ \\
\phantom{přitom} \quad $J_{1z}\ket{j_1j_2;m_1m_2} = \hbar m_1\ket{j_1j_2m_1m_2}$ \hfill + podob.\ $\vec J_2^{\,2}, J_{2z}$

\kw{alternativní ÚSKO:} \quad $\vec J^{\,2}, \vec J_1^{\,2}, \vec J_2^{\,2}, J_z$ \quad
\dots\ důvod \dots\ $\vec J$ je gener.\ rotací \\
\hspace*{3cm} \dots\ často $[H,\vec J] = 0$, ale $[H,\vec J_1] \neq 0$, \ $[H,\vec J_2] \neq 0$
\end{minipage}\right.$

\begin{enumerate}[label=\textcircled{\arabic*},start=2]
  \item opět platí: \quad $\vec J^{\,2} = \vec J_1^{\,2} + \vec J_2^{\,2} + 2J_{1z}J_{2z} + J_{1+}J_{2-} + J_{1-}J_{2+}$
  \item $\Rightarrow\ [J^2, J_1^2] = 0$ \quad (neboť $J_1^2$ komut.\ se $\forall$ 2-oper.\ a $[J_1^2, J_{1\pm}] = 0$) \\
        \phantom{$\Rightarrow$} $[J_1^2, J_z]$ jednoduchý $= [J_1^2, J_{1z} + J_{2z}] = [J_1^2, J_{1z}] = 0$
\end{enumerate}

tj.\ komutují \dots\ značení $\ket{j_1j_2;jm}$:

\begin{minipage}{0.66\textwidth}
\begin{align*}
  J_1^2\ket{j_1j_2;j,m} &= \hbar^2 j_1(j_1+1)\ket{j_1j_2;jm} \quad\text{+ podob.\ } J_2^2 \\
  J^2\ket{j_1j_2;j,m} &= \hbar^2 j(j+1)\ket{j_1j_2;jm} \\
  J_z\ket{j_1j_2;j,m} &= \hbar m\ket{j_1j_2;jm}
\end{align*}
\end{minipage}
\begin{minipage}{0.32\textwidth}
{\color{defbarva}\small tohle jediné nenulové, tj.\ nepřevede na netriv.\ transformaci}

\begin{tikzpicture}[defbarva,scale=0.85,every node/.style={font=\small}]
  \draw (0,0.5) ellipse (0.62 and 0.3); \node at (0,0.5) {$J_1^2, J_2^2$};
  \draw (0,-0.2) ellipse (0.62 and 0.3); \node at (0,-0.2) {$J_1^2 J_2^2$};
  \draw (0.75,0.85) rectangle (2.4,-0.55);
  \draw (0.75,0.15) -- (2.4,0.15);
  \draw (1.5,0.85) -- (1.5,-0.55);
  \node at (1.12,0.5) {$J_{1z}$}; \node at (1.95,0.5) {$J_{2z}$};
  \node at (1.12,-0.2) {$J^2$}; \node at (1.95,-0.2) {$J_z$};
  \draw[<->,thick] (1.5,0.6) -- (1.5,-0.25);
\end{tikzpicture}
\end{minipage}

Pozor, ačkoli $[J^2, J_z] = 0$, je \quad $[J^2, J_{1z}] \neq 0$ \quad a \quad $[J^2, J_{2z}] \neq 0$ \\
\quad tj.\ $\ket{j_1j_2;jm}$ nejsou totožné s $\ket{j_1j_2;m_1m_2}$, ale
\[
  \vzorec{\ket{j_1j_2;jm} = \sum_{m_1m_2}\ket{j_1j_2;m_1m_2}\underbrace{\braket{j_1j_2;m_1m_2}{j_1j_2;jm}}_{\text{Clebsch--Gordan coefficients}}}
\]
\quad (obecně \dots\ báze v invar.\ podprostorech kart.\ součinu grup.)

\begin{veta}[mnohé koef.\ jsou nulové]
\begin{enumerate}[label=\textcircled{\arabic*}]
  \item musí být \quad $m = m_1 + m_2$, \quad neboť
  \[
    \bra{j_1j_2;m_1m_2}\underbrace{J_z - J_{1z} - J_{2z}}_{0}\ket{j_1j_2;jm} = (m - m_1 - m_2)\braket{j_1j_2;m_1m_2}{j_1j_2;jm} = 0
  \]
  \item $\Delta$-nerovnost \dots\ klasicky:
  \begin{tikzpicture}[baseline=0.2cm,scale=0.9,>=Stealth]
    \draw[->] (0,0) -- (1.6,0) node[midway,below] {\small $j_2$};
    \draw[->] (0,0) -- (0.5,0.65) node[midway,left] {\small $j_1$};
    \draw[->,thick] (0,0) -- (2.1,0.65) node[above] {\small $j$};
    \draw[dashed] (0.5,0.65) -- (2.1,0.65) (1.6,0) -- (2.1,0.65);
  \end{tikzpicture}
  \quad \dots\ $|\vec j| < |j_1| + |j_2|$ \ a \ $j_1 < j + j_2$
  \[
    \Rightarrow\quad \vzorec{|j_1 - j_2| \leq j \leq j_1 + j_2}
    \qquad\dots\ \text{pokud není splň.:}\ \langle\ |\ \rangle = 0
  \]
  {\color{defbarva}opět jediné netriviální kvůli $[J_{1z}, J^2] \neq 0$; $[J_{2z}, J^2] \neq 0$}
\end{enumerate}
\end{veta}

%----------------------------------------------------------------------
\strana{AA-4}
je konzistentní s dimenzionalitou problému:

\dots\ dimenze lin.\ obalu $\{\ket{j_1j_2;m_1m_2}\}$ \quad \dots\ $N = (2j_1+1)(2j_2+1)$

dimenze $\{\ket{j_1j_2;j,m}\}$ \quad \dots\ předp., že $\forall j$ právě jeden a $j_1 > j_2$ \dots\ BÚNO
\begin{align*}
  N = \sum_{j=j_1-j_2}^{j_1+j_2}(2j+1)
  &= \tfrac12\underbrace{\big[2(j_1+j_2+j_1-j_2) + 2\big]}_{1.\ +\ \text{poslední}}\underbrace{\big[j_1+j_2+1-j_1+j_2\big]}_{\text{počet členů řady}} \\
  &= \tfrac12\,[4j_1+2]\,[2j_2+1] \quad\dots\ \text{stejné}
\end{align*}

{\color{defbarva}
\noindent\rule{\linewidth}{0.4pt}
\underline{pozn:} podrobnější konstrukce (opakování postupu ze $s_1 = \frac12 + s_2 = \frac12$) \hfill viz \sout{též} FORMÁNEK \\
\hspace*{9.5cm} (prozatím vynechávám $\ket{j_1j_2;\dots}$ ve značení)

\begin{center}
\begin{tikzpicture}[defbarva,>=Stealth,every node/.style={font=\small}]
  \node[anchor=west] (a) at (0,0) {$\ket{m_1 = j_1, m_2 = j_2}$ je identický s $\ket{j = j_1+j_2,\ m = j_1+j_2}$};
  \node[anchor=west] at (1,-0.45) {(ověření: $\vec J^{\,2} = \dots$)};
  \draw[->] (0.4,-0.25) -- (0.4,-1.15) node[midway,right] {$J_- = J_{1-} + J_{2-}$};
  \node[anchor=west] (b) at (-0.6,-1.45) {$\sqrt{2j}\,\ket{j,m=j-1} = \sqrt{2j_1}\,\ket{m_1 = j_1-1, m_2 = j_2} + \sqrt{2j_2}\,\ket{m_1 = j_1, m_2 = j_2 - 1}$};
  \draw[->] (0.4,-1.75) -- (0.4,-2.45) node[midway,right] {$J_-$};
  \node[anchor=west] at (-0.1,-2.7) {$\ket{j,m=j-2}$};
  \node at (0.4,-3.2) {$\vdots$};
  \draw[->] (2.6,-1.75) -- (2.6,-2.0) -- (3,-2.0) node[right] {ortogonalizace: \ \dots\ ověření, že $J^2$ odpov.\ $j-1$ pomocí $J_+$};
  \node[anchor=west] at (3,-2.55) {$\ket{j-1,m-1} = \sqrt{\tfrac{2j_2}{2j}}\ket{m_1-1,m_2} - \sqrt{\tfrac{2j_1}{2j}}\ket{m_1,m_2-1}$};
  \draw[->] (6.2,-2.9) -- (6.2,-3.5) node[midway,right] {$J_-$};
  \node at (6.2,-3.8) {$\ket{j-1,m-2}$};
  \draw (3.6,-3.8) ellipse (4.6 and 0.4);
  \draw[->] (8.4,-3.7) -- (9.1,-3.7) node[midway,above] {ortogonalizace};
  \node[anchor=west] at (9.1,-3.8) {$\ket{j-2,m-2} \longrightarrow\ \dots$};
  \draw[->] (10.3,-4.1) -- (10.3,-4.6);
  \node at (10.3,-4.8) {$\vdots$};
  \node at (6.2,-4.6) {$\vdots$};
\end{tikzpicture}
\end{center}
diskuse $N$ výše nás přesvědčí, že tak vyčerpáme celý prostor $\ket{m_1,m_2}$\par
\noindent\rule{\linewidth}{0.4pt}
}

\subsection*{Vlastnosti Clebsch--Gordanových koef.}

\begin{minipage}{0.62\textwidth}
\textcircled{1} $\braket{j_1j_2;m_1m_2}{j_1j_2;jm}$ tvoří \kw{unitární matici}, \\
\quad neboť jde o přechod mezi bázemi $m_1m_2 \to jm$ uvnitř vl.\ podprostoru $j_1, j_2$
\end{minipage}
\begin{minipage}{0.36\textwidth}
\begin{tikzpicture}[every node/.style={font=\small}]
  \node at (0,0) {$\left(\rule{0pt}{1.1cm}\ \braket{j_1j_2;m_1m_2}{j_1j_2;jm}\ \right)$};
  \node at (0,1.35) {$\cdots\ jm\ \cdots$};
  \node[rotate=90] at (-2.25,0) {$\cdots\ m_1m_2\ \cdots$};
  \node[defbarva,rotate=90,align=center,font=\scriptsize] at (2.75,0) {ve vhodném uspoř.\\ blok.\ diagonální\\ (neboť $m = m_1+m_2$)};
\end{tikzpicture}
\end{minipage}

\quad $\uparrow$\ \sout{con.}\ definice nefixuje fázi koef.\ \dots\ konvence \dots\ \kw{Condon--Shortley}
\begin{enumerate}[label=\arabic*),nosep]
  \item požadavek, aby fungovala konstrukce $\uparrow\uparrow$ \ fixuje spoustu fází
  \item požadavek reálnosti \ $\braket{j_1j_2;m_1m_2}{j_1j_2;jm}\in\mathbb{R}$ \ fixuje vše až na znaménka
\end{enumerate}

\underline{důsledky:} \quad $\braket{j_1j_2;m_1m_2}{j_1j_2;jm} = \braket{j_1j_2;jm}{j_1j_2;m_1m_2}$ \\
reálná unit.\ matice je ortogonální, tj.\ \ $UU^T = U^TU = I$ \\
tj.\ její řádky tak sloupce tvoří ON bázi v $\mathbb{R}^N$, \quad neboť:
\begin{align*}
  \text{a)}\quad &\sum_{jm}\braket{j_1j_2;m_1m_2}{j_1j_2;jm}\braket{j_1j_2;m_1'm_2'}{j_1j_2;jm} = \delta_{m_1m_1'}\delta_{m_2m_2'} \\
  \text{b)}\quad &\sum_{m_1m_2}\braket{j_1j_2;m_1m_2}{j_1j_2;jm}\braket{j_1j_2;m_1m_2}{j_1j_2;j'm'} = \delta_{jj'}\delta_{mm'}
\end{align*}

občas se setkáme s tzv.\ \kw{Wignerův 3j symbol}:
\[
  \vzorec{\braket{j_1j_2;m_1m_2}{j_1j_2;jm} \equiv (-1)^{j_1-j_2+m}\sqrt{2j+1}\begin{pmatrix} j_1 & j_2 & j \\ m_1 & m_2 & -m\end{pmatrix}}
\]
který má více symetrií, např.\ se nemění při sudé permut.\ sloupců a lichá vyrobí znaménko $(-1)^{j_1+j_2+j_3}$ \\
více viz Edmonds: \emph{Angular Momentum in QM} \hfill \sout{více ve}

%----------------------------------------------------------------------
\strana{AA-5}
\subsection*{Konstrukce C-G koeficientů}

\begin{veta}[Rekurentní relace]
{\color{defbarva}(pro stručnost vynechávám $j_1j_2$ \dots\ automat.\ pracuje v přísl.\ podprost.)}
\begin{align*}
  \sqrt{(j\mp m)(j\pm m+1)}\,\braket{m_1m_2}{j,m\pm1}
  &= \sqrt{(j_1\mp m_1+1)(j_1\pm m_1)}\,\braket{m_1\mp1,m_2}{jm} \tag{$R_\pm$} \\
  &\quad + \sqrt{(j_2\mp m_2+1)(j_2\pm m_2)}\,\braket{m_1,m_2\mp1}{jm}
\end{align*}
\end{veta}
\pozn{v originále ve 2.\ odmocnině $(j_2\pm m_1)$}

\underline{DK:} \quad $J_\pm\ket{jm} = (J_{1\pm} + J_{2\pm})\sum_{m_1m_2}\ket{m_1m_2}\braket{m_1m_2}{jm}$
\begin{align*}
  \sqrt{(j\mp m)(j\pm m+1)}\,\ket{j,m\pm1} = \sum_{m_1'm_2'}\Big\{&\sqrt{(j_1\mp m_1')(j_1\pm m_1'+1)}\,\ket{m_1'\pm1,m_2'} \\
  &+ \sqrt{(j_2\mp m_2')(j_2\pm m_2'+1)}\,\ket{m_1',m_2'\pm1}\Big\}\braket{m_1'm_2'}{jm}
\end{align*}
$\downarrow$ vynásobení $\bra{m_1m_2}$ + užití ortogonality \dots\ první člen $m_1 = m_1'\pm1$, $m_2 = m_2'$; \\
\hspace*{9.2cm} druhý: $m_1 = m_1'$, $m_2 = m_2'\pm1$ \\
$\Rightarrow$ C.B.D.

\begin{minipage}{0.25\textwidth}
Použití relací:

\medskip
$j$ se nemění \\
\dots\ nebo opakovat
\end{minipage}
\begin{minipage}{0.74\textwidth}
\begin{tikzpicture}[every node/.style={font=\small}]
  \fill (0,1.3) circle (2pt) node[above] {$(m_1-1,m_2)$};
  \fill (2,1.3) circle (2pt) node[above] {$(m_1,m_2)$};
  \fill (2,0) circle (2pt) node[below] {$(m_1,m_2-1)$};
  \draw[dashed] (0,1.3) -- (2,1.3) -- (2,0) -- cycle;
  \node at (1.4,0.85) {$J_+$};
  \begin{scope}[xshift=5.5cm]
    \fill (0,1.3) circle (2pt) node[above] {$(m_1,m_2+1)$};
    \fill (0,0) circle (2pt) node[below] {$(m_1,m_2)$};
    \fill (2,0) circle (2pt) node[below right] {$(m_1+1,m_2)$};
    \draw[dashed] (0,1.3) -- (0,0) -- (2,0) -- cycle;
    \node at (0.6,0.45) {$J_-$};
  \end{scope}
\end{tikzpicture}
\end{minipage}

neznáme jedno číslo $\forall j$ \\
$\to$ z normalizace a \kw{Condon--Shortley} konvence:
\[
  \braket{j_1j_2}{j_1+j_2,\,j_1+j_2} = 1, \qquad \braket{j_1,\,j-j_1}{j,j} > 0
\]

\textbf{SCHÉMA KONSTRUKCE:} \quad ($\forall j$ fixní $= j_1+j_2,\ j_1+j_2-1,\ \dots,\ 0$)

\begin{center}
\begin{tikzpicture}[defbarva,scale=0.85,every node/.style={font=\small}]
  % levý diagram
  \begin{scope}
    \fill[pattern=north east lines,pattern color=defbarva!60] (-2,-0.5) -- (-2,2) -- (0.5,2) -- (2,0.5) -- (2,-2) -- (-0.5,-2) -- cycle;
    \draw[->] (-3.4,0) -- (3.4,0) node[below] {$m_1$};
    \draw[->] (0,-3.2) -- (0,3.4) node[right] {$m_2$};
    \draw (-3.4,2) -- (3.2,2) node[above,pos=0.12] {$m_2 = j_2$};
    \draw (-3.4,-2) -- (3.2,-2) node[below,pos=0.12] {$m_2 = -j_2$};
    \draw (-2,-3.2) -- (-2,3.2);
    \node[left] at (-2,-0.4) {$-m_1 = j_1$};
    \draw (2,-3.2) -- (2,3.2);
    \node[right] at (2,-2.7) {$m_1 = j_1$};
    \draw (-0.8,3.3) -- (3.2,-0.7) node[below] {$m_1+m_2 = j$};
    \draw (-3.2,0.7) -- (0.8,-3.3) node[right] {$m_1+m_2 = -j$};
    \node[fill=white,fill opacity=0.8,text opacity=1,inner sep=1pt,align=center] at (-0.7,-0.3) {OBLAST\\ NENULOVÝCH\\ CG\\ (pro fix $j$)};
    \draw[prbarva,thick] (0.2,0.2) rectangle (2.6,2.6);
    \draw[prbarva,->,thick] (2.6,2.3) to[bend left=10] node[above] {Detail} (5.4,2.3);
  \end{scope}
  % detail
  \begin{scope}[xshift=6.3cm,yshift=-1.6cm,scale=0.95]
    \foreach \y in {0,1,2,3} \draw (-0.4,\y) -- (3.6,\y);
    \draw (0,-0.6) -- (0,4.4);
    \draw (1,-0.6) -- (1,3.6);
    \draw[very thick] (2,-0.6) -- (2,3);
    \draw (-0.2,4.5) -- (1,3) -- (2.9,1.1);
    \foreach \x/\y in {0/4,0/3,0/2,0/0,1/3,1/2,1/1,1/0,2/2,2/1,2/0} \fill (\x,\y) circle (2.5pt);
    \draw (0,1) circle (2.5pt);
    \foreach \x/\y in {3.3/2,3.3/1,3.3/0,2/3.3} \node at (\x,\y) {$0$};
    \node[right] at (2.6,3.6) {CG $= 0$};
    \node[above right] at (2,2) {A \dots\ zvolím};
    \node[above left] at (1,2) {C};
    \node[below left] at (1,1) {D};
    \node[below right] at (2,1) {B};
    \node[below right] at (2,0) {E};
    \begin{scope}[prbarva,thick]
      \draw (1.1,1.9) -- (1.9,1.9) -- (1.9,1.1) -- cycle;
      \node at (1.65,1.65) {\scriptsize $J_+$};
      \draw (1.1,1.85) -- (1.1,1.1) -- (1.85,1.1) -- cycle;
      \node at (1.35,1.35) {\scriptsize $J_-$};
      \draw (2.1,1.9) -- (2.1,1.1) -- (2.9,1.1) -- cycle;
      \node at (2.35,1.35) {\scriptsize $J_\mp$};
      \draw (1.1,0.9) -- (1.9,0.9) -- (1.9,0.1) -- cycle;
      \node at (1.65,0.65) {\scriptsize $J_+$};
    \end{scope}
    \foreach \y in {-0.4,0.6,1.6} \draw[defbarva!60] (-0.3,\y) -- (0.8,\y+1.1);
  \end{scope}
\end{tikzpicture}
\end{center}
\hfill postupná konstrukce z A, např.\ v pořadí ABCDE \dots

\begin{priklad}[PŘÍKLAD: částice se spinem $\frac12$ ve sfér.\ poli, $\hat{\vec L}$ komutuje s $H$]
$\mathcal{H} = \mathcal{H}_{\text{orbit}}\otimes\mathcal{H}_{\text{spin}}$ \qquad $\hat{\vec J} = \hat{\vec L} + \hat{\vec S}$

\dots\ tj.\ \quad
$\begin{aligned} j_1 &= l & m_1 &= m_l \\ j_2 &= s = \tfrac12 & m_2 &= m_s = \pm\tfrac12 \end{aligned}$
\qquad $\Delta$-nerovnost povoluje jen $j = l\pm\frac12$ \\
\hspace*{11.2cm} (pro $l = 0$ jen $j = \frac12$)
\end{priklad}

%----------------------------------------------------------------------
\strana{AA-6}
\subsubsection*{Konstrukce C-G pro $\vec J = \vec L + \vec S$:}

\begin{center}
\begin{tikzpicture}[defbarva,x=0.9cm,y=1cm,every node/.style={font=\small}]
  \draw[->] (-4.5,0) -- (7,0) node[right] {$m_l$};
  \draw[->] (0,-1) -- (0,1.2) node[right] {$m_s$};
  \node at (-4.8,0.5) {$+\frac12$};
  \node at (-4.8,-0.5) {$-\frac12$};
  % řádek +1/2
  \foreach \x in {-3.5,-2} \draw (\x,0.5) circle (2pt);
  \foreach \x in {0,1,2,3,4} \fill (\x,0.5) circle (2.5pt);
  % řádek -1/2
  \foreach \x in {-3.5,-2} \draw (\x,-0.5) circle (2pt);
  \foreach \x in {0,1,2,3,4} \fill (\x,-0.5) circle (2.5pt);
  % nuly nad
  \foreach \x in {2,3,4} \draw (\x,1.2) circle (2pt);
  \node[right] at (4.6,0.45) {$\leftarrow m_s = \frac12$};
  \node[above right] at (4,0) {$+l$};
  % diagonály
  \draw (1.2,1) -- (3.3,-1.1) node[below left] {$j$};
  \draw (2.2,1) -- (4.2,-1) node[below] {$j = l-\frac12$};
  \draw (3.3,1.2) -- (5.4,-1) node[below right] {$j = l+\frac12$};
  % J- trojúhelníky
  \begin{scope}[prbarva,thick]
    \draw (2,1.2) -- (2,0.5) -- (3,0.5) -- cycle; \node at (2.3,0.75) {\scriptsize $J_-$};
    \draw (3,1.2) -- (3,0.5) -- (4,0.5) -- cycle; \node at (3.3,0.75) {\scriptsize $J_-$};
  \end{scope}
\end{tikzpicture}
\end{center}

pro $j = l + \frac12$ \ ($R_-$):
\begin{multline*}
  \sqrt{\big(l+\tfrac12+m+1\big)\big(l+\tfrac12-m\big)}\,\braket{m-\tfrac12,\tfrac12}{l+\tfrac12,m} \\
  = \sqrt{\big(l+m+\tfrac12\big)\big(l-m+\tfrac12\big)}\,\braket{m+\tfrac12,\tfrac12}{l+\tfrac12,m+1}
\end{multline*}
\hspace*{7.5cm} $\nwarrow$ pro $m = l+\frac12,\ l-\frac12,\ \dots,\ -l-\frac12$

přičemž konvence diktuje $\braket{l,\tfrac12}{l+\tfrac12,\,l+\tfrac12} = 1$ \\
$\to$ iterace
\begin{align*}
  &\braket{m-\tfrac12,\tfrac12}{l+\tfrac12,m} \\
  &\quad= \sqrt{\frac{(l+m+\frac12)\,{\color{gray}(l-m+\frac12)}}{(l+m+\frac32)\,{\color{gray}(l-m+\frac12)}}}\ \braket{m+\tfrac12,\tfrac12}{l+\tfrac12,m+1}
     \quad \pozn{krácení} \\
  &\quad= \sqrt{\frac{l+m+\frac12}{l+m+\frac32}}\sqrt{\frac{l+m+\frac32}{l+m+\frac52}}\ \braket{m+\tfrac32,\tfrac12}{l+\tfrac12,m+2}
     \quad \pozn{$m\to m+1$} \\
  &\quad\ \ \vdots \\
  &\quad= \sqrt{\frac{l+m+\frac12}{2l+1}}\ \underbrace{\braket{l,\tfrac12}{l+\tfrac12,\,l+\tfrac12}}_{=1}
\end{align*}
tj.\ \quad
$\vzorec{\braket{m-\tfrac12,\tfrac12}{l+\tfrac12,m} = \sqrt{\frac{l+m+\frac12}{2l+1}}}$
\quad \dots\ jen horní řádek pro $j = \frac12 + l$; \\
\hspace*{8.6cm} ještě dolní řádek a hodně $j = l - \frac12$

C-G matice je blok.\ diagonální \dots\ nenulové jen bloky $m = m_1 + m_2$ \quad $\swarrow$ matice pro fixní $l, s$

\begin{minipage}{0.5\textwidth}
\begin{tikzpicture}[scale=0.85,every node/.style={font=\scriptsize}]
  \draw (-0.2,2.2) to[bend right=15] (-0.2,-1.2);
  \draw (4.6,2.2) to[bend left=15] (4.6,-1.2);
  \draw (0,1.2) -- (4.4,1.2) (0,0) -- (4.4,0);
  \draw (1.8,2.4) -- (1.8,-1.4) (3.1,2.4) -- (3.1,-1.4);
  \node at (2.45,0.9) {$\cos\phi\ \ \sin\phi$};
  \node at (2.45,0.35) {$-\sin\phi\ \ \cos\phi$};
  \node[defbarva,above] at (2.45,2.4) {$j = l+\frac12,\ l-\frac12$};
  \draw[defbarva,->,thick] (2.3,2.4) -- (2.3,1.6);
  \node[left,align=right] at (-0.2,0.6) {$m_l\,m_s$\quad $m-\frac12,\ +\frac12$\\ $m+\frac12,\ -\frac12$};
  \draw[defbarva,->,thick] (0,0.9) -- (1.6,0.9);
  \draw[decorate,decoration={brace,amplitude=3pt}] (4.7,1.2) -- (4.7,0) node[midway,right=3pt] {blok $m$};
  \draw[decorate,decoration={brace,mirror,amplitude=3pt}] (1.8,-1.5) -- (3.1,-1.5) node[midway,below=3pt] {blok $m$};
\end{tikzpicture}
\end{minipage}
\begin{minipage}{0.46\textwidth}
\dots\ + \fcolorbox{defbarva}{white}{UNITARITA}
\[
  \cos\phi = \sqrt{\frac{l+m+\frac12}{2l+1}}
\]
\begin{align*}
  \sin\phi &= \sqrt{1-\cos^2\phi} \\
  &= \sqrt{\frac{l-m+\frac12}{2l+1}} = \braket{m-\tfrac12,\tfrac12}{\ }
\end{align*}
\end{minipage}

\noindent\rule{\linewidth}{0.4pt}

\begin{definice}[spin-úhlové funkce \pozn{spin-orbitaly}]
\begin{align*}
  \mathcal{Y}_l^{\,j=l\pm\frac12,m}
  &= \pm\sqrt{\frac{l\pm m+\frac12}{2l+1}}\,Y_l^{m-\frac12}(\theta,\varphi)\,\chi_+
     + \sqrt{\frac{l\mp m+\frac12}{2l+1}}\,Y_l^{m+\frac12}(\theta,\varphi)\,\chi_- \\
  &= \frac{1}{\sqrt{2l+1}}\begin{pmatrix} \pm\sqrt{l\pm m+\frac12}\ Y_l^{m-1/2}(\theta,\varphi) \\ \sqrt{l\mp m+\frac12}\ Y_l^{m+1/2}(\theta,\varphi)\end{pmatrix}
  \qquad\dots\ \text{vlast.\ fce } L^2, S^2, J^2 \text{ a } J_z
\end{align*}
\end{definice}

$\Rightarrow$ také vl.\ fce \quad
$\vec L\cdot\vec S = \frac12\big(J^2 - L^2 - S^2\big) = \frac{\hbar^2}{2}\Big(j(j+1) - l(l+1) - \frac34\Big)$
$\begin{cases}
  \oplus:\ \ \frac{l}{2}\hbar^2 & \text{pro } j = l+\frac12 \\[2pt]
  \ominus:\ -\frac{l+1}{2}\hbar^2 & j = l-\frac12
\end{cases}$


\newpage
%======================================================================
% Poznamky_7_identicke_castice.pdf
%======================================================================
\section{Systémy několika identických částic -- fyzika}
\strana{I1}

\kw{identické částice} $\to$ všechny vlastnosti stejné
\begin{itemize}[nosep]
  \item[--] Hamiltonián symetrický vůči záměně $i\leftrightarrow j$
\end{itemize}
např.:
\[
  H = \frac{p_1^2}{2m} + \frac{p_2^2}{2m} + V_{\text{int}}\big(|\vec r_1 - \vec r_2|\big) + V_{\text{ext}}(\vec r_1) + V_{\text{ext}}(\vec r_2)
\]
v QM hlubší důsledky než v klas. Např.\ částice nelze identifikovat podle trajektorie.

\kw{nerozlišitelnost}: žádným měřením nelze rozhodnout, o kterou částici jde

\subsection*{Výměnná (permutační) symetrie}
operace symetrie -- záměna $i\leftrightarrow j$ $\to$ unitární operátor \\
\kw{permutační operátor} $P_{ij}$

\begin{priklad}[Př: dvě částice]
$\mathcal{H} = \mathcal{H}_1\otimes\mathcal{H}_2$ \\
nerozlišitelnost \dots\ izomorfie $\mathcal{H}_1\equiv\mathcal{H}_2\equiv\mathcal{H}$ \\
nechť: částice 1 je ve stavu $\ket{\alpha}\in\mathcal{H}$ \\
\phantom{nechť:} částice 2 \phantom{je} ve stavu $\ket{\beta}\in\mathcal{H}$ \quad \dots\ $\ket{\psi} = \ket{\alpha}\ket{\beta}$

def: \quad $P_{12}\ket{\alpha}\ket{\beta} = \ket{\beta}\ket{\alpha}$ \\
podob.\ na složené stavy \quad
$P_{12}\sum_{\alpha\beta}c_{\alpha\beta}\ket{\alpha}\ket{\beta} = \sum_{\alpha\beta}c_{\alpha\beta}\ket{\beta}\ket{\alpha} = \sum_{\alpha\beta}c_{\beta\alpha}\ket{\alpha}\ket{\beta}$

$x$-reprezentace:
\begin{align*}
  P_{12}\int\psi(x_1,x_2)\ket{x_1}\ket{x_2}\,\dd x_1\dd x_2 &= \int\psi(x_1,x_2)\ket{x_2}\ket{x_1}\,\dd x_1\dd x_2 \\
  \dots\ \text{přeznačení integ.\ prom.\ } x_1\leftrightarrow x_2 &= \int\psi(x_2,x_1)\ket{x_1}\ket{x_2}\,\dd x_1\dd x_2
\end{align*}
tj.\ \quad $P_{12}\psi(x_1,x_2) = \psi(x_2,x_1)$
\end{priklad}

\underline{vl.\ vektory $P_{12}$:} \quad $P_{12}^2 = I$ \quad $\Rightarrow$ \quad $\lambda = \pm1$ \dots\ vl.\ č.
\begin{align*}
  \lambda = 1 &\quad\dots\quad \psi(x_1,x_2) = \psi(x_2,x_1) \quad\dots\ \text{symetrie } c_{\alpha\beta} = c_{\beta\alpha} \\
  \lambda = -1 &\quad\dots\quad \psi(x_1,x_2) = -\psi(x_2,x_1) \quad\dots\ \text{antisymetrie } c_{\alpha\beta} = -c_{\beta\alpha}
\end{align*}

\underline{symetrie $H$:} \quad $P_{12}HP_{12} = H$ \quad tj.\ $[P_{12},H] = 0$ \dots\ společné vl.\ v. \\
\quad tj.\ $\forall$ vl.\ stav $H$ lze volit buď sym.\ nebo antisym. \\
\quad (nemusí být degenerace)

\underline{více částic:} \quad problém \quad $[P_{12},P_{23}]\neq 0$ \quad atd.
\quad \pozn{$P_{12}P_{23}\ket{123} = \ket{312}$, \ $P_{23}P_{12}\ket{123} = \ket{231}$} \\
\dots\ nelze volit společnou úplnou množinu \sout{vektorů} vl.\ vektorů, \\
ale invariantní podprostory:

%----------------------------------------------------------------------
\strana{I2}
\begin{priklad}[Př: invariantní podprostory pro 3 částice]
\begin{itemize}
  \item totál.\ symetrický (1D) / antisymetrický
  \[
    \tfrac{1}{\sqrt6}\big\{\ket{\alpha\beta\gamma} + \ket{\beta\gamma\alpha} + \ket{\gamma\alpha\beta} \pm \ket{\gamma\beta\alpha} \pm \ket{\beta\alpha\gamma} \pm \ket{\alpha\gamma\beta}\big\}
    \qquad \pozn{$P_\forall\ (\pm)$}
  \]
  \item 2D:
  \[
  \begin{array}{r@{\,}c@{\ }c@{\ }c@{\ }c@{\ }c@{\ }c@{\ }l}
    \frac{1}{\sqrt{12}}\big\{ & 2\ket{\alpha\beta\gamma} & +2\ket{\beta\alpha\gamma} & -\ket{\alpha\gamma\beta} & -\ket{\gamma\alpha\beta} & -\ket{\beta\gamma\alpha} & -\ket{\gamma\beta\alpha} & \big\}\quad P_{12}\,(+) \\[3pt]
    \frac12\big\{ & 0 & 0 & +\ket{\alpha\gamma\beta} & -\ket{\gamma\alpha\beta} & +\ket{\beta\gamma\alpha} & -\ket{\gamma\beta\alpha} & \big\}\quad P_{12}\,(-)
  \end{array}
  \]
  \item 2D:
  \[
  \begin{array}{r@{\,}c@{\ }c@{\ }c@{\ }c@{\ }c@{\ }c@{\ }l}
    \frac{1}{\sqrt{12}}\big\{ & 2\ket{\alpha\beta\gamma} & -2\ket{\beta\alpha\gamma} & +\ket{\alpha\gamma\beta} & -\ket{\gamma\alpha\beta} & +\ket{\gamma\beta\alpha} & -\ket{\beta\gamma\alpha} & \big\}\quad P_{12}\,(-) \\[3pt]
    \frac12\big\{ & 0 & 0 & + & + & - & - & \big\}\quad P_{12}\,(+)
  \end{array}
  \]
\end{itemize}
\end{priklad}

Permutační operátory míchají vektory z invar.\ prostoru jen mezi sebou, \\
navíc $\forall$ komutují s $H$ \quad $\to$ vl.\ v.\ $H$ lze konstr.\ z vl.\ podprostorů, \\
\hspace*{5.5cm} navíc čas.\ vývoj ponechává vl.\ podprostory

\underline{naštěstí:} uvidíme, že jen tot.\ sym.\ / tot.\ antisym.\ jsou povoleny

\sout{Nerozlišitelnost}: \underline{Ostatní pozorovatelné}: \quad
(příklad: pozitronium $e^-e^+$, $H$ invar.\ $e^+\leftrightarrow e^-$, ale rozliš.) \\
nerozlišitelnost $\to$ žádné měření nerozliší $\ket{\psi}$ a $P_{ij}\ket{\psi}$ \\
\quad tj.\ \quad $\mel{\psi}{P_{ij}AP_{ij}}{\psi} = \mel{\psi}{A}{\psi}\ \ \forall\ket{\psi}$ \quad $\Rightarrow$ \quad $[A,P_{ij}] = 0$ \\
$\to$ $\forall$ pozorovatelné jsou invar.\ vůči permutacím

\pozor{vlastní podprostory se nemíchají!} \\
např.: \quad $P_{ij}\ket{s} = \ket{s}$, \quad $P_{ij}\ket{a} = -\ket{a}$ \quad $\Rightarrow$ \quad $\mel{s}{A}{a} = -\mel{s}{A}{a} = 0$ \\
\quad tj.\ pro $\ket{\psi} = \alpha\ket{s} + \beta\ket{a}$ je $\mel{\psi}{A}{\psi} = |\alpha|^2\mel{s}{A}{s} + |\beta|^2\mel{a}{A}{a}$ \\
$\to$ interference není pozorovatelná \quad \dots\ \kw{superselekční pravidlo} \\
\hspace*{6cm} (výběrová prav.\ je pro $H$; proto super)

\begin{veta}[Symetrizační postulát]
V přírodě nastává jen jeden z těchto případů:
\begin{enumerate}[label=(\alph*)]
  \item částice má celočíselný spin a její mnohočásticové stavy jsou totálně symetrické \dots\ \kw{BOSON} \dots\ Bose--Einstein statistika
  \item částice má poločíselný spin a její mnohočásticové stavy jsou totálně antisymetrické \dots\ \kw{FERMION} \dots\ Fermi--Dirac st.
\end{enumerate}
\end{veta}

\begin{veta}[Důsledek: Pauliho vylučovací princip]
dva nerozlišitelné fermiony nemohou být ve stejném stavu.
\end{veta}

%----------------------------------------------------------------------
\strana{I3}
\subsection*{Příklad: dvouelektronové systémy}
Předp., že $[H,S] = 0$, kde $S = S_1 + S_2$ je celkový spin

Potom vl.\ stavy $H$ separované: \quad
$\ket{\psi} = \underbrace{\ket{\phi}}_{\text{prostor.\ část}}\underbrace{\ket{\chi}}_{\text{spin.\ část}}$

kde $\ket{\chi}$ lze volit jako vl.\ stav $S^2, S_z$, tj.\ (viz skládání m.\ hybn.)
\[
  \ket{\chi} = \begin{cases}
    S = 1: & \left.\begin{array}{l}
               \ket{++} \\
               \frac{1}{\sqrt2}\big(\ket{+-} + \ket{-+}\big) \\
               \ket{--}
             \end{array}\right\}
             \ \text{\uv{triplet}}\dots\text{symetrický}\quad P_{12}\underbrace{\ket{\chi}}_{\ket{1m}} = \ket{\chi} \\[18pt]
    S = 0: & \ \frac{1}{\sqrt2}\big(\ket{+-} - \ket{-+}\big)
             \quad\dots\ \text{\uv{singlet}}\dots\text{antisym.}\quad P_{12}\ket{\chi} = -\ket{\chi}
  \end{cases}
\]

\begin{itemize}
  \item triplet je symetrický, neboť $S_\pm = S_{1\pm} + S_{2\pm}$ komutují s $P_{12}$ a stav $\ket{++}$ je symetrický.
  \item permutační operátor lze faktorizovat \quad $P_{12} = P_{12}^{(x)}P_{12}^{(s)}$, kde \\
        spinová část \quad $P_{12}^{(s)} = \frac12\Big(I + \frac{4}{\hbar^2}\vec S_1\cdot\vec S_2\Big)$, \quad
        neboť vl.\ č.\ $S_1\cdot S_2 = \begin{cases} \frac{\hbar^2}{4} & \text{triplet} \\ -\frac34\hbar^2 & \text{singlet}\end{cases}$
\end{itemize}
\[
  P_{12}\ket{\phi}\ket{\chi} = P_{12}^{(x)}\ket{\phi}\,P_{12}^{(s)}\ket{\chi}
\]
$\ket{\psi}$ musí být totálně antisym. \quad $\Rightarrow$ \quad
$\begin{cases} \ket{\phi}\ \text{sym.\ pro singlet} \\ \ket{\phi}\ \text{antisym.\ pro triplet}\end{cases}$

\begin{minipage}{0.35\textwidth}
\underline{Důsledky:} dva atomy
\end{minipage}
\begin{minipage}{0.6\textwidth}
\begin{tikzpicture}[gray]
  \draw (0,0) circle (0.75);
  \node at (-0.35,0.05) {$A$}; \node at (0.05,0.3) {$e^-$};
  \node[right] at (0.8,-0.2) {\small $x_1\in\Omega_A$};
  \draw (4,0) circle (0.75);
  \node at (3.65,0.05) {$B$}; \node at (4.05,0.3) {$e^-$};
  \node[right] at (4.8,0) {\small $x_2\in\Omega_B$};
\end{tikzpicture}
\end{minipage}

vln.\ fce atomu $A$: \ $\phi_A(\vec x_1)$ \hfill vl.\ fce atomu $B$: \ $\phi_B(\vec x_2)$

vln.\ fce složeného systému: \quad
$\phi(x_1,x_2) = \frac{1}{\sqrt2}\big[\phi_A(x_1)\phi_B(x_2) \pm \phi_A(x_2)\phi_B(x_1)\big]$ \\
\hspace*{9cm} pro singlet / triplet.

\dots\ \pozor{MRZUTÉ!} \dots\ atom v laboratoři je promíchaný s atomem na Měsíci!? \\
naštěstí nemá pozorovatelné důsl., pokud $\Omega_A\cap\Omega_B = \emptyset$

Např.: hustota pravd.\ nalezení $e^-$ v $A$ a druhého v $B$: \ $x_1\in\Omega_A$, $x_2\in\Omega_B$
\begin{align*}
  p(x_1,x_2) = \tfrac12\Big\{&|\phi_A(x_1)|^2\cdot|\phi_B(x_2)|^2 + |\phi_A(x_2)|^2\cdot|\phi_B(x_1)|^2 \\
  &\pm \underbrace{2\,\mathrm{Re}\,\phi_A(x_1)\phi_B(x_2)\phi_A^*(x_2)\phi_B^*(x_1)}_{\text{výměnný člen}}\Big\}
\end{align*}
\quad výměnný člen \dots\ udělá pro $x_1 = x_2$ \quad hustotu \quad
$\begin{cases} 0 & \text{pro Fermiony} \\ 2\times\text{větší} & \text{pro Bosony}\end{cases}$

ale pro $x_1\in\Omega_A$, $x_2\in\Omega_B$ zbude jen $p_A\cdot p_B$

%----------------------------------------------------------------------
\strana{I4}
\subsection*{Atom Helia}
\[
  H = \frac{p_1^2}{2m} + \frac{p_2^2}{2m} - \frac{Ze^2}{r_1} - \frac{Ze^2}{r_2} + \boxed{\frac{e^2}{r_{12}}}
\]
řešení jako atom vodíku, ale $Z = 1\to Z = 2$

0.\ přiblížení: zanedbáme $e^-$--$e^-$ interakci $\uparrow$:
\[
  \phi(\vec x_1,\vec x_2) = \frac{1}{\sqrt2}\Big[\psi_{nlm}(\vec x_1)\psi_{n'l'm'}(\vec x_2) \pm \psi_{nlm}(\vec x_2)\psi_{n'l'm'}(\vec x_1)\Big]
  \qquad \pozn{$+$ pro singlet, $-$ pro triplet}
\]
speciálně základní stav \quad $n = n' = 1$, $l = l' = 0$, $m = m' = 0$ \dots\ dovoleno jen $+$
\[
  \phi_0(\vec x_1,\vec x_2) = \frac{Z^3}{\pi a_0^3}e^{-Z(r_1+r_2)/a_0}\cdot\chi_{\text{singlet}}
  \qquad \pozn{$\chi_{\text{singlet}} = \frac{1}{\sqrt2}\big(\ket{\uparrow\downarrow} - \ket{\downarrow\uparrow}\big)$}
\]
\hfill $E_0 = 2\times\Big(-\dfrac{4e^2}{2a_0}\Big) = -8\cdot\dfrac{e^2}{2a_0}$ \quad \pozn{$\sim -108.84$\,eV}

\begin{itemize}
  \item poruchová teorie \dots\ 1.\ korekce:
  \[
    \Delta_{(1s)^2} = \Big\langle\frac{e^2}{r_{12}}\Big\rangle_{\phi_0} = \iint\frac{Z^6}{\pi^2a_0^6}e^{-2Z(r_1+r_2)/a_0}\frac{e^2}{r_{12}}\,\dd^3x_1\,\dd^3x_2
  \]
  + použití multipól.\ rozvoje: \quad
  $\dfrac{1}{r_{12}} = \displaystyle\sum_{lm}\frac{4\pi}{2l+1}\frac{r_<^l}{r_>^{l+1}}Y_{lm}^*\Big(\frac{x_1}{r_1}\Big)Y_{lm}\Big(\frac{x_2}{r_2}\Big)$ \\
  + $\displaystyle\int\dd\Omega_{x_1}Y_{lm}\Big(\frac{x_1}{r_1}\Big) \sim \delta_{l0}\delta_{m0}$ -- zruší sumu \\
  ve shodě je třeba rozdělit $\int_0^\infty\!\int_0^\infty\dd r_1\dd r_2$ na $\int_0^\infty\Big[\int_0^{r_1} + \int_{r_1}^\infty\Big]$

  \dots\ $\Delta E = \frac52\Big(\dfrac{e^2}{2a_0}\Big)$ \quad \dots \quad
  $E = \Big(-8 + \frac52\Big)\Big(\dfrac{e^2}{2a_0}\Big) \doteq -74.8$\,eV

  \item Variační metoda: \\
  nejjednodušší vzít stejnou funkci, ale proměnné $Z'$ (screening) \\
  tj.\ \quad $\phi_v(\vec x_1,\vec x_2) = \dfrac{Z'^3}{\pi a_0^3}e^{-Z_{\text{eff}}(r_1+r_2)/a_0}$

  výpočet $E(Z') = \langle H\rangle_{\phi_v}$ probíhá jako $\uparrow$
  \[
    E(Z') = \langle T_0\rangle + \langle V_0\rangle + \Big\langle\frac{e^2}{r_{12}}\Big\rangle
    = \Big(2\frac{Z'^2}{2} - 2ZZ' + \frac58Z'\Big)\Big(\frac{e^2}{a_0}\Big)
  \]
  $\to$ minimum \quad $Z' = 2 - \frac{5}{16}$ \quad $\to$ \quad $\vzorec{E = -77.5\,\text{eV}}$ \\
  \hspace*{6.3cm} $E_{\text{exp}} \doteq -78.8$\,eV
\end{itemize}

\underline{Excit.\ stavy:} \quad $n'l'm' \neq (1,0,0)$
\[ E = E_{100} + E_{nlm} + \Delta E \]
v 1.\ řádu: \quad $\Delta E = \Big\langle\dfrac{e^2}{r_{12}}\Big\rangle = I \pm J$
\begin{align*}
  \text{přímá interakce}\quad I &= \int\dd^3x_1\int\dd^3x_2\,|\psi_{100}(x_1)|^2|\psi_{nlm}(x_2)|^2\frac{e^2}{r_{12}} \\
  \text{\kw{výměnný integrál}}\quad J &= \int\dd^3x_1\,\dd^3x_2\,\psi_{100}^*(\vec x_1)\psi_{nlm}^*(\vec x_2)\frac{e^2}{r_{12}}\psi_{100}(\vec x_2)\psi_{nlm}(\vec x_1) \\
  &\qquad\text{(single center)}
\end{align*}

%======================================================================
\strana{I5}
\section{Systémy identických částic -- rozvinutí formalismu}
\hfill (literatura: Formánek, Ballentine)

\underline{opakování} (shrnutí):
\begin{itemize}[nosep]
  \item[--] permutační (výměnná) symetrie
  \item[--] rozpad $\forall$ vl.\ podprostorů podle IR permutační grupy
  \item[--] symetrizační postulát $\to$ jen totálně symetrická / antisymetrická 1-D reprezentace se realizuje v přír.
\end{itemize}

\underline{závěr:} $\mathcal{H} = \mathcal{H}_1\otimes\mathcal{H}_2\otimes\dots\otimes\mathcal{H}_N$ není správným stav.\ prostorem pro systém ident.\ částic, ale jen jeho podprostor $\mathcal{H}_S/\mathcal{H}_A$

dva způsoby popisu \dots\ \textcircled{A} uvnitř $\mathcal{H}$ jako podprostor \\
\hspace*{3.8cm} \textcircled{B} elegantní formalismus \kw{Fockova prostoru}

\textcircled{A} Pro popis stavů používáme jen totálně symetrické (bosony) / totálně antisymetrické (fermiony) funkce z $\mathcal{H}$ \\
Popis měření \dots\ operátory na $\mathcal{H}$, ale správná veličina $[A,P_{ij}] = 0\ \ \forall i,j$

\underline{popis vektorů z $\mathcal{H}_S/\mathcal{H}_A$}

\dots\ úplná množina pozorovatelných na $\mathcal{H}_1\equiv\mathcal{H}_2\equiv\dots\equiv\mathcal{H}_N$ \\
\quad $\to$ \quad $\hat C = \sum_n c_n\ket{n}\bra{n}$ \quad \dots\ $\hat C\ket{n} = c_n\ket{n}$ \\
stavy z $\mathcal{H}$ \dots\ $\ket{n_1n_2\dots n_N} \equiv \ket{n_1}\otimes\ket{n_2}\otimes\dots\otimes\ket{n_N}$

permutace $p = \begin{pmatrix} 1 & 2 & \dots & n \\ p_1 & p_2 & \dots & p_n\end{pmatrix}$ \quad \dots\ operátor \quad
$\hat P_p\ket{n_1\dots n_N} = \ket{n_{p_1}\dots n_{p_N}}$ \\
\pozn{dosud jsme měli jen $P_{ij}$ \dots\ $\begin{pmatrix}1&2&\dots&i&\dots&j&\dots&n\\1&2&\dots&j&\dots&i&\dots&n\end{pmatrix} \equiv p_{ij}$}

totálně sym./antisym.\ vektory \quad $\hat P_p\ket{\psi} = \pm\ket{\psi}$ \quad $\leftarrow$ pro liché permutace

\begin{veta}[Projektory na $\mathcal{H}_S/\mathcal{H}_A$]
\[
  \hat S = \frac{1}{N!}\sum_p\hat P_p, \qquad
  \hat A = \frac{1}{N!}\sum_p(-1)^p\hat P_p \qquad (-1)^p\ \dots\ \text{znaménko permutace}
\]
\end{veta}

DK, že projektory \dots\ $\forall$ permutace lze složit z transpozic $\to$ $\hat P_p = \hat P_{ij}\dots\hat P_{i'j'}$ \\
\quad pro $P_{ij}$ víme \dots\ $P_{ij}^2 = I$, \quad $P_{ij}^\dagger = P_{ij}$ \quad $\Rightarrow$ \quad $S^\dagger = S$; \ $A^\dagger = A$ \\
\hspace*{9.5cm} (suma probíhá přes $p^{-1}$)

+ platí \quad
$\begin{aligned}[t]
  \hat P_{jk}\hat S &= \hat S\hat P_{jk} = \hat S \\
  \hat P_{jk}\hat A &= \hat A\hat P_{jk} = -\hat A
\end{aligned}$
\quad (suma přes $p_{jk}p$ \dots\ probíhá $\forall$ permutace)

$\Rightarrow$ \quad 1) \ $P_{jk}\ket{\phi} = +\ket{\phi}$ \ pro $\ket{\phi} = S\ket{\psi}$ \\
\hspace*{1.95cm} $= -\ket{\phi}$ \ \dots\ $\ket{\phi} = A\ket{\psi}$ \\[3pt]
\quad\quad 2) \ $\hat S^2 = \frac{1}{N!}\sum_p\hat P\hat S = \frac{1}{N!}N!\,\hat S$, \qquad
$\hat A^2 = \frac{1}{N!}\sum_p(-1)^p\hat P\hat A = \frac{1}{N!}N!\,(-1)^{2p}\hat A$

%----------------------------------------------------------------------
\strana{I6}
působení v $x$-reprezentaci \dots\ $\ket{\psi} = \ket{\phi_1}\ket{\phi_2}\dots\ket{\phi_N}$ \\
\hspace*{3.5cm} $\to$ \quad $\psi(x_1\dots x_N) \equiv \braket{x_1\dots x_N}{\psi} = \phi_1(x_1)\dots\phi_N(x_N)$

\quad 1) \quad $\hat P\psi = \phi_1(x_{p_1})\dots\phi_N(x_{p_N})$

\kw{Slaterův determinant}: \\
\quad 2) \quad
\[
  \mathcal{N}\cdot\hat A\psi \equiv \mathcal{N}\cdot\hat A\,\phi_1(x_1)\dots\phi_N(x_N)
  = \vzorec{\frac{1}{\sqrt{N!}}
  \begin{vmatrix}
    \phi_1(x_1) & \phi_1(x_2) & \dots & \phi_1(x_N) \\
    \phi_2(x_1) & & & \vdots \\
    \vdots & \vdots & & \\
    \phi_N(x_1) & \phi_N(x_2) & \dots & \phi_N(x_N)
  \end{vmatrix}}
\]
\quad $\mathcal{N}$ \dots\ normalizační koef.\ $\sqrt{N!}$

\subsection*{Reprezentace obsazovacích čísel}
problém \dots\ konstrukce báze \ $S\ket{n_1n_2\dots n_N}$ \quad \dots\ $\mathcal{H}\to\mathcal{H}_S$ \\
\hspace*{4.4cm} $\downarrow\ \ \dots\ \ \downarrow$ \\
\quad při procházení $\forall$ hodnot \quad $c_{n_1}\ \dots\ c_{n_N}$ \ operátoru $\hat C$ \\
\quad dostanu vektory \sout{s $n_1 = n_2$} vícekrát

např.\ \quad $S\ket{n_1\dots n_i\dots n_j\dots n_N} = S\ket{n_1\dots n_j\dots n_i\dots n_N}$

$S\ket{n_1\dots n_N}$ lze jednoznačně charakterizovat tzv.\ \kw{obsaz.\ čísly} \\
\quad $N_1, N_2,\dots$ \quad $N_1 + N_2 + \dots = N$ \quad \dots\ $N_i$ udává, kolikrát se vyskytuje \\
\hspace*{3cm} v $\ket{n_1\dots n_N}$ číslo $n_i$, tj.\ hodnota $c_{n_i}$ \\
\hspace*{3cm} \dots\ kolik částic má je ve stavu $\ket{n_i}$ \dots

\begin{definice}[obsazovací čísla]
tj.\ def \quad $\ket{N_1N_2\dots} = \mathcal{N}\hat S\ket{n_1\dots n_N}$, \quad kde číslo $n_i$ se vyskytuje $N_{n_i}$ krát
\end{definice}

např.\ \quad $\ket{10201\dots} = \mathcal{N}\hat S\ket{1335}$ \quad \pozn{+ rozvést \dots\ několik 1.\ stavů v obou značeních}

normalizační koeficient $\mathcal{N}$? \quad \dots \quad $\vzorec{\mathcal{N} = \sqrt{\frac{N!}{N_1!N_2!\dots}}}$

\underline{DK:} \quad
$\hat S\ket{n_1\dots n_N} = \Big(\dfrac{N!}{N_1!N_2!\dots}\ \text{navz.\ OG členů}\Big)\times\big(\text{faktor } N_1!N_2!\dots\big)\cdot\underbrace{\dfrac{1}{N!}}_{\text{z def } \hat S}$
\[
  \Rightarrow\quad \mathcal{N} = \sqrt{\frac{N_1!N_2!\dots}{N!}\cdot\frac{N!}{N_1!N_2!\dots}} = \ \nearrow
\]

relace úplnosti na $\mathcal{H}_S$ \dots \quad
$I\,(=\hat S) = \displaystyle\sum_{\substack{N_1N_2\dots \\ \sum N_k = N}}\ket{N_1N_2\dots}\bra{N_1N_2\dots}$

Fermiony \dots\ jediný rozdíl, že $N_i\in\{0,1\}$, jinak $\hat A\ket{n_1\dots n_N} = 0$
\[
  \Rightarrow\quad \ket{N_1N_2\dots} = \sqrt{N!}\,A\ket{n_1\dots n_N}
\]
\hrule\smallskip

pozn: relace úplnosti na $\mathcal{H}$: \quad
$\displaystyle\sum_{n_1n_2\dots n_N}\ket{n_1n_2\dots n_N}\bra{n_1\dots n_N} = 1_{\mathcal{H}}$ \quad $\Big\downarrow A(\cdot)A$
\[
  \sum_{n_1\dots n_N}\underbrace{A\ket{n_1\dots n_N}}_{\frac{1}{\sqrt{N!}}\ket{N_1N_2\dots}}\bra{n_1\dots n_N}A = A = 1_{\mathcal{H}_A}
\]
$\sum_{N_1N_2\dots}$, ale $\forall$ člen je tam $N!\times$ \quad $\to$ \quad
$N!\displaystyle\sum_{N_1N_2\dots}\Big(\frac{1}{\sqrt{N!}}\Big)^2\ket{N_1N_2\dots}\bra{N_1N_2\dots} = 1_{\mathcal{H}_A}$

podob.\ $\mathcal{H}_S$

%----------------------------------------------------------------------
\strana{I7}
\subsection*{Operátory ve Fokově (Fockově) prostoru}
\begin{minipage}{0.55\textwidth}
\underline{pozn:} Vladimir Alexandrovič Fok \pozn{v originále azbukou} \\
přepis do češtiny? FOK / FOCK?
\end{minipage}
\begin{minipage}{0.43\textwidth}
\underline{Google} (2011): \\
Hartree--Fock \dots\ cca 500\,000 \\
Hartree--Fok \dots\ cca 500 \\
Hartree--Fork \dots\ cca 3000
\end{minipage}

\begin{definice}[Fockův prostor]
\[
  \mathcal{H}_F \equiv \bigoplus_{N=0}^{\infty}\mathcal{H}_S^{(N)}, \quad\text{kde } \mathcal{H}_S^{(N)} \text{ je tot.\ sym.\ prostor } N \text{ částic}
\]
podob.\ \quad $\bigoplus_{N=0}^{\infty}\mathcal{H}_A^{(N)}$ \quad pro Fermiony
\end{definice}

\textcircled{1} \underline{případ Bosonů}

definitoricky $\mathcal{H}_S^{(0)}$ je 1D prostor s bází $\ket{0}$ \dots\ \kw{vakuum} \dots\ žádná částice \\
\hspace*{2.3cm} $\mathcal{H}_S^{(N)}$ \dots\ tot.\ sym.\ prostor s bází $\ket{N_1N_2\dots}$, \quad $N_1 + N_2 + \dots = N$

hodí se pro popis otevřených systémů nebo kreaci/anihilaci částic, ale výhody i pro fixní $N$ \dots\ algebraická struktura \\
\hspace*{3cm} $N_n$ \dots\ shodné spektrum jako nezáv.\ harm.\ oscil.

\begin{definice}[kreační operátor]
\kw{kreační operátor} $\hat a_n^\dagger$ částice ve stavu $\ket{n}$ \\
\hspace*{3cm} (nezapomínat \dots\ záv.\ na bázi $\hat C\ket{n} = c_n\ket{n}$) \\
$\hookrightarrow$ operátor na $\mathcal{H}_F$ \quad $a_n^\dagger: \mathcal{H}_S^{(N)}\to\mathcal{H}_S^{(N+1)}$
\[
  \vzorec{a_n^\dagger\ket{N_1\dots N_n\dots} = \sqrt{N_n+1}\,\ket{N_1\dots N_n+1\dots}}
  \qquad\dots\ \text{přidá částici do stavu } \ket{n}
\]
\end{definice}

\underline{vlastnosti:}
\begin{itemize}
  \item $\vzorec{a_n\ket{N_1\dots N_n\dots} = \sqrt{N_n}\,\ket{N_1\dots N_n-1\dots}}$ \quad \dots\ \kw{anihilační operátor} \\
        speciálně \quad $a_n\ket{0} = 0$ \\
        \pozn{DK: $\mel{N_1'\dots N_n'}{a_n^\dagger}{N_1\dots N_n\dots} = \mel{N_1\dots N_n\dots}{a_n}{N_1'\dots N_n'\dots}^*$}
  \item $\ket{N_1,N_2\dots} = \dfrac{1}{\sqrt{N_1!}}(a_1^\dagger)^{N_1}\dfrac{1}{\sqrt{N_2!}}(a_2^\dagger)^{N_2}\dots\ket{0}$ \\
        neboli \quad $S\ket{n_1}\dots\ket{n_N} = \dfrac{1}{\sqrt{N!}}a_{n_1}^\dagger\dots a_{n_N}^\dagger\ket{0}$ \quad
        \pozn{$= \sqrt{\frac{N_1!N_2!\dots}{N!}}\ket{N_1N_2\dots}$}
  \item $[a_i^\dagger,a_j^\dagger] = [a_i,a_j] = 0$ \qquad $\vzorec{[a_i,a_j^\dagger] = \delta_{ij}}$ \\
        \pozn{DK: $\uparrow$ je vidět z def.} \qquad
        \pozn{rovnež z def.\ -- pro $i = j$: \ $a_ia_i^\dagger\ket{N_1\dots N_i\dots} = (N_i+1)\ket{N_1\dots N_i\dots}$, \ $a_i^\dagger a_i\ket{N_1\dots N_i\dots} = N_i\ket{N_1\dots N_i\dots}$ \ $\forall N_1,N_2$}
  \item operátor počtu částic ve stavu $i$: \quad $\hat N_i = a_i^\dagger a_i$
  \item operátor celk.\ počtu částic \dots\ $\hat N = \sum_i\hat N_i$
\end{itemize}

%----------------------------------------------------------------------
\strana{I8}
\underline{změna báze:} \quad $\hat C\ket{n} = c_n\ket{n}$ \quad $\to$ \quad $\hat B\ket{b{:}n} = b_n\ket{b{:}n}$ \\
kreační operátory \quad $\ket{c{:}n} = \hat a_n^\dagger(c)\ket{0}$ \qquad $\ket{b{:}n} = \hat a_n^\dagger(b)\ket{0}$ \\
transformace báze \quad $\ket{b{:}n} = \sum_m\ket{c{:}m}\underbrace{\braket{c{:}m}{b{:}n}}_{\text{koef.\ rozvoje}}$
\[
  \Rightarrow\quad
  \left.\begin{aligned}
    \hat a_n^\dagger(b) &= \sum_m\hat a_m^\dagger(c)\braket{c{:}m}{b{:}n} \\
    \hat a_n(b) &= \sum_m\hat a_m(c)\braket{b{:}n}{c{:}m}
  \end{aligned}\right\}
  \ \to\ [\hat a_n(b),\hat a_m^\dagger(b)] = \delta_{nm}\ \ \text{díky unitaritě } \braket{c{:}m}{b{:}n}
\]

\underline{spojitá báze:} \quad $\braket{x}{n} = \phi_n(x)$ \quad \dots\ def.\ kreač.\ oper.\ $\psi^\dagger(x)\ket{0} = \ket{x}$
\[
  \Rightarrow\quad
  \left.\begin{aligned}
    \hat\psi^\dagger(x) &= \sum_n\hat a_n^\dagger\braket{n}{x} = \sum_n\phi_n(x)^*\hat a_n^\dagger \\
    \hat\psi(x) &= \sum_n\phi_n(x)\hat a_n
  \end{aligned}\right\}
  \quad\text{\uv{\kw{field operators}}}
\]
Podobně lze zahrnout spin \quad $\hat\psi_m^\dagger(x)\ket{0} = \ket{x}\ket{s,m}$ \\
operátor hustoty počtu částic v místě $x$: \quad $\hat N(x) = \hat\psi^\dagger(x)\hat\psi(x)$

\subsection*{Vyjádření operátorů}
\uv{\kw{Jednočásticový operátor}} na $\mathcal{H} = \mathcal{H}_1\otimes\dots\otimes\mathcal{H}_N$

příklad \quad $\hat{\mathbb{C}} = \sum_{i=1}^{N}\hat c^{(i)}$, \quad kde $\hat c^{(i)} \equiv \underset{1}{\hat 1}\otimes\dots\otimes\underset{i}{\hat c}\otimes\dots\otimes\underset{N}{\hat 1}$ \\
\quad $\uparrow$ ale jen \dots\ nemá smysl se ptát, jaké $\hat c$ naměřím na částici 1, (nerozlišitelnost)
\[
  \hat{\mathbb{C}}\ket{n_1}\ket{n_2}\dots\ket{n_N} = \underbrace{(c_{n_1} + c_{n_2} + \dots + c_{n_N})}_{\text{hodnota } c_{n_i} \text{ je zde } N_{n_i} \text{ krát}}\ket{n_1}\dots\ket{n_N}
  = \underbrace{(c_1N_1 + c_2N_2 + \dots)}_{\sum_ic_i\underbrace{\scriptstyle\hat a_i^\dagger\hat a_i}_{N_i}}\ket{n_1}\dots\ket{n_N}
\]
Symetrizace $\mathcal{N}\hat S\cdot$ \quad $[\hat{\mathbb{C}},\hat S] = 0$ \quad $\to$ \quad
$\hat{\mathbb{C}}\ket{N_1N_2\dots} = \sum_ic_i\hat a_i^\dagger\hat a_i\ket{N_1N_2\dots}$

tj.\ \quad $\hat{\mathbb{C}} = \sum_ic_i\hat a_i^\dagger\hat a_i$ \quad + přechod k jiné bázi \quad
$\hat a_i^\dagger(c) = \sum_j\hat a_j^\dagger(b)\braket{b{:}j}{c{:}i}$
\[
  \hat{\mathbb{C}} = \sum_{jj'}\underbrace{\sum_i\braket{b{:}j}{c{:}i}c_i\braket{c{:}i}{b{:}j'}}_{\mel{b{:}j}{\hat c}{b{:}j'}\equiv c_{jj'}}\hat a_j^\dagger\hat a_{j'}(b)
  \qquad\leftarrow\ \text{vyjádření nezáv.\ na bázi}
\]

\begin{veta}[Jednočásticový operátor ve Fockově prostoru]
tj.\ $\forall$ operátor $\hat{\mathbb{R}} \equiv \sum_i\hat r^{(i)}$: \quad
\[
  \vzorec{\hat{\mathbb{R}} = \sum_{nn'}\underbrace{\mel{n}{\hat r}{n'}}_{R_{nn'}}\hat a_n^\dagger\hat a_{n'}}
\]
$\uparrow$ vyjádření, které platí pro 1-čás.\ operátor bez ohledu na počet částic, tj.\ $\hat{\mathbb{R}} = \sum_{i=1}^{N}\hat r^{(i)}$ na $\forall\mathcal{H}_S^{(N)}\subset\mathcal{H}_F$
\end{veta}

%----------------------------------------------------------------------
\strana{I9}
Podobně \kw{dvoučásticový operátor}
\[
  \hat{\mathbb{R}} \equiv \sum_{i<j}\hat r^{(ij)}
  \quad\dots\ \text{ukážeme za chvíli pro fermiony}\quad
  \longrightarrow\ \frac12\sum_{klmn}\mel{kl}{\hat r^{(ij)}}{mn}\,a_k^\dagger a_l^\dagger a_n a_m
\]
\pozn{šipky: pořadí $m,n$ u anihilačních operátorů je obrácené}

\textcircled{2} \underline{případ Fermionů}
\[
  \mathcal{H}_F \equiv \bigoplus_{N=0}^{\infty}\mathcal{H}_A^{(N)}
\]
na $\forall\mathcal{H}_A^{(N)}$ báze (ortonormální) \quad $\ket{N_1,N_2\dots} \equiv \sqrt{N!}\,\hat A\ket{n_1n_2\dots n_N}$, \\
\hspace*{7.5cm} kde $N_1 + N_2 + \dots = N$ \\
pokud nepořadujeme $\nearrow$ máme bázi na celém $\mathcal{H}_F$

Pauliho princip \dots\ $N_i\in\{0,1\}$ \quad (jinak $\hat A\ket{n_1\dots n_N} = 0$ \dots\ pro $n_i = n_j$)

\begin{definice}[kreační operátor (fermiony)]
\[
  \left(\begin{aligned}
    b_n^\dagger\ket{N_1\dots N_n\dots} &= \sqrt{N_n}\,\ket{N_1\dots N_n+1\dots}\ ? \\
    b_n^\dagger\ket{N_1\dots N_n{=}0\dots} &= \ket{N_1\dots N_n{=}1\dots}\ ?
  \end{aligned}\right)\ \pozn{?}
\]
$\to$ ano, ale jen $N_n = 0$, tj.
\[
  \hat A\ket{n_1\dots n_N} = \frac{1}{\sqrt{N!}}b_{n_1}^\dagger\dots b_{n_N}^\dagger\ket{0} \qquad (*)
\]
\end{definice}
\[
  \hat A\ket{n_1\dots n_i\dots n_j\dots n_N} = -\hat A\ket{n_1\dots n_j\dots n_i\dots n_N}
  \qquad \boxed{\text{vzpomeň: } \hat P_{ij}\hat A = -\hat A}
\]
spec.\ pro dvoučástic.\ stavy: \quad
$\frac{1}{\sqrt2}b_n^\dagger b_m^\dagger\ket{0} = -\frac{1}{\sqrt2}b_m^\dagger b_n^\dagger\ket{0}$ \\
\hspace*{6cm} $\uparrow$ nemůžou komutovat
\begin{itemize}
  \item \dots\ relace \quad $\vzorec{\{b_i^\dagger,b_j^\dagger\} = 0}$ \quad $\leftarrow$ automaticky dá i splnění $\uparrow$ \\
        \hspace*{2.2cm} $\downarrow$
  \item \dots\ hermit.\ sdruž.\ \quad $\{b_i,b_j\} = 0$
\end{itemize}

\underline{pozn:}
\begin{itemize}[nosep]
  \item[--] relace $\{b_i^\dagger,b_j^\dagger\} = 0$ v sobě obsahuje $b_j^{\dagger 2} = 0$ \dots\ tj.\ $N_j < 2$
  \item[--] díky antikom.\ relaci je \quad
        $b_n^\dagger\ket{N_1\dots N_n{=}0\dots} = (-1)^{N_1+\dots+N_{n-1}}\ket{N_1\dots N_n{=}1\dots}$ \quad \pozn{$(*)$}
  \item[--] z $(*)$ plyne, že správně norm.\ vektor je prostě $b_{n_1}^\dagger\dots b_{n_N}^\dagger\ket{0}$, \\
        tj.\ vektory $\forall$ možné obsazení $n_1\dots n_N$ tvoří bázi v $\mathcal{H}_F$
\end{itemize}

\begin{itemize}
  \item \dots\ platí antikom.\ relace \quad $\vzorec{\{b_i,b_j^\dagger\} = \delta_{ij}}$
\end{itemize}
DK: hermit.\ sdruž.\ operátoru splní $(*)$:
\[
  b_n\ket{N_1\dots N_n\dots} = \begin{cases}
    0 & \dots\ \text{pro } N_n = 0 \\
    (-1)^{N_1+\dots+N_{n-1}}\ket{N_1\dots\underbrace{N_n-1}_{0}\dots}
  \end{cases}
\]
\[
  \Rightarrow\quad
  \left.\begin{aligned}
    b_nb_m^\dagger &= \pm\ket{\dots N_m{=}1\dots N_n{=}0\dots}\ \text{nebo } 0 \\
    b_m^\dagger b_n &= \mp\ket{\dots N_m{=}1\dots N_n{=}0\dots}\ \text{nebo } 0
  \end{aligned}\right\}
  \ \{b_n,b_m^\dagger\} = 0 \text{ pro } n\neq m
\]

%----------------------------------------------------------------------
\strana{I10}
zvlášť $n = m$:
\begin{align*}
  b_nb_n^\dagger\ket{N_1\dots N_n\dots} &= \begin{cases}
    0 & \text{pro } N_n = 1 \\
    \underbrace{(-1)^{2(N_1+\dots+N_{n-1})}}_{1}\ket{N_1\dots N_n\dots} & \text{pro } N_n = 0
  \end{cases} \\
  b_n^\dagger b_n\ket{N_1\dots N_n\dots} &= +1\ket{N_1\dots\overset{1}{N_n}\dots}\ \ \text{pro } N_n = 1 \\
  &= 0 \hspace{3.3cm} \text{pro } N_n = 0
\end{align*}
\[ \Rightarrow\quad \{b_n,b_m^\dagger\} = \delta_{nm} \]

\underline{pozn:} $\ket{N_1\dots N_n\dots}$ je vl.\ vektor $b_nb_n^\dagger$ s vl.\ hodn.\ $1 - N_n$ \\
\hspace*{3.6cm} a vl.\ vektor $b_n^\dagger b_n$ s vl.\ hodn.\ $N_n$ \\
\quad tj.\ operátor počtu částic opět \quad $\hat N = \sum_nb_n^\dagger b_n$

\underline{záměna báze} \dots\ projde stejně jako u bosonů: \quad $\tilde b_m^\dagger = \sum_nb_n^\dagger\braket{n}{\tilde m}$ \\
\quad kde 1-částic.\ báze $\hat C\ket{n} = c_n\ket{n}$ se nahradí $\hat{\tilde C}\ket{\tilde m} = \tilde c_m\ket{\tilde m}$

speciálně \uv{polní operátory} \quad $\hat\psi^\dagger(x) = \sum_nb_n^\dagger\braket{n}{x} = \sum_nb_n^\dagger\phi_n(x)^*$

komut.\ relace $\Rightarrow$ \quad
$\{\hat\psi^\dagger(x),\hat\psi^\dagger(x')\} = \{\hat\psi(x),\hat\psi(x')\} = 0$ \\
\hspace*{3.85cm} $\vzorec{\{\hat\psi(x),\hat\psi^\dagger(x')\} = \delta(x - x')}$

\underline{operátory:} \quad 1-částicový \quad
$\hat{\mathbb{R}} = \sum_i\hat r^{(i)} = \sum_{nn'}\underbrace{\mel{n}{\hat r}{n'}}_{R_{nn'}}\hat b_n^\dagger\hat b_{n'}$ \\
\hspace*{2.4cm} spočítáno v $\mathcal{H}^{(1)}$ \dots\ stejné $\forall\hat r^{(i)}$ $\nearrow$

2-částicová interakce: \quad např.\ coulomb.\ repulze \ $\sum_{i<j}\dfrac{e^2}{|r_i - r_j|}$ \\
obecně \quad $\hat{\mathbb{R}} = \sum_{i<j}\hat r^{(ij)}$; \quad kde $r^{(ij)}$ působí jen na částici $i$ a $j$

\begin{align*}
  \hat{\mathbb{R}}\ket{n\dots n_i\dots n_j\dots} &= \sum_{i<j}\mel{n_in_j}{\hat r^{(ij)}}{n_in_j}N_iN_j\ket{n\dots n_i\dots n_j\dots} \\
  \text{\pozn{přepis do báz.\ obsaz.\ čísel, vysvětlit detailně}}\quad
  &= \sum_{k<l}\mel{kl}{\hat r}{kl}\hat N_k\hat N_l\ket{\dots N_k\dots N_l\dots} \\
  &= \frac12\sum_{kl}\mel{kl}{\hat r}{kl}\,b_k^\dagger b_l^\dagger b_lb_k\ket{\dots}
\end{align*}
+ záměna báze \quad $b_k^\dagger = \sum_m\tilde b_m^\dagger\braket{m}{k}$ \\
\hspace*{3cm} $\ket{kl} = \sum_{m_1m_2}\braket{m_1}{k}\braket{m_2}{l}\ket{m_1m_2}$
\[
  \to\quad \hat{\mathbb{R}} = \frac12\sum_{\substack{m_1m_2\\ m_1'm_2'}}
  \underbrace{\Big(\sum_{kl}\braket{m_1}{k}\braket{m_2}{l}\mel{kl}{\hat r}{kl}\braket{k}{m_1'}\braket{l}{m_2'}\Big)}_{\mel{m_1m_2}{\hat r}{m_1'm_2'}}
  \tilde b_{m_1}^\dagger\tilde b_{m_2}^\dagger\tilde b_{m_2'}\tilde b_{m_1'}
\]

%----------------------------------------------------------------------
\strana{I11}
\begin{veta}[Dvoučásticový operátor ve Fockově prostoru]
tj.\ obecně \quad $\hat{\mathbb{R}} = \sum_{i<j}\hat r^{(ij)}$, \quad pak
\[
  \vzorec{\hat{\mathbb{R}} = \frac12\sum_{klmn}\mel{kl}{\hat r}{mn}\,\hat b_k^\dagger\hat b_l^\dagger\hat b_n\hat b_m}
\]
\pozn{pozn:} \quad $\longrightarrow\ ;\ \longrightarrow\ \longleftarrow$ \quad \pozor{pořadí operátorů!}
\end{veta}

\subsection*{Rovnice Hartree--Fock-ovy}
systém $N$ identických fermionů s Hamiltonianem
\[
  \hat H = \hat H_1 + \hat H_2 = \sum_{mn}h_{mn}b_m^\dagger b_n + \frac12\sum_{klmn}v_{kl,mn}b_k^\dagger b_l^\dagger b_nb_m
\]
\underline{Idea:} hledání nejlepší vln.\ funkce ve tvaru
$\underbrace{\overset{1}{\ket{\phi_1}}\overset{2}{\ket{\phi_2}}\dots\overset{N}{\ket{\phi_N}}}_{\ket{\psi}}$ \\
\quad + variační metoda \quad $\longrightarrow$ \ \kw{Hartree} \\
\quad \dots\ nezávislé částice i pro interagující systém, ale vidí \uv{\kw{střední pole}} \\
\quad \dots\ neobsahuje \uv{korelace}, aby $j$-tý elektron \uv{věděl}, kde je $i$-tý \dots\ $\phi_i(x_i)$ \\
\hspace*{9cm} muselo by $\phi_j(x_j,x_i)$

problém \dots\ $\ket{\psi}$ není správná fce pro fermiony, ale:
\[
  \ket{\psi} = \sqrt{N!}\,A\overset{1}{\ket{\phi_1}}\dots\overset{N}{\ket{\phi_N}} = b_{\phi_1}^\dagger\dots b_{\phi_N}^\dagger\ket{0}
  = \frac{1}{\sqrt{N!}}\begin{vmatrix}\phi_1(x_1) & \dots & \phi_1(x_N) \\ \vdots & & \vdots \\ \phi_N(x_1) & \dots & \phi_N(x_N)\end{vmatrix}
\]
tj.\ hledáme nejlepší $\ket{\psi}$ ve tvaru $\uparrow$ ve smyslu Ritz.\ var.\ principu.

abychom mohli užívat formalismus vybudovaný výše, musí navíc
\begin{align*}
  \int\phi_n^*\phi_m \equiv \braket{n}{m} = \delta_{nm} &\quad\dots\ \text{vedlejší podmínka} \\
  \equiv \braket{\phi_n}{\phi_m} &\quad\dots\ \text{minimalizace funkcionálu s vazbou}
\end{align*}

\textcircled{1} \underline{Funkcionál energie}
\[
  E[\phi_1,\phi_2,\dots] = \mel{\psi}{H}{\psi} = \sum_n\big[\mel{\psi}{H_1}{\psi} + \mel{\psi}{H_2}{\psi}\big]
\]
\begin{align*}
  \mel{\psi}{H_1}{\psi} &= \bra{0}b_N\dots b_1\Big(\sum_{mn}^{\infty}h_{mn}b_m^\dagger b_n\Big)b_1^\dagger\dots b_N^\dagger\ket{0} \\
  &= \bra{0}\sum_{m,n=0}^{N}b_m\,h_{mn}b_m^\dagger b_n\,b_n^\dagger\ket{0} \\
  &= \sum_{n=0}^{N}h_{nn}\bra{0}b_n\underbrace{b_n^\dagger b_n}_{N_n}b_n^\dagger\ket{0} = \sum_{n=0}^{N}h_{nn} \equiv \sum_n\mel{\phi_n}{\hat h}{\phi_n}
\end{align*}
podob.
\begin{align*}
  \mel{\psi}{H_2}{\psi} &= \frac12\sum_{klmn}v_{kl,mn}\bra{0}b_N\dots b_1\big(b_k^\dagger b_l^\dagger b_nb_m\big)b_1^\dagger\dots b_N^\dagger\ket{0}
  \quad\text{musí } n,m\in\{1\dots N\} \\
  \text{+ musí být vypárované:}\quad
  &\frac12\sum_{klmn}v_{kl,mn}\big\{\delta_{ln}\delta_{km} + \delta_{lm}\delta_{kn}\big\}\bra{0}b_kb_lb_k^\dagger b_l^\dagger b_nb_mb_l^\dagger b_k^\dagger\ket{0}
\end{align*}

%----------------------------------------------------------------------
\strana{I12}
\begin{align*}
  \text{a)}\quad &\bra{0}b_kb_lb_k^\dagger b_l^\dagger\,\textcolor{red}{\underline{b_lb_k}}\,b_l^\dagger b_k^\dagger\ket{0} = \\
  &\qquad \underbrace{(\delta_{kl} - b_l^\dagger b_k)b_k^\dagger\ket{0} = \delta_{kl}b_k^\dagger\ket{0} - b_l^\dagger\ket{0}}_{\delta_{kl}\ket{0} - \ket{0}\ =\ \ket{0}(\delta_{kl} - 1)} \\
  &= (\delta_{kl} - 1)\bra{0}b_kb_lb_k^\dagger b_l^\dagger\ket{0} = (1 - \delta_{kl}) \\[4pt]
  \text{b)}\quad &\bra{0}b_kb_lb_k^\dagger b_l^\dagger\,\textcolor{red}{\underline{b_kb_l}}\,b_l^\dagger b_k^\dagger\ket{0} = (\delta_{kl} - 1)
  \qquad \textcolor{red}{\nwarrow\ \text{liší se jen pořadí}}
\end{align*}
\[
  \text{tj.:}\quad \mel{\psi}{H_2}{\psi} = \frac12\sum_{kl}\big(v_{kl,kl} - v_{kl,lk}\big)
  = \frac12\sum_{kl}\bra{\phi_k\phi_l}\hat v\big(\ket{\phi_k}\ket{\phi_l} - \ket{\phi_l}\ket{\phi_k}\big)
\]

vedlejší podmínka $\braket{\phi_k}{\phi_l} = \delta_{kl}$ \quad $\to$ Lagrangeovy multiplik.\ $\varepsilon_{kl}$ \\
$\to$ minimalizace funkcionálu $E[\phi_1\dots\phi_N] = \mel{\psi}{H}{\psi}$ za podmínky $\uparrow$

\pozn{mezivýsl.\ -- funkcionál}
\[
  \to\quad \delta\Big[\underbrace{\sum_n^N\mel{\phi_n}{\hat h}{\phi_n}}_{\textcolor{OliveGreen}{\textcircled{\scriptsize 1}}}
  + \underbrace{\frac12\sum_{kl}^N\bra{\phi_k}\bra{\phi_l}\hat v\big(\ket{\phi_k}\ket{\phi_l} - \ket{\phi_l}\ket{\phi_k}\big)}_{\textcolor{OliveGreen}{\textcircled{\scriptsize 2}}}
  - \underbrace{\sum_{kl}^N\varepsilon_{kl}\braket{\phi_k}{\phi_l}}_{\textcolor{OliveGreen}{\textcircled{\scriptsize 3}}}\Big] = 0
\]
změna do prvního řádu při záměně $\ket{\phi_n}\to\ket{\phi_n} + \ket{\delta\phi_n}$:

\textcolor{OliveGreen}{\textcircled{1}} \quad $\delta\sum_n\mel{\phi_n}{\hat h}{\phi_n} = \sum_n\mel{\delta\phi_n}{\hat h}{\phi_n} + \text{c.c.}$

\textcolor{OliveGreen}{\textcircled{2}}
\begin{align*}
  \delta\,\frac12\sum_{kl}\bra{\phi_k}\bra{\phi_l}\hat v\big(\ket{\phi_k}\ket{\phi_l} - \ket{\phi_l}\ket{\phi_k}\big)
  &= \frac12\sum_{kl}\bra{\delta\phi_k}\bra{\phi_l}\hat v\big(\ket{\phi_k\phi_l} - \ket{\phi_l\phi_k}\big) + \text{c.c.} \\
  &\quad + \frac12\sum_{kl}\bra{\phi_k}\bra{\delta\phi_l}\hat v\big(\ket{\phi_k\phi_l} - \ket{\phi_l\phi_k}\big)
\end{align*}
+ symetrie výměnná: \quad
$\bra{\phi}\bra{\psi}\hat v\hat P_{12}\ket{\phi'}\ket{\psi'} = \bra{\phi}\bra{\psi}\hat v\ket{\psi'}\ket{\phi'}$ \\
\hspace*{4.25cm} $= \bra{\phi}\bra{\psi}\hat P_{12}\hat v\ket{\phi'}\ket{\psi'} = \bra{\psi}\bra{\phi}\hat v\ket{\phi'}\ket{\psi'}$ \\
\hspace*{9cm} tj.\ možno současně vyměnit bra i kety \\
+ přeznačení sčít.\ proměnné
\[
  = \sum_{kl}\bra{\delta\phi_k}\bra{\phi_l}\hat v\big(\ket{\phi_k}\ket{\phi_l} - \ket{\phi_l}\ket{\phi_k}\big) + \text{c.c.}
\]

\textcolor{OliveGreen}{\textcircled{3}} \quad člen $\sum_{kl}^N\varepsilon_{kl}\braket{\phi_k}{\phi_l}$ \dots\ jistě máme volbu, $\sum_lU_{kl}\ket{\phi_l}$ změní \\
\quad bázi, ale ta nezmění determinant $A\ket{\phi_1}\dots\ket{\phi_N}$ \\
\quad tj.\ $\exists$ báze, kde \quad $\longrightarrow$ \quad $\sum_k\varepsilon_k\braket{\phi_k}{\phi_k}$

%----------------------------------------------------------------------
\strana{I13}
tj.\ \quad $\delta\sum_k\varepsilon_k\braket{\phi_k}{\phi_k} = \sum_k\varepsilon_k\braket{\delta\phi_k}{\phi_k} + \text{c.c.}$

tj.\ DOHROMADY \quad $\delta(\dots) = 0$ \quad \dots\ \textcolor{OliveGreen}{\textcircled{1}+\textcircled{2}+\textcircled{3}} $\to$

\begin{veta}[Rovnice Hartree--Foka (Hartree--Fock equation)]
\[
  \vzorec{\hat h\ket{\phi_k} + \sum_{\substack{l=1\\ l\neq k}}^{N}\bra{\phi_l^{(2)}}\hat v^{(12)}\Big(\ket{\phi_k^{(1)}}\ket{\phi_l^{(2)}} - \ket{\phi_l^{(1)}}\ket{\phi_k^{(2)}}\Big) = \epsilon_k\ket{\phi_k}}
\]
\textcolor{OliveGreen}{\small $\uparrow$ čísla označují přiřazení 2-částic.\ operátoru $v^{(12)}$}

lze interpretovat jako \fbox{$\hat H_{HF}\ket{\phi_k} = \epsilon_k\ket{\phi_k}$}, ale lépe uvidíme ze zápisu v konkrétní bázi:
\end{veta}

původní 1-částicová báze \dots\ $\hat c\ket{n} = c_n\ket{n}$ \dots\ kreač.\ oper.\ $\bar b_n^\dagger$ \\
nová báze v 1-částic.\ prostoru \ $\hat H_{HF}\ket{\phi_k} = \epsilon_k\ket{\phi_k}$ \dots\ \hspace{1.2cm} $b_k^\dagger$

\underline{připomenout:} H--F rovnice \dots\ podmínka, aby $b_1^\dagger\dots b_N^\dagger\ket{0}$ bylo nejlepší \\
\hspace*{2.7cm} řešení mnohočást.\ S.R.\ ve smyslu var.\ principu

zápis HF rovnic v pův.\ bázi: \quad
$\hat H_1 = \displaystyle\sum_{mn}^{\infty}\bar h_{mn}\bar b_m^\dagger\bar b_n$ \\
\hspace*{5cm} $\hat H_2 = \displaystyle\sum_{mnm'n'}^{\infty}\bar v_{mn,m'n'}\bar b_m^\dagger\bar b_n^\dagger\bar b_{n'}\bar b_{m'}$
\[
  \ket{\phi_k} = \sum_m^{\infty}\ket{m}\braket{m}{\phi_k}
\]
\begin{align*}
  \Rightarrow\quad H_{mn}^{(HF)} \equiv \mel{m}{\hat H_{HF}}{n} = \bar h_{mn}
  &+ \underbrace{\sum_{l=1}^{N}\sum_{m_1m_2}^{\infty}\bar v_{mm_1,nm_2}\braket{m_2}{\phi_l}\braket{\phi_l}{m_1}}_{\textcolor{OliveGreen}{\text{přímá interakce}}} \\
  &- \underbrace{\sum_{l=1}^{N}\sum_{m_1m_2}^{\infty}\bar v_{mm_1,m_2n}\braket{m_2}{\phi_l}\braket{\phi_l}{m_1}}_{\textcolor{OliveGreen}{\text{výměnná interakce}}}
\end{align*}
\hspace*{3.5cm} $\uparrow$ mohu psát $l$ jen $l\neq k$

\textcolor{OliveGreen}{operátor 1-částicové hustoty:} \quad $\hat\rho = \sum_l\ket{\phi_l}\bra{\phi_l}$ \quad \textcolor{OliveGreen}{$\to$ (výskyt výše)}

\underline{interpretace H-F rovnic:} \quad
$\hat H_{HF}\underbrace{\ket{\phi_k}}_{\textcolor{OliveGreen}{\text{orbitaly}}} = \underbrace{\epsilon_k}_{\textcolor{OliveGreen}{\text{1-čás.\ energie}}}\ket{\phi_k}$ \\
\hspace*{3.9cm} \textcolor{OliveGreen}{1-částicové vln.\ fce $\equiv$ orbitaly}

\pozor{pozor:} variační ($\equiv$ celková) energie: \\
\dots\ hezčí v bázi $\ket{\phi_k}$ \dots\ \quad
$v_{kl,k'l'} \equiv \sum\braket{\phi_k}{m}\braket{\phi_l}{n}\bar v_{mn,m'n'}\braket{m'}{\phi_{k'}}\braket{n'}{\phi_{l'}}$ \\
viz předch.\ str.\ \dots \quad
$\vzorec{E = \sum_{n=1}^{N}\epsilon_n - \frac12\sum_{nl=1}^{N}\big(v_{nl,nl} - v_{nl,ln}\big)}$ \\
\hspace*{6.5cm} \textcolor{OliveGreen}{$\uparrow$ korekce na část započtenou $2\times$}

$\hat H_{HF}[\hat\rho]$ \dots\ obsahuje 1-část.\ hamiltonián $\hat h$ \\
\quad \dots\ a \kw{střední pole} skládající se z přímé a výměn.\ části

problém: $\hat H_{HF}$ závisí na $\hat\rho \equiv \sum_l\ket{\phi_l}\bra{\phi_l}$, ale $\ket{\phi_l}$ řešením pro $H_{HF}$ \\
\quad \dots\ \kw{selfkonsistentní úloha} \quad \dots\ iterace

startovací hodnota \dots\ např.\ řešení jen pro $\hat H_1$ \\
\textcircled{1} najít $\forall\ket{\phi_k}$, $\epsilon_k$ pro dané $H_{HF}[\rho]$ \quad
\textcircled{2} považovat nejnižší $\epsilon_k$ za obsazené a spočíst nové $H_{HF}[\rho]$ a pokračovat na \textcircled{1}, dokud nezkonverguje

%----------------------------------------------------------------------
\strana{I14}
\subsection*{HF-rovnice v $x$-reprezentaci}
částice se spinem $\frac12$ \dots\ báze $\ket{x\sigma}$; \ $\sigma = \pm\frac12$ \\
výsledné jednočásticové stavy \dots\ $\phi_{k\sigma}(x) \equiv \braket{x\sigma}{\phi_{k\sigma}}$ \\
pro jednoduchost uvažujme spin.\ nezáv.\ interakci
\[
  \left.\begin{aligned}
    \hat h &= \frac{\hat p^2}{2m} + V(\hat x) \\
    \hat v^{(12)} &= v\big(\hat x_1^{(1)} - \hat x_2^{(2)}\big)
  \end{aligned}\right\}
  \ \Rightarrow\ \mel{x}{H^{(HF)}}{\phi_{k\sigma}}:
\]
\begin{align*}
  \Big[-\frac{\hbar^2}{2m}\Delta_x + V(x) + \int v(x - x')\rho(x')\,\dd x'\Big]\phi_{k\sigma}(x)& \\
  - \sum_l\phi_{l\sigma}(x)\int\phi_{l\sigma}^*(x')v(x - x')\phi_{k\sigma}(x')\,\dd x' &= \epsilon_k\phi_{k\sigma}(x)
\end{align*}
kde \quad $\rho(x) \equiv \sum_{l\sigma}^{N}|\phi_{l\sigma}(x)|^2$ \quad \dots\ jednočásticová hustota

\pozn{pozn: \ $v_{mm_1,nm_2} \equiv \bra{x\sigma}\bra{x_1\sigma_1}v\big(\hat x^{(1)} - \hat x^{(2)}\big)\ket{x'\sigma'}\ket{x_2\sigma_2} = \delta(x - x')\delta_{\sigma\sigma'}\delta(x_1 - x_2)\delta_{\sigma_1\sigma_2}$}

\underline{interpretace:}
\begin{itemize}[nosep]
  \item přímý člen má význam interakce s oblakem okolních nábojů s hustotou $\rho(x)$
  \item výměnný člen $\leftarrow$ přispívají jen $\phi_{l\sigma}$ se stejným spinem \\
        člen je \pozor{nelokální v $x$!}
\end{itemize}

\underline{Poznámky:} \quad metody typu \uv{\kw{meanfield}} \quad \dots\ teorie středního pole
\begin{itemize}
  \item stejný přístup se dá uplatnit pro
  \begin{enumerate}[label=(\alph*)]
    \item rozlišitelné částice \dots\ \kw{Hartree-ho rovnice} \\
          -- stejná rovnice, bez výměnného členu
    \item nerozlišitelné bosony \dots\ \underline{\kw{Gross--Pitaevskii} rovnice} \\
          speciálně pro základní stav obsazena jen nejnižší hladina \dots\ $\forall N$ elektrony v jediném $\psi_0(x)$, tj.\ výměnný člen nepřispívá (měl by znaménko $+$)
  \end{enumerate}
\end{itemize}
\[
  \vzorec{\Big[-\frac{\hbar^2}{2m}\Delta_x + V(x)\Big]\psi_0(x) + \Big[\int|\psi_0(x')|^2v(x - x')\,\dd x'\Big]\psi_0(x) = \epsilon_0\psi_0(x)}
\]
\hspace*{5cm} \textcolor{OliveGreen}{$\uparrow$ \uv{nelineární Schrödingerova rovnice}}

$\psi_0(x)$ \dots\ \uv{\kw{makroskopická vln.\ fce}} \quad \dots\ $\forall N$ částic ve stej.\ stavu \\
\hspace*{9.3cm} (v principu mnoho)



\end{document}
